% concatenated from main.tex
\documentclass[11pt]{amsart}
\usepackage[
  left=2.5cm,
  right=2.5cm,
  top=3cm,
  bottom=3cm
]{geometry}
% =================================================
%                 ENCODING & LANGUAGE
% =================================================
\usepackage[utf8]{inputenc}
\usepackage[T1]{fontenc}
\usepackage[british]{babel}
\usepackage[numbers,sort&compress]{natbib}

% =================================================
%                 MATHEMATICS
% =================================================
\usepackage{amsmath,amssymb,amsfonts,amsthm}


\def\smallsection#1{\smallskip\noindent\textbf{#1}.}
\usepackage{mathtools}
\usepackage{mathrsfs}
\usepackage{bbm}
\usepackage{dsfont}
\usepackage{stmaryrd}
\usepackage{bm}
\usepackage{braket}


% =================================================
%                 GRAPHICS & TIKZ
% =================================================
\usepackage{graphicx}
\usepackage{pgf,tikz}
\usetikzlibrary{positioning,fit,calc,arrows.meta}

% =================================================
%                 TEXT / STRUCTURE
% =================================================
\usepackage{paralist}
\usepackage{comment}
\usepackage{soul}
\usepackage{xcolor}
\usepackage{framed}

% =================================================
%                 BIBLIOGRAPHY
% =================================================
\usepackage{natbib}

% =================================================
%                 HYPERREF (LOAD LAST)
% =================================================
\usepackage{hyperref}
\usepackage{cleveref}

\definecolor{navyblue}{rgb}{0.0,0.0,0.5}
\hypersetup{
  colorlinks=true,
  linkcolor=black,
  citecolor=blue,
  urlcolor=navyblue,
  anchorcolor=blue
}

% =================================================
%                 MACROS
% =================================================
\newcommand{\nn}{\nonumber}
\newcommand{\1}{\mathbbm{1}}
\newcommand{\Tr}{\operatorname{Tr}}
\newcommand{\tr}[1]{\operatorname{Tr}\!\left(#1\right)}
\newcommand{\norm}[1]{\left\lVert#1\right\rVert}
\newcommand{\ceil}[1]{\left\lceil {#1} \right\rceil}

\newcommand{\ketbra}[2]{|#1\rangle\langle#2|}
\newcommand{\kb}[1]{|#1\rangle\!\langle#1|}
\newcommand{\scalar}[2]{\langle#1|#2\rangle}

\setlength{\itemsep}{0pt}
\setlength{\parskip}{0pt}
\setlength{\parsep}{0pt}

% three special one-off environments
\newtheorem{conditionA}{Condition}
\renewcommand{\theconditionA}{A}

\newtheorem{conditionB}{Condition}
\renewcommand{\theconditionB}{B}

\newtheorem{conditionC}{Condition}
\renewcommand{\theconditionC}{C}

\newcommand{\id}{\mathbb{I}}
\newcommand{\R}{\mathbb{R}}
\newcommand{\C}{\mathbb{C}}
\newcommand{\N}{\mathbb{N}}
\newcommand{\Z}{\mathbb{Z}}

\newcommand{\cH}{\mathcal{H}}
\newcommand{\cB}{\mathcal{B}}
\newcommand{\cL}{\mathcal{L}}
\newcommand{\cE}{\mathcal{E}}

\newcommand{\Var}{\mathrm{Var}}
\newcommand{\E}{\mathbb{E}}
\newcommand{\Inf}{\operatorname{Inf}}
\newcommand{\spec}{\operatorname{Sp}}

\newcommand{\eps}{\varepsilon}
\newcommand{\NA}{N_{\mathcal{A}}}
\renewcommand{\epsilon}{\varepsilon}
\renewcommand{\phi}{\varphi}
\renewcommand{\d}{d}
\renewcommand{\L}{\mathcal{L}}
\newcommand{\cA}{\mathcal{A}}
\newcommand{\hatfM}{\widehat{f}_{\!\scalebox{0.6}{$\mathscr{M}$}}
}
\newcommand{\hatfMdelta}{\widehat{f}_{\!\scalebox{0.6}{$\mathscr{M}$}_\delta}}
\newcommand{\fM}{f_{\!\scalebox{0.6}{$\mathscr{M}$}}
}
\newcommand{\fMdelta}{f_{\!\scalebox{0.6}{$\mathscr{M}$}_\delta}}
\newcommand{\fMdeltakappa}{f_{\!\scalebox{0.6}{$\mathscr{M}$}_{\delta,\kappa}}}

\newcommand{\Ntot}{N_{\operatorname{tot}}}
\newcommand{\bin}[1]{}
\newcommand{\red}[1]{\textcolor{red}{#1}}
\newcommand{\rob}[1]{\textcolor{cyan}{[rob: #1]}}
\newcommand{\robi}[1]{\textcolor{cyan}{#1}}

% =================================================
%                 THEOREMS
% =================================================
\theoremstyle{plain}
\newtheorem{theorem}{Theorem}[section]
\newtheorem{lemma}[theorem]{Lemma}
\newtheorem{proposition}[theorem]{Proposition}
\newtheorem{corollary}[theorem]{Corollary}
\newtheorem{conjecture}[theorem]{Conjecture}
\newtheorem{hypothesis}[theorem]{Hypothesis}
\newtheorem{remark}[theorem]{Remark}

\theoremstyle{definition}
\newtheorem{definition}[theorem]{Definition}
\newtheorem{assumption}[theorem]{Assumption}
\newtheorem{notation}[theorem]{Notation}
\newtheorem{condition}[theorem]{Condition}
\newtheorem{problem}[theorem]{Problem}

\theoremstyle{remark}
\newtheorem{example}[theorem]{Example}
\newtheorem{claim}{Claim}[theorem]

\numberwithin{equation}{section}


% =================================================
%                 DOCUMENT
% =================================================
\begin{document}

\title[Infinite-dimensional Quantum Gibbs sampling]{Quantum Gibbs Sampling in infinite dimensions:\\
\tiny{Generation, mixing times and circuit implementation}}


\author{Simon Becker}
\address{Bocconi University, Milan, Italy}
\email{simon.becker@unibocconi.it}



\author{Cambyse Rouz\'e}
\address{Inria, T\'el\'ecom Paris -- LTCI,  
Institut Polytechnique de Paris,  
91120 Palaiseau, France}
\email{cambyse.rouze@inria.fr}

\author{Robert Salzmann}
\address{RWTH Aachen, Department of Physics,  
Otto-Blumenthal-Strasse 20,  
52074 Aachen, Germany\\
Inria, ENS de Lyon, Université Claude Bernard Lyon 1, LIP, 69342, Lyon cedex 07, France}
\email{robert.salzmann@rwth-aachen.de}

\maketitle
\vspace{-0.3cm}

\begin{abstract}
\noindent 
We develop a rigorous and implementable framework for Gibbs sampling of infinite-dimensional quantum systems governed by unbounded Hamiltonians. Extending dissipative Gibbs samplers beyond finite dimensions raises fundamental obstacles, including ill-defined generators, the absence of spectral gaps on natural Banach spaces, and tensions between implementability and convergence guarantees.
We overcome these issues by constructing KMS-symmetric quantum Markov semigroups on separable Hilbert spaces that are both well-posed and efficiently implementable on qubit hardware. Our generation theory is based on the abstract framework of Dirichlet forms, adapted here to the case of algebras of bounded operators over separable Hilbert spaces.
Leveraging the spectral properties of our self-adjoint generators, we establish quantitative convergence results in trace distance, including regimes of fast thermalization.  In contrast, we also identify Hamiltonians for which a naive choice of generators guaranteeing implementability generally comes at the cost of losing convergence of the associated evolutions, thereby establishing a strong trade-off between implementability and convergence.
 Our framework applies to a wide class of models—including Schrödinger operators, Gaussian systems, and Bose–Hubbard Hamiltonians—and provides a unified approach linking rigorous infinite-dimensional analysis with algorithmic Gibbs state preparation.

\end{abstract}
\vspace{-0.3cm}
\tableofcontents





\section{Introduction}
 


Davies semigroups \cite{Davies1974} are among the most prominently used models for the thermalization of quantum systems into their Gibbs states. Formally, for a system at inverse temperature $\beta$ with an associated finite-dimensional Hilbert space $\mathcal{H}$ and Hamiltonian $H$, and a finite set of so-called \emph{bare jump} operators $\{A^\alpha\}_{\alpha\in\mathcal{A}}$ on $\mathcal{H}$ that is closed under taking adjoints, i.e.~$\{A^\alpha\}_{\alpha\in\mathcal{A}}=\{(A^\alpha)^\dagger\}_{\alpha\in\mathcal{A}}$ with trivial commutant, the generator in the Schr\"{o}dinger picture is defined as
\begin{align}\label{eq:Davies}
\mathcal{L}_{\operatorname{D}}(\rho)
&\!= \!\!\!\!\!\sum_{\substack{\alpha\in\mathcal{A}\\\omega\in B(H)}} \!\!\!\!\Upsilon(\omega)
   \Big( A^\alpha_\omega \rho (A^\alpha_\omega)^\dagger
   - \tfrac{1}{2}\{(A^\alpha_\omega)^\dagger A^\alpha_\omega, \rho\} \Big).
\end{align}
Here $B(H)=\operatorname{Sp}(H)-\operatorname{Sp}(H)$ denotes the set of Bohr frequencies of $H$. Given the spectral decomposition $H=\sum_{E\in\operatorname{Sp}(H)}EP_E$, the jump operators take the form
\begin{align*}
A^\alpha_\omega
&:= \sum_{\substack{E,E'\in\operatorname{Sp}(H)\\ E-E'=\omega}}
    P_{E} \, A^\alpha \, P_{E'}.
\end{align*}
The function $\Upsilon:\mathbb{R}\to\mathbb{R}_+$ encodes the rate at which each jump occurs, and in order for the evolution to fix the Gibbs state $\sigma_\beta := \frac{e^{-\beta H}}{\Tr(e^{-\beta H})}$, it satisfies the KMS symmetry condition $\Upsilon(-\omega) = e^{\beta \omega}\,\Upsilon(\omega)$. A well-known difficulty with Davies generators, already present in finite dimensions, is that the jump operators depend on the generally unknown spectral decomposition of $H$. This dependence renders both their circuit implementation and general proofs of convergence for the resulting dynamics particularly delicate. Recently, several alternatives to Davies generators have been proposed whose definitions avoid the spectral decomposition of $H$, making them appealing for quantum simulation tasks \cite{wocjan2023szegedy,rall2023thermal,chen2023quantum}. Previous methods generally achieved only approximate preparation of the target Gibbs state, and only under restrictive assumptions. By contrast, \cite{chen2023efficient} provided the first exact sampler based on Lindbladian dynamics. This result was later complemented by simpler Lindbladian constructions \cite{ding2025efficient,gilyen2024quantum} that require only finitely many jumps.
In \cite{ding2025efficient}, the jump operators are defined as matrix-valued integrals of the form
\begin{align}
L^\alpha
&:= \int e^{itH} A^\alpha e^{-itH}\, {f}(t)\, dt = \sum_{E,E'\in\operatorname{Sp}(H)} \,\widehat{f}(E{-}E')\, P_E A^\alpha P_{E'} = \sum_{\nu\in B(H)}\widehat{f}(\nu)\, A^\alpha_\nu,\label{def:Lalpha}
\end{align}
with a smooth and sufficiently fast decaying \textit{filter function} $f\in L^1(\mathbb{R})$. Note that we take the slightly unconventional definition $\widehat{f}(\nu) =\int_{\mathbb R} f(t) e^{i\nu t} \ dt$ from \cite{ding2025efficient} above. The integral formulation \eqref{def:Lalpha} permits implementation via oracle access to block encodings of the Hamiltonian evolution and bare jumps $A^\alpha$, after time-discretization. The associated Lindblad generator takes the GKLS form

\vspace{- 0.1 in}

\begin{align}\label{eq:GKLS}
\!\!\!\mathcal{L}_{\widehat{f},H}(\rho)
\!&=\! \!-i[B,\rho]
 + \!\!\sum_{\alpha\in\mathcal{A}}\!\! \Big( L^\alpha \rho(L^\alpha)^\dagger
 \!-\! \tfrac{1}{2}\{(L^\alpha)^\dagger L^\alpha, \rho\} \Big),
\end{align}
where the Hermitian operator $B$ is carefully chosen so that the Gibbs state remains a fixed point of the evolution, $\mathcal{L}_{\widehat{f},H}(\sigma_\beta)=0$.


It was recently shown that the convergence properties of the evolution generated by~$\mathcal{L}_{\widehat{f},H}$ are strictly weaker than those of~$\mathcal{L}_{\mathrm D}$ \cite{slezak2026polynomial}, revealing a fundamental tension between \emph{implementability} and \emph{efficiency}. This tension becomes even more pronounced for very large systems, for which the choice of the filter function $f$ becomes key \cite{rouze2024efficient}. This leads to the central question addressed in this paper:
\begin{center}
\textit{Is there a viable way to reconcile \textbf{efficiency} and \textbf{implementability} \\for the dissipative preparation of Gibbs states of infinite-dimensional systems?}
\end{center}

\noindent
To tackle this question, we develop a rigorous framework for Gibbs samplers of infinite-dimensional quantum systems that simultaneously ensures
\begin{itemize}
\item well-posedness of the dynamics, as in infinite dimensions the Lindblad generator may fail to generate a trace-preserving semigroup,
\item spectral convergence guarantees, as convergence properties change dramatically due to the absence of a spectral gap on the trace class operators in this setting, and
\item efficient implementation on finite-dimensional, qubit-base hardware.
\end{itemize}


\smallskip

 \subsection{Gibbs-preserving Markovian dynamics}\label{sec:GibbsMarkovintro}
 
\noindent  It is well-known that unbounded operators that formally satisfy a GKLS-type equation on a natural domain may fail to generate legitimate quantum Markovian dynamics; for instance, the two-photon pure birth process defined with a vanishing Hamiltonian and jump operator $L=(a^\dagger)^2$, where $a^\dagger$ denotes the creation operator over $L^2(\mathbb{R})$, does not preserve the trace \cite[Example~3.3]{chebotarev2000sufficient}. Related pathologies were subsequently identified by Fagnola et al., who proposed a resolution to the problem by imposing additional structural and domain conditions on the generators \cite{fagnola1999quantum,fagnola2000characterization,fagnola2001generators}.
Other approaches to the generation problem include the seminal works of Davies \cite{Davies1977,davies1976quantum,davies1977generator} and Holevo \cite{holevo1993generators,holevo1996quantum}, who established abstract sufficient conditions for unbounded generators of QMSs, albeit those are often difficult to verify in concrete many-body models. More recently, simpler and more explicit sufficient conditions for generation were obtained in \cite{gondolf2024energy} for classes of generators whose jump operators are polynomials in creation and annihilation operators. While well suited to a broad class of continuous-variable models, these results do not apply to the type of generators \eqref{eq:GKLS} considered in the present work. 





Here, instead, we make use of the abstract theory of KMS-symmetric quantum Markov semigroups developed in \cite{Albeverio1977,Davies1992,Goldstein1995} in order to derive our generation theorem. Although the generator we consider is formally identical to that introduced in \cite{ding2025efficient}, additional compatibility conditions relating the Hamiltonian $H$, the jump operators, and the filter function are required in order to establish the well-posedness of the master equation: 
For the moment, we pick an arbitrary filter function $\widehat{f}:\mathbb{R}\to\mathbb{C}$ which simply needs to satisfy the symmetry condition
\begin{align}\label{eq:symmetryintro}
\overline{\widehat{f}(\nu)}=\widehat{f}(-\nu)\,e^{-\beta\nu/2}
\end{align}
and boundedness assumption 
\begin{align}
\sup_\nu\,|\widehat{f}(\nu)|,\   \sup_{\nu}e^{\frac{\beta\nu}{2}}|\widehat{f}(\nu)| <\infty.
\end{align}
In Proposition \ref{propDirichlettoSchro}, we show that these choices lead to a well-defined Lindbladian $\mathcal{L}$ generating a semigroup of quantum channels over a separable Hilbert space $\mathcal{H}$ formally defined as on the right-hand side of \Cref{eq:GKLS},
and with a coherent term $B$ satisfying the following:
for any $E,E'\in \operatorname{Sp}(H)$ with corresponding eigenstates $|E\rangle$, $|E'\rangle$,
\begin{equation}
\begin{split}\label{eq:Bintro} \langle E'&|B|E\rangle=\frac{i}{2}\operatorname{tanh}\Big(\tfrac{\beta(E'-E)}{4}\Big)\sum_{\alpha\in\mathcal{A}}\sum_{\substack{\nu_1,\nu_2\in B(H)\\\nu_2-\nu_1=E'-E}}\overline{\widehat{f}(\nu_1)}\widehat{f}(\nu_2)\,\langle E'|(A^\alpha)^\dagger P_{E+\nu_2}A^\alpha |E\rangle. 
\end{split}
\end{equation}
\smallskip

\noindent At this stage, the generator $\mathcal{L}_{\widehat{f},H}$ still seems to depend on the spectral decomposition of $H$. To circumvent this, we need representations of the jumps $L^\alpha$ and of the coherent term $B$ that are independent of the spectral decomposition of $H$. For this, we first 
 assume that $\widehat{f}$ is Schwartz, so that it has a smooth and rapidly decaying Fourier transform, giving the jumps a potentially discretizable integral formula as in the first identity of \Cref{def:Lalpha}. Next, we find a spectrally agnostic representation of the coherent part $B$. In finite dimensions, \cite{ding2025efficient} considered the drift part $G\rho+\rho \,G^\dagger$, with 
\begin{align*}
G&:=-iB-\frac{1}{2}\sum_\alpha (L^\alpha)^\dagger L^\alpha  =\sum_{\alpha\in\mathcal{A}}\sum_{\nu\in B(H)}g(\nu)\big((L^\alpha)^\dagger L^\alpha\big)_\nu,
\end{align*}
with $\widehat{g}(\nu):=-\frac{1}{2}\big(\operatorname{tanh}(-\beta \nu/4)+1\big) $. As such, since the function $\widehat{g}$ does not possess a smooth Fourier transform, $G$ does not seem to admit an integral representation.
In \cite{ding2025efficient}, the issue was resolved by introducing a compactly supported cut-off function $\kappa(\nu)=1$ for $\nu\in [-2\|H\|,2\|H\|]$, and denoting 
$
\widehat{g}_\kappa(\nu):=\widehat{g}(\nu)\,\kappa(\nu).
$
This way, $G$ coincides with the operator defined by replacing $\widehat{g}$ with $\widehat{g}_\kappa$. Moreover, since $\widehat{g}_\kappa$ admits a Fourier transform $g_\kappa$, we can write
\begin{align*}
G=\sum_{\alpha\in \mathcal{A}}\int_{\mathbb{R}} g_\kappa(t) \,e^{iHt}\big((L^\alpha)^\dagger L^\alpha\big) e^{-itH}\,dt.
\end{align*}
When $H$ is unbounded, a sharp cutoff function $\kappa$ regularizing the function $\widehat{g}_\kappa$, without modifying the generator $G$, no longer exists. Moreover, imposing a cut-off would eventually lead to a loss of the Gibbs-preservation property of the evolution generated by $\mathcal{L}_{\widehat{f},H}$. 

\smallskip
To resolve these conflicting constraints, we consider the following Gaussian weighted version of the generator $\mathcal{L}_{\widehat{f},H}$, which has already been considered in finite dimensions in \cite{gilyen2024quantum,slezak2026polynomial}: for a given width $\sigma_E\ge 0$,
\begin{align}
%\label{eq:DefLsigma_E}
&\mathcal{L}_{\sigma_E,\widehat{f},H}(\rho)\!:=\!\sum_\alpha\sum_{\nu_1,\nu_2}\!e^{-\frac{(\nu_1-\nu_2)^2}{8\sigma_E^2}}\,\overline{\widehat{f}(\nu_1)}\widehat{f}(\nu_2) \Big(\!\!-i[B^\alpha_{\nu_1,\nu_2},\rho] +A^\alpha_{\nu_2}\rho (A^\alpha_{\nu_1})^\dagger\!-\!\frac{1}{2}\{(A^\alpha_{\nu_1})^\dagger A^\alpha_{\nu_2},\rho\}\Big),\nonumber
\end{align}
where 
\begin{align*}
B^\alpha_{\nu_1\nu_2}
&=\frac{i}{2}\tanh(\beta(\nu_1-\nu_2)/4)\,(A^\alpha_{\nu_1})^\dagger A^\alpha_{\nu_2}.
\end{align*}
This formal definition is shown to generate a semigroup of quantum channels with a unique fixed state $\sigma_{\beta}$ in \Cref{monotonicitygap}. The semigroup is, in fact, KMS-symmetric with respect to $\sigma_\beta$, which means that for any $t$, the Heisenberg-dual, weak$^*$-continuous semigroup generated by $\mathcal{L}_{\smash{\sigma_E,\widehat{f},H}}^\dagger$, is self-adjoint with respect to the so-called KMS scalar product on the bounded operators on $\mathcal{H}$
\begin{align}\label{eq:KMSsymmetry}
\langle \bullet,\bullet \rangle_{\sigma_\beta}:(X,Y)\mapsto \Tr\Big(\sigma_\beta^{\frac{1}{2}}X^\dagger \sigma_\beta^{\frac{1}{2}}Y\Big).
\end{align}
Moreover, $\mathcal{L}_{\smash{\sigma_E,\widehat{f},H}}$ interpolates between the generator $\mathcal{L}_{\smash{\widehat{f},H}}$ at  $\sigma_E=\infty$ and the Davies generator $\mathcal{L}_{\operatorname{D}}$ associated with $H$ at a rate of $\Upsilon=|\widehat{h}|^2$ at $\sigma_E=0$. Most importantly, the added Gaussian envelope helps in regularizing the drift: in Proposition \ref{integralrep}, we show that
\begin{align}\label{eq:LsigmaEintro}
\mathcal{L}_{\sigma_E,\widehat{f},H}(\rho)=G_{\sigma_E}\,\rho+\rho\, G_{\sigma_E}^\dagger+\Phi_{\sigma_E,\widehat{f},H}(\rho),
\end{align}
with 
\begin{align*}
 G_{\sigma_E}:=-\sum_{\alpha\in\mathcal{A}}\int_{-\infty}^\infty g(t)\,e^{-itH}((L^{\alpha})^\dagger L^\alpha)e^{itH}dt,
\end{align*}
with 
$g(t)=\frac{1}{2\pi}\int_{-\infty}^\infty \frac{e^{-\nu^2/8\sigma_E^2}}{1+e^{\beta\nu/2}} e^{-i\nu t}d\nu$, $X^{\alpha}_s:=e^{isH}L^\alpha e^{-isH}$, and a CP map 
\begin{align*}
\Phi_{\sigma_E,\widehat{f},H}(\rho):=\sigma_E\sqrt{\frac{2}{\pi}}\sum_{\alpha}\int_{\mathbb{R}}e^{-2\sigma_E^2 s^2}\, \, X_s^{\alpha} \rho (X_s^\alpha)^\dagger \,ds.
\end{align*}
In summary, whenever the function $f$ is Schwartz, we can construct a spectral-agnostic generator that exactly fixes the Gibbs state of $H$ at the inverse temperature $\beta$. 

 % \subsection{Efficient implementation}\label{subsec:efficientimplem}


\smallskip

 \subsection{Convergence guarantees via spectral analysis}\label{sec.fastconv}

\noindent Next, we argue that the evolution generated by $\mathcal{L}_{\sigma_E,\widehat{f},H}$ converges to its fixed point $\sigma_\beta$.
Given a subset $\mathscr{S}$ of input quantum states, we define the mixing time by
\[
t_{\operatorname{mix}}(\epsilon,\mathscr{S})
\!:=\!\inf\Bigl\{ t \ge 0 \Big| \| \rho_t - \sigma_\beta \|_1 \le \epsilon \,\forall\rho \in \mathscr{S}\! \Bigr\}
\]
where $\rho_t$ denotes the state $e^{t\mathcal{L}_{\sigma_E,\widehat{f},H}}(\rho)$ evolved according to the semigroup generated by $\mathcal{L}_{\sigma_E,\widehat{f},H}$, and $\sigma_\beta$ is its unique stationary state. 
Establishing polynomial-time convergence for Gibbs sampling dynamics is notoriously challenging, as the mixing time depends delicately on both the Hamiltonian $H$ and the inverse temperature $\beta$. In \cite{rouze2024efficient}, it was first proved that high-temperature Gibbs states of geometrically local Hamiltonians with locally finite-dimensional constituents can be prepared in time $t_{\mathrm{mix}}=\mathcal{O}(n)$; this was later improved in \cite{rouze2411optimal} to the optimal bound $t_{\mathrm{mix}}=\mathcal{O}(\log n)$. These mixing results were subsequently extended to several low-temperature regimes, including spin chains \cite{https://doi.org/10.48550/arxiv.2510.08533}, certain CSS codes above critical thermodynamic temperatures \cite{stengele2025modified}, and perturbations of Gaussian fermionic models \cite{vsmid2025rapid,vsmid2025polynomial,tong2025fast}. In contrast, exponential lower bounds below critical temperatures were also obtained through quantum extensions of the bottleneck lemma \cite{gamarnik2024slow}. Nevertheless, all such existing approaches currently rely fundamentally on the boundedness of the generator.


\smallskip
First, we contemplate the possibility of getting uniform convergence of the evolution over the entire set $\mathscr{S}\equiv \mathscr{S}(\mathcal{H})$ of states: given a strongly continuous semigroup over a Banach space $\mathscr{B}$, a necessary condition for it to converge in norm is that the spectrum of its generator $A$ is gapped, meaning that there exists a positive constant $\lambda_0>0$ such that for all $\lambda\in\operatorname{Sp}(A)\backslash \{0\}$ we have $\operatorname{Re}(\lambda)\le -\lambda_0$ \cite[Corollary 4.1.2]{salzmann2024robustness}.
 A natural Banach space on which one would want to study the dynamics generated by $\mathcal{L}_{\sigma_E,\widehat{f},H}$ is the space of trace-class operators over $\mathcal{H}$, and the presence of a gap would thus imply that the dynamics converges to $\sigma_\beta$ uniformly over the set of all input states. Unfortunately, this condition often fails, even in the simplest settings (see Proposition \ref{thm.qOUnogap}). This first observation, in sharp contrast with the finite-dimensional setting, strongly suggests the need to restrict the set $\mathscr{S}$ of input states.

\smallskip

Here again, the imposed KMS-symmetry condition \eqref{eq:KMSsymmetry} can be used to our advantage. Indeed, this condition implies that we can associate with the generator $\mathcal{L}_{\sigma_E,\widehat{f},H}$ a self-adjoint operator $L_{\sigma_E,\widehat{f},H}$ on the Hilbert space $\mathscr{T}_2$ of Hilbert-Schmidt operators on $\mathcal{H}$, via the defining property that for any $x\in \mathscr{T}_2$, 
\begin{align*}
\sigma_\beta^{\frac{1}{4}}e^{tL_{\sigma_E,\widehat{f},H}}(x)\sigma_\beta^{\frac{1}{4}}=e^{t\mathcal{L}_{\sigma_E,\widehat{f},H}}(\sigma_\beta^{\frac{1}{4}}x\sigma_\beta^{\frac{1}{4}}).
\end{align*}
By a standard use of H\"{o}lder's inequality, we conclude that for input states of the form $\rho=\sigma_{\beta}^{\frac{1}{4}}x\sigma_\beta^{\frac{1}{4}}$,
\begin{align*}
\big\|e^{t\mathcal{L}_{\sigma_E,\widehat{f},H
}}(\rho)-\sigma_{\beta}\big\|_1\le \|e^{tL_{\sigma_E,\widehat{f},H}}(x)-\sigma_{\beta}^{\frac{1}{2}}\|_2.
\end{align*}
Moreover, for a gapped self-adjoint semigroup, the above norm decays exponentially in $t$ with a rate $\lambda_2\equiv \operatorname{gap}(L_{\sigma_E,\widehat{f},H})$ given by the largest value for which $L_{\sigma_E,\widehat{f},H}$ satisfies the condition that $\operatorname{Re}(\lambda)\le -\lambda_2$ for all $\lambda\in \operatorname{Sp}(L_{\sigma_E,\widehat{f},H})\backslash \{0\}$, so we obtain the convergence
\begin{align*}
\big\|e^{t\mathcal{L}_{\sigma_E,\widehat{f},H}}(\rho)-\sigma_{\beta}\big\|_1\!\le\! e^{-\lambda_2t}\|x-\sigma_\beta^{\frac{1}{2}}\|_2\le e^{-\lambda_2 t}\|x\|_2,
\end{align*}
where in the last bound we also used that $x-\sqrt{\sigma_\beta}$ is orthogonal to $\sqrt{\sigma_\beta}$ due to the normalization of the state $\rho$. Thus, on the set $\mathscr{S}\equiv \mathscr{S}_2$ of states of the form $\rho={\sigma_\beta}^{1/4}x{\sigma_\beta}^{1/4}$ with $x\in \mathscr{T}_2$, the dynamics converges exponentially fast to $\sigma_\beta$. Interestingly, in \Cref{lem.gapmonoton} we also find that the spectral gap is monotonically decreasing with $\sigma_E$ (see also \cite{slezak2026polynomial}):
\begin{align*}
\operatorname{gap}(L_{\sigma_E,\widehat{f},H})\nearrow\qquad 
\text{ as } \qquad \sigma_E\searrow \,,
\end{align*}
which justifies the intuition that the Davies dynamics converges faster than that of \cite{ding2025efficient}.
In Theorem \ref{thm:fromL2toLp}, we push the analysis and prove that for any $p\in (1,\infty)$, there exists a constant $M_p<\infty$ such that on the set $\mathscr{S}\equiv \mathscr{S}_p$ of states of the form $\rho=\sigma_{\beta}^{1/2\hat{p}}x\sigma_{\beta}^{1/2\hat{p}}$  for some operator $x$ in the Schatten $p$ class $\mathscr{T}_p$ and $\hat{p}$ the H\"{o}lder conjugate of $p$  
\begin{align}\label{eq:Tpconvergence}
\|e^{t\mathcal{L}_{\sigma_E,\widehat{f},H}}(\rho)-\sigma_\beta\|_1\le M_p\, e^{-\lambda_2 t}\|x\|_{p}.
\end{align}
Thus, while we cannot prove uniform convergence for $p=1$, we can get arbitrarily close to it. The constraint that $\|x\|_p<\infty$ coincides with the condition that the Sandwiched Rényi divergence $\widehat{D}_p(\rho\|\sigma_\beta)<\infty$, and \eqref{eq:Tpconvergence} is equivalent to
\begin{align}\label{convgapDp}
\|e^{t\mathcal{L}_{\sigma_E,\widehat{f},H}}(\rho)-\sigma_\beta\|_1\le M_p\, e^{-\lambda_2 t}\,e^{\hat{p}\widehat{D}_p(\rho\|\sigma_\beta)}.
\end{align}
In \Cref{sec:spectralgaps}, we show that the generator $L_{\sigma_E,\widehat{f},H}$ is gapped for various models: first in \Cref{sec:gaussexamples}, we consider Gaussian models with quadratic Hamiltonians. All these results are valid for a Schwartz filter function $f$. In contrast, we show in \Cref{prop.no-gogap} that, beyond quadratic models and perturbations thereof, the gap closes for such $f$. This is generally a consequence of the induced decay of $\widehat{f}$ as $\nu\to-\infty$. Instead, in \Cref{sec:GeneralNumberviaBirthDeath}, we choose the Metropolis-type filter function already considered in \cite{ding2025efficient}
\begin{align}
\label{eq:filterFunction}
    \hatfM(\nu) = \exp\left(-\frac{\sqrt{1+(\beta\nu)^2}+ \beta\nu}{4}\right).
\end{align}


\noindent For this function, we show in Theorem \ref{thm:SpectralGapGeneralh(N)} that in the case of a single-mode bosonic system with the associated total photon number observable $N$, given $H=h(N)$ for some eventually non-decreasing function $h$ with large enough energy differences, the generator $L_{\sigma_E,\hatfM,H}$ remains gapped. 
In the companion paper \cite{BeckerRouzeSalzmannBose}, we derive analogous results for Bose--Hubbard models.



\bin{\subsection{Efficient implementations}


\noindent Our discussion thus far clearly illustrates the tension between the rate at which the sampler generated by
$\mathcal{L}_{\sigma_E,\widehat{f},H}$ converges to its equilibrium Gibbs state and the amount of resources required for its implementation as a quantum circuit: while smooth, rapidly decaying functions are required to justify the integral representation \eqref{def:Lalpha} of the jumps \(L^\alpha\)—with faster decay leading to shorter circuit implementations—such conditions on \(f\) result in gapless generators, a signature of non-convergent dynamics. }








\bin{ the present continuous variable setting, the bare jumps $a_i,a_i^\dagger$ are unbounded, which renders the question of their block encoding non-trivial.  }



\subsection{Efficient implementation}

In order to implement the sampler generated by the unbounded Lindbladian  $\cL_{\sigma_E,\widehat f,H}$ we consider in Section~\ref{sec:FiniteTruncBareJumps} a finite-dimensional approximation scheme which we then use to obtain an efficient circuit implementation on a qubit-based quantum computer in Section~\ref{sec:FullFiniteDimPipeline}. For that we consider for each truncation level $M\in\N$ a finite rank projection  $P_M$ and for each $\alpha\in\cA$ an $(M+1)$-rank projection $\pi^\alpha_M.$ Here, the truncated subspace $\operatorname{im}(P_M)$ serves as the \emph{system register} of the quantum device on which we aim to implement the Gibbs sampler approximately. For instance, in many-body or multi-mode systems, the number of bare jumps $|\mathcal{A}|$ is typically proportional to the number of particles or modes. Accordingly, one usually considers truncations for which the local register space is associated with the image of $\pi^\alpha_M$ and the dimension of the full system register satisfies
$
\log\left( \dim\bigl(\operatorname{im}(P_M)\bigr)\right)
= \mathcal{O}\bigl(|\mathcal{A}|\log(M)\bigr),
$ although, at this stage, we are in principle free to leave the dimension of the system register unspecified.

Using the projections $\pi^\alpha_M$ and $P_M$ we can define finite-dimensional truncations of the bare jumps and the Hamiltonian as 
\begin{align}
\left(A^\alpha\right)^{\le M} := \pi^\alpha_M A^\alpha\pi^\alpha_M\qquad\text{and}\qquad  H_{\le M} := P_M H P_M.
\end{align}
This allows us to define the finite-dimensional Lindblad generator
$\mathcal{L}^{\smash{\leq M}}_{\smash{\sigma_E,\widehat f,H_{\leq M}}}$
by replacing the bare jump operators and the Hamiltonian in the unbounded generator
$\mathcal{L}_{\smash{\sigma_E,\widehat f,H}}$
with $\bigl(A^\alpha\bigr)^{\smash{\leq M}}$ and $H_{\smash{\leq M}}$, respectively. If, in addition, the compatibility condition
$
\bigl(A^\alpha\bigr)^{\smash{\leq M}}\operatorname{im}(P_M)\subseteq \operatorname{im}(P_M)$
is satisfied, then
$\mathcal{L}^{\smash{\leq M}}_{\smash{\sigma_E,\widehat f,H_{\leq M}}}$
generates a quantum Markov semigroup on the finite-dimensional system register.

In Section~\ref{sec:FiniteTruncBareJumps}, we show that this finite-dimensional generator provides a good approximation to the target generator $\mathcal{L}_{\sigma_E,\widehat f,H}$ for large truncation parameter $M$. Since $\mathcal{L}_{\smash{\sigma_E,\widehat f,H}}$ is typically unbounded, the approximation has to be understood pointwise on suitably energy-constrained states. To make this precise, we introduce a self-adjoint, positive semidefinite energy observable $\NA$ and consider states whose expectation values with respect to suitable exponential weights in $\NA$ are finite. We then assume that both the truncated bare jumps and the truncated Hamiltonian approximate their original counterparts well on such inputs, with errors that are exponentially small in $M$ up to polynomial prefactors, see~\eqref{eq:AbstractFiniteTruncBareJumps} and~\eqref{eq:HamiltonianLeakTrunc}. We further require that the Hamiltonian evolution is compatible with the same energy constraint, in the sense that it drives low-energy states into higher-energy sectors only at a controlled exponential rate, c.f.~\eqref{eq:ExpBound}. As discussed in Sections~\ref{sec:TruncateJumps} and \ref{sec:FinitDimGenerator} all of these assumptions are naturally satisfied for multi-mode bosonic systems with bare jumps being the creation and annihilation operators, certain Hamiltonians constructed as bounded degree polynomials in $a_i$ and $a^\dagger_i$ and for the choice $\NA \equiv \sum_{i} a^\dagger_ia_i.$   

Under these assumptions, we show in Theorem~\ref{thm:SchwartzGibbsDynamics} for Schwartz filter functions $\widehat{f}$ and input states $\rho$ satisfying
\begin{align}
\label{eq:rhocIntro}
    \rho \le \mathfrak{c} \,\sigma_\beta
\end{align}
for some $\mathfrak{c}\ge 1$ that for evolution time $t\ge 0$ and accuracy $\eps>0$ we can achieve\footnote{Here, the $\widetilde{\mathcal{O}}$ notation hides constants independent of the displayed parameters and additionally suppresses subdominant $\operatorname{poly}\log\log$ factors.}  \begin{align}\label{eq:MetalIntro}
       &\left\|\left(e^{t\cL_{\sigma_E,\widehat f,H}}-e^{t\cL^{\le M}_{\sigma_E,\widehat f,H_{\le M}}}\right)(\rho)\right\|_1\le \eps,\qquad\text{with}\qquad M = \widetilde{\mathcal{O}} \left(\operatorname{poly}\left(\log\left(\frac{t\,\mathfrak{c}\,E_{\operatorname{Gibbs}}|\cA|}{ \eps}\right)\right)\right).
    \end{align}
Here, $E_{\operatorname{Gibbs}}$ denotes the expectation value of the mentioned exponential energy observable involving $\NA$ with respect to the Gibbs state.

In Section~\ref{sec:FullFiniteDimPipeline}, we then combine this result with the work of \cite{chen2023quantum,gilyen2024quantum,chen2023efficient,ding2025efficient}, in particular \cite[Theorem 18]{ding2025efficient}, which provide efficient circuit implementations of such finite dimensional Lindbladian dynamics by approximating the involved time integrals via linear combinations of unitaries \cite{LiWang23} given oracle access to the Hamiltonian evolution $e^{\smash{-isH_{\le M}}}$  and block encoding of the bare jumps of the truncated bare jumps $\left(A\right)^{\smash{\le M}}$. In particular, provided a state preparation circuit for input state $\rho$ in \eqref{eq:rhocIntro}, we find in Theorem~\ref{thm:CircuitDynamicsSchwartz} that
$e^{\smash{t\cL_{\sigma_E,\widehat f, H}}}(\rho)$ can be prepared on the finite dimensional system register within $\eps$-trace distance with order
\begin{align*}
\widetilde{\mathcal{O}}\left( \,t\,\operatorname{poly}\left(|\mathcal{A}|\,,\, \log\left(\frac{\mathfrak{c}\,E_{\operatorname{Gibbs}}}{\epsilon}\right)\right)\right) \text{ total Hamiltonian simulation time corresponding to }H_{\le M}.
\end{align*}
Therefore, given positivity of spectral gap, $\lambda_2\equiv \operatorname{gap}(L_{\sigma_E,\widehat f, H})>0,$ we show in Corollary~\ref{cor:GibbsCircuitSchwartz} that the Gibbs state of the Hamiltonian $H$ can be prepared via a finite dimensional circuit using order
\begin{align*}
\widetilde{\mathcal{O}}\left(\frac{1}{\lambda_2}\operatorname{poly}\left(|\mathcal{A}|\,,\,\log\left(\frac{\mathfrak{c}E_{\operatorname{Gibbs}}}{\epsilon}\right)\right)\right)\text{ Hamiltonian simulation time with respect to } H_{\le M}.
\end{align*}
As discussed above, for many Hamiltonians of interest it is necessary to move beyond Schwartz filter functions to obtain a positive spectral gap and consider instead filter functions such as  \eqref{eq:filterFunction}.
To extend our implementation theory to this choice, we consider in Section~\ref{sec:singular} for parameter $\delta>0$ the regularisation
\begin{align}\label{eq:hatfMdeltaeqIntro}
    \hatfMdelta(\nu) := \hatfM(\nu)e^{-\delta\eta_{2,\theta}(\nu)}\qquad \text{with} \qquad  \eta_{2,\theta}(\nu) := e^{(1+(\beta\nu)^2)^\theta} \qquad \text{ and }\qquad \theta\in(0,1/2).
\end{align}
 By a simple continuity bound argument, we show in Proposition~\ref{prop:DeltaTo0} closeness of the unbounded generators $\cL_{\smash{\sigma_E,\hatfM,H}}$ and $\cL_{\smash{\sigma_E,\hatfMdelta,H}}$ on certain energy-constrained states for small $\delta.$ As $\hatfMdelta$ is again a Schwartz function, we can combine this result with the finite-dimensional approximation scheme outlined above. In this way, Theorem~\ref{thm:SingularGibbsdynamics} shows that, given a state-preparation circuit for the state $\rho$ in~\eqref{eq:rhocIntro}, the state
$e^{\smash{t\mathcal{L}_{\sigma_E,\hatfMdelta, H}}}(\rho)$
can be prepared on the finite-dimensional system register with essentially the same resource requirements as those described above. In the case of positive spectral gap $\lambda_2\equiv \operatorname{gap}(L_{\smash{\sigma_E,\hatfM,H}})>0$, we analogously find in Corollary~\ref{cor:GibbsCircuitSingular} that the Gibbs state of the Hamiltonian $H$ can be prepared by a finite dimensional circuit given the same resource requirements as in the case for Schwartz filter functions.  

For many-body continuous-variable quantum systems of interest, both the quantity $\mathfrak{c}$ in~\eqref{eq:rhocIntro} and the Gibbs energy $E_{\operatorname{Gibbs}}$ typically scale exponentially with the number of particles or modes, much like the partition function of the Gibbs state. Owing to the logarithmic dependence of the above complexity bounds on these quantities, which stems from the exponential energy constraint underlying the finite-dimensional approximation scheme, the resulting implementations of the Lindblad dynamics and Gibbs state preparation remain efficient, with resource requirements that scale polynomially in the number of particles or modes.



\bigskip

\smallsection{Acknowledgement}
SB would like to thank Lin Lin for fruitful discussions on the Gibbs sampling of powers of the number operator. He would also like to thank Jeff Galkowski and Maciej Zworski for bringing the issue of defining Gibbs dynamics for Schr\"odinger operators to his attention. This led to Theorem \ref{thm:Schroedinger_operators}.
SB would also like to acknowledge support from the SNF Grant PZ00P2\_216019. 
CR is supported by France 2030 under the
French National Research Agency award number ''ANR-22-EXES-0013''. RS acknowledges support by the European Research Council (ERC Grant Agreement No.~948139 and ERC Grant AlgoQIP, Agreement No. 851716), from the Excellence Cluster Matter and Light for Quantum Computing (ML4Q-2), from the QuantERA II Programme of the
European Union’s Horizon 2020 research and innovation programme under Grant Agreement No
101017733 (VERIqTAS) as well as the government grant managed by the Agence Nationale de la
Recherche under the Plan France 2030 with the reference ANR-22-PETQ-0007. 



\section{KMS-symmetric generators in infinite dimensions}\label{sec:generalframe}


\noindent In this section, we construct a family of quantum Gibbs samplers for infinite-dimensional quantum systems. Given a densely defined self-adjoint operator $(H,D(H))$ on a separable Hilbert space $\mathcal{H}$ with $H\ge -h_0 I$, $h_0\ge 0$, and inverse temperature $\beta>0$, we aim to prepare the corresponding Gibbs state of the form
\begin{align*}
\sigma_\beta:=\frac{e^{-\beta H}}{\mathcal{Z}(\beta)},\qquad \text{ with }\qquad \mathcal{Z}(\beta):=\Tr(e^{-\beta H})<\infty,
\end{align*}
for Hamiltonians for which the trace in the previous line is finite.
We recall that the finiteness of the partition function $\mathcal{Z(\beta)}$ for all $\beta>0$, also known as the \textit{Gibbs hypothesis}, directly implies that the spectrum is unbounded above, that the spectrum is discrete, and that the energy levels cannot become \emph{too dense} with increasing energy values. In contrast to the finite-dimensional setting, where generators of quantum dynamical semigroups have long been fully characterized \cite{Lindblad1976,Gorini1976}, extensions to unbounded jumps are more intricate. Here, we leverage a certain detailed balance condition that will allow us to define our evolutions following the abstract theory of KMS-symmetric quantum Markov semigroups developed in \cite{Albeverio1977,Davies1992,Goldstein1995,Cipriani1997}.

We consider a finite set $\{A^\alpha\}_{\alpha\in\mathcal{A}}$ of closed, densely defined jump operators $A^\alpha$ with a common domain $D\subseteq \cH$ that is invariant under taking adjoints, i.e. $\{A^\alpha\}_{\alpha\in \mathcal{A}}=\{(A^\alpha)^\dagger\}_{\alpha\in \mathcal{A}}$. We also assume that $D$ includes all the eigenstates of $H$, and we require the following condition throughout the paper
% \begin{framed}
% \begin{conditionA}\label{eq:condAalphas}
% There exist some moments $0\le\gamma,  \mu,\zeta$, with $\gamma\le \mu$ and $\gamma+\mu\le \zeta$, as well as a constant $C>0$ such that, denoting $\widetilde{H}:=H+(h_0+1)I$,
% \begin{align}
% \label{eq:BoundBareJumpsWithHam}
% \|A^\alpha \widetilde{H}^{-\gamma}\|,\quad \|\widetilde{H}^\gamma\,A^\alpha \widetilde{H}^{-\mu}\|,\quad 
%  \|\widetilde H^{\mu+\gamma}A^\alpha\widetilde  H^{-\zeta}\|\le C.
% \end{align}
% Next, we consider a function $\widehat{f}:\mathbb{R}\to\mathbb{C}$ with 
% \begin{align}
% \overline{\widehat{f}(\nu)}=\widehat{f}(-\nu)\,e^{-\beta\nu/2}\qquad \forall\nu\in\mathbb{R}.
% \end{align}



% Throughout this work, we also assume that the function $\widehat{f}$ satisfies these boundedness conditions:
% \begin{align}
% \sup_\nu\,|\widehat{f}(\nu)|,\quad \sup_{\nu}e^{\frac{\beta\nu}{4}}|\widehat{f}(\nu)|,\quad  \sup_{\nu}e^{\frac{\beta\nu}{2}}|\widehat{f}(\nu)| \le C'
% \end{align}
%  for some constant $C'>0$. Finally, we assume the following summability condition on the eigenvalues of $H$:
%  \begin{align}
% \sum_{E\in\operatorname{Sp}(H)} (E+h_0+1)^{-\gamma}<\infty\,.
%  \end{align}
 
% \end{conditionA}
% \end{framed}

\begin{framed}
\begin{conditionA}\label{eq:condAalphas}
There exist some moments $0\le\gamma,  \mu$, with $\gamma\le \mu$, as well as a constant $C>0$ such that, denoting $\widetilde{H}:=H+(h_0+1)I$,
\begin{align}
\label{eq:BoundBareJumpsWithHam}
\|A^\alpha \widetilde{H}^{-\gamma}\|,\quad \|\widetilde{H}^\gamma\,A^\alpha \widetilde{H}^{-\mu}\|\le C.
\end{align}
Next, we consider a function $\widehat{f}:\mathbb{R}\to\mathbb{C}$ with 
\begin{align}
\label{eq:KMSFilterFunction}
\overline{\widehat{f}(\nu)}=\widehat{f}(-\nu)\,e^{-\beta\nu/2}\qquad \forall\nu\in\mathbb{R}.
\end{align}
We also assume there is a constant $C'>0$ such that
\begin{align}
\sup_\nu\,|\widehat{f}(\nu)|,\quad  \sup_{\nu}e^{\frac{\beta\nu}{2}}|\widehat{f}(\nu)| \le C'.
\end{align}
%  Furthermore, we assume that for all $E\in\spec(H)$ we have
%  \begin{align}
%  \label{eq:SumsFin1}
% F_1(E) := \sum_{E'\in\spec(H)} |\widehat f(E'-E)| < \infty\quad\text{and}\quad \sum_{E'\in\spec(H)} |\widehat f(E^{\textcolor{red}{'}}-E)| e^{\beta(E'-E)/4}<\infty  
%  \end{align}
%  and furthermore
%  \begin{align}
%  \label{eq:SumsFin2}
%  \sum_{E\in\spec(H)} F_1(E)e^{-\beta E/4},\quad  \sum_{E\in\spec(H)} e^{-\beta E/4},\quad \sum_{E\in\spec(H)} e^{-\beta E/2}<\infty.
%  \end{align}
\end{conditionA}
\end{framed}







\noindent Writing the spectral and eigenvalue decompositions of $H$ as $H=\sum_{E\in\operatorname{Sp}(H)}EP_E=\sum_{i}E_i|E_i\rangle\langle E_i|$, with $E_0\le E_1\le \cdots $, where $\{\ket{E_i}\}_{i\in\mathbb{N}_0}$ denotes the energy eigenbasis of $H$ by slight abuse of notation ($\ket{E_i}\ne \ket{E_{i+1}}$, even if $E_i=E_{i+1}$), we will extensively make use of the subspaces
\begin{align*}
\mathcal{F}:=\operatorname{span}\{\ket{E_i}\}_{E_i\in\operatorname{Sp}(H)}\,,\qquad \mathscr{F}:=\operatorname{span}\{|E_i\rangle\langle E_j|\}_{E_i,E_j\in\operatorname{Sp}(H)} 
\end{align*}
of the Hilbert space $\mathcal{H}$, resp.~of the space $\mathscr{T}_1(\mathcal{H})$ of trace-class operators over $\mathcal{H}$. The next claim is standard. 
\begin{claim}\label{densityFFs}
The space $\mathcal{F}$ is dense in $\mathcal{H}$, while $\mathscr{F}$ is dense in $\mathscr{T}_1(\mathcal{H})$.
\end{claim} 

For Schr\"odinger operators, we then have the following theorem that shows that, for the set of bare jumps given by the creation and annihilation operators, i.e., $\{A^\alpha\}_{\alpha\in\cA}\equiv \{a_j,a^\dagger_j\}_{j=1}^m,$ Condition \theconditionA \ is satisfied under very general assumptions. 
A different perspective from our Dirichlet form approach, by directly verifying Davies' conditions to obtain a semigroup in the space of trace-class operators, has been pursued in \cite{GalkowskiZworskiDavies}.    
\begin{theorem}
\label{thm:Schroedinger_operators}
    We consider the Schr\"odinger operator $H=-\Delta+V.$ 
    Let $V$ be real-valued and satisfy 
\begin{itemize}
\item $V \in C^{\infty}(\mathbb R^d)$ with $
V(x)\ge c\langle x\rangle^r-c_0,$ with $ c>0,c_0 \ge 0, r\ge 1,$
as well as
\[
|\partial_x^\alpha V(x)|\le C_\alpha \langle x\rangle^{r-|\alpha|},
\qquad \alpha\in\mathbb N^d.
\]
Choose $\lambda>C_0+1$ and set 
\[
\widetilde H:=H+\lambda.
\]
Let $a_j,a_j^{\dagger}$ be the annihilation and creation operators.
Then for every $n\in\mathbb N$,
\[
\widetilde H^n a_j \widetilde H^{-n-1},\widetilde H^n a_j^{\dagger} \widetilde H^{-n-1}\in \mathcal B(L^2(\mathbb R^d)),
\]
i.e. $\gamma=1,\mu=2$ are admissible in Condition \theconditionA \  for the set of bare jumps being $\{A^\alpha\}_{\alpha\in\cA}\equiv \{a_j,a^\dagger_j\}_{j=1}^d.$
\item $V=\nu |x|^2 + W$ for $\nu>0$ with $W \in  L^{\max\{2,d/2\}}(\mathbb R^d)+ \langle x \rangle^{\alpha} L^{\infty}(\mathbb R^d)\text{ with } \alpha<2$, then $\gamma=1/2$ and $\mu=1$ are admissible in Condition \theconditionA \ for the set of bare jumps being $\{A^\alpha\}_{\alpha\in\cA}\equiv \{a_j,a^\dagger_j\}_{j=1}^d.$
\end{itemize}
\end{theorem}
We stress that the first case includes trapping potentials to very high order and the second case includes singular potentials such as Coulomb interactions. 
The proof of this theorem is given in Appendix \ref{appendixC}.














\subsection{Dirichlet forms and Gibbs generators on Hilbert-Schmidt operators}


\noindent 
Next, for any Bohr frequency $\nu\in B(H)=\operatorname{Sp}(H)-\operatorname{Sp}(H)$ and label $\alpha\in\mathcal{A}$, we introduce the energy jump operators $(A^\alpha_\nu,\mathcal{F})$ by
\begin{align*}
A^\alpha_\nu\ket{E_i}:=P_{E_i+\nu}A^\alpha\ket{E_i}.
\end{align*}
Above, we also write $P_E=0$ whenever $E$ is not an eigenvalue of $H$.
With slight abuse of notation, we also denote $(A^\alpha_\nu)^\dagger=\big((A^\alpha)^\dagger\big)_{-\nu}$, where we recall that $\{A^\alpha\}_{\alpha\in\mathcal{A}}$ is invariant under adjoints. Similarly, we formally define the operators $(A^\alpha_{\nu_1})^\dagger A^\alpha_{\nu_2}$ on $\mathcal{F}$, whose actions on an energy eigenstate $\ket{E_i}$ are
\begin{align*}
(A^\alpha_{\nu_1})^\dagger A^\alpha_{\nu_2}\ket{E_i}=P_{E_i+\nu_2-\nu_1}(A^\alpha)^\dagger P_{E_i+\nu_2}A^\alpha\ket{E_i}.
\end{align*}
Next, we formally introduce the operator
\begin{align}\label{def:L}
L_{\widehat{f},H}(\lambda x+\mu\sqrt{\sigma_\beta})\!:=-\lambda \sum_{\substack{\alpha\in\mathcal{A}\\\nu_1,\nu_2\in B(H)}}\frac{\overline{\widehat{f}(\nu_1)}\widehat{f}(\nu_2)\,e^{\beta(\nu_1+\nu_2)/4}}{2\cosh((\nu_1-\nu_2)\beta/4)}\,(\delta^\alpha_{\nu_1})^{\dagger}\delta^\alpha_{\nu_2}(x)
\end{align}
where 
\begin{equation}
\label{eq:delta}\delta^\alpha_\nu(x):=e^{-\frac{\beta \nu}{4}}A^\alpha_\nu x-e^{\frac{\beta\nu}{4}}xA^\alpha_\nu
\end{equation}
and $(\delta_\nu^\alpha)^{\dagger}$ denotes the formal adjoint of $\delta^\alpha_\nu$ in the Hilbert space $\mathscr{T}_2(\mathcal{H})$ of Schatten-2 operators over $\cH$ endowed with the Hilbert-Schmidt inner product $\langle A,B\rangle:=\Tr(A^\dagger B)$, i.e.
\begin{align*}
(\delta^\alpha_{\nu})^\dagger(x)=e^{-\frac{\beta\nu}{4}}(A^\alpha_\nu)^\dagger x-e^{\frac{\beta\nu}{4}}x(A^\alpha_\nu)^\dagger.
\end{align*}


 \begin{lemma}\label{lem:definL}
 Under \Cref{eq:condAalphas}, the expression \eqref{def:L} defines a densely defined, negative semidefinite symmetric operator $(L_{\widehat{f},H},\operatorname{span}(\mathscr{F}\cup \{\sqrt{\sigma_\beta}\}))$ over $\mathscr{T}_2(\mathcal{H})$, with $\sqrt{\sigma_\beta}\in \operatorname{Ker}(L_{\widehat{f},H})$.
 \end{lemma}
 
 \begin{proof}
We consider $x=|E_i\rangle\langle E_j|$,
\begin{align}
\label{eq:deltaTriangle}
\nn\|(\delta^\alpha_{\nu_1})^\dagger\delta^\alpha_{\nu_2}x\|_2&\le \!e^{-\frac{\beta(\nu_1+\nu_2)}{4}}\|(A^\alpha_{\nu_1} )^\dagger A^\alpha_{\nu_2}x\|_2+e^{\frac{\beta(\nu_1+\nu_2)}{4}}\|xA^\alpha_{\nu_2}(A^\alpha_{\nu_1})^\dagger\|_2\\
&\quad +e^{\frac{\beta(\nu_2-\nu_1)}{4}}\|(A^\alpha_{\nu_1})^\dagger x A^\alpha_{\nu_2}\|_2+e^{\frac{\beta(\nu_1-\nu_2)}{4}}\|A^\alpha_{\nu_2}x(A^\alpha_{\nu_1})^\dagger\|_2.
\end{align}
% \red{Each of the four terms above can be controlled using \Cref{eq:condAalphas} on the jumps $A^\alpha$. For instance
% \begin{align*}
% (A^\alpha_{\nu_1})^\dagger A^\alpha_{\nu_2}|E_i\rangle\langle E_j|&=P_{E_i+\nu_2-\nu_1}(A^\alpha )^\dagger P_{E_i+\nu_2}A^\alpha  |E_i\rangle\langle E_j|
% \end{align*}
% which implies that, denoting $\widetilde{E}:=(E+h_0+1)$.
% \begin{align*}
% \|(A^\alpha_{\nu_1})^\dagger A^\alpha_{\nu_2}|E_i\rangle\langle E_j|\|_2\le \mathbf{1}_{\{\nu_2\ge -h_0-E_i\}}\mathbf{1}_{\{E_i+\nu_2-\nu_1\ge -h_0\}}  C^2\,\widetilde{E_i}^{{\zeta}}(\widetilde{E_i}+\nu_2)^{-\gamma} (\widetilde{E_i}+\nu_2-\nu_1)^{-\gamma} 
% \end{align*}
% since the operator vanishes whenever $\nu_2<-h_0-E_i$. For the second term, we get:
% \begin{align*}
% |E_i\rangle\langle E_j|A^\alpha_{\nu_2}(A^\alpha_{\nu_1})^\dagger=|E_i\rangle\langle E_j| A^\alpha P_{E_j-\nu_2}(A^\alpha)^\dagger P_{E_j-\nu_2+\nu_1}.
% \end{align*}
% Again, the above term vanishes whenever $\nu_2>E_j-h_0$. Hence,
% \begin{align*}
% \|\ket{E_i}\bra{ E_j}A^\alpha_{\nu_2}(A^\alpha_{\nu_1})^\dagger\|_2\le \mathbf{1}_{\{E_j-\nu_2\ge -h_0\}}\mathbf{1}_{\{E_j-\nu_2+\nu_1\ge-h_0\}}C^2\widetilde{E_j}^{{\zeta}}(\widetilde{E_j}-\nu_2)^{{-\gamma}}  (\widetilde{E_j}-\nu_2+\nu_1)^{-\gamma}.
% \end{align*}
% More straightforwardly, we also get that
% \begin{align*}
% &\|(A^\alpha_{\nu_1})^\dagger x A^\alpha_{\nu_2}\|_2\le C^2\widetilde{E_i}^{{\mu}}\widetilde{E_j}^{{\mu}}(\widetilde{E_i}-\nu_1)^{-\gamma}(\widetilde{E_j}-\nu_2)^{-\gamma}\mathbf{1}_{\{E_j-\nu_2\ge -h_0\}}\mathbf{1}_{\{E_i-\nu_1\ge - h_0\}}\\
% &\|A^\alpha_{\nu_2}x(A^\alpha_{\nu_1})^\dagger\|_2\le {C^2\widetilde{E_i}^{{\mu}}\widetilde{E_j}^{{\mu}}}(\widetilde{E_i}+\nu_2)^{-\gamma}(\widetilde{E_j}+\nu_1)^{-\gamma}\mathbf{1}_{\{E_j+\nu_1\ge -h_0\}}\mathbf{1}_{\{E_i+\nu_2\ge -h_0\}}.
% \end{align*}
% Bringing all the estimates together and only keeping track of the dependencies on $\nu_1,\nu_2$, we have, using the notations $(\widetilde{E}+\nu)_{\ge -h_0}:=(\widetilde{E}+\nu)\mathbf{1}_{E+\nu\ge -h_0}$,
% \begin{align*}
% \frac{\|(\delta^\alpha_{\nu_1})^\dagger\delta^\alpha_{\nu_2}x\|_2}{\cosh((\nu_1-\nu_2)\beta/4)}&\lesssim (\widetilde{E_i}+\nu_2)_{\ge -h_0}^{-\gamma}(\widetilde{E_j}+\nu_1)_{\ge -h_0}^{-\gamma}+(\widetilde{E_i}-\nu_1)_{\ge -h_0}^{-\gamma}(\widetilde{E_j}-\nu_2)_{\ge -h_0}^{-\gamma}\! \!\\
% &+\frac{ e^{-\frac{\beta(\nu_1+\nu_2)}{4}}}{\cosh((\nu_1-\nu_2)\beta/4)}(\widetilde{E_i}+\nu_2)_{\ge -h_0}^{-\gamma} (\widetilde{E_i}+\nu_2-\nu_1)_{\ge -h_0}^{-\gamma}\\
% &+\frac{\!e^{\frac{\beta(\nu_1+\nu_2)}{4}}}{\cosh((\nu_1-\nu_2)\beta/4)}(\widetilde{E_j}-\nu_2)_{\ge -h_0}^{{-\gamma}}  (\widetilde{E_j}-\nu_2+\nu_1)_{\ge -h_0}^{-\gamma}\\
% &\equiv T_1(\nu_1,\nu_2)+T_2(\nu_1,\nu_2)+T_3(\nu_1,\nu_2)+T_4(\nu_1,\nu_2).
% \end{align*}
% Next, we verify the summability of each term resulting from the above bound: first
% \begin{align*}
% \sum_{\nu_1,\nu_2}e^{\frac{\beta(
% \nu_1+\nu_2)}{4}}|\widehat{f}(\nu_1)\widehat{f}(\nu_2)|T_1(\nu_1,\nu_2)&=\sum_{\nu_1,\nu_2}e^{\frac{\beta(
% \nu_1+\nu_2)}{4}}|\widehat{f}(\nu_1)\widehat{f}(\nu_2)|(\widetilde{E_i}+\nu_2)_{\ge -h_0}^{-\gamma}(\widetilde{E_j}+\nu_1)_{\ge -h_0}^{-\gamma}\\
% &\le (C')^2 \left(\sum_{E\ge -h_0}\widetilde{E}^{-\gamma}\right)^2
% \end{align*}
% Similarly, using \Cref{eq:condAalphas}, we can show that the three remaining terms satisfy the same bound, and hence
% \begin{align*}
% &{\|L_{\widehat{f},H}x\|_2}\lesssim \left(\sum_{E\ge -h_0}\widetilde{E}^{-\gamma}\right)^2<\infty.
% \end{align*}}


We treat each term in \eqref{eq:deltaTriangle} one by one: For the first one we see 
\begin{align*}
&\sum_{\nu_1,\nu_2} \frac{|\widehat{f}(\nu_1)\widehat{f}(\nu_2)|\,}{2\cosh((\nu_1-\nu_2)\beta/4)}\|(A^\alpha_{\nu_1} )^\dagger A^\alpha_{\nu_2}|E_i\rangle\langle E_j|\|_2 \\
&\le \sum_{E',E\in\spec(H)} \frac{|\widehat{f}(E'-E)\widehat{f}(E'-E_i)|\,}{2\cosh((E_i-E)\beta/4)}\|P_E(A^\alpha)^\dagger P_{E'}A^\alpha|E_i\rangle\langle E_j|\|_2  \\
&\le \|\widehat f\|_\infty \sum_{E,E'\in\spec(H)} \frac{|\widehat{f}(E'-E)|}{2\cosh((E_i-E)\beta/4)}\|P_E(A^\alpha)^\dagger P_{E'}A^\alpha|E_i\rangle\langle E_j|\|_2 \\
&\le\frac{e^{\beta E_i/4}}{2}\|\widehat f\|_\infty\|\widetilde H^{-\gamma}A^\alpha\| \|\widetilde H^{\gamma}A^\alpha\widetilde H^{-\mu}\| (E_i+1+h_0)^\mu\sum_{E,E'\in\spec(H)}|\widehat{f}(E'-E)| e^{-\beta E/4}.
\end{align*}
It remains to argue that the last sum above converges. Indeed, by the  boundedness conditions of $\widehat{f}$ in \Cref{eq:condAalphas}: 
\begin{align*}
\sum_{E,E'\in\spec(H)}|\widehat{f}(E'-E)| e^{-\beta E/4}&=\sum_{E,E'\in\spec(H)}|\widehat{f}(E'-E)|^{3/4} |\widehat{f}(E'-E)|^{1/4}e^{\beta(E'-E)/8}e^{\beta(E-E')/8}e^{-\beta E/4}\\
&=C' \sum_{E,E'\in\spec(H)}e^{-\beta(E+E')/8}<\infty
\end{align*}
where the convergence is ensured by the Gibbs hypothesis.
For the second term in \eqref{eq:deltaTriangle} we see
\begin{align*}
    &\sum_{\nu_1,\nu_2} \frac{|\widehat{f}(\nu_1)\widehat{f}(\nu_2)|\,}{2\cosh((\nu_1-\nu_2)\beta/4)}e^{\frac{\beta(\nu_1+\nu_2)}{2}}\||E_i\rangle\langle E_j|A^\alpha_{\nu_2}(A^\alpha_{\nu_1})^\dagger\|_2\\&\le\sup_{\nu_1\in\R} \left(e^{\beta \nu_1/2}|\widehat f(\nu_1)| \right)\|\widehat f\|_\infty\sum_{E',E\in\spec(H)} \frac{e^{\frac{\beta(E_j-E)}{2}}}{2\cosh((E'-E_j)\beta/4)}\||E_i\rangle\langle E_j|A^\alpha P_E(A^\alpha)^\dagger P_{E'}\|_2 \\&\le \frac{e^{\beta E_j/4}}{2}\sup_{\nu_1\in\R} \left(e^{\beta \nu_1/2}|\widehat f(\nu_1)| \right)\|\widehat f\|_\infty \| A^\alpha \widetilde H^{-\gamma}\| \|\widetilde H^{-\mu}A^\alpha \widetilde H^{\gamma}\| e^{\beta E_j/2}(E_j+h_0+1)^\mu \times \\&\qquad\times \left(\sum_{E\in\spec(H)}e^{-\beta E/2}\right)\left(\sum_{E'\in\spec(H)}e^{-\beta E'/4}\right)<\infty\,,
\end{align*}
 For the third term in \eqref{eq:deltaTriangle} we argue as
\begin{align*}
    &\sum_{\nu_1,\nu_2} \frac{|\widehat{f}(\nu_1)\widehat{f}(\nu_2)|\,}{2\cosh((\nu_1-\nu_2)\beta/4)}
e^{\frac{\beta\nu_2}{2}}\|(A^\alpha_{\nu_1})^\dagger \ket{E_i}\!\bra{E_j}A^\alpha_{\nu_2}\|_2\\&\le e^{\frac{\beta E_j}{2}} \|\widehat f\|^2_\infty \sum_{E,E''\in\spec(H)} \frac{e^{\frac{-\beta E''}{2}}}{2\cosh((E''+E_i-E-E_j)\beta/4)}
 \|P_E(A^\alpha)^\dagger \ket{E_i}\!\bra{E_j}A^\alpha P_{E''}\|_2 \\&\le \frac{e^{\frac{\beta (E_j+E_i)}{4}}}{2} \|\widehat f\|^2_\infty\|\widetilde H^{-\gamma}A^\alpha\|^2(E_i+h_0+1)^{\gamma}(E_j+h_0+1)^{\gamma}  \left(\sum_{E\in\spec(H)} e^{-\beta E/4}\right)^2<\infty.
\end{align*}
Finally, for the fourth term in \eqref{eq:deltaTriangle} we argue as
\begin{align*}
    &\sum_{\nu_1,\nu_2} \frac{|\widehat{f}(\nu_1)\widehat{f}(\nu_2)|\,}{2\cosh((\nu_1-\nu_2)\beta/4)}
e^{\frac{\beta\nu_1}{2}}\|A^\alpha_{\nu_2} \ket{E_i}\!\bra{E_j}(A^\alpha_{\nu_1})^\dagger\|_2\\&\le  \left(\sup_{\nu_1\in\R}e^{\beta\nu_1/2}|\widehat f(\nu_1)|\right)\sum_{E',E'''\in\spec(H)} \frac{|\widehat{f}(E'''-E_i)|\,}{2\cosh((E'+E_i-E'''-E_j)\beta/4)}
\|P_{E'''}A^\alpha \ket{E_i}\!\bra{E_j}(A^\alpha)^\dagger P_E'\|_2 \\&\lesssim \frac{e^{\beta E_j/4}}{2}\|A^\alpha \widetilde H^{-\gamma}\|^2\sum_{E',E'''\in\spec(H)} |\widehat f(E'''-E_i)| e^{\beta (E'''-E_i)/4} e^{-\beta E'/4}\\
&\le C'\frac{e^{\beta E_j/4}}{2}\|A^\alpha \widetilde H^{-\gamma}\|^2\sum_{E',E'''\in\spec(H)}  e^{-\beta (E'''-E_i)/8} e^{-\beta E'/4}<\infty
\end{align*}
where the last bound is once again a consequence of the Gibbs hypothesis. From this we can conclude that
\begin{align*}
    \|L_{\widehat f,H}\ket{E_i}\!\bra{E_j}\|_2<\infty.
\end{align*}
% The first term above is summable since $\nu\mapsto\widehat{f}(\nu)e^{\smash{\frac{\beta\nu}{4}}}$ is summable. For the second term, we lower bound the hyperbolic cosine by $e^{\smash{\beta(\nu_2-\nu_1)/4}}$, so that the series converges as long as $\nu\mapsto\widehat{f}(\nu)e^{\smash{\frac{\beta\nu}{4}}}$ and $\nu\mapsto\widehat{f}(\nu)e^{\smash{-\frac{\beta\nu}{4}}}\mathbf{1}_{\smash{\{\nu\ge -h_0-E\}}}$ are summable. Finally, in the last term, we choose to lower bound the hyperbolic cosine by $e^{\smash{\beta(\nu_1-\nu_2)/4}}$, so that the series converges as long as $\nu\mapsto\widehat{f}(\nu)e^{\smash{\frac{\beta\nu}{4}}}$ and $\nu\mapsto\widehat{f}(\nu)e^{\smash{\frac{3\beta\nu}{4}}}\mathbf{1}_{\smash{\{\nu\le E'-h_0\}}}$ are summable. All these conditions hold, and therefore $\smash\|L_{\smash{\widehat{f},H}}x\|_2<\infty$, whenever \Cref{eq:condAalphas} is satisfied.
Moreover, for any $x\in \mathscr{F}$, denoting $x_\nu:=\widehat{f}(\nu)e^{\smash{\frac{\beta\nu}{4}}}\delta_\nu^\alpha x$ and $x_t:=\sum_{\smash{\nu\in B(H)}}e^{\smash{it\nu}}x_\nu$, 
\begin{align}
\mathcal{E}_{\widehat{f},H}(x)&:=-\langle x,L_{\widehat{f},H}x\rangle=\sum_{\substack{\alpha\in\mathcal{A}\\\nu_1,\nu_2\in B(H)}}\frac{1}{2\cosh((\nu_1-\nu_2)\beta/4)}\,\langle x_{\nu_1},x_{\nu_2}\rangle\label{eq:Dirichletformdef} \\
& =\sum_{\substack{\alpha\in\mathcal{A}\\\nu_1,\nu_2\in B(H)}}\!\!\!\int_{-\infty}^\infty \frac{1}{\beta \cosh(2\pi t/\beta)}\,\langle e^{it\nu_1}x_{\nu_1},e^{it\nu_2}x_{\nu_2}\rangle dt\nonumber
 =\sum_{\alpha\in\mathcal{A}}\int_{-\infty}^\infty \frac{1}{\beta \cosh(2\pi t/\beta)}\,\|x_t\|_2^2 \,dt,\nonumber 
\end{align}
where the above manipulations are justified by Fubini's theorem, together with
\begin{align*}
\|x_t\|_2\le \sum_{\nu}|\widehat{f}(\nu)|e^{\frac{\beta\nu}{4}}\|\delta^\alpha_\nu(x)\|_2<\infty,
\end{align*}
where we used $\beta^{-1}\int_{-\infty}^\infty\cosh(2\pi t/\beta)^{-1}dt=\frac{1}{2}<\infty$ together with \Cref{eq:condAalphas} so that
\begin{align}
&\sum_{\nu}|\widehat{f}(\nu)|\,e^{\frac{\beta\nu}{4}}\,\|\delta^\alpha_\nu(x)\|_2\\
&\qquad\le \sum_{\nu}|\widehat{f}(\nu)|\|A^\alpha_\nu\ket{E_i}\bra{E_j}\|_2+\sum_{\nu}|\widehat{f}(\nu)|e^{\frac{\beta\nu}{2}}\|\ket{E_i}\bra{E}_jA^\alpha_\nu\|_2\nonumber\\
&\qquad= \sum_{E}|\widehat{f}(E-E_i)|\|P_EA^\alpha \ket{E_i}\bra{E_j}\|_2+\sum_{E'}|\widehat{f}(E_j-E')|e^{\frac{\beta(E_j-E')}{2}}\|\ket{E_i}\bra{E_j}A^\alpha P_{E'}\|_2\nonumber\nonumber\\
&\qquad= \sum_{E}|\widehat{f}(E-E_i)|\|P_EA^\alpha \ket{E_i}\bra{E_j}\|_2+\sum_{E'}|\widehat{f}(E'-E_j)|\|\ket{E_i}\bra{E_j}A^\alpha P_{E'}\|_2\nonumber\nonumber\\
&\qquad<\infty.\label{eq.deltaEEbraket}
\end{align}
Thus, $L_{\widehat{f},H}$ is negative semidefinite on $\mathscr{F}$. It remains to verify that for any $x\in \mathscr{F}$, $\langle \sqrt{\sigma_\beta},x\rangle=0$. For this, it suffices to compute, for any two Bohr frequencies $\nu_1,\nu_2\in B(H)$ and $x=\ket{E_i}\bra{E_j}$, the inner products
\begin{align*}
\sqrt{\mathcal{Z}(\beta)}\langle \sqrt{\sigma_\beta},(\delta^\alpha_{\nu_1})^\dagger\delta^\alpha_{\nu_2}(x)\rangle&=e^{-\frac{\beta(\nu_1+\nu_2)}{4}}e^{-\frac{\beta E_j}{2}}\bra{ E_j}(A^\alpha_{\nu_1})^\dagger A^\alpha_{\nu_2}\ket{E_i}\\
&+e^{\frac{\beta(\nu_1+\nu_2)}{4}}e^{-\frac{\beta E_i}{2}}\bra{E_j}A^\alpha_{\nu_2}(A^\alpha_{\nu_1})^\dagger \ket{E_i}\\
&-\sqrt{\mathcal{Z}(\beta)}\,e^{\frac{\beta(\nu_1-\nu_2)}{4}}\bra{E_j}(A^\alpha_{\nu_1})^\dagger \sqrt{\sigma_\beta }A^\alpha_{\nu_2}\ket{E_i}\\
&-\sqrt{\mathcal{Z}(\beta)}\,e^{\frac{\beta(\nu_2-\nu_1)}{4}}\bra{E_j}A^\alpha_{\nu_2}\sqrt{\sigma_\beta}(A^\alpha_{\nu_1})^\dagger\ket{E_i} \\
&=\delta_{E_i+\nu_2,E_j+\nu_1}e^{-\frac{\beta(\nu_1+\nu_2)}{4}}e^{-\frac{\beta E_j}{2}}\bra{E_j}(A^\alpha_{\nu_1})^\dagger P_{E_i+\nu_2}A^\alpha_{\nu_2}\ket{E_i}\\
&+\delta_{E_j-\nu_2,E_i-\nu_1}e^{\frac{\beta(\nu_1+\nu_2)}{4}}e^{-\frac{\beta E_i}{2}}\bra{E_j}A^\alpha_{\nu_2}P_{E_i-\nu_1}(A^\alpha_{\nu_1})^\dagger\ket{E_i}\\
&-\delta_{E_i+\nu_2,E_j+\nu_1}e^{\frac{\beta(\nu_1-\nu_2)}{4}}e^{-\frac{\beta(E_i+\nu_2)}{2}}\bra{E_j}(A^\alpha_{\nu_1})^\dagger P_{E_i+\nu_2}A^\alpha_{\nu_2}\ket{E_i}\\
&-\delta_{E_i-\nu_1,E_j-\nu_2}e^{\frac{\beta(\nu_2-\nu_1)}{4}}e^{\frac{-\beta(E_i-\nu_1)}{2}}\bra{E_j}A^\alpha_{\nu_2}P_{E_i-\nu_1}(A^\alpha_{\nu_1})^\dagger\ket{E_i}=0.
\end{align*}
Therefore, $\langle\sqrt{\sigma_\beta},L_{\widehat{f},H}(x)\rangle=0$ for all $x\in\mathscr{F}$ by the definition of $L_{\widehat{f},H}(x)$. Thus, $L_{\widehat{f},H}$ remains symmetric over $\operatorname{span}(\mathscr{F}\cup\{\sqrt{\sigma_{\beta}}\})$, and negativity follows directly from the observation that $\mathcal{E}_{\widehat{f},H}(\lambda x+\mu\sqrt{\sigma_\beta})=|\lambda|^2\mathcal{E}_{\widehat{f},H}(x)$ for all $x\in\mathscr{F}$ and all $\lambda,\mu\in\mathbb{C}$.
\end{proof}




\noindent In what follows, we denote by $\mathscr{F}_{\sigma_\beta}:=\operatorname{span}(\mathscr{F}\cup\{\sqrt{\sigma_\beta}\})$. Next, by the Friedrichs extension recalled below, the quadratic form $\mathcal{E}_{\widehat{f},H}$ over $\mathscr{F}_{\sigma_\beta}\times \mathscr{F}_{\sigma_\beta}$ defined as $\mathcal{E}_{\widehat{f},H}(x,y):=-\langle x,L_{\widehat{f},H}y\rangle$ induces the generator of a strongly continuous semigroup
\cite[Lemma 10.16, Theorem 10.7]{Schmdgen2012}: let $\mathcal{E}$ be a quadratic form on $\mathscr{T}_2(\mathcal{H})_{\mathbb{R}}$, namely a non-negative definite, symmetric bilinear form $D(\mathcal{E})\times D(\mathcal{E})\to\mathbb{R}$, where $D(\mathcal{E})$ is a dense subspace of $\mathscr{T}_2(\mathcal{H})_{\mathbb{R}}$. The form $\mathcal{E}$ is said to be closed if $D(\mathcal{E})$ equipped with the norm $\|x\|_{D(\mathcal{E})}:=\sqrt{\|x\|_2^2+\mathcal{E}(x)}$ is a Hilbert space. $\mathcal{E}$ is said to be closable if it admits a closed extension, i.e., there exists a closed quadratic form $\mathcal{E}'$ such that $D(\mathcal{E})\subset D(\mathcal{E}')$ and $\mathcal{E}'$ coincide with $\mathcal{E}$ on $D(\mathcal{E})\times D(\mathcal{E})$.
In the case of a closed form $(\mathcal{E},D(\mathcal{E}))$, we extend it to a form $ \mathcal{E}':\mathscr{T}_2(\mathcal{H})\to \mathbb{R}_+\cup \{+\infty\} $ by setting $\mathcal{E}'(x)=+\infty$ whenever $x\notin D(\mathcal{E})$. In that case, $\mathcal{E}'$ is lower semicontinuous on $\mathscr{T}_2(\mathcal{H})$ \cite[Proposition 10.1]{Schmdgen2012}: for any sequence $x_n\to x$ in $\mathscr{T}_2(\mathcal{H})$,
\begin{align}\label{eq.lsc}
\mathcal{E}'\big(\lim_{n\to\infty} x_n\big)\le \liminf_{n\to\infty}\mathcal{E}'\big(x_n\big).
\end{align}
In what follows, we do not distinguish between $\mathcal{E}$ and $\mathcal{E}'$. Finally, given a densely defined, closed, non-negative symmetric linear form $\mathcal{E}$ over $\mathscr{T}_2(\mathcal{H})$, a subspace $\mathscr{D}\subseteq D(\mathcal{E})$ is called a form core for $\mathcal{E}$ if for every $x\in D(\mathcal{E})$, there exists a sequence $(x_n)\subseteq \mathscr{D}$ such that $\|x-x_n\|_{D(\mathcal{E})}\to 0$ as $n\to\infty$.

\begin{lemma}[Friedrichs extension]\label{lem:Friedrichs}
Let $L$ be a densely defined negative semidefinite  symmetric operator and define $\mathcal{E}(x,y):=-\langle L x,y \rangle$ with $x,y \in D(L)$. Then $\mathcal{E}$ is closable, and $D(L)$ is a form core for its closure.
Conversely, any closed quadratic form $\mathcal{E}$ admits a unique densely defined non-positive self-adjoint operator $L$ on  the real Hilbert space $\mathscr{T}_2(\mathcal{H})_{\mathbb{R}}$ of Hermitian elements in $\mathscr{T}_2$, defined on the set of elements $x\in D(\mathcal{E})$ for which there exists $y\in \mathscr{T}_2(\mathcal{H})_{\mathbb{R}}$ such that, for all $ z\in D(\mathcal{E})$, $\mathcal{E}(x,z)=-\Tr(zy)$. In that case, $L(x)=y$. 
\end{lemma} 
\noindent The above construction readily implies that $\mathcal{E}_{\widehat{f},H}$ is closable. For simplicity of notation, we continue to denote its closure by $\mathcal{E}_{\widehat{f},H}$ over the domain $D(\mathcal{E}_{\widehat{f},H})\times D(\mathcal{E}_{\widehat{f},H})$, so that $\mathscr{F}_{\sigma_\beta}$ is a form core for $\mathcal{E}_{\widehat{f},H}$. By Friedrichs extension, it also admits a unique densely defined non-positive self-adjoint operator on $\mathscr{T}_2(\mathcal{H})_{\mathbb{R}}$, called the generator of $\mathcal{E}_{\widehat{f},H}$. Since the latter extends $L_{\widehat{f},H}$, we also identify the two by abuse of notation. In other words, $L_{\widehat{f},H}$ generates a strongly continuous, symmetric semigroup of contractions over $\mathscr{T}_2(\cH)$ \cite[Proposition 6.14]{Schmdgen2012}, which we denote by $\{e^{\smash{t L_{\widehat{f},H}}}\}_{t\ge 0}$.



\smallskip

Next, we argue that the semigroup $\{e^{t L_{\widehat{f},H}}\}_{t\ge 0}$ induces a semigroup of quantum channels over the trace-class operators $\mathscr{T}_1(\mathcal{H})$. This is achieved by leveraging the connection between semigroups of completely positive, trace preserving maps over $\mathscr{T}_1(\mathcal{H})$ and completely Markov semigroups over $\mathscr{T}_2(\mathcal{H})$: given a faithful quantum state $\omega$ on $\mathcal{H}$, we say that a densely defined, closed, non-negative definite, symmetric bilinear form $\mathcal{E}$ is a Dirichlet form (with respect to $\omega$) if it satisfies the following condition: for any self-adjoint $x\in D(\mathcal{E})$, denoting by $x_+$ the positive part of $x$ and $x_\wedge:=x-(x-\sqrt{\omega})_+$, it holds that $x_+,x_\wedge\in D(\mathcal{E})$ and $
\mathcal{E}(x_+),\mathcal{E}(x_\wedge)\le \mathcal{E}(x)$.
A strongly continuous, symmetric semigroup $\{T_t\}_{t\ge 0}$ contraction is said to be Markov (with respect to $\omega$) if and only if $0\le x\le \sqrt{\omega}$ implies that $0\le T_t(x)\le \sqrt{\omega}$ for all $t\ge0$. More generally, let ${\mathscr{M}}_n$ denote the algebra of $n\times n$ matrices acting on $\mathbb{C}^n$. We let $\mathcal{E}^{(n)}$ denote the quadratic form on $\mathscr{T}_2(\mathbb{C}^n\otimes \mathcal{H})_{\mathbb{R}}$, where the state on ${\mathscr{M}}_n$ is taken as the usual trace, given by $D(\mathcal{E}^{(n)}):={\mathscr{M}}_n\otimes D(\mathcal{E})$, with $\mathcal{E}^{(n)}(\{x_{ij}\})=\sum_{i,j=1}^n\mathcal{E}(x_{ij})$.
Then the form $\mathcal{E}$ is said to be $n$-Dirichlet if the form $\mathcal{E}^{(n)}$ is Dirichlet for the underlying state $\omega^{(n)}:=\tau_n\otimes \omega$, where $\tau_n$ denotes the maximally mixed state in ${\mathscr{M}}_n$, and completely Dirichlet if it is $n$-Dirichlet for all $n$. Similarly, the semigroup $\{T_t\}_{t\ge 0}$ is $n$-Markov if for all $0\le x\le \sqrt{\tau_n\otimes \omega}$, $0\le \operatorname{id}_n\otimes T_t(x)\le \sqrt{\tau_n\otimes \omega}$, and completely Markov if it is $n$-Markov for all $n$. Next, we make use of the equivalence between completely Dirichlet forms and completely Markov semigroups \cite[Theorem 5.7]{Goldstein1995} to prove the complete Markov property of the semigroup generated by $L_{\widehat{f},H}$:
\begin{proposition}\label{GibbsDirichlet}
The closed quadratic form $\mathcal{E}_{\widehat{f},H}$ is completely Dirichlet. Moreover, the strongly continuous, symmetric semigroup $\{e^{\smash{tL_{\widehat{f},H}}}\}_{t\ge 0}$ is completely Markov with respect to the state $\sigma_\beta$.
\end{proposition}





\begin{proof}
From the above discussion, it suffices to show that the 
form $\mathcal{E}_{\widehat{f},H}$ is completely Dirichlet. First, the invariance under taking adjoints directly follows from the invariance of the set of jumps under adjoint taking. Next, let $\mathbf{x}\equiv \{x_{ij}\}\in {\mathscr{M}}_n\otimes \mathscr{F}_{\sigma_\beta}$ be a self-adjoint operator with decomposition $\mathbf{x}=\mathbf{x}_+-\mathbf{x}_-$ into positive and negative parts. Decomposing $L_{\widehat{f},H}=K+T$ with, given $x\in\mathscr{F}$,



\begin{align*}
&K(x):=-\frac{1}{2}\sum_{\alpha\in\mathcal{A}}\sum_{\nu_1,\nu_2\in B(H)}\frac{\overline{\widehat{f}(\nu_1)}\widehat{f}(\nu_2)\,e^{\beta(\nu_1+\nu_2)/4}}{\cosh((\nu_1-\nu_2)\beta/4)}  \Big((A^\alpha_{\nu_1})^\dagger A^\alpha_{\nu_2}x+x A^\alpha_{\nu_2}(A^\alpha_{\nu_1})^\dagger\Big)\,,
\end{align*}
and 
\begin{align}
T(x)&:=\frac{1}{2}\sum_{\alpha\in\mathcal{A}}\sum_{\nu_1,\nu_2\in B(H)}\frac{\overline{\widehat{f}(\nu_1)}\widehat{f}(\nu_2)\,e^{\beta(\nu_1+\nu_2)/4}}{\cosh((\nu_1-\nu_2)\beta/4)} \Big((A^\alpha_{\nu_1})^\dagger x A^\alpha_{\nu_2}+A^\alpha_{\nu_2}x(A^\alpha_{\nu_1})^\dagger \Big)\nonumber \\
&= \frac{1}{2}\sum_{\alpha\in\mathcal{A}}  \int_{-\infty}^\infty \frac{1}{\beta\operatorname{cosh}(2\pi t/\beta)}\,\Big((R^\alpha_t)^\dagger x R^\alpha_t+R^\alpha_tx(R^\alpha_t)^\dagger \Big) \,dt\label{eq.FourierintsechTmap} \,,
\end{align}
with $R^\alpha_t:=\sum_{\nu\in B(H)}e^{it\nu+\beta\nu/4}A^\alpha_\nu\widehat{f}(\nu)$, we have that
\begin{align*}
\mathcal{E}_{\widehat{f},H}^{(n)}(\mathbf{x})-\mathcal{E}_{\widehat{f},H}^{(n)}(\mathbf{x}_+)&=-\langle \mathbf{x}_-,(\operatorname{id}_n\otimes L)(\mathbf{x}_-)\rangle  +\langle \mathbf{x}_+,(\operatorname{id}_n\otimes L)(\mathbf{x}_-)\rangle+\langle \mathbf{x}_-,(\operatorname{id}_n\otimes L)(\mathbf{x}_+)\rangle\\
&\ge \langle \mathbf{x}_+,(\operatorname{id}_n\otimes L)(\mathbf{x}_-)\rangle+\langle \mathbf{x}_-,(\operatorname{id}_n\otimes L)(\mathbf{x}_+)\rangle\\
&=\langle \mathbf{x}_+,({\operatorname{id}}_n\otimes T)(\mathbf{x}_-)\rangle+\langle \mathbf{x}_-,({\operatorname{id}}_n\otimes T)(\mathbf{x}_+)\rangle\ge 0\,.
\end{align*}
In the penultimate line, we use the fact that the terms involving the map $K$ are all identically zero due to the cyclicity of the trace, the definition of $K$ as a sum of left and right multiplication operators, and the fact that $\mathbf{x}_+\mathbf{x}_-=\mathbf{x}_-\mathbf{x}_+=0$. The last line follows from the complete positivity of $T$, which is apparent from the Fourier integral representation of $\operatorname{sech}$ (cf. Eq. \eqref{eq.FourierintsechTmap}), and since $\mathbf{x}_-,\mathbf{x}_+$ are positive semidefinite. We conclude using the form core property of $\mathscr{M}_n\otimes \mathscr{F}_{\sigma_\beta}$ with respect to the form norm $\|\cdot \|_{D(\mathcal{E})}$
 and lower semicontinuity of $\mathcal{E}$ (cf.~\Cref{eq.lsc}). Similarly, we can prove that $\mathcal{E}_{\widehat{f},H}^{(n)}(\mathbf{x}_\wedge)\le \mathcal{E}_{\widehat{f},H}^{(n)}(\mathbf{x})$, simply from $\mathbf{x}_\wedge,\mathbf{x}-\mathbf{x}_\wedge\ge 0$ and $\mathbf{x}_\wedge (\mathbf{x}-\mathbf{x}_\wedge)=(\mathbf{x}-\sigma_\beta^{\smash{(n)\frac{1}{2}}})_+(\mathbf{x}-\sigma_\beta^{\smash{(n)\frac{1}{2}}})_-=0$.

\end{proof}





\subsection{From Hilbert to Schr\"{o}dinger}\label{appendixHilberttoSchro}


Next, we show how to define a semigroup of quantum channels on the space $\mathscr{T}_1(\mathcal{H})$ of trace-class operators from the completely Markov semigroup constructed in \Cref{GibbsDirichlet}. We start by deriving a simple density argument: in what follows, given a faithful quantum state $\omega$ on $\mathcal{H}$ and two real numbers $\lambda\le  \mu$, we denote by $B_{[\lambda,\mu]}(\omega)$ the set of self-adjoint, trace-class operators $x$ with $\lambda \,\omega \le x\le \mu\,\omega$.

\begin{lemma}\label{densityT1T2}
Let $\omega$ be a faithful quantum state on $\mathcal{H}$. Then the set of positive semidefinite, trace-class operators $\mathscr{T}_1(\mathcal{H})_+$ satisfies
\begin{align*}
\mathscr{T}_1(\mathcal{H})_+=\overline{\bigcup_{\lambda\in\mathbb{R}_+}B_{[0,\lambda]}(\omega)}^{\|\cdot \|_1}\,.
\end{align*}
\end{lemma}

\begin{proof}
 
 Denote by $P_K$ the projection onto the subspace corresponding to the first $K$ largest eigenvalues of $\omega$. Then, for any positive semidefinite, trace-class operator $a$, $0\le P_ka P_k\le \|a\|P_K\le \frac{\|a\|}{\lambda_K}\omega$, where $\lambda_K$ denotes the $K$-th largest eigenvalue of $\omega$. Moreover, we can clearly see that $\|P_Ka P_K-a\|_1\to 0$ as $K\to\infty$. The result follows.
\end{proof}



\noindent  Next, we consider the completely Dirichlet form $\mathcal{E}_{\widehat{f},H}$ over $\mathscr{T}_2(\mathcal{H})_{\mathbb{R}}$ with the associated generator $L_{\widehat{f},H}$.
Using the non-commutative Radon-Nikodym theorem (see \cite[Lemma 2.2c]{schmitt1986radon} or \cite[Lemma 1.5]{Goldstein1995}), for any $0\le y\le \sigma_\beta$, there exists $X\in \mathscr{B}(\mathcal{H})_+$ with $\|X\|\le 1$ and $y=\sigma_\beta^{\smash{\scriptstyle 1/2}}X\sigma_\beta^{\smash{\scriptstyle 1/2}}$. We denote $x=\sigma_\beta^{\smash{\scriptstyle 1/4}}X\sigma_\beta^{\smash{\scriptstyle 1/4}}\in\mathscr{T}_2(\mathcal{H})_+$, so that $y=\sigma_{\beta}^{\smash{1/4}}x\sigma_\beta^{\smash{1/4}}$. Uniqueness of $X$ (and thus $x$) follows from the fact that $\sigma_\beta^{\smash{1/2}}X\sigma_\beta^{\smash{1/2}}=\sigma_\beta^{\smash{1/2}}Y\sigma_\beta^{\smash{1/2}}\Rightarrow X=Y$ by the faithfulness of $\sigma_\beta$. Next, given the continuous embedding 
\begin{align*}
\iota_2:\mathscr{T}_2(\mathcal{H})\to\mathscr{T}_1(\mathcal{H})\,,\qquad \iota_2(x)=\sigma_\beta^{\frac{1}{4}}x\sigma_\beta^{\frac{1}{4}},
\end{align*}
we can define induced maps on the trace-class operators as
\begin{align}\label{eqcLtoL}
\Phi_t\circ \iota_2(x)=\iota_2\circ e^{tL_{\widehat{f},H}}(x).
\end{align}
By the complete Markov property, 
\begin{align*}
0\le x\le \sqrt{ \sigma_\beta}\quad \Rightarrow\quad  0\le e^{tL_{\widehat{f},H}}(x)\le \sqrt{ \sigma_\beta}\quad \Rightarrow \quad 0\le  \Phi_t(y)=  \iota_2\Big( e^{tL_{\widehat{f},H}}(x) \Big)\le \sigma_\beta.
\end{align*}

\begin{lemma}\label{T1semigroupgeneration}
The maps $\Phi_t$ defined above can be extended by linearity to a strongly continuous semigroup of uniformly bounded, completely positive, trace preserving linear maps on $\mathscr{T}_1(\mathcal{H})_{\mathbb{R}}$ with $\sup_{t\ge 0}\|\Phi_t\|_{1\to 1}\le 1$ and $\Phi_t(\sigma_\beta)=\sigma_\beta$ for all $t\ge 0$.
\end{lemma}
\begin{proof}
Consider $ y\in B_{[0,\lambda]}(\sigma_\beta)$, $\lambda\in\mathbb{R}_+$. Therefore $\|y\|_1\le \lambda$. By order preservation, we have $0\le \Phi_t(y)\le \lambda \sigma_\beta$. Hence $\|\Phi_t(y)\|_1\le \lambda$. By the density of these elements $y$ in the set $\mathscr{T}_1(\mathcal{H})_+$, as seen in Lemma \ref{densityT1T2}, the maps $\Phi_t$ can be extended to bounded maps on $\mathscr{T}_1(\mathcal{H})_+$. Finally, since any $y\in \mathscr{T}_1(\mathcal{H})_{\mathbb{R}}$ can be decomposed as $y=y_+-y_-$, where both $y_+,y_-\in\mathscr{T}_1(\mathcal{H})_+$ and $y_+y_-=0$, we get
\begin{align*}
\|\Phi_t(y)\|_1\le \|\Phi_t(y_+)\|_1+\|\Phi_t(y_-)\|_1
\le \|y_+\|_1+\|y_-\|_1
=\|y\|_1
\end{align*}
and uniform boundedness follows. The semigroup property and strong continuity can be easily shown via density and uniform boundedness of the maps. Complete positivity follows from the completeness of the Dirichlet form. Next, $\sigma_\beta$ is a fixed point of $\Phi_t$, since 
\begin{align*}
\Phi_t(\sigma_\beta)=\iota_2\Big(e^{tL_{\widehat{f},H}}(\sqrt{\sigma_\beta})\Big)=\sigma_\beta,
\end{align*}
where we used that $L_{\smash{\widehat{f},H}}(\sqrt{\sigma_\beta})=0$. Finally, $\Phi_t$ is also trace preserving for all $t\ge 0$. Indeed, we have that the Heisenberg dual map $\Phi^{\smash{\dagger}}_t$ is contractive, which implies that $\Phi^{\smash{\dagger}}_t(I)\le I$. Assuming that the previous inequality is strict, then 
\begin{align*}
1=\Tr(\sigma_\beta I)=\Tr(\Phi_t(\sigma_\beta)I)=\Tr(\sigma_\beta \Phi_t^{\dagger}(I))<1
\end{align*}
where the last strict inequality holds due to the faithfulness of $\sigma_\beta$. This gives a contradiction. Thus $\Phi_t^{\dagger}$ is unital, which implies that $\Phi_t$ is trace preserving. 
\end{proof}





 
 \smallskip
 
\noindent Back to the generator $L_{\smash{\widehat{f},H}}$ introduced in \Cref{def:L}, we aim to obtain a Lindblad form for the generator of the associated semigroup $\{\Phi_t\}_{t\ge 0}$. For this, we formally introduce the jumps and drift operators:
\begin{align}
\label{eq:LalphainGenerationTheory}
L^\alpha:&=\sum_{\nu\in B(H)}\widehat{f}(\nu)A^\alpha_\nu=\sum_{E,E'\in\spec(H) }\widehat{f}(E'-E)P_{E'}A^\alpha P_{E},\qquad \text{for }\alpha\in\mathcal{A},\\
\nn G:&=-\sum_{\substack{\alpha\in\mathcal{A}\\\nu_1,\nu_2\in B(H)}}\, \frac{\overline{\widehat{f}(\nu_1)}\widehat{f}(\nu_2)\,e^{\frac{\beta(\nu_1-\nu_2)}{4}}}{2\cosh((\nu_1-\nu_2)\beta/4)}\,(A^\alpha_{\nu_1})^\dagger A^\alpha_{\nu_2} \\\label{eq:GinGenerationTheore}&=-\sum_{\substack{\alpha\in \mathcal{A}\\E,E',E''\in\spec(H)}}  \frac{\overline{\widehat{f}(E'-E'')}\widehat{f}(E'-E)\,e^{\frac{\beta(E-E'')}{4}}}{2\cosh((E-E'')\beta/4)}P_{E''}\left(A^\alpha\right)^\dagger P_{E'} A^\alpha P_E .
\end{align}
It is clear that these operators are well-defined on the domain $\mathcal{F}.$ In practice, they can, in fact, however, directly be defined on the larger domain $D(\widetilde H^\kappa)$ for certain $\kappa \ge 0$, as we show in Lemma~\ref{lem.LalphaGwelldefinedF} below.
To see this, we introduce the functions\footnote{From \Cref{eq:condAalphas} we know that $\widehat f$ is uniformly bounded. In this case, the function $F_2$ can be bounded by the expression $F_2(E)\lesssim F_1(E)\sum_{E''\in\spec(H)}\frac{1}{1+e^{\beta(E''-E)/2}},$, which is often easier to analyze in practice.}
\begin{align}
\label{eq:defF_1F_2}
    F_1(E):= \sum_{E'\in\spec(H)} |\widehat f(E'-E)|,\quad\text{and}\quad F_2(E):=\sum_{E',E''\in\spec(H)}  \frac{|\widehat{f}(E'-E'')\widehat{f}(E'-E)|}{1+e^{\beta(E''-E)/2}}.
\end{align}
for $E\in\spec(H).$ Note that from \Cref{eq:condAalphas} we easily see that $F_1(E)<\infty$ for all $E\in\spec(H).$ In practice, these are well-behaved and growing polynomially in $E$.
\begin{framed}
\begin{conditionB}\label{eq:F_1Kill}
There exists $\kappa_1>\gamma$ such that the following sums converge:
\begin{align}
\label{eq:kappa_1}
\sum_{E\in\spec(H)} F_1(E)(1+h_0+E)^{-(\kappa_1-\gamma)}&<\infty\, \text{ and }\,  \sum_{E\in\spec(H)} F_2(E) \left(1+h_0+E\right)^{-(\kappa_1-\gamma)}<\infty.
\end{align}
\end{conditionB}
\end{framed}
\begin{lemma}
\label{lem.LalphaGwelldefinedF}
For $\mu\ge \gamma\ge 0$ being defined in \Cref{eq:condAalphas},
let $\kappa_1> \gamma$ be such that \Cref{eq:F_1Kill} holds. Then the operators $L^\alpha$ are well-defined on the dense domain $D(\widetilde H^{\kappa_1})$ for all $\alpha\in\mathcal{A}$ and relatively $\widetilde H^{\kappa_1}$-bounded. Furthermore,  $G$ is well-defined on the dense domain $D(\widetilde H^{\kappa_1+\mu-\gamma})$ and relatively $\widetilde H^{\kappa_1+\mu-\gamma}$-bounded.
\end{lemma}
\begin{proof}
As by \Cref{eq:condAalphas} we have that $P_E'A^\alpha P_E$ is bounded and, in particular, well-defined on $D(\widetilde H^{\kappa_1})$, it remains to show that the sum on the right hand side of the definition $L^\alpha$ in \eqref{eq:LalphainGenerationTheory} when applied on $\ket{\psi}\in D(\widetilde H^{\kappa_1})$ has finite Hilbert space norm\footnote{More precisely, we show absolute convergence of $\sum_{E,E'\in\spec{H}}|\widehat f(E'-E)|\|P_{E'}A^\alpha P_{E}\ket{\psi}\|.$ Hence, the sequence $\ket{\phi_{\tilde E}} := \sum_{\substack{E,E'\in\spec(H)\\E,E'\le \tilde E}} \widehat{f}(E'-E)P_{E'}A^\alpha P_{E}\ket{\psi}$ converges in $\cH$ and we define $L^\alpha\ket{\psi}:=\lim_{\tilde E}\ket{\phi_{\tilde E}}.$}.
For $\ket{\psi} \in D(\widetilde H^{\kappa_1})$ we see
\begin{align*}
    \left\|L^\alpha \ket{\psi}\right\|&\le \sum_{E,E'\in\spec{H}}|\widehat f(E'-E)|\|P_{E'}A^\alpha P_{E}\ket{\psi}\|  \\&\le \|A^{\alpha}\widetilde H^{-\gamma}\| \left\|\widetilde H^{\kappa_1}\ket{\psi}\right\|\sum_{E\in\spec(H)}F_1(E) \left(1+h_0+E\right)^{-(\kappa_1-\gamma)} \lesssim \left\|\widetilde H^{\kappa_1}\ket{\psi}\right\|
\end{align*}
where we used \Cref{eq:condAalphas} on the jumps $A^\alpha$ and \Cref{eq:F_1Kill}.
We can argue similarly for the $G$ operator as for $\ket{\psi}\in D(\widetilde H^{\kappa_1+\mu-\gamma})$ we have
\begin{align*}
\|G\ket{\psi}\|&\le \sum_{\substack{\alpha\in \mathcal{A}\\E,E',E''\in\spec(H)}}  \frac{|\widehat{f}(E'-E'')\widehat{f}(E'-E)|\,e^{\frac{\beta(E-E'')}{4}}}{2\cosh((E-E'')\beta/4)}\left\|P_{E''}\left(A^\alpha\right)^\dagger P_{E'} A^\alpha P_E\ket{\psi}\right\| \\&\le |\mathcal{A}|\left\|\widetilde H^{-\gamma}A^\alpha\right\|\left\|\widetilde H^{\gamma}A^\alpha \widetilde H^{-\mu}\right\| \left\|\widetilde H^{\kappa_1+\mu-\gamma}\ket{\psi}\right\|\sum_{E\in\spec(H)} F_2(E) \left(1+h_0+E\right)^{-(\kappa_1-\gamma)}\\&\lesssim \left\|\widetilde H^{\kappa_1+\mu-\gamma}\ket{\psi}\right\|.
\end{align*}
\end{proof}




% \textcolor{red}{Suggestion: remove all this, extend to finite moments using integral reps in the section on $\cL_{\sigma_E}'s$:
% \begin{lemma}
% \label{lem.LalphaGwelldefinedF} \textcolor{red}{need to change this} Let $\zeta\ge\mu\ge \gamma\ge 0$ being defined in \Cref{eq:condAalphas}. Then the operators $L^\alpha$ are well-defined on the dense domain $D(\widetilde H^{\mu+\gamma})$ for all $\alpha\in\mathcal{A}$ and relatively $\widetilde H^{\mu+\gamma}$-bounded. Furthermore,  $G$ is well-defined on the dense domain $D(\widetilde H^{\zeta+\gamma})$ and relatively $\widetilde H^{\zeta+\gamma}$-bounded.
% \end{lemma}
% \begin{proof}
% As by \Cref{eq:condAalphas}, we have that $P_E'A^\alpha P_E$ is bounded and, in particular, well-defined on $D(\widetilde H^{\mu+\gamma})$, it remains to show that the sum on the right hand side of the definition $L^\alpha$ in \eqref{eq:LalphainGenerationTheory} when applied to $\ket{\psi}\in D(\widetilde H^{\mu+\gamma})$ has a finite Hilbert space norm\footnote{More precisely, we show absolute convergence of $\sum_{E,E'\in\spec{H}}|\widehat f(E'-E)|\|P_{E'}A^\alpha P_{E}\ket{\psi}\|.$ Hence, the sequence $\ket{\phi_{\tilde E}} := \sum_{\substack{E,E'\in\spec(H)\\E,E'\le \tilde E}} \widehat{f}(E'-E)P_{E'}A^\alpha P_{E}\ket{\psi}$ converges in $\cH$ and we define $L^\alpha\ket{\psi}:=\lim_{\tilde E}\ket{\phi_{\tilde E}}.$}.
%     For $\ket{\psi} \in D(\widetilde H^{\mu+\gamma}),$ we see
% \begin{align*}
%     \left\|L^\alpha \ket{\psi}\right\|&\le \sum_{E,E'\in\spec{H}}|\widehat f(E'-E)|\|P_{E'}A^\alpha P_{E}\ket{\psi}\|  \\&\lesssim \left\|\widetilde H^{\gamma}A^{\alpha}\widetilde H^{-\mu}\right\| \left\|\widetilde H^{\mu+\gamma}\ket{\psi}\right\|\left(\sum_{E\in\spec(H)} \left(1+h_0+E\right)^{-\gamma} \right)^2\lesssim \left\|\widetilde H^{\mu+\gamma}\ket{\psi}\right\|
% \end{align*}
% where we used \Cref{eq:condAalphas}.
% We can argue similarly for the $G$ operator as for $\ket{\psi}\in D(\widetilde H^{\zeta+\gamma})$ we have
% \begin{align*}
%         \|G\ket{\psi}\|&\le \sum_{\substack{\alpha\in \mathcal{A}\\E,E',E''\in\spec(H)}}  \left\|P_{E''}\left(A^\alpha\right)^\dagger P_{E'} A^\alpha P_E\ket{\psi}\right\| \\&\lesssim |\mathcal{A}|\left\| \widetilde H^{-\mu}A^\alpha\widetilde H^{\gamma}\right\|\left\|\widetilde H^{\mu+\gamma}A^\alpha \widetilde H^{-\zeta}\right\| \left\|\widetilde H^{\zeta+\gamma}\ket{\psi}\right\|
%     \left(\sum_{E\in\spec(H)}\left(1+h_0+E\right)^{-\gamma}\right)^3\\&\lesssim \left\|\widetilde H^{\zeta+\gamma}\ket{\psi}\right\|,
% \end{align*}
% where we used \Cref{eq:condAalphas} and the fact that $|\frac{e^{\beta(\nu_1-\nu_2)/4}}{2\cosh((\nu_1-\nu_2)\beta/4)}| = \frac{1}{1+e^{\beta(\nu_2-\nu_1)/2}}\le 1.$
% \bin{\begin{align*}
% &\left\|\sum_{\substack{\alpha\in\mathcal{A}\\\nu_1,\nu_2\in B(H)\\E_i+\nu_2\ge -h_0\\E_i+\nu_2-\nu_1\ge-h_0}}\frac{\overline{\widehat{f}(\nu_1)}\widehat{f}(\nu_2)e^{\frac{\beta(\nu_1-\nu_2)\beta}{4}}}{2\cosh((\nu_1-\nu_2)\beta/4)}(A^\alpha)^\dagger P_{E_i+\nu_2} A^\alpha \ket{E_i}\right\|\\
% &\qquad\qquad\qquad\qquad \le C^2(1+h_0+E_i)^\mu\sum_{\substack{\alpha,\nu_1,\nu_2\\E_i+\nu_2\ge -h_0\\E_i+\nu_2-\nu_1\ge -h_0}}(1+h_0+E_i+\nu_2)^\gamma\,{|{\widehat{f}(\nu_1)}\widehat{f}(\nu_2)|e^{\frac{\beta(\nu_1-\nu_2)\beta}{4}}}\\
% &\qquad\qquad\qquad\qquad \lesssim (C')^2C^2(1+h_0+E_i)^\mu. 
% \end{align*}
% Above, we used that, since the sum over $\nu_2$ is bounded from below and for large $\nu_2$, the exponentially decaying function $e^{-\beta\nu_2/4}$ tames the polynomial growth $(1+h_0+E_i+\nu_2)^\gamma$, so that the sum over $\nu_2$ can be controlled in terms of $C'$. Similarly, the sum over $\nu_1$ is controlled by $C'$. Next, we argue that $L^\alpha$ and $G$ can be extended to $\widetilde{H}^\gamma$-bounded, resp.~$\widetilde{H}^\mu$-bounded operators on $\mathcal{H}$. First, for any $\ket{\psi}\in\mathcal{F}$, we have by the exact same computation as above that \textcolor{red}{needs correction}
% \begin{align*}
% \|L^\alpha\widetilde{H}^{-\gamma}\ket{\psi}\|\lesssim CC' \|\ket{\psi}\|.
% \end{align*}
% Now, since $\mathcal{F}$ is dense in $\mathcal{H}$, we get that $L^\alpha\widetilde{H}^{-\gamma}$ extends to a bounded operator over $\mathcal{H}$. The same holds for the drift $G$.}
% \end{proof}}

% \textcolor{blue}{Remove this if not needed any more
% \begin{lemma}
% \label{lem.LalphaGwelldefinedF} Assume \Cref{eq:condAalphas} holds. Then the operators $L^\alpha$ are well-defined on the dense domain $D(e^{\delta H})$, for any $\delta<2\beta$, for all $\alpha\in\mathcal{A}$ and relatively $e^{\delta H}$-bounded. Furthermore,  $G$ is well-defined on the dense domain $D(???)$ and relatively $\widetilde ???$-bounded.
% \end{lemma}
% \begin{proof}
% As by \Cref{eq:condAalphas}, we have that $P_E'A^\alpha P_E$ is bounded and, in particular, well-defined on $D(e^{\delta H})$, it remains to show that the sum on the right hand side of the definition $L^\alpha$ in \eqref{eq:LalphainGenerationTheory} when applied to $\ket{\psi}\in D(e^{\delta H})$ has a finite Hilbert space norm\footnote{More precisely, we show absolute convergence of $\sum_{E,E'\in\spec{H}}|\widehat f(E'-E)|\|P_{E'}A^\alpha P_{E}\ket{\psi}\|.$ Hence, the sequence $\ket{\phi_{\tilde E}} := \sum_{\substack{E,E'\in\spec(H)\\E,E'\le \tilde E}} \widehat{f}(E'-E)P_{E'}A^\alpha P_{E}\ket{\psi}$ converges in $\cH$ and we define $L^\alpha\ket{\psi}:=\lim_{\tilde E}\ket{\phi_{\tilde E}}.$}.
%     For $\ket{\psi} \in D( e^{\delta H}),$ we see
% \begin{align*}
%     \left\|L^\alpha \ket{\psi}\right\|&\le \sum_{E,E'\in\spec{H}}|\widehat f(E'-E)|\|P_{E'}A^\alpha P_{E}\ket{\psi}\|  \\
%     &\le \sum_{E,E'\in\operatorname{Sp}(H)} \, |\widehat{f}(E'-E)|e^{\delta(E'-E)/4}\,e^{\delta(E-E')/4}e^{-\delta E/2}\|A^\alpha e^{-\delta H/2}\|\,\|e^{\delta H}\ket{\psi}\|\\
%     &\le \sup_\nu \left(|\widehat{f}(\nu)|\,e^{\beta\nu/2}\right)^{\frac{\delta}{2\beta}}\, \sup_\nu |\widehat{f}(\nu)|^{1-\frac{\delta}{2\beta}}\, \|A^\alpha e^{-\delta H/2}\|\,\|e^{\delta H}\ket{\psi}\|\,\sum_{E,E'\in\operatorname{Sp}(H)}\, e^{- \delta (E+E')/4}<\infty
% \end{align*}
% where we used \Cref{eq:condAalphas}.
% We can argue similarly for the $G$ operator as for $\ket{\psi}\in D(\widetilde H^{\zeta+\gamma})$ we have
% \begin{align*}
%         \|G\ket{\psi}\|&\le \sum_{\substack{\alpha\in \mathcal{A}\\E,E',E''\in\spec(H)}}  \left\|P_{E''}\left(A^\alpha\right)^\dagger P_{E'} A^\alpha P_E\ket{\psi}\right\| \\&\lesssim |\mathcal{A}|\left\| \widetilde H^{-\mu}A^\alpha\widetilde H^{\gamma}\right\|\left\|\widetilde H^{\mu+\gamma}A^\alpha \widetilde H^{-\zeta}\right\| \left\|\widetilde H^{\zeta+\gamma}\ket{\psi}\right\|
%     \left(\sum_{E\in\spec(H)}\left(1+h_0+E\right)^{-\gamma}\right)^3\\&\lesssim \left\|\widetilde H^{\zeta+\gamma}\ket{\psi}\right\|,
% \end{align*}
% where we used \Cref{eq:condAalphas} and the fact that $|\frac{e^{\beta(\nu_1-\nu_2)/4}}{2\cosh((\nu_1-\nu_2)\beta/4)}| = \frac{1}{1+e^{\beta(\nu_2-\nu_1)/2}}\le 1.$
% \bin{\begin{align*}
% &\left\|\sum_{\substack{\alpha\in\mathcal{A}\\\nu_1,\nu_2\in B(H)\\E_i+\nu_2\ge -h_0\\E_i+\nu_2-\nu_1\ge-h_0}}\frac{\overline{\widehat{f}(\nu_1)}\widehat{f}(\nu_2)e^{\frac{\beta(\nu_1-\nu_2)\beta}{4}}}{2\cosh((\nu_1-\nu_2)\beta/4)}(A^\alpha)^\dagger P_{E_i+\nu_2} A^\alpha \ket{E_i}\right\|\\
% &\qquad\qquad\qquad\qquad \le C^2(1+h_0+E_i)^\mu\sum_{\substack{\alpha,\nu_1,\nu_2\\E_i+\nu_2\ge -h_0\\E_i+\nu_2-\nu_1\ge -h_0}}(1+h_0+E_i+\nu_2)^\gamma\,{|{\widehat{f}(\nu_1)}\widehat{f}(\nu_2)|e^{\frac{\beta(\nu_1-\nu_2)\beta}{4}}}\\
% &\qquad\qquad\qquad\qquad \lesssim (C')^2C^2(1+h_0+E_i)^\mu. 
% \end{align*}
% Above, we used that, since the sum over $\nu_2$ is bounded from below and for large $\nu_2$, the exponentially decaying function $e^{-\beta\nu_2/4}$ tames the polynomial growth $(1+h_0+E_i+\nu_2)^\gamma$, so that the sum over $\nu_2$ can be controlled in terms of $C'$. Similarly, the sum over $\nu_1$ is controlled by $C'$. Next, we argue that $L^\alpha$ and $G$ can be extended to $\widetilde{H}^\gamma$-bounded, resp.~$\widetilde{H}^\mu$-bounded operators on $\mathcal{H}$. First, for any $\ket{\psi}\in\mathcal{F}$, we have by the exact same computation as above that \textcolor{red}{needs correction}
% \begin{align*}
% \|L^\alpha\widetilde{H}^{-\gamma}\ket{\psi}\|\lesssim CC' \|\ket{\psi}\|.
% \end{align*}
% Now, since $\mathcal{F}$ is dense in $\mathcal{H}$, we get that $L^\alpha\widetilde{H}^{-\gamma}$ extends to a bounded operator over $\mathcal{H}$. The same holds for the drift $G$.}
% \end{proof}}

\noindent Next, we provide an expression for the generator $\cL_{\widehat{f},H}$ of the semigroup $\{\Phi_t\}_{t\ge 0}$ on $\mathscr{F}$.
\begin{proposition}\label{propDirichlettoSchro}
The generator $\mathcal{L}_{\widehat{f},H}$ of the semigroup $\{\Phi_t\}_{t\ge 0}$ on $\mathscr{T}_1(\mathcal{H})$ satisfies  $\mathscr{F}\subset D(\mathcal{L}_{\widehat{f},H})$. Moreover, given any two energy eigenstates $\ket{E_i},\ket{E_j}$, the generator $\mathcal{L}_{\smash{\widehat{f},H}}$ associated with the semigroup $\{\Phi_t\}_{t\ge 0}$ evaluated at $\ket{E_i}\bra{E_j}$ takes the form
\begin{align}\label{eq:cLonF}
\mathcal{L}_{\widehat{f},H}(\ket{E_i}\bra{E_j})\!=\sum_{\alpha\in\mathcal{A}}L^\alpha \ket{E_i}(L^\alpha\ket{E_j})^\dagger+G\ket{E_i}\bra{E_j}+\ket{E_i}(G\ket{E_j})^\dagger. 
\end{align}
Following standard notations, we denote $\Phi_t:=e^{t\mathcal{L}_{\widehat{f},H}}$. 
\end{proposition}






\begin{proof}
In order to derive \eqref{eq:cLonF}, it suffices to consider $y=\ket{E_j}\bra{E_i}$, for any two energies $E_i,E_j\in\operatorname{Sp}(H)$. It is clear that there exists a unique $x=\lambda_{i,j}y\in\mathscr{T}_2(\mathcal{H})$, with $\lambda_{i,j}=\mathcal{Z}(\beta)^{1/2}e^{\smash{\beta(E_j-E_i)/4}}$, such that $y=\iota_2(x)$. Therefore, since $x\in D(L_{\widehat{f},H})$, we have that $y\in D(\mathcal{L}_{\widehat{f},H})$ and, by the continuity of the embedding $\iota_2$,
\begin{align*}
\mathcal{L}_{\widehat{f},H}(y)=\lim_{t\to 0}^{\mathscr{T}_1(\mathcal{H})}\frac{\Phi_t(y)-y}{t}=\iota_2\Big(\lim^{\mathscr{T}_2(\mathcal{H})}_{t\to 0}\frac{e^{tL_{\widehat{f},H}}(x)-x}{t}\Big)=\iota_2\circ L_{\widehat{f},H}(x).
\end{align*}
Equation \eqref{eq:cLonF} follows by the definition \eqref{def:L} of $L_{\widehat{f},H}$ and direct computation: first, we directly get that
\begin{align*}
\iota_2\circ L_{\widehat{f},H}(x)=-\sum_{\substack{\nu_1,\nu_2\in B(H)\\\alpha\in\mathcal{A}}}\frac{\overline{\widehat{f}(\nu_1)}\widehat{f}(\nu_2)e^{\frac{\beta(\nu_1+\nu_2)}{4}}}{2\cosh((\nu_1-\nu_2)\beta/4)}\,x^\alpha_{\nu_1,\nu_2}\equiv L_1+L_2+L_3+L_4,
\end{align*}
where each $L_j$ is a sum corresponding to one of the four elements in the decomposition of $x^\alpha_{\nu_1,\nu_2}$ below:
\begin{align}\label{eq:4termstomassage}
x^\alpha_{\nu_1,\nu_2}:=e^{-\frac{\beta\nu_2}{2}}(A^\alpha_{\nu_1})^\dagger A^\alpha_{\nu_2}y+e^{\frac{\beta\nu_2}{2}}yA^\alpha_{\nu_2}(A^\alpha_{\nu_1})^\dagger-e^{-\frac{\beta\nu_2}{2}}A^\alpha_{\nu_2}y(A^\alpha_{\nu_1})^\dagger-e^{\frac{\beta\nu_2}{2}}(A^\alpha_{\nu_1})^\dagger yA^\alpha_{\nu_2}.
\end{align}
Although this is obvious by construction, by an almost identical analysis to the one done in the proof \Cref{lem:definL}, we can also verify that each of the four sums above defines an element in $\mathscr{T}_1(\mathcal{H})$ by hand. Next, we consider the sums associated with the second and fourth terms in \eqref{eq:4termstomassage}: 
\begin{align*}
L_2=-\sum_{\substack{\nu_1,\nu_2\in B(H)\\\alpha\in\mathcal{A}}}\,\frac{\overline{\widehat{f}(\nu_1)}\widehat{f}(\nu_2)e^{\frac{\beta(\nu_1+\nu_2)}{4}}}{2\cosh((\nu_1-\nu_2)\beta/4)}\,e^{\frac{\beta\nu_2}{2}}yA^\alpha_{\nu_2}(A^\alpha_{\nu_1})^\dagger.
\end{align*}
\Cref{eq:condAalphas} implies that $\overline{\widehat{f}(-\nu_2)}\widehat{f}(-\nu_1)e^{\frac{-\beta(\nu_1+\nu_2)}{4}}=\widehat{f}(\nu_2)\overline{\widehat{f}(\nu_1)}e^{\frac{\beta(\nu_1+\nu_2)}{4}}$. Thus, by change of variables $\nu_1\leftarrow -\nu_2$, $\nu_2\leftarrow -\nu_1$, the invariance of $\{A^\alpha\}_{\alpha\in\mathcal{A}}$ under adjoints further implies that 
\begin{align*}
L_2=-\sum_{\substack{\nu_1,\nu_2\in B(H)\\\alpha\in\mathcal{A}}}\,\frac{\overline{\widehat{f}(\nu_1)}\widehat{f}(\nu_2)e^{-\frac{\beta(\nu_1+\nu_2)}{4}}}{2\cosh((\nu_1-\nu_2)\beta/4)}\,e^{-\frac{\beta\nu_1}{2}}y(A^\alpha_{\nu_1})^\dagger A^\alpha_{\nu_2}.
\end{align*}
Similarly, we get that 
\begin{align*}
L_4=\sum_{\substack{\nu_1,\nu_2\in B(H)\\\alpha\in\mathcal{A}}}\frac{\overline{\widehat{f}(\nu_1)}\widehat{f}(\nu_2)e^{\frac{\beta(\nu_1+\nu_2)}{4}}}{2\cosh((\nu_1-\nu_2)\beta/4)}\,e^{-\frac{\beta\nu_1}{2}}A^\alpha_{\nu_2} y(A^\alpha_{\nu_1})^\dagger.
\end{align*}
With these observations, we can rewrite the generator in the Schr\"{o}dinger picture as
\begin{align*}
\mathcal{L}_{\widehat{f},H}(y)=\sum_{\substack{\nu_1,\nu_2\\\alpha\in\mathcal{A}}}\overline{\widehat{f}(\nu_1)}\widehat{f}(\nu_2)\,\left(A^\alpha_{\nu_2} y(A^\alpha_{\nu_1})^\dagger-\frac{e^{\frac{\beta(\nu_1-\nu_2)}{4}}(A^\alpha_{\nu_1})^\dagger A^\alpha_{\nu_2}y+e^{\frac{\beta(\nu_2-\nu_1)}{4}}y(A^\alpha_{\nu_1})^\dagger A^\alpha_{\nu_2}}{2\cosh((\nu_1-\nu_2)\beta/4)}\right),
\end{align*}
as claimed in \eqref{eq:cLonF}. 




\end{proof}




\noindent Our next goal is to show that the domain of $\mathcal{L}_{\widehat{f},H}$ contains a larger set of energy constrained quantum states. For this, as well as in order to streamline the analysis of the next sections, we make use of a simple tool defined in \cite{gondolf2024energy}: given $\delta_1,\delta_2\ge 0$, we introduce the quantum Sobolev spaces defined on 
\begin{align*}
D(\mathcal{W}_H^{\delta_1,\delta_2}):=\Big\{\widetilde{H}^{-\delta_1}a \widetilde{H}^{-\delta_2}\,\Big|\,a\in\mathscr{T}_1(\mathcal{H})\Big\}\text{ with norms }
\|x\|_{\mathcal{W}_H^{\delta_1,\delta_2}}:=\|\mathcal{W}_H^{\delta_1,\delta_2}(x)\|_1,
\end{align*}
where $\mathcal{W}^{\delta_1,\delta_2}_H(x):=\widetilde{H}^{\delta_1} x \widetilde{H}^{\delta_2}$. Since the inverse $(\mathcal{W}_H^{\delta_1,\delta_2})^{-1}$ is bounded, $\mathcal{W}^{\delta_1,\delta_2}_H$ is closed and $(D(\mathcal{W}_H^{\delta_1,\delta_2}),\|\cdot \|_{\smash{\mathcal{W}_H^{\delta_1,\delta_2}}})$ is a Banach space equivalent to the domain $D(\mathcal{W}^{\delta_1,\delta_2}_{\smash{H}})$ endowed with the graph norm of $\mathcal{W}^{\delta_1,\delta_2}_{\smash{H}}$
\cite[Theorem 2.2]{gondolf2024energy}. Moreover, for any $E\ge 0$
\[
\mathscr{S}_E(\widetilde{H}^{2\delta})
:=\bigl\{\, \rho \in \mathscr{S}(\mathcal{H}) \;\big|\; \Tr(\widetilde{H}^{2\delta}\rho) \le E \,\bigr\}\subset \mathscr{S}(\mathcal{H})\cap D(\mathcal{W}_H^{\delta,\delta}).
\]
\begin{lemma}\label{mathFcoreW}
For any $\delta_1,\delta_2>0$, $\mathscr{F}$ is a core for $\mathcal{W}_H^{\delta_1,\delta_2}$.
\end{lemma}
\begin{proof}
For any $a\in D(\mathcal{W}_H^{\delta_1,\delta_2})$, there is $x\in\mathscr{T}_1(\mathcal{H})$ such that $a:=(\mathcal{W}_H^{\delta_1,\delta_2})^{-1}(x)$. We truncate both $x$ and $a$ by the projection $P_n$ onto the subspace spanned by the eigenstates $\ket{E_0},\cdots,\ket{E_{n}}$: $x_n=P_nxP_n$ and $a_n=P_naP_n$. We have by construction that $a_n=(\mathcal{W}_H^{\delta_1,\delta_2})^{-1}(x_n)$. By the proof of \Cref{densityFFs}, we have that $a_n\to a$ and $x_n\to x$. 
\end{proof}
Next, for any two operators $(L_1,D(L_1))$, $(L_2,D(L_2))$ that are relatively $\widetilde{H}^{\delta_1}$-bounded and relatively $\widetilde{H}^{\delta_2}$-bounded, the operators $L_1\widetilde{H}^{-\delta_1}$ and $L_2\widetilde{H}^{-\delta_2}$ can be extended into bounded operators, such that for any $x\in D(\mathcal{W}_H^{\delta_1,\delta_2})$, the trace-class operator 
\begin{align*}
L_1\cdot  x\cdot L_2^\dagger:=L_1\widetilde{H}^{-\delta_1}(L_2\widetilde{H}^{-\delta_2}(\mathcal{W}_H^{\delta_1,\delta_2}(x))^\dagger)^\dagger
\end{align*}
satisfies
\begin{align*}
\|L_1\cdot x\cdot L_2^\dagger\|_1\le \|L_1\widetilde{H}^{-\delta_1}\|\,\|L_2\widetilde{H}^{-\delta_2}\|\,\|x\|_{\mathcal{W}^{\delta_1,\delta_2}_H}. 
\end{align*}
In what follows, we also denote $I\cdot x\cdot L^\dagger_2\equiv x\cdot L^\dagger_2$ and $L_1\cdot x\cdot I\equiv  L_1\cdot x$.
\begin{proposition}\label{prop.generationtheorem} Under \Cref{eq:condAalphas} and \Cref{eq:F_1Kill}, we have $ D(\mathcal{W}_H^{\kappa_1,\kappa_1})\cap D(\mathcal{W}_H^{\kappa_1+\mu-\gamma,0})\cap D(\mathcal{W}_H^{0,\kappa_1+\mu-\gamma}) \subseteq D(\mathcal{L}_{\widehat{f},H})$, and for any $x\in D(\mathcal{W}_H^{\kappa_1,\kappa_1})\cap D(\mathcal{W}_H^{\kappa_1+\mu-\gamma,0})\cap D(\mathcal{W}_H^{0,\kappa_1+\mu-\gamma})$,
\begin{align}\label{eq.LindbladformL}
\mathcal{L}_{\widehat{f},H}(x)=\sum_{\alpha\in\mathcal{A}}L^\alpha\cdot x\cdot (L^\alpha)^\dagger+G\cdot x+x\cdot G^\dagger.
\end{align}
\end{proposition}
\begin{proof}
We use that $\mathcal{L}_{\smash{\widehat{f},H}}$ is closed in $\mathscr{T}_{\smash{1}}(\mathcal{H})$ by the property of being the generator of the strongly continuous semigroup $\{\Phi_t\}_{t\ge 0}$. Therefore, it suffices to show that for any $x\in D(\mathcal{W}_H^{\kappa_1,\kappa_1})\cap D(\mathcal{W}_H^{\kappa_1+\mu-\gamma,0})\cap D(\mathcal{W}_H^{0,\kappa_1+\mu-\gamma})$ there is a sequence $x_n\to x$ in $\mathscr{T}_1(\mathcal{H})$ and $y\in \mathscr{T}_1(\mathcal{H})$ with $\mathcal{L}_{\widehat{f},H}(x_n)\to y$ in $\mathscr{T}_1(\mathcal{H})$. First, we observe that the expression \eqref{eq.LindbladformL} coincides with the definition \eqref{eq:cLonF}
on $\mathscr{F}$. Moreover, for any $x\in D(\mathcal{W}_H^{\kappa_1,\kappa_1})\cap D(\mathcal{W}_H^{\kappa_1+\mu-\gamma,0})\cap D(\mathcal{W}_H^{0,\kappa_1+\mu-\gamma})$, 
there is a sequence $\{x_n\}_{n}$ in $\mathscr{F}$ such that $x_n\to x$ and $\mathcal{W}_H^{\kappa_1,\kappa_1}(x_n)\to \mathcal{W}_H^{\kappa_1,\kappa_1}(x),\,\mathcal{W}_H^{\kappa_1+\mu-\gamma,0}(x_n)\to \mathcal{W}_H^{\kappa_1+\mu-\gamma,0}(x) $ and $\mathcal{W}_H^{0,\kappa_1+\mu-\gamma}(x_n)\to\mathcal{W}_H^{0,\kappa_1+\mu-\gamma}(x)$ by the core property of $\mathscr{F}$ (cf.~Lemma \ref{mathFcoreW}). Fixing $y\equiv \sum_{\alpha\in\mathcal{A}}L^\alpha\cdot x \cdot (L^\alpha)^\dagger+G\cdot x+x\cdot G$, it suffices to show that for each $\alpha$, $L^\alpha\cdot x_n \cdot (L^\alpha)^\dagger\to L^\alpha\cdot x\cdot (L^\alpha)^\dagger$, $G\cdot x_n\to G\cdot x$ and $x_n\cdot G\to x\cdot G$. This we show using Lemma~\ref{lem.LalphaGwelldefinedF} and explicitly only verify the first convergence, as the others follow similarly:
\begin{align*}
\|L^\alpha\cdot x_n\cdot (L^\alpha)^\dagger- L^\alpha\cdot x\cdot (L^\alpha)^\dagger\|_1&=\|L^\alpha\widetilde{H}^{-\kappa_1}(L^\alpha\widetilde{H}^{-\kappa_1}(\mathcal{W}_H^{\kappa_1,\kappa_1}(x_n-x))^\dagger)^\dagger\|_1\\
&\le \|L^\alpha \widetilde{H}^{-\kappa_1}\|^2\,\|x_n-x\|_{\mathcal{W}_H^{\kappa_1,\kappa_1}}\to 0.
\end{align*}
% This can be done by decomposing the sums over Bohr frequencies into sums over differences of energies: defining $a^\alpha_{i,j}=\bra{E_i}A^\alpha\ket{E_j}$ and $\rho_{i,j}:=\bra{E_i}\rho\ket{E_j}$, we have
% \begin{align*}
% \mathcal{L}_{\widehat{f},H}(\rho)=&\sum_{\substack{E_i,E_j,E_k,E_l\in \operatorname{Sp}(H)\\\alpha\in\mathcal{A}}}\,\overline{\widehat{f}(E_l-E_j)}\widehat{f}(E_k-E_i)\rho_{i,j}a^\alpha_{k,i}\overline{a^\alpha_{l,j}}\ket{E_k}\bra{E_l}\\
% &-\sum_{\substack{E_i,E_j,E_k,E_l\in \operatorname{Sp}(H)\\\alpha\in\mathcal{A}}}\,\overline{\widehat{f}(E_k-E_l)}\widehat{f}(E_k-E_i)\frac{e^{\frac{\beta(E_i-E_l)}{4}}}{2\cosh((E_i-E_l)\beta/4)}\rho_{i,j}a^\alpha_{k,i}\overline{a^\alpha_{k,l}}\ket{E_l}\bra{E_j}\\
% &-\sum_{\substack{E_i,E_j,E_k,E_l\in \operatorname{Sp}(H)\\\alpha\in\mathcal{A}}}\,\overline{\widehat{f}(E_k-E_j)}\widehat{f}(E_k-E_l)\frac{e^{\frac{\beta(E_j-E_l)}{4}}}{2\cosh((E_j-E_l)\beta/4)}\rho_{i,j}a^\alpha_{k,l}\overline{a^\alpha_{k,j}}\ket{E_i}\bra{E_l}.
% \end{align*}
\end{proof}

\subsection{Uniqueness of ground state}


The uniqueness of the ground state can be established under some technical assumptions on the Hamiltonian which includes finite rank perturbations of integer powers of the total number operator and the kinetic operator on the torus. 
For this, we require the following simple Lemma on Bohr frequencies.
\begin{lemma}
\label{lemm:spectrum}
Let $\Sigma\subset\mathbb R$ be a discrete set of eigenvalues such that $\Sigma_0\subset \omega\mathbb Z$ for some $\omega>0$, and suppose that $\Sigma$ is obtained from $\Sigma_0$ by changing only finitely many eigenvalues. Then the set of Bohr frequencies
\[
\Lambda:=\Sigma-\Sigma=\{E-E':E,E'\in \Sigma\}
\]
is contained in a finite union of translates of $\omega\mathbb Z$, i.e., there exist $M\in\mathbb N$ and $a_1,\dots,a_M\in\mathbb R$ such that
\[
\Lambda\subset \bigcup_{r=1}^M (a_r+\omega\mathbb Z).
\]
\end{lemma}

\begin{proof}
Since only finitely many eigenvalues are changed, we may write $\Sigma=A\cup F$, where $A\subset \omega\mathbb Z$ and $F=\{f_1,\dots,f_N\}$ is finite. Then
\[
\Lambda=\Sigma-\Sigma\subset (A-A)\cup (A-F)\cup (F-A)\cup (F-F).
\]
Clearly $A-A\subset \omega\mathbb Z$. Moreover, for each $j$ we have $A-f_j\subset (-f_j)+\omega\mathbb Z$ and $f_j-A\subset f_j+\omega\mathbb Z$, hence
\[
A-F\subset \bigcup_{j=1}^N((-f_j)+\omega\mathbb Z),\qquad
F-A\subset \bigcup_{j=1}^N(f_j+\omega\mathbb Z).
\]
Finally, $F-F=\{f_i-f_j:1\le i,j\le N\}$ is finite. Collecting these inclusions,
\[
\Lambda\subset \omega\mathbb Z
\cup \bigcup_{j=1}^N((-f_j)+\omega\mathbb Z)
\cup \bigcup_{j=1}^N(f_j+\omega\mathbb Z)
\cup (F-F).
\]
Since any finite set $C\subset\mathbb R$ satisfies $C\subset \bigcup_{c\in C}(c+\omega\mathbb Z)$, the last term can be absorbed into finitely many translates of $\omega\mathbb Z$, yielding the claim.
\end{proof}

We then have that for Bohr frequencies satisfying the conclusion of the previous Lemma, we can perform a stable phase retrieval as the next Lemma shows.
\begin{lemma}[Stable phase retrieval]\label{lem:Riesz_finite_lattices}
Let
\[
w(t):=\frac{1}{\beta\cosh(2\pi t/\beta)}, 
\qquad
k(s):=\int_{\mathbb R} w(t)e^{its}\,dt
= \frac{1}{2\cosh(\beta s/4)}.
\]
Let $\omega>0$ and assume that the frequency set $\Lambda\subset\mathbb R$ satisfies
\[
\Lambda \subset \bigcup_{r=1}^M (a_r+\omega\mathbb Z),
\]
for some $M\in\mathbb N$ and pairwise distinct $a_1,\dots,a_M\in[0,\omega)$.

Then there exists a constant $A>0$ such that for every finitely supported family $(c_\nu)_{\nu\in\Lambda}$ in a Hilbert space $\mathcal H$,
\begin{equation}\label{eq:Riesz_lower}
A \sum_{\nu\in\Lambda} \|c_\nu\|_{\mathcal H}^2
\;\le\;
\int_{\mathbb R} w(t)\Big\|\sum_{\nu\in\Lambda} c_\nu e^{it\nu}\Big\|_{\mathcal H}^2\,dt.
\end{equation}
In addition, if the non-zero Bohr frequencies are at least $\delta>0$, then for $\beta>0$ large enough, then $\mathcal E(x)=0$ implies that $x_{\nu}=0.$
\end{lemma}

\begin{proof}
For each $r=1,\dots,M$, set
\[
\Lambda_r := \Lambda \cap (a_r+\omega\mathbb Z),
\qquad
c_{r,n} := c_{a_r+n\omega},
\]
with the convention that $c_{r,n}=0$ if $a_r+n\omega\notin\Lambda$.
Then
\[
\sum_{\nu\in\Lambda} c_\nu e^{it\nu}
=
\sum_{r=1}^M e^{ita_r}\sum_{n\in\mathbb Z} c_{r,n} e^{in\omega t}.
\]
Define
\[
F_r(t):=\sum_{n\in\mathbb Z} c_{r,n} e^{in\omega t},
\qquad
F(t):=\sum_{r=1}^M e^{ita_r}F_r(t).
\]

Expanding the square and using the definition of $k$, we obtain
\begin{align*}
\int_{\mathbb R} w(t)\|F(t)\|_{\mathcal H}^2\,dt
&=
\sum_{r,s=1}^M \sum_{n,m\in\mathbb Z}
\langle c_{r,n},c_{s,m}\rangle_{\mathcal H}
\int_{\mathbb R} w(t)e^{it[(a_r-a_s)+(n-m)\omega]}\,dt \\
&=
\sum_{r,s=1}^M \sum_{n,m\in\mathbb Z}
k\big((a_r-a_s)+(n-m)\omega\big)\,
\langle c_{r,n},c_{s,m}\rangle_{\mathcal H}.
\end{align*}

For $\theta\in[-\pi,\pi]$, define
\[
\widehat c_r(\theta):=\sum_{n\in\mathbb Z} c_{r,n}e^{in\theta},
\]
and define the scalar symbols
\[
m_{rs}(\theta)
:=
\sum_{\ell\in\mathbb Z}
k\big((a_r-a_s)+\ell\omega\big)e^{-i\ell\theta}.
\]
Since $k$ decays exponentially, the series defining $m_{rs}$ converges absolutely and uniformly, hence $m_{rs}$ is continuous. Moreover,
\[
m_{sr}(\theta)=\overline{m_{rs}(\theta)},
\]
so $M(\theta):=(m_{rs}(\theta))_{r,s=1}^M$ is self-adjoint.

By setting $\ell=n-m$, we obtain
\[
\sum_{n,m\in\mathbb Z}
k\bigl((a_r-a_s)+(n-m)\omega\bigr)\,
\langle c_{r,n},c_{s,m}\rangle_{\mathcal H}
= \sum_{\ell\in\mathbb Z}
k\bigl((a_r-a_s)+\ell\omega\bigr)
\sum_{m\in\mathbb Z}
\langle c_{r,m+\ell},c_{s,m}\rangle_{\mathcal H}.
\]
By Plancherel on $\ell^2(\mathbb Z)$,
\[
\sum_{m\in\mathbb Z}
\langle c_{r,m+\ell},c_{s,m}\rangle_{\mathcal H}
=
\int_{-\pi}^{\pi}
e^{-i\ell\theta}\,
\langle \widehat c_r(\theta),\widehat c_s(\theta)\rangle_{\mathcal H}
\,\frac{d\theta}{2\pi}.
\]
Hence
\[
\sum_{\ell\in\mathbb Z}
k\bigl((a_r-a_s)+\ell\omega\bigr)
\sum_{m\in\mathbb Z}
\langle c_{r,m+\ell},c_{s,m}\rangle_{\mathcal H}= \int_{-\pi}^{\pi}
\Bigl(
\sum_{\ell\in\mathbb Z}
k\bigl((a_r-a_s)+\ell\omega\bigr)e^{-i\ell\theta}
\Bigr)
\langle \widehat c_r(\theta),\widehat c_s(\theta)\rangle_{\mathcal H}
\frac{d\theta}{2\pi}.
\]
This means that
\[
\int_{\mathbb R} w(t)\|F(t)\|_{\mathcal H}^2\,dt
=
\int_{-\pi}^{\pi}
\sum_{r,s=1}^M
m_{rs}(\theta)\,
\langle \widehat c_r(\theta),\widehat c_s(\theta)\rangle_{\mathcal H}
\,\frac{d\theta}{2\pi}.
\]

We now show that $M(\theta)$ is strictly positive definite for each $\theta$.
By Poisson summation,
\[
m_{rs}(\theta)
=
\frac{2\pi}{\omega}\sum_{j\in\mathbb Z}
w\!\left(\frac{\theta+2\pi j}{\omega}\right)
e^{\,i(a_r-a_s)(\theta+2\pi j)/\omega}.
\]
Hence, for any $z=(z_1,\dots,z_M)\in\mathbb C^M$,
\begin{align*}
\sum_{r,s=1}^M m_{rs}(\theta) z_r\overline{z_s}
&=
\frac{2\pi}{\omega}\sum_{j\in\mathbb Z}
w\!\left(\frac{\theta+2\pi j}{\omega}\right)
\sum_{r,s=1}^M
z_r\overline{z_s}
e^{\,i(a_r-a_s)(\theta+2\pi j)/\omega} \\
&=
\frac{2\pi}{\omega}\sum_{j\in\mathbb Z}
w\!\left(\frac{\theta+2\pi j}{\omega}\right)
\Big|\sum_{r=1}^M z_r e^{ia_r(\theta+2\pi j)/\omega}\Big|^2.
\end{align*}
This shows that $M(\theta)$ is positive semidefinite.

Assume now that
\[
\sum_{r,s=1}^M m_{rs}(\theta) z_r\overline{z_s}=0.
\]
Since $w>0$, every term in the above sum must vanish, hence
\[
\sum_{r=1}^M z_r e^{ia_r(\theta+2\pi j)/\omega}=0
\qquad\text{for every }j\in\mathbb Z.
\]
Set
\[
\lambda_r:=e^{2\pi i a_r/\omega},
\qquad
c_r:=z_r e^{ia_r\theta/\omega}.
\]
Then the preceding identities become
\[
\sum_{r=1}^M c_r \lambda_r^j = 0
\qquad\text{for every }j\in\mathbb Z.
\]
Since the $a_r$ are distinct modulo $\omega$, the numbers $\lambda_r$ are pairwise distinct. Taking $j=0,\dots,M-1$, we obtain a Vandermonde system, hence
\[
c_r=0 \quad\text{for all }r.
\]
Therefore $z_r=0$ for all $r$, and so $M(\theta)$ is strictly positive definite.

Let $\lambda_{\min}(\theta)$ denote the smallest eigenvalue of $M(\theta)$. Since $M(\theta)$ depends continuously on $\theta$ and is strictly positive definite for every $\theta$, compactness of $[-\pi,\pi]$ yields
\[
A:=\min_{\theta\in[-\pi,\pi]} \lambda_{\min}(\theta)>0.
\]
Therefore
\[
\sum_{r,s=1}^M
m_{rs}(\theta)\,
\langle u_r,u_s\rangle_{\mathcal H}
\ge
A\sum_{r=1}^M \|u_r\|_{\mathcal H}^2
\qquad
\text{for all }u_1,\dots,u_M\in\mathcal H.
\]
Applying this with $u_r=\widehat c_r(\theta)$ and integrating, we get
\[
\int_{\mathbb R} w(t)\|F(t)\|_{\mathcal H}^2\,dt
\ge
A\int_{-\pi}^{\pi}\sum_{r=1}^M \|\widehat c_r(\theta)\|_{\mathcal H}^2\,\frac{d\theta}{2\pi}.
\]
Finally, Parseval gives
\[
\int_{-\pi}^{\pi}\|\widehat c_r(\theta)\|_{\mathcal H}^2\,\frac{d\theta}{2\pi}
=
\sum_{n\in\mathbb Z}\|c_{r,n}\|_{\mathcal H}^2,
\]
and so
\[
\int_{\mathbb R} w(t)\|F(t)\|_{\mathcal H}^2\,dt
\ge
A\sum_{r=1}^M\sum_{n\in\mathbb Z}\|c_{r,n}\|_{\mathcal H}^2
=
A\sum_{\nu\in\Lambda}\|c_\nu\|_{\mathcal H}^2.
\]
This proves \eqref{eq:Riesz_lower}.
\end{proof}

If the different Bohr frequencies are separated by a distance of at least $\delta>0$, it is still possible to recover the coefficients for low enough temperatures, i.e. $\beta \gg 1$ as in this case the integral kernel is sufficiently diagonal.

\begin{lemma}[Coercive form]
\label{lemm:coercive}
Let \(B\subset \mathbb R\) be countable, and for \(\beta>0\) define
\[
K_{\nu\mu}:=\frac{1}{2\cosh((\nu-\mu)\beta/4)},
\qquad \nu,\mu\in B.
\]
Let \(x=(x_\nu)_{\nu\in B}\) be a countable family in a Hilbert space \(\mathcal H\), and assume that
\[
\mathcal E_\beta(x):=\sum_{\nu,\mu\in B} K_{\nu\mu}\,\langle x_\nu,x_\mu\rangle
\]
is absolutely convergent. Set $M_\beta:=\sup_{\nu\in B}\sum_{\mu\neq \nu} K_{\nu\mu}.$
Then, if $
M_\beta<\frac12,$ we find
\[
\mathcal E_\beta(x)\ge \Big(\frac12-M_\beta\Big)\sum_{\nu\in B}\|x_\nu\|^2.
\]
In particular, assume that \(B\) is uniformly separated, i.e. there exists \(\delta>0\) such that
\[
|\nu-\mu|\ge \delta
\qquad\text{for all }\nu\neq\mu\text{ in }B.
\]
Then
\[
\mathcal E_\beta(x)\ge c_\beta \sum_{\nu\in B}\|x_\nu\|^2,
\]
where for $\beta$ large enough
\[
c_\beta:=\frac12-2\sum_{m=1}^\infty \frac{1}{2\cosh(m\delta\beta/4)}>0
\]
\end{lemma}

\begin{proof}
Since $K_{\nu\nu}=\frac{1}{2\cosh(0)}=\frac12,$
we decompose
\[
\mathcal E_\beta(x)
=
\frac12\sum_{\nu\in B}\|x_\nu\|^2
+
\sum_{\nu\neq\mu}K_{\nu\mu}\langle x_\nu,x_\mu\rangle.
\]
By the Cauchy-Schwarz inequality, we have $\Re\langle x_\nu,x_\mu\rangle\ge -\|x_\nu\|\,\|x_\mu\|,$
and thus get
\[
\mathcal E_\beta(x)
\ge
\frac12\sum_{\nu\in B}\|x_\nu\|^2
-
\sum_{\nu\neq\mu}K_{\nu\mu}\|x_\nu\|\,\|x_\mu\|.
\]
Now put \(u_\nu:=\|x_\nu\|\). Since \(K_{\nu\mu}\ge 0\), Schur's test yields
\[
\sum_{\nu\neq\mu}K_{\nu\mu}u_\nu u_\mu
\le
\Big(\sup_{\nu\in B}\sum_{\mu\neq\nu}K_{\nu\mu}\Big)\sum_{\nu\in B}u_\nu^2
=
M_\beta\sum_{\nu\in B}\|x_\nu\|^2.
\]
Therefore
\[
\mathcal E_\beta(x)
\ge
\Big(\frac12-M_\beta\Big)\sum_{\nu\in B}\|x_\nu\|^2,
\]
which proves the claim.
For the second part we fix \(\nu\in B\). By the separation assumption, for each \(m\ge1\) there is at most one point of \(B\) in each interval
\[
[\nu+m\delta,\nu+(m+1)\delta)
\quad\text{and}\quad
(\nu-(m+1)\delta,\nu-m\delta].
\]
Therefore
\[
\sum_{\mu\neq\nu} K_{\nu\mu}
\le
2\sum_{m=1}^\infty \frac{1}{2\cosh(m\delta\beta/4)}.
\]
Setting then
$S_\beta:=2\sum_{m=1}^\infty \frac{1}{2\cosh(m\delta\beta/4)},$ then implies
\[
\sup_{\nu\in B}\sum_{\mu\neq\nu}K_{\nu\mu}\le S_\beta \xrightarrow[\beta \to \infty]{}0.
\]
\end{proof}



We can now use this to conclude the existence of a unique invariant state of the Lindbladian for certain Hamiltonians that exhibit a coercive form. The coercivity is satisfied for Hamiltonians with a spectrum as in Lemma \ref{lemm:spectrum} or Bohr frequencies as in Lemma \ref{lemm:coercive}.  First, we conclude from $x_n \to x$ and $\mathcal E(x_n) \to \mathcal E(x)=0$ that for all $\nu\in B(H)$ and $ \alpha\in\mathcal A$ by Lemma \ref{lem:Riesz_finite_lattices}
\[
\lim_{n \to \infty} (x_n)_\nu^\alpha=0.
\]
By the definition of $x_\nu^\alpha$ and the assumption $\widehat f(\nu)\neq 0$, it follows that for all $\nu\in B(H)$ and $ \alpha\in\mathcal A$
\[
\lim_{n \to \infty} \delta_\nu^\alpha(x_n)=0.
\]
We recall that 
\[
\langle \varphi, \delta_\nu^\alpha(x) \psi\rangle=\langle \varphi, (e^{-\beta \nu/4}A_\nu^\alpha x-e^{\beta \nu/4}xA_\nu^\alpha) \psi\rangle.
\]
We have, by continuity, that for suitable $\varphi,\psi \in D(e^{\beta H/4})$ with $\varphi \in D(e^{\beta H/4} (A_\nu^\alpha)^*)$ and $\psi \in D(e^{\beta H/4} A_\nu^\alpha )$
\begin{equation}
\label{eq:limit}
\begin{split}
0&=\lim_{n \rightarrow \infty} \langle e^{\beta H/4} \varphi, \delta_\nu^\alpha(x_n) e^{\beta H/4}\psi\rangle\\
&=\lim_{n \rightarrow \infty}\langle e^{\beta H/4} \varphi, (e^{-\beta \nu/4}A_\nu^\alpha x_n-e^{\beta \nu/4}x_nA_\nu^\alpha) e^{\beta H/4} \psi\rangle \\
&=\lim_{n \rightarrow \infty}\langle e^{\beta H/4} (A_\nu^\alpha)^* \varphi,  x_ne^{\beta H/4} \psi\rangle-\langle e^{\beta H/4} \varphi, x_n e^{\beta H/4} A_\nu^\alpha \psi\rangle\\
&=\langle e^{\beta H/4} (A_\nu^\alpha)^* \varphi,  xe^{\beta H/4} \psi\rangle-\langle e^{\beta H/4} \varphi, x e^{\beta H/4} A_\nu^\alpha \psi\rangle.
\end{split}
\end{equation}
We then have that by summing $A_\nu^\alpha$ over $\nu$
\[ \sum_{\nu \in B(H)} u_{\nu}= A^{\alpha}\psi. \] 
In addition, we have 
\begin{equation} 
\label{eq:closedness}
\sum_{\nu \in B(H)} e^{\beta H/4}u_{\nu} =  \sum_{\nu \in B(H)}e^{\beta H/4}A_\nu^\alpha  \psi = \sum_{\nu \in B(H)} e^{\beta \nu /4}A_\nu^\alpha  e^{\beta H/4}\psi. 
\end{equation}

If this sum converges, then since $e^{\beta H/4}$ is closed, it follows that 
\[ \sum_{\nu \in B(H)} e^{\beta H/4}u_{\nu} =  e^{\beta H/4} A^{\alpha} \psi. \] 
For states $\varphi,\psi$ for which \eqref{eq:limit} holds and for which \eqref{eq:closedness} is finite we have 
\begin{equation}
\label{eq:comm_relations}
    \langle e^{\beta H/4} (A^\alpha)^* \varphi,  xe^{\beta H/4} \psi\rangle=\langle e^{\beta H/4} \varphi, x e^{\beta H/4} A^\alpha \psi\rangle.
\end{equation}

To simplify the presentation, we focus on the case of one mode in the following result showing that the nullspace of the quadratic form is spanned by the Gibbs state 
\[ \mathcal E(x)=0 \Rightarrow x \propto e^{-\beta H/2}.\]
\begin{theorem}[Uniqueness of the invariant state]
Let either $\mathcal H=L^2(\mathbb R)$ with $N=a^\dagger a$ or $\mathcal H=L^2(\mathbb R/(2\pi\mathbb Z))$ with $N=-\Delta=D_x^2$. Consider Hamiltonians
\[
H=N^m+K,\qquad m\in\mathbb N,
\]
where $K\in\mathcal B(\mathcal H)$ is finite rank and satisfies $K=1_{[0,n]}(N)K1_{[0,n]}(N)$ for some $n\in\mathbb N$. In the first case choose jump operators $\{a,a^\dagger\}$, and in the second case $\{D_x,E,E^\dagger\}$ with $E=e^{ix}$. Then
\[
\mathcal E_{\hat f,H}(x)=0 \ \Longrightarrow\ x\propto e^{-\beta H/2}.
\]
\end{theorem}

\begin{proof}
By Lemma~\ref{lem:Riesz_finite_lattices}, the finite-energy space
\[
D_{\mathrm{fin}}:=\bigoplus_{\lambda\in\operatorname{Spec}(N)}\ker(N-\lambda)
\]
is contained in $D(e^{\beta H/4})$, and the weak commutation relations \eqref{eq:comm_relations} hold on $D_{\mathrm{fin}}$. Hence we may define the sesquilinear form
\[
B(\varphi,\psi):=\langle e^{\beta H/4}\varphi,\, x\,e^{\beta H/4}\psi\rangle,\qquad \varphi,\psi\in D_{\mathrm{fin}}.
\]

\textit{Case 1: $\mathcal H=L^2(\mathbb R)$, $N=a^\dagger a$.}
The weak relations read
\[
B(a^\dagger\varphi,\psi)=B(\varphi,a\psi),\qquad B(a\varphi,\psi)=B(\varphi,a^\dagger\psi).
\]
Let $(e_n)_{n\ge0}$ be the Hermite basis and set $b_{mn}:=B(e_m,e_n)$. Using $ae_n=\sqrt n\,e_{n-1}$ and $a^\dagger e_n=\sqrt{n+1}\,e_{n+1}$, we obtain
\[
\sqrt{m+1}\,b_{m+1,n}=\sqrt n\,b_{m,n-1},\qquad \sqrt m\,b_{m-1,n}=\sqrt{n+1}\,b_{m,n+1}.
\]
Setting $n=0$ gives $b_{p0}=0$ for $p\ge1$, while $m=0$ gives $b_{0q}=0$ for $q\ge1$. Iterating the recursions yields $b_{mn}=0$ for $m\neq n$. For the diagonal, taking $n=m+1$ shows $b_{m+1,m+1}=b_{mm}$, hence $b_{nn}=\lambda$ for some $\lambda\in\mathbb C$. Thus $B(e_m,e_n)=\lambda\delta_{mn}=\lambda\langle e_m,e_n\rangle$, and by sesquilinearity $B(\varphi,\psi)=\lambda\langle\varphi,\psi\rangle$ on $D_{\mathrm{fin}}$.

\textit{Case 2: $\mathcal H=L^2(\mathbb R/(2\pi\mathbb Z))$, $N=D_x^2$.}
Let $e_k(x)=(2\pi)^{-1/2}e^{ikx}$, $k\in\mathbb Z$. Then $D_x e_k=ke_k$, $Ee_k=e_{k+1}$, $E^\dagger e_k=e_{k-1}$, and the weak relations give
\[
B(D_x\varphi,\psi)=B(\varphi,D_x\psi),\quad B(E\varphi,\psi)=B(\varphi,E^\dagger\psi),\quad B(E^\dagger\varphi,\psi)=B(\varphi,E\psi).
\]
Setting $b_{jk}:=B(e_j,e_k)$, the $D_x$-relation yields $j\,b_{jk}=k\,b_{jk}$, hence $b_{jk}=0$ for $j\neq k$. The $E$-relation gives $b_{j+1,k}=b_{j,k-1}$, and setting $k=j+1$ yields $b_{j+1,j+1}=b_{jj}$, so $b_{jj}=\lambda$. Thus again $B(e_j,e_k)=\lambda\delta_{jk}$ and hence $B(\varphi,\psi)=\lambda\langle\varphi,\psi\rangle$.

In both cases,
\[
B(\varphi,\psi)=\lambda\langle\varphi,\psi\rangle,\qquad \varphi,\psi\in D_{\mathrm{fin}}.
\]
On the other hand,
\[
\langle e^{\beta H/4}\varphi,\, \lambda e^{-\beta H/2}e^{\beta H/4}\psi\rangle=\lambda\langle\varphi,\psi\rangle,
\]
so for all $\varphi,\psi\in D_{\mathrm{fin}}$,
\[
\langle e^{\beta H/4}\varphi,\, x\,e^{\beta H/4}\psi\rangle=\langle e^{\beta H/4}\varphi,\, \lambda e^{-\beta H/2}e^{\beta H/4}\psi\rangle.
\]
Since $D_{\mathrm{fin}}$ is a form core for $e^{\beta H/4}$, it follows that $x=\lambda e^{-\beta H/2}$, proving the claim.
\end{proof}

\subsection{Davies generators}\label{sec:Davies}

\noindent We end this first section by briefly sketching a generation theory for Davies generators, which parallels the previous construction. We start by introducing the symmetric operator on $\operatorname{span}(\mathscr{F}\cup \{\sqrt{\sigma_{\beta}}\})$
\begin{align}
L_{0,\widehat{f},H}(\lambda x+\mu\sqrt{\sigma_\beta})\!:=-\lambda \sum_{\substack{\alpha\in\mathcal{A}\\ \nu\in B(H)}}\frac{|\widehat{f}(\nu)|^2\,e^{\beta\nu/2}}{2}\,(\delta^\alpha_{\nu})^{\dagger}\delta^\alpha_{\nu}(x)
\end{align}
Compared with the generator $L_{\widehat{f},H}$ constructed in \eqref{def:L}, $L_{0,\widehat{f},H}$ corresponds to the sum over Bohr frequencies $\nu_1=\nu_2$. By the same reasoning as done in \Cref{propDirichlettoSchro}, we get that the associated Lindbladian in the Schr\"{o}dinger picture takes the following form when evaluated on $x\in\mathscr{F}$
\begin{align}
\label{eq:DaviesGenerator}
    \mathcal{L}_{0,\widehat{f},H}(x)&=\sum_{\substack{\nu\in B(H)\\\alpha\in\mathcal{A}}}{|\widehat{f}(\nu)|^2}\,\left(A^\alpha_{\nu} \cdot x\cdot (A^\alpha_{\nu})^\dagger-\frac{(A^\alpha_{\nu})^\dagger A^\alpha_{\nu}\cdot x+x\cdot (A^\alpha_{\nu})^\dagger A^\alpha_{\nu}}{2}\right).
\end{align}
Next, we aim to extend the domain of definition of $\mathcal{L}_{0,\widehat{f},H}$ to a quantum Sobolev space, thereby including certain finite moment states in Proposition~\ref{prop:DaviesFiniteMoment} below. For that, we proceed similarly as for the proof of Lemma~\ref{lem.LalphaGwelldefinedF} and Proposition~\ref{prop.generationtheorem} and define the function
\begin{align}\label{def:F3E}
F(E):=\sum_{E'\in\spec(H)}|\widehat f(E'-E)|^2
\end{align}
for $E\in\spec(H).$
\begin{proposition}
\label{prop:DaviesFiniteMoment}
    For $\mu\ge\gamma\ge 0$ being defined in \Cref{eq:condAalphas}, let $\kappa_2,\kappa_3\ge\gamma$ be such that
    \begin{align}
    \sum_{E\in \spec(H)} F(E)(1+h_0+E)^{-(\kappa_2-\gamma)} &<\infty,\quad \sum_{E\in \spec(H)} (1+h_0+E)^{-(\kappa_3-\gamma)} <\infty.
\end{align}
Then $D(\mathcal{W}_H^{\kappa_2,\kappa_3})\cap D(\mathcal{W}_H^{\kappa_2+\mu-\gamma,0})\cap D(\mathcal{W}_H^{0,\kappa_2+\mu-\gamma})\subseteq D(\mathcal{L}_{0,\widehat{f},H})$ and for $x\in D(\mathcal{W}_H^{\kappa_2,\kappa_3})\cap D(\mathcal{W}_H^{\kappa_2+\mu-\gamma,0})\cap D(\mathcal{W}_H^{0,\kappa_2+\mu-\gamma})$ the action of $\mathcal{L}_{0,\widehat{f},H}$ on $x$ is explicitly given by \eqref{eq:DaviesGenerator}.

\end{proposition}

\begin{proof}
\bin{\begin{align*}
\mathcal{L}_{0,\widehat{f},H}(x)&=\sum_{\substack{E,E',E'',E'''\in\spec(H)\\ E'-E=E'''-E''}}|\widehat f(E'-E)|^2
P_{E'}A^{\alpha}P_ExP_{E''} (A^{\alpha})^{\dagger}P_{E'''}\\\quad&-\sum_{E,E'\in\spec(H)}|\widehat f(E'-E)|^2\frac{P_{E}(A^{\alpha})^{\dagger}P_{E'}A^{\alpha}P_Ex+xP_{E}(A^{\alpha})^{\dagger}P_{E'}A^{\alpha}P_E}{2}.
\end{align*}}
For the CP-term of the Davies generator in \eqref{eq:DaviesGenerator} and $x\in D(\mathcal{W}^{\kappa_2,\kappa_3}_H)$ we see
\begin{align*}
&\sum_{\nu\in B(H)}|\widehat f(\nu)|^2 \left\|A^{\alpha}_{\nu}\cdot x\cdot \left(A^{\alpha}_\nu\right)^{\dagger}\right\|_1 \le 
\sum_{\substack{E,E',E'',E'''\in\spec(H)\\ E'-E=E'''-E''}}|\widehat f(E'-E)|^2
\left\|P_{E'}A^{\alpha}P_ExP_{E''} (A^{\alpha})^{\dagger}P_{E'''} \right\|_1\\& \le  \sum_{\substack{E,E',E''\in\spec(H)}}|\widehat f(E'-E)|^2
\left\|A^{\alpha}P_ExP_{E''} (A^{\alpha})^{\dagger} \right\|_1\\&\le\|A^\alpha \widetilde H^{-\gamma}\|^2\|\widetilde H^{\kappa_3}x\widetilde H^{\kappa_4}\|_1 \sum_{E,E''\in \spec(H)} F(E) \left(1+h_0+E\right)^{-(\kappa_2-\gamma)} \left(1+h_0+E''\right)^{-(\kappa_3-\gamma)} \\
&\lesssim \|x\|_{\mathcal{W}^{\kappa_2,\kappa_3}_H}.
\end{align*}
Furthermore, for the first term in the anticommutator part and $x\in D(\mathcal{W}^{\kappa_2+\mu-\gamma,0}_H)$ we get
\begin{align*}
    &\sum_{\nu\in B(H)}|\widehat f(\nu)|^2\left\|(A_\nu^{\alpha})^{\dagger}A_\nu^{\alpha}\cdot x\right\|_1\le \sum_{E,E'\in \spec(H)}|\widehat f(E'-E)|^2\left\|P_{E}(A^{\alpha})^{\dagger}P_{E'}A^{\alpha}P_Ex\right\|_1 \\&\le\|\widetilde H^{-\gamma}A^{\alpha}\| \|\widetilde H^{\gamma}A^{\alpha}\widetilde H^{-\mu}\|\|\widetilde H^{\kappa_2+\mu-\gamma}x\|_1 \sum_{E\in\spec(H)} F(E) \left(1+h_0+E\right)^{-(\kappa_2-\gamma)} \lesssim\|x\|_{\mathcal{W}^{\kappa_2+\mu-\gamma,0}_H}.
\end{align*}
The other term in the anticommutator part can be treated similarly. 

Using that by construction $\mathcal{L}_{0,\widehat f,H}$ is closed and employing Lemma~\ref{mathFcoreW} again, we see from the above inequalities and the same argument as in  Proposition~\ref{prop.generationtheorem} that $D(\mathcal{W}_H^{\kappa_2,\kappa_3})\cap D(\mathcal{W}_H^{\kappa_2+\mu-\gamma,0})\cap D(\mathcal{W}_H^{0,\kappa_2+\mu-\gamma})\subseteq D(\mathcal{L}_{0,\widehat{f},H}).$  
\end{proof}





\section{Spectral gap}\label{sec:spectralgaps}

Now that we have rigorously constructed our semigroups of quantum channels that will serve as our Gibbs samplers, we need to develop a framework to study their convergence towards their fixed point $\sigma_\beta$. In this section, we utilize the spectral properties of the operator $L_{\widehat{f},H}$ in order to control the mixing time
\[
t_{\operatorname{mix}}(\epsilon,\mathscr{S}_{D_p(\sigma_\beta)})
\!:=\!\inf\Bigl\{ t \ge 0 \Big| \| \rho_t - \sigma_\beta \|_1 \le \epsilon \,\forall\rho \in \mathscr{S}_{D_p(\sigma_\beta)}\! \Bigr\}
\]
for a given subset $\mathscr{S}_{D_p(\sigma_{\beta})}$ of input quantum states $\rho$ with $\widehat{D}_p(\rho\|\sigma_\beta)\le D_p(\sigma_\beta)<\infty$, $p>1$. As explained in Section \ref{sec.fastconv}, given $\lambda_2\equiv \operatorname{gap}(L_{\widehat{f},H})$, 
\begin{align*}
\big\|e^{t\mathcal{L}_{\sigma_E,\widehat{f},H}}(\rho)-\sigma_{\beta}\big\|_1\le e^{-\lambda_2 t+2{D}_2(\sigma_\beta)}
\quad \Longrightarrow \quad t_{\operatorname{mix}}(\epsilon,\mathscr{S}_{D_p(\sigma_\beta)})\le  \frac{1}{\lambda_2}\Big(2D_2(\sigma_\beta)+\log\Big(\frac{1}{\epsilon}\Big)\Big).
\end{align*}
Thus, it suffices to control the spectral gap in order to control the mixing time. 

\subsection{The harmonic oscillator}\label{sec:gaussexamples}


In order to build some intuition, we first consider the simple case of a Gaussian thermal state over a single-mode quantum bosonic system. Here, $\mathcal{H} = L^2(\mathbb{R})$ and we denote by $a$ and $a^\dagger$ the annihilation and creation operators, defined on a common dense domain $\mathcal{S}(\mathbb{R})$ of Schwartz functions, where they satisfy the canonical commutation relation
\[
[a,a^\dagger] = I.
\]
The associated number operator is given by
\begin{align}\label{def:Numberop}
N := a^\dagger a = \sum_{n \in \mathbb{N}} n\, |n\rangle\langle n|,
\end{align}
where $\{ |n\rangle \}_{n \in \mathbb{N}}$ denotes the Fock basis. We recall that these operators act on the Fock basis as
\[
a|n\rangle = \sqrt{n}\,|n-1\rangle,
\qquad
a^\dagger |n\rangle = \sqrt{n+1}\,|n+1\rangle.
\]
We start by choosing our bare jumps as $\{a,a^\dagger\}$ and considering the Hamiltonian $H=\gamma N$. Then, given a function $\widehat{f}$ satisfying the conditions of Section \ref{sec:generalframe}, we conclude that
\begin{align*}
L^+:=\sum_\nu \widehat{f}(\gamma\nu)\, (a^\dagger)_\nu =\sum_{k\in\mathbb{Z}}\sum_{n\in\mathbb{N}}\widehat{f}(\gamma k)|n+k\rangle\langle n+k|a^\dagger |n\rangle\langle n|
=\sum_{n\in\mathbb{N}}\, \widehat{f}(\gamma)\,\sqrt{n+1}|n+1\rangle\langle n|=\widehat{f}(\gamma)a^\dagger.
\end{align*}
Similarly,
\begin{align*}
L^-:=\sum_\nu \widehat{f}(\gamma\nu)\, a_\nu
=\sum_{k\in\mathbb{Z}}\sum_{n\in\mathbb{N}}\widehat{f}(\gamma k)|n+k\rangle\langle n+k|a |n\rangle\langle n|
=\sum_{n\in\mathbb{N}}\, \widehat{f}(-\gamma)\,\sqrt{n}|n-1\rangle\langle n|=\widehat{f}(-\gamma)a.
\end{align*}
Computing the last two terms of \eqref{eq:cLonF} similarly, we easily derive
\begin{align}\label{eq:qOUgen}
\mathcal{L}_{\widehat{f},\gamma N}(\rho)=&|\widehat{f}(\gamma)|^2\,\Big(a^\dagger \rho a-\frac{1}{2}\{aa^\dagger,\rho\}\Big)+|\widehat{f}(-\gamma)|^2\,\Big(a\rho a^\dagger-\frac{1}{2}\{a^\dagger a,\rho\}\Big)\,.
\end{align}
This coincides with the generator of the quantum Ornstein-Uhlenbeck (qOU) semigroup \cite{OU}, with a birth rate $\nu_+:=|\widehat{f}(\gamma)|^2$ and a death rate $\nu_-:=|\widehat{f}(-\gamma)|^2$. Moreover, since $\overline{\widehat{f}(\gamma)}=\widehat{f}(-\gamma)e^{-\beta \gamma/2}$, it is clear that 
$\nu_+=e^{-\beta\gamma }\nu_-<\nu_-$. 

\smallskip
 
 Next, we choose $\gamma=1$ and now consider the generator $L_{\widehat{f},\gamma N}$ on Hilbert-Schmidt operators, as defined in Equation \eqref{def:L}. A direct computation yields \cite{OU}
\begin{align}
  \label{eq:HSgenerator}
L_{\widehat{f},N}(x) =& -\left( \frac{\nu_- + \nu_+}{2}(Nx + xN) + \nu_+ x \right)  + \sqrt{\nu_+ \nu_-}( a x a^{\dagger} +  a^{\dagger} x a) .
\end{align}
Cipriani, Fagnola, and Lindsay \cite[Thm.~7.2]{OU} showed that the spectrum of \( L_{\widehat{f},N} \) is
\[
\operatorname{Sp}(L_{\widehat{f},N}) = -\left\{ n\left( \frac{\nu_- - \nu_+}{2} \right) \;\middle|\; n \in \mathbb{N}_0 \right\}.
\]
In the limiting case \( \nu_- = \nu_+ \), the spectrum becomes continuous and fills the half-line \( [0, \infty) \), corresponding to the so-called quantum Brownian motion \cite[Thm. 8.1]{OU}.


\smallskip

\paragraph{Why Work in the Hilbert-Schmidt Space?}
One might ask why we study \( L_{\widehat{f},N} \) on the Hilbert-Schmidt space instead of working with \( \mathcal{L}_{\widehat{f},N} \) directly on the set \( \mathscr{T}_1(\mathcal{H}) \) of trace-class operators. The reason is that \( \mathcal{L}_{\widehat{f},N} \) lacks a spectral gap in the space of trace-class operators, as the next theorem demonstrates. We recall that, for a linear operator $(A,D(A))$ on a Banach space $\mathscr{B}$, its resolvent $\rho(A)$ is the set of complex numbers $\lambda$ such that $A-\lambda I$ has a bounded inverse $B:\mathscr{B}\to D(A)$, i.e., $(A-\lambda I)B=I$ and $B(A-\lambda I)=I_{D(A)}$. The spectrum $\operatorname{Sp}(A)$ is the complementary set of $\rho(A)$ in $\mathbb{C}$. The spectrum of a closed operator is closed, and that of the generator of a strongly continuous semigroup is included in $\mathbb{C}_-$. A necessary condition for the uniform exponential convergence of the strongly continuous contraction semigroup associated with $\mathcal{L}_{\widehat{f},N}$ is the existence of a spectral gap, namely a constant $\lambda_0<0$ such that for all $\lambda\in \operatorname{Sp}(\mathcal{L}_{\widehat{f},H})\backslash \{0\}$, $\operatorname{Re}(\lambda)\le \lambda_0$ \cite[Corollary 4.1.2]{salzmann2024robustness}.


\begin{proposition}\label{thm.qOUnogap}
The spectrum of the OU generator \( \mathcal{L}_{\widehat{f},N} \) on the space of trace-class operators is the entire closed left complex half-plane \( \mathbb{C}_- \).
\end{proposition}

\begin{proof}
We recall that the set $\sigma_{\mathrm{p}}(A)$ of eigenvalues $\lambda$ of a linear map $(A,D(A))$ is included in $\operatorname{Sp}(A)$. The formal adjoint of the generator of the classical Ornstein-Uhlenbeck on $L^\infty(\mathbb{R})$ 
\[
G_{\mathrm{OU}} \varphi(x) = q \varphi''(x) - b x \varphi'(x),
\]
 is
\[
G_{\mathrm{OU}}' \varphi(x) = q \varphi''(x) + b x \varphi'(x) + b\varphi(x).
\]
To find its spectrum, we solve the eigenvalue problem
\[
G_{\mathrm{OU}}' \varphi = \lambda \varphi \quad \Rightarrow \quad q \varphi''(x) + b x \varphi'(x) - \mu \varphi(x) = 0,
\]
where \( \mu = \lambda - b \). 
Taking the Fourier transform of this equation yields
\[
- q \xi^2 \widehat{\varphi}(\xi) - b\left( \widehat{\varphi}(\xi) + \xi \widehat{\varphi}'(\xi) \right) - \mu \widehat{\varphi}(\xi) = 0.
\]
Solving this differential equation in Fourier space gives
\[
\widehat{\varphi}(\xi) = c e^{ -\frac{q \xi^2}{2b} } |\xi|^{ -\frac{\lambda}{b} },
\]
for some \( c \neq 0 \). Taking the inverse Fourier transform yields
\[
\varphi(x) = \tfrac{c}{2} \left( \tfrac{2b}{q} \right)^{\frac{b-\lambda}{2b}} \Gamma\left( \tfrac{b-\lambda}{2b}\right) \, {}_1F_1\left( \tfrac{b-\lambda}{2b}, \tfrac{1}{2}, -\tfrac{b x^2}{2q} \right),
\]
where $F_1$ denotes the confluent hypergeometric function. Asymptotically, we have
\[ \vert \varphi(x)\vert = \mathcal O(x^{-1+\frac{\lambda}{b}}).\]
 Thus, $\varphi \in L^1(\mathbb R)$ for $\Re(\lambda)<0$. This shows that $\mathbb C_- \subset \operatorname{Sp}(G'_{\operatorname{OU}})= \operatorname{Sp}(G_{\operatorname{OU}})$. Since $\mathcal L_{\widehat{f},N}$ generates a contraction semigroup, we have $\operatorname{Sp}(\mathcal L_{\widehat{f},N}) \subset \mathbb C_-.$ 
Finally, it was shown in \cite{OU} that, given the multiplication operator \( M_\varphi \) defined by \( M_\varphi(x) = \varphi(x) \),
 the generator $\mathcal{L}_{\widehat{f},N}$ satisfies
\[
G_{\mathrm{OU}} \varphi(x) := \mathcal{L}_{\widehat{f},N}(M_\varphi)(x) = q \varphi''(x) - b x \varphi'(x),
\]
with
\[
q := \frac{\nu_- + \nu_+}{4} > 0, \quad b := \frac{\nu_- - \nu_+}{2} > 0.
\]
  We conclude that
 \[ \mathbb C_- \subset \operatorname{Sp}(G_{\operatorname{OU}}) \subset \operatorname{Sp}(\mathcal L_{\widehat{f},N}) \subset \mathbb C_- \Longrightarrow \operatorname{Sp}(\mathcal L_{\widehat{f},N})=\mathbb C_-. \]
 
\end{proof}





























































\subsection{Independence of spectra}\label{Lpsemigroupsindependence}


\noindent Next, we show how to relax the condition that input states are in $\mathscr{S}_{D_2(\sigma_\beta)}$ to $\mathscr{S}_{D_p(\sigma_\beta)}$ for $1<p<2$, while keeping the same convergence rate given by the spectral gap of $L_{\widehat{f},H}$. As $p$ approaches $1$, this allows us to consider increasingly larger sets of initial states, although still under an exponential moment constraint. For our analysis, we shall use the following family of interpolating Banach spaces:
Let \(\omega\in\mathscr B(\mathcal H)\) be a faithful state (\(\omega>0\), \(\Tr\omega=1\)).
For \(1\le p<\infty\) define on \(\mathscr{B}(\mathcal{H})\)
\begin{equation}
\label{eq:norms}
\begin{split}
  \|x\|_{p,\omega}
  := \Big(\Tr\,\big|\,\omega^{1/(2p)}\,x\,\omega^{1/(2p)}\big|^{\,p}\Big)^{1/p} =\Big\|\omega^{1/(2p)}\,x\,\omega^{1/(2p)}\Big\|_{p}.
  \end{split}
\end{equation}
The non\mbox{-}commutative \(L^p\) space \(L^p(\omega)\) is the completion of \(\mathscr{B}(\mathcal H)\)
with respect to \(\|\cdot\|_{p,\omega}\). 
Set \(L^\infty(\omega):=\mathscr{B}(\mathcal H)\) with the operator norm \(\|\cdot\|_\infty\).
If \(1/p+1/q=1\), the duality pairing between \(L^p(\omega)\) and \(L^q(\omega)\) is
\[
  \langle x, y\rangle_\omega
  := \Tr\!\left(\omega^{1/(2p)}\,x\,\omega^{1/(2q)}\,y^{\dagger}\right),
 \ x\in L^p(\omega),\ y\in L^q(\omega),
\]
so that \(L^q(\omega)\cong L^p(\omega)^{*}\).
In particular, for \(p=2\),
\(\langle x,y\rangle_{2,\omega}=\Tr(\omega^{1/4}x^{*} \omega^{1/4}y)\) and \(L^2(\omega)\) is a Hilbert space.
 Next, we build a family of semigroups $\{T^{(p)}_t\}_{t\ge 0}$ on $L^p(\sigma_{\beta})$ for each $p\ge 1$ from the $\mathscr{T}_2(\cH)$ semigroup $\{e^{\smash{tL_{\widehat{f},H}}}\}_{\smash{t\ge 0}}$. For this, we first introduce the isometry $\eta_2: L^2(\sigma_\beta)\to \mathscr{T}_2(\cH)$, with $\eta_2(x)=\sigma_\beta^{\smash{\frac{1}{4}}}x\sigma_\beta^{\smash{\frac{1}{4}}}$. Note that, formally, $\eta_2$ coincides with $\iota_2$. Then, for each $t\ge0$, we define the map $T_t^{(2)}$ as 
\[T^{(2)}_t=\eta_2^\dagger\, e^{tL_{\widehat{f},H}}\,\eta_2.\]
Next, we start from the $L^2(\sigma_\beta)$ semigroup with spectral resolution
\[ T^{(2)}_t = \sum_{k=0}^{\infty} e^{-\lambda_k t} P_k \,, \]
where $\{\lambda_k\}_k$ form a discrete set of eigenvalues, tending to infinity, and $P_k$ the associated spectral projections. 
Then $\{T^{(2)}_t\}_{t>0}$ is a family of compact operators on $L^2(\sigma_\beta)$, since it is the norm limit of finite rank operators $T^{(2,n)}_t:=\sum_{k=0}^{n} e^{-\lambda_kt} P_k.$ By \cite[Theorem 1.4.1]{Davies_1989}, it follows that, since $T^{(2)}_t$ is a symmetric Markov semigroup on $L^2(\sigma_\beta)$, it can be extended from $L^1(\sigma_\beta) \cap L^\infty(\sigma_\beta)$ to a positive one-parameter contraction
semigroup $T^{(p)}_t$ on $L^p(\sigma_\beta)$ for all $1 \le p \le \infty$. These semigroups
are strongly continuous if $1 \le p < \infty$, and are consistent in the
sense that
\[
    T^{(p)}_t(x) = T^{(q)}_t(x)
\]
if $x \in L^p(\sigma_\beta) \cap L^q(\sigma_\beta) = L^{\text{max}\{p,q\}}(\sigma_\beta)$. They are self-adjoint in the sense that
\begin{equation}
\label{eq:duality}
    T^{(p)\dagger}_t = T^{(q)}_t\text{ if }1 \le p < \infty\text{ and }p^{-1} + q^{-1} = 1.
\end{equation}
\noindent By interpolation, since $T^{(2)}_t$ is an analytic semigroup which follows for instance from \cite[Section 4, Theorem 4.6]{engel1999one}, all semigroups $T^{(p)}_t$, for $p \in (1,\infty)$ are analytic, too. See \cite[Theorem 6.6]{Lunardi2018} for a thorough discussion using Stein's interpolation theorem. 
It also follows from interpolation, see \cite[Theorem 3.1]{COBOS1990351}, that since $T^{(2)}_t$ with $t>0$ is compact on $L^2(\sigma_\beta)$, $T^{(p)}_t$ is also compact on $L^p(\sigma_\beta)$ for $p \in (1,\infty)$. 

Denoting by $\mathcal{L}^{(p)}_{\widehat{f},H}$ the generators of the semigroups $\{T^{(p)}_t\}_{t\ge 0}$, by the Laplace transform, we have \[ (\lambda-\cL^{(p)}_{\widehat{f},H})^{-1}  =\int_0^{\infty} e^{-\lambda t} T^{(p)}_t  \ dt.\]
This implies that the resolvent is also compact on all $L^p(\sigma_\beta)$ spaces for $p\in (1,\infty)$, and thus the spectrum of $\mathcal{L}^{(p)}_{\widehat{f},H}$, the generator of $\{T^{(p)}_t\}_{t\ge 0}$, is discrete. Every $L^2(\sigma_\beta)$ eigenfunction is automatically also a $L^p(\sigma_\beta)$ eigenfunction for $1\le p<2$, since these are weaker norms. On the other hand, every $L^p(\sigma_\beta)$ eigenfunction for $p>2$ is automatically an $L^2(\sigma_\beta)$ eigenfunction. We have thus shown using \eqref{eq:duality} the following:
\begin{theorem}
    The spectrum of $\mathcal{L}^{(p)}_{\widehat{f},H}$ is independent of $p \in (1,\infty).$
\end{theorem}

\noindent An important property of analytic semigroups $\{T_t\}_{t\ge 0}$, with generator $A$, is that their growth bound
\[\omega(A):=\inf_{t>0} \frac{1}{t} \log \Vert T(t)\Vert\]
is equal to their spectral bound 
\[ s(A):=\sup\{\Re(\lambda); \lambda \in \operatorname{Sp}(A) \}; \]
and the spectral mapping theorem holds 
\[ \operatorname{Sp}(T(t)) \setminus \{0\} = e^{\overline{t \operatorname{Sp}(A)}} \setminus \{0\} \] 
see \cite[\SS 3, Corr. 3.12]{engel1999one}. For contraction semigroups as above on spaces $L^p(\sigma_\beta)$, we may study the semigroup on $(I-P)L^p(\sigma_\beta),$ where $P$ is the projection onto the $0$ eigenspace of the generator. By the existence of a spectral gap, we thus have that there exists $\lambda>0$, equal to the spectral gap, and $M_p>0$ such that
\[\Vert (T^{(p)}_t-P)x \Vert_{p,\sigma_\beta} \le M_p e^{-\lambda t} \Vert x \Vert_{p,\sigma_\beta}. \]
 We have thus shown the following in our setting:
\begin{proposition}\label{thm:fromL2toLp}
    Assume that $L_{\widehat{f},H}$ has a compact resolvent, and let $\operatorname{gap}(L_{\smash{\widehat{f},H}})>0$ be the spectral gap between eigenvalue 0 and its next largest eigenvalue. Then, there exists $M_p>0$ such that for all $p \in (1,\infty)$ and $P_p$, the spectral projection associated with the $0$ eigenvalue in $L^p(\sigma_\beta)$,
    \[\Vert (e^{t\cL_{\widehat{f},H}^{(p)}}-P_p)x \Vert_{p,\sigma_\beta} \le M_p\, e^{-\operatorname{gap}(L_{\widehat{f},H}) t} \Vert x \Vert_{p,\sigma_\beta} ~~\text{ for all }~~ x \in L^p(\sigma_\beta). \]
\end{proposition}











































\subsection{Absence of gap for decaying filter functions}
\label{sec:AbsenceOfGap}



An important class of models consists of Hamiltonians $H$ that can be formally expressed as polynomials in creation and annihilation operators. Among these, the Bose--Hubbard model is particularly prominent and is analysed in our companion paper \cite{BeckerRouzeSalzmannBose}. As expected, the difficulty when analyzing such Hamiltonians increases significantly when the degree of the polynomial exceeds $2$. In this section, we demonstrate that the generator $L_{\smash{\widehat{f},N^2}}$ is gapless whenever $\widehat{f}$ is excessively regular, thereby revealing a first explicit tension with the implementability of the corresponding evolution $e^{t\mathcal{L}_{\widehat{f},N^2}}$.
More broadly, we consider the properties of Hamiltonians 
of the form 
\begin{align}
\label{eq:NumPresHamiltonian}
    H = h(N) = \sum_{n=0}^\infty h(n) \kb{n}.
\end{align}
where $N$ is the number operator defined in \eqref{def:Numberop}. Here, $\left\{\ket{n}\right\}_{n\in\N_0}$ denotes the Fock basis, and we consider $h:\N_0\to \R$ to be non-decreasing for $n\ge n_0$ for some $n_0\in\N_0$, and such that $H$ has a well-defined Gibbs state. In the following proposition we consider for $\sigma_E\in(0,\infty]$ the generator $\cL_{\sigma_E,\widehat f,H}$ defined and studied in Section~\ref{monotonicitygap} and \ref{integralrep} below and which generalises $\cL_{\widehat f,H}$  as $$\cL_{\infty,\widehat f,H}=\cL_{\widehat f,H}.$$ Note, as established in Proposition~\ref{lem.gapmonoton} below, we have that the spectral gap of the corresponding generator on the space of Hilbert-Schmidt operators, $L_{\sigma_E,\hat f,H},$ is non-increasing in the parameter $\sigma_E,$  i.e.
$\operatorname{gap}(L_{\sigma_E,\widehat{f},H})\ge \operatorname{gap}(L_{\sigma'_E,\widehat{f},H})$  for any  $0< \sigma_E\le \sigma'_E\le \infty.$

\begin{proposition}\label{prop.no-gogap}
    Let the filter function $\widehat{f}\in\mathcal{S}(\R)$ be such that
    \[\lim_{m \to \infty} \widehat{f}(\pm (h(m+1)-h(m))) \sqrt{m}=0,\]
    then the generator $\mathcal{L}_{\sigma_E,\widehat{f},H}$ as in \eqref{eq:LsigmaEintro}  with bare jumps $\{a,a^\dagger\}$ is compact for $\sigma_E \in (0,\infty]$. In particular, $0$ is part of the essential spectrum.

    
      If the filter function $\widehat{f}$ satisfies 
    \[ \lim_{m \to \infty} e^{\pm \beta (h(m)-h(m+1))/4}\widehat{f}(\pm ( h(m)-h(m+1))) \sqrt{|m|}=0,\]
    where we consider all possible combinations of $\pm$, then $L_{\sigma_E,\widehat{f},H}$ is compact for $\sigma_E \in (0,\infty]$ as well.
 This holds, for instance, if $\widehat{f}$ has Gaussian decay. 
\end{proposition}
\begin{proof}
Denoting $L^a$ the operator $L^\alpha$ corresponding to the bare jump $a$ (cf.~\eqref{eq:LalphainGenerationTheory}), we have 
\[L^a  =\sum_{n,n'}\ket{n}\bra{n}a\ket{n'}\bra{n'}\widehat{f}(h(n)-h(n'))=\widehat{f}(h(N)-h(N+1))\, a.  \]
We can write the annihilation operator as $a=\sqrt{N+1} U$ with $U\vert n\rangle := \vert n-1 \rangle$, where $U$ is a (bounded) shift operator.
Thus, under the assumption that 
\[ \lim_{m \to \infty} \widehat{f}(h(m)-h(m+1)) \sqrt{m+1}=0, \]
it follows that the operator $L^a$ is a compact operator, as $L^a$ can be approximated by finite rank operators 
\[\lim_{m \to \infty} 1_{[0,m]}(N)\widehat{f}(h(N)-h(N+1)) a =\widehat{f}(h(N)-h(N+1)) a.\]
This implies by \cite[Theorem C.7]{engel1999one} that  \begin{align*}
 G_{\sigma_E}:=-\sum_{\alpha\in\mathcal{A}}\int_{-\infty}^\infty g(t)\,e^{-itH}((L^{\alpha})^\dagger L^\alpha)e^{itH}dt,
\end{align*}
with 
$g(t)=\frac{1}{2\pi}\int_{-\infty}^\infty \frac{e^{-\nu^2/8\sigma_E^2}}{1+e^{\beta\nu/2}} e^{-i\nu t}d\nu$, $X^{\alpha}_s:=e^{isH}L^\alpha e^{-isH}$ are compact, and a CP map 
\begin{align*}
\Phi_{\sigma_E,\widehat{f},h(n)}(\rho):=\sigma_E\sqrt{\frac{2}{\pi}}\sum_{\alpha}\int_{\mathbb{R}}e^{-2\sigma_E^2 s^2}\, \, X_s^{\alpha} \rho (X_s^\alpha)^\dagger \,ds
\end{align*}
are compact operators, too. In the case of $\Phi_{\sigma_E,\hat{f},h(n)},$ this argument applies up to $\sigma_E=\infty.$ For $G$ and $\sigma_E=\infty$, we use the representation
\begin{equation}
    \begin{split}
G:&=-\sum_{\substack{\alpha\in \mathcal{A}\\E,E',E''\in\spec(H)}}  \frac{\overline{\widehat{f}(E'-E'')}\widehat{f}(E'-E)\,e^{\frac{\beta(E-E'')}{4}}}{2\cosh((E-E'')\beta/4)}P_{E''}\left(A^\alpha\right)^\dagger P_{E'} A^\alpha P_E \\
&=\sum_{\substack{\alpha\in \mathcal{A}\\n \in \mathbb N_0}}  \frac{\overline{\widehat{f}(E'-E'')}\widehat{f}(E'-E)\,e^{\frac{\beta(E-E'')}{4}}}{2\cosh((E-E'')\beta/4)}P_{h(n)}\left(A^\alpha\right)^\dagger P_{h(n\pm 1)} A^\alpha P_{h(n)} \\
&=\sum_{\substack{\alpha\in \mathcal{A}\\n \in \mathbb N_0}}  \vert\widehat{f}(h(n)-h(n\pm 1))\vert^2P_{h(n)}\left(A^\alpha\right)^\dagger P_{h(n\pm 1)} A^\alpha P_{h(n)}
    \end{split}
\end{equation}
which shows that this operator $G$ is a uniform limit of the finite rank approximations for $m \in \mathbb N$
\begin{equation}
    \begin{split}
G_m:
&=\sum_{\substack{\alpha\in \mathcal{A}\\n \in [m]}}  \vert\widehat{f}(h(n)-h(n-1))\vert^2P_{h(n)}\left(A^\alpha\right)^\dagger P_{h(n-1)} A^\alpha P_{h(n)}
    \end{split}
\end{equation}
and thus compact.

Doing a similar computation for $a^{\dagger}$, we find that under the condition that \[ \lim_{m \to \infty} \widehat{f}(\pm (h(m+1)-h(m))) \sqrt{m}=0, \]
the generator $\cL_{\sigma_E,\widehat{f},h(N)}$ is a compact operator, too.

\smallskip

This can be extended to the generator $L_{\sigma_E,\widehat{f},h(N)}$ as well. In this case, one has to conjugate $L^a$ by appropriate Gibbs states to find 
\[ \Delta^{\pm 1/4}(L^a)=\widehat{f}(h(N)-h(N+1)) e^{\pm \beta (h(N)-h(N+1))/4} a \]
and 
\[ \Delta^{\pm 1/4}(L^a)=\widehat{f}(h(N-1)-h(N)) e^{\pm \beta (h(N-1)-h(N))/4} a^{\dagger} \]
with $ \Delta(\cdot )=\sigma_\beta[\cdot ]\sigma_\beta^{-1}.$
Thus, let $\widehat{f}$ be such that $e^{\pm \beta (h(m)-h(m+1))/4}\widehat{f}(\pm  h(m)-h(m+1)) \sqrt{|m|}=0$; then the $\mathscr T^2$ generator $L_{\sigma_E,\hat f,h(N)}$ is a compact operator as well.
\end{proof}
\noindent Since compact operators always have $0$ in their essential spectrum, the generator does not have a spectral gap. In other words, the compactness of $L_{\sigma_E,\widehat{f},h(N)}$ implies, by the spectral theorem and the dominated convergence theorem, that we still have pointwise convergence 
\[\begin{split} \lim_{t \to \infty} \left\lVert e^{L_{\sigma_E,\widehat{f},h(N)}t} (x)-\langle \sqrt{\sigma_\beta},x\rangle\,\sqrt{\sigma_\beta}\right\rVert_2^2&=\lim_{t\to\infty}\left\lVert  \sum_{\lambda \in \operatorname{Spec}(L_{\sigma_E,\widehat{f},h(N)})} e^{\lambda t} \langle x,y_{\lambda} \rangle y_{\lambda} - \langle \sqrt{\sigma_\beta},x\rangle\,\sqrt{\sigma_\beta} \right\rVert_2^2\\
&=\sum_{\lambda \in \operatorname{Spec}(L_{\sigma_E,\widehat{f},h(N)});\lambda<0} \lim_{t \to \infty} e^{2\lambda t} \vert \langle x,y_{\lambda} \rangle \vert^2=0,  
\end{split}\]
where $y_\lambda$ corresponds to an eigenvector associated to eigenvalue $\lambda$, but there is no uniform convergence unless $\widehat{f}$ is suitably chosen. In Section \ref{sec:GeneralNumberviaBirthDeath} below, we make a special choice of function $\widehat{f}$, illustrated in Figure \ref{fig:placeholder}, which violates the compactness condition above and allows for the existence of a spectral gap.


\subsection{Single-mode number preserving Hamiltonians}
\label{sec:GeneralNumberviaBirthDeath}



\noindent
In the previous section, we saw that choosing an excessively regular function $\widehat{f}$ may lead to gapless generators $L_{\widehat{f},H}$. Here instead, we consider a function which only decays in one direction, akin to the classical Metropolis-Hastings rate function. In the following, we denote $E_n= h(n)$ for the eigenvalues of $H$, consider the bare jump operators $\{a,a^\dagger\}$ and the function for $\beta>0$
\begin{align}
\label{eq:Function}
    \hatfM(\nu) = \exp\left(-\frac{\sqrt{1+(\beta\nu)^2}+ \beta\nu}{4}\right).
\end{align}


 \begin{figure}[h!]
     \centering
     \includegraphics[width=0.3\linewidth]{hatf.pdf}
     \caption{Metropolis-type filter $\hatfM$ in \eqref{eq:Function} for $\beta=1$ smoothly approximating a step function.}
     \label{fig:placeholder}
 \end{figure}

\noindent With that, the corresponding jump operators from Proposition~\ref{prop.generationtheorem} are given by
\begin{align*}
    &L^+ =\sum_{n,m=0}^{\infty} \hatfM(E_n - E_m) \kb{n}a^\dagger\kb{m} =a^\dagger\,\hat f(h(N+1)-h(N))\\
    &L^-   =  \sum_{n}^{\infty} \hatfM(E_{n-1} - E_n)\sqrt{n} \ket{n-1}\!\bra{n} = a\hatfM(h(N-1)-h(N)).
\end{align*}
From this, we can formally write the corresponding generator from Proposition~\ref{prop.generationtheorem} as


\begin{align}
  \label{eq:LindbladNumPresHamiltonian} \cL_{\hatfM,h(N)}(x) &= L^+ x(L^+)^\dagger + L^- x (L^-)^\dagger - \frac{1}{2}\Big\{\left((L^+)^\dagger L^+ + (L^-)^\dagger L^-\right),x\Big\} \\
  &=\nn a^\dagger\,\hatfM(h(N+1)-h(N))x \hatfM(h(N+1)-h(N))\,a  \\
&\quad+ a\,\hatfM(h(N-1)-h(N))x \hatfM(h(N-1)-h(N))\, a^\dagger\nn\\
  &\quad-\frac{1}{2} \Big\{(N+1)|\hatfM(h(N+1)-h(N))|^2,x\Big\}-\frac{1}{2}\Big\{ N|\hatfM(h(N-1)-h(N))|^2,x\Big\}.\nn
\end{align}


\noindent Next, we consider the spectral gap of the associated generator $L_{\hatfM,h(N)}$ on the space of Hilbert-Schmidt operators $\mathscr{T}_2(\mathcal{H})$ through the relation \eqref{eqcLtoL}.
A direct computation shows that, e.g.~on $\mathscr{F}$,
\begin{equation}
\label{eq:L22}
\begin{split}
  L_{\hatfM,h(N)}(y) &= a^\dagger\, g_+(N) y g_+(N)\,a  + a\,g_-(N) y g_-(N)\, a^\dagger -\frac{1}{2} \Big\{(N+1)|\widehat f(h(N+1)-h(N))|^2, y\Big\}\\&\quad-\frac{1}{2}\Big\{ N|\widehat f(h(N-1)-h(N))|^2,y\Big\},
\end{split}
\end{equation}
where we denote the functions
\begin{equation}
\begin{split}
    g_{+}(n) &= \hatfM(h(n+1)-h(n))e^{\frac{\beta(h(n+1)-h(n))}{4}}\text{ and }
    g_{-}(n) = \hatfM(h(n-1)-h(n))e^{-\frac{\beta(h(n)-h(n-1))}{4}}.
    \end{split}
\end{equation}

\begin{theorem}[Spectral gap of $L_{\hatfM,h(N)}$]
\label{thm:SpectralGapGeneralh(N)}
    Let $n_0\in\N_0$ and $H$ be of the form \eqref{eq:NumPresHamiltonian} with eigenvalues $E_n$ being non-decreasing for all $n\ge n_0.$ Assume that there exists $\delta>0$ and $s\in\N$ such that
\begin{align}
\label{eq:LowerboundEn+s-En}
    E_{m+s} -E_m \ge \delta
    \end{align}
    for all $m\ge n_0.$
    Then for all $\beta>0$ we have $\operatorname{gap}(L_{\hatfM,h(N)})>0.$ Furthermore, we can lower bound \begin{align}
    \label{eq:UniformGapBound}
    \operatorname{gap}(L_{\hatfM,h(N)}) \ge \kappa(\beta,n_0,\delta,s,\Delta_E) >0
\end{align}
for a constant $\kappa(\beta,n_0,\delta,s,\Delta_E)>0$ only depending on $\beta,n_0,\delta,s$ and $\Delta_E :=\max_{0\le j,m\le 2n_0} |E_j-E_m|.$
    \bin{Furthermore, for $n_0=0,$ i.e. $\left(E_n\right)_{n\in\N_0}$ being a non-decreasing sequence we have that $\operatorname{gap}(L)=\Omega(1)$ in the limit of large $\beta,$ i.e. for all $\beta'>0$ we have $\operatorname{gap}(L_{\widehat{f},h(N)})\ge C_{\beta'}$ for some $C_{\beta'}>0$ and all $\beta\in[\beta',\infty).$}
\end{theorem}

\noindent In order to prove  Theorem~\ref{thm:SpectralGapGeneralh(N)}, we first notice 
that the generator $\mathcal{L}_{\hatfM,h(N)}$ in \eqref{eq:LindbladNumPresHamiltonian} can be seen as a quantum birth and death generator, as studied in \cite{carboneFagnola_exponentiall2-convergence_2000}, with square roots of the birth and death rates given as\footnote{To see the scaling relations captured by $\sim$ in \eqref{eq:ScaleBirthing} and \eqref{eq:ScaleDying}, we first note that since the $E_n$ are non-decreasing for $n\ge n_0$, $\inf_{n\ge0}(E_{n+1}-E_n) = \min_{0\le n\le n_0}\{E_{n+1}-E_n,0\} \ge -\max_{0\le j,m\le n_0} |E_j-E_m|$ and $\sup_{n\ge0}(E_{n-1}-E_n) = \max_{0\le n\le n_0}\{E_{n-1}-E_n,0\} \le \max_{0\le j,m\le n_0} |E_j-E_m|.$ \eqref{eq:ScaleBirthing} and \eqref{eq:ScaleDying} then follow by noting
that for $C\in\R$ and $\nu\ge C$
, we have $
e^{-\frac{\sqrt{1+(\beta C)^2}-\beta C}{4}}\,
e^{-\beta \nu/2}
\;\le\;
\hatfM(\nu)
\;\le\;
e^{-\beta \nu/2}$ 
and $e^{-\frac{\sqrt{1+(\beta C)^2}-\beta C}{4}}>0$; further, for $c\in\R$ and $\nu\le c$, we have
$
0<\hatfM(c)\le \hatfM(\nu)\le 1.$
}
\begin{align}
   & \mu^+_n\! =\! \sqrt{n+1}\hatfM(E_{n+1}-E_n) 
   \!=\!\sqrt{n+1} \,e^{-\frac{\sqrt{1+(\beta(E_{n+1}-E_n))^2}+\beta(E_{n+1}-E_n)}{4}}\!\!\sim\!\label{eq:ScaleBirthing} \sqrt{n+1}\,e^{-\beta\frac{E_{n+1}-E_n}{2}}\!\!,\\
    & \mu^-_n = \sqrt{n}\,\hatfM(E_{n-1}-E_n) =\sqrt{n} \label{eq:ScaleDying} e^{-\frac{\sqrt{1+(\beta(E_{n-1}-E_n))^2}+\beta(E_{n-1}-E_n)}{4}} \sim \sqrt{n}.
\end{align}
In  \cite[Theorem 4.2]{carboneFagnola_exponentiall2-convergence_2000}, Carbone and Fagnola found a general set of conditions to ensure the positivity of the spectral gap of such quantum birth and death generators. In particular, they proved positivity of spectral gap under the condition
\begin{align}
\label{eq:FagnolaGapCondition0}
    \inf_{n\ge 1}\left((\mu^-_n)^2 + (\mu^+_n-\mu^+_0)^2\right)>0
\end{align}
and assuming that there exists $\gamma\in(0,1)$ and $c,d\ge 0$ such that the following inequalities hold for all $m,k\in\N$:
\begin{align}
    &\sum_{j>m} \sqrt{\bra{j}\sigma_\beta\ket{j}\bra{j+k}\sigma_\beta\ket{j+k}} \le \label{eq:FagnolaGapCondition1} c\, \mu^+_{m}\mu^+_{m+k} \sqrt{\bra{m}\sigma_\beta\ket{m}\bra{m+k}\sigma_\beta\ket{m+k}}
\end{align}
and
\begin{align}
   \sum_{j>m} \!\gamma^{-j}\mu^+_j\mu^+_{j+k}\sqrt{\bra{j}\!\sigma_\beta\!\ket{j}\bra{j+k}\!\sigma_\beta\!\ket{j+k}}\!\le\! \label{eq:FagnolaGapCondition2} d\, \gamma^{-m}\mu^+_{m}\mu^+_{m+k} \sqrt{\bra{m}\!\sigma_\beta\!\ket{m}\bra{m+k}\!\sigma_\beta\!\ket{m+k}}.
\end{align}
 Under these assumptions the spectral gap of $L_{\hatfM,h(N)}$ satisfies the following lower bound
\begin{equation}
\label{eq:LowerBoundGapBirthDeath}
\mathrm{gap}(L_{\hatfM,h(N)}) \ge 
\min \left\{\frac{1}{c} 
\left( d + 1 + \frac{\gamma}{1 - \gamma} \right)
^{-1},
\;
\inf_{n > 0} 
\frac{(\mu^-_n)^2 + (\mu^+_n - \mu^+_0)^2}
{1 + c \mu^+_0 \mu^+_n \left( 1 + \frac{d + 1}{\gamma} \right)}
\right\}>0.
\end{equation}

% For the proof of Theorem~\ref{thm:SpectralGapGeneralh(N)} we verify the conditions \eqref{eq:FagnolaGapCondition0},
% \eqref{eq:FagnolaGapCondition1} and \eqref{eq:FagnolaGapCondition2} to employ \cite[Theorem 4.2]{carboneFagnola_exponentiall2-convergence_2000}. The main work here lies in showing that the assumption \eqref{eq:LowerboundEn+s-En} implies \eqref{eq:FagnolaGapCondition1} and \eqref{eq:FagnolaGapCondition2}.





\begin{proof}[Proof of Theorem~\ref{thm:SpectralGapGeneralh(N)}]
We prove Theorem~\ref{thm:SpectralGapGeneralh(N)} by verifying \eqref{eq:FagnolaGapCondition0}, \eqref{eq:FagnolaGapCondition1} and \eqref{eq:FagnolaGapCondition2} and conclude positivity of the spectral gap by \cite[Theorem 4.2]{carboneFagnola_exponentiall2-convergence_2000}. The condition \eqref{eq:FagnolaGapCondition0} follows directly using \eqref{eq:ScaleBirthing} and \eqref{eq:ScaleDying}.
Hence, we focus in the following on verifying \eqref{eq:FagnolaGapCondition1} and \eqref{eq:FagnolaGapCondition2}: Using \eqref{eq:LowerboundEn+s-En} and Lemma~\ref{lem:EnEquivalences}, in particular Equation \eqref{eq:ExplicitBoundGamma} for $\gamma=1,$ we see 
\bin{\begin{align*}
   \sup_{\substack{m,k\in\N_0\\m+k\ge n_0}}\sum_{j>m} e^{-\beta(E_{j+k}-E_{m+k+1}+E_j -E_{m+1})/2}   \le \sup_{m\in\N_0}\sum_{j\ge m+1} e^{-\beta(E_j -E_{m+1})/2} <\infty
\end{align*}
and therefore, using that also $$\max_{\substack{m,k\in\N_0\\m+k<n_0}} \sum_{j>m} e^{-\beta(E_{j+k}-E_{m+k+1}+E_j -E_{m+1})/2}<\infty $$  we have}
\begin{align}
\label{eq:tildec}
 & \sup_{m,k\in\N_0}\sum_{j>m} e^{-\beta(E_{j+k}-E_{m+k+1}+E_j -E_{m+1})/2}  \le (n_0+1) e^{\beta\Delta_E} \frac{e^{\beta\delta}}{1-e^{-\beta\delta/s}}  =:\tilde c <\infty,
\end{align} 
where we denoted $\Delta_E :=\max_{0\le j,m\le 2n_0} |E_j-E_m|.$
Combining this with \eqref{eq:ScaleBirthing}, we see that for all $k,m\in\N_0$
\begin{align*}
    \sum_{j>m} \!e^{-\frac{\beta(E_{j+k}+E_j)}{2}} &=\! e^{-\frac{\beta(E_{m+k}+E_m)}{2}-\frac{\beta(E_{m+1}-E_m)}{2}-\frac{\beta(E_{m+k+1}-E_{m+k})}{2}}  \sum_{j>m} e^{-\frac{\beta(E_{j+k}-E_{m+k+1}+E_j -E_{m+1})}{2}} \\&\lesssim \tilde c  \,e^{-\beta(E_{m+k}+E_m)/2}\mu^+_m\mu^+_{m+k}, 
\end{align*}
for all $m,k\in\N,$ 
which shows that condition \eqref{eq:FagnolaGapCondition1} is satisfied. Similarly, using again \eqref{eq:LowerboundEn+s-En}
and Lemma~\ref{lem:EnEquivalences}, in particular \eqref{eq:ExplicitBoundGamma}, for $\gamma,\tilde\gamma,\in(e^{-\beta\delta/s},1)\neq \emptyset$ with $\tilde \gamma>\gamma$ and $q= \tilde\gamma^{-1}e^{-\beta\delta/s}<1$ we have

\begin{align} 
\label{eq:tilded}\sup_{\substack{m,k\in\N_0}}\sum_{j>m}  \gamma^{-(j-m)}  &e^{-\beta(E_{j+k+1}-E_{m+k+1}+E_{j+1} -E_{m+1})/2} \sqrt{\tfrac{(j+1)(j+1+k)}{(m+1)(m+k+1)}}\\
&\le C_{\gamma,\tilde \gamma}\sup_{\substack{m,k}}\!
\sum_{j>m}\! \! \tilde \gamma^{-(j-m)}  
e^{-\beta(E_{j+k+1}-E_{m+k+1}+E_{j+1} -E_{m+1})/2} 
\nn\\
&\nn\le C_{\gamma,\tilde \gamma}\,(n_0+1)\tilde\gamma^{-n_0}e^{\beta\Delta_E}\frac{e^{\beta\delta}}{1-q}=:\tilde{d} < \infty,
\end{align}
where we denoted $$C_{\gamma,\tilde\gamma}:=\sup _{\substack{k,m\ge 0\\j>m}}\left(\frac{\gamma}{\tilde\gamma}\right)^{j-m}\sqrt{\tfrac{(j+1)(j+k+1)}{(m+1)(m+k+1)}}$$ which is finite since $\gamma/\tilde\gamma < 1.$
\bin{Since further also \begin{align}
    &\max_{m+k<n_0}\sum_{j>m}  \gamma^{-(j-m)}  e^{-\beta(E_{j+k+1}-E_{m+k+1}+E_{j+1} -E_{m+1})/2} \times\\&\quad\quad\nn\sqrt{\tfrac{(j+1)(j+1+k)}{(m+1(m+k+1)}} \\ &\le\sup_{m,k}\sum_{j\ge m+1}  \gamma^{-(j-m)}  e^{-\beta(E_{j} -E_{m+1})/2} \sqrt{\tfrac{(j+1)(j+k+1)}{(m+1)(m+k+1)}} <\infty
\end{align}
we see using \eqref{eq:tilded} that
\begin{align*} 
\nn\tilde d&:=\sup_{m,k}\sum_{j>m}  \gamma^{-(j-m)}  e^{-\beta(E_{j+k+1}-E_{m+k+1}+E_{j+1} -E_{m+1})/2} \times\\&\quad\quad\nn\times \sqrt{\tfrac{(j+1)(j+1+k)}{(m+1(m+k+1)}}<\infty.
\end{align*}}
Combining \eqref{eq:tilded} with \eqref{eq:ScaleBirthing}, we get
\begin{align*}
    \sum_{j>m} \gamma^{-j}e^{-\beta(E_{j+k}+E_j)/2} \mu^+_j\mu^+_{j+k}&\sim \sum_{j>m} \gamma^{-j}e^{-\beta(E_{j+k+1}+E_{j+1})/2} \sqrt{(j+1)(j+k+1)}\\& = \gamma^{-m} e^{-\tfrac{\beta(E_{m+k}+E_m)}{2}} e^{-\tfrac{\beta(E_{m+1}-E_m)}{2}} e^{-\tfrac{\beta(E_{m+k+1}-E_{m+k})}{2}}\\
    &\sum_{j>m}  e^{-\tfrac{\beta(E_{j+k+1}-E_{m+k+1}+E_{j+1} -E_{m+1})}{2}} \! \frac{\sqrt{(j+1)(j+k+1)}}{\gamma^{j-m}}  \\ 
    &\sim \gamma^{-m}e^{-\beta(E_{m+k}+E_m)/2}\mu^+_{m}\mu^+_{m+k}  \sum_{j>m}  \gamma^{-(j-m)} \times \\
    &\qquad e^{-\beta(E_{j+k+1}-E_{m+k+1}+E_{j+1} -E_{m+1})/2} \sqrt{\tfrac{(j+1)(j+k+1)}{(m+1)(m+k+1)}} \\&\le \tilde d\, \gamma^{-m}e^{-\beta(E_{m+k}+E_m)/2}\mu^+_{m}\mu^+_{m+k}
\end{align*}
for all $m,k\in\mathbb{N},$ which shows that condition \eqref{eq:FagnolaGapCondition2} is satisfied and therefore $\operatorname{gap}(L_{\smash{\hatfM,h(N)}})>0.$
Furthermore, from \eqref{eq:tildec} and \eqref{eq:tilded} we see that we can pick constants $\gamma,c$ and $d$ for which \eqref{eq:FagnolaGapCondition1} and \eqref{eq:FagnolaGapCondition2} are satisfied
and 
which only depend on $\beta,n_0,\delta, s$
and $\Delta_E$ but are, apart from that, independent of the specific sequence $\left(E_n\right)_{n\in\N_0}.$ Furthermore, for $\gamma\in(0,1)$ and $c,d\ge 0$ fixed
we use  \eqref{eq:ScaleBirthing} and \eqref{eq:ScaleDying}  to see 
\begin{align*}
\inf_{n > 0} 
\frac{(\mu^-_n)^2 + (\mu^+_n - \mu^+_0)^2}
{1 + c \mu^+_0 \mu^+_n \left( 1 + \frac{d + 1}{\gamma} \right)} \ge \inf_{n > 0} 
\frac{(\mu^-_n)^2}
{1 + c \mu^+_0 \mu^+_n \left( 1 + \frac{d + 1}{\gamma} \right)} &
\gtrsim \inf_{n > 0} 
\frac{n}
{1 + c \sqrt{n+1} \left( 1 + \frac{d + 1}{\gamma} \right)} \\ 
&=  \frac{1}
{1 + \sqrt{2}\, c   \left( 1 + \frac{d + 1}{\gamma} \right)}
\end{align*}
where we used
 $\mu^+_0 = \hatfM(E_1 -E_0)\le 1. $
Using \eqref{eq:LowerBoundGapBirthDeath} we therefore see that we can lower bound the spectral gap as \begin{align}
    \operatorname{gap}(L_{\hatfM,h(N)}) \ge \kappa(\beta,n_0,\delta,s,\Delta_E) >0
\end{align}
where $\kappa(\beta,n_0,\delta,s,\Delta_E)>0$ is some constant that only depends on $\beta,n_0,\delta,s$ and $\Delta_E.$ 

\bin{Lastly, note that for $n_0=0,$ we have that for $\gamma\in(0,1)$ fixed the constants $\tilde c$ and $\tilde d$ in \eqref{eq:tildec} and \eqref{eq:tilded} are decreasing for increasing $\beta.$ Hence, from \eqref{eq:LowerBoundGapBirthDeath} we see that for all $\beta'>0$ the spectral gap satisfies $\operatorname{gap}(L)\ge C_{\beta'}$ for some $C_{\beta'}>0$ and all $\beta\in[\beta',\infty).$}
\end{proof}








\section{Efficient implementation}\label{sec:implementations}




\noindent In Proposition \ref{propDirichlettoSchro}, we constructed the generator $\cL_{\smash{\widehat{f},H}}$ of a quantum dynamical semigroup associated with a function $\widehat{f}$ that satisfies minimal summability and symmetry assumptions in \Cref{eq:condAalphas}. In Section \ref{sec:spectralgaps}, we argued that the spectral gap of the corresponding evolution—which governs its mixing time—crucially depends on the choice of $\widehat{f}$. In particular, choosing $\widehat{f}\equiv 
    \hatfM(\nu)$ as defined in \eqref{eq:Function}, we found in Theorem \ref{thm:SpectralGapGeneralh(N)} that the Gibbs sampler corresponding to $H = N^2$ is gapped; however, this is no longer true for rapidly decaying functions, cf.~\Cref{prop.no-gogap}. However, the particular choice of \eqref{eq:Function} appears to hinder algorithmic implementations of the dynamics, as the lack of decay of $\hatfM$ as $\nu\to-\infty$ implies singularity of its Fourier transform $\fM$ at $0$. Another serious issue compared to the finite-dimensional setting concerns the possible unboundedness of the bare Hamiltonian $H$. 
We first address this second problem and leave the issue of the regularity of $\widehat{f}$ to Section \ref{sec:singular}.

 We consider an alternative family of generators for which the efficiency results established in Section \ref{sec:gaussexamples} remain valid, and we show that these generators can be implemented using infinite-dimensional extensions of the LCU technique. Our starting point is the following generator, which is well defined on $\mathscr{F}$ for any $\sigma_E\ge 0$, under \Cref{eq:condAalphas}:
\begin{align}
\label{eq:DefLsigma_E}
\mathcal{L}_{\sigma_E,\widehat{f},H}(\rho)\!:=\!\sum_\alpha&\sum_{\nu_1,\nu_2}e^{-\frac{(\nu_1-\nu_2)^2}{8\sigma_E^2}}\,\overline{\widehat{f}(\nu_1)}\widehat{f}(\nu_2)A^\alpha_{\nu_2}\rho (A^\alpha_{\nu_1})^\dagger+G_{\sigma_E}\rho+\rho G_{\sigma_E}^\dagger\,,
\end{align}
where 
\begin{align*}
 G_{\sigma_E}
&=-\sum_{\substack{\alpha\in\mathcal{A}\\ \nu_1,\nu_2}}\!e^{-\frac{(\nu_1-\nu_2)^2}{8\sigma_E^2}}\!\!\! \frac{1}{1+e^{\frac{\beta(\nu_2-\nu_1)}{2}}}\overline{\widehat{f}(\nu_1)}\widehat{f}(\nu_2)(A^\alpha_{\nu_1})^\dagger A^\alpha_{\nu_2}.
\end{align*}
Formally, we retrieve the generator $\mathcal{L}_{\widehat{f},H}$ of Proposition \ref{propDirichlettoSchro} by setting $\sigma_E=\infty$. It is not hard to see, by adapting the generation results of Section \ref{sec:generalframe}, that the above expression gives rise to a KMS-symmetric quantum Markov semigroup for any $\sigma_E\ge0$, whose generator in the Hilbert-Schmidt setting we denote by $L_{\sigma_E,\widehat{f},H}$. On the other hand, taking the limit $\sigma_E\to 0$, since $B^\alpha_{\nu\nu}=0$,  we obtain a Davies-type generator $\mathcal{L}_{0,\widehat{f},H}\equiv \mathcal{L}_{\operatorname{D}}$, with
\begin{align}
\label{eq:DefDavies}
\mathcal{L}_{\operatorname{D}}(\rho)=\sum_{\substack{\alpha \\ \nu\in B(H)}}|\widehat{f}(\nu)|^2\,\big(A^\alpha_\nu\rho (A^\alpha_\nu)^\dagger-\frac{1}{2}\{(A^\alpha_\nu)^\dagger A^\alpha_\nu,\rho\}\big),
\end{align}
that is, the rates $\Upsilon$ introduced in \eqref{eq:Davies} all coincide with $|\widehat{f}|^2$. In other words, the family $\{\mathcal{L}_{\sigma_E,\widehat{f},H}\}_{\sigma_E\ge 0}$ interpolates between $\mathcal{L}_{\widehat{f},H}$ and $\mathcal{L}_{\operatorname{D}}$. The next sections are organized as follows: first, in Section \ref{monotonicitygap}, we briefly justify the well-posedness of the generators $\cL_{\sigma_E,\widehat{f},H}$ and $L_{\sigma_E,\widehat{f},H}$, and argue that the gap can only increase when introducing the Gaussian envelope, which implies that the gap lower bounds of Section \ref{sec:spectralgaps} derived for $\mathcal{L}_{\smash{\widehat{f},H}}$ directly hold for $\mathcal{L}_{\smash{\sigma_E,\widehat{f},H}}$ for any $\sigma_E\ge 0$. In Section \ref{integralrepfschwartz}, we derive an integral formulation of the generator $\mathcal{L}_{\smash{\sigma_E,\widehat{f},H}}$ whenever $\widehat{f}$ can be assumed to be Schwartz. The rest of this Section is devoted to the derivation of an implementation scheme for the evolution generated by $\cL_{\smash{\sigma_E,\widehat{f},H}}$.

\subsection{Gaussian-convoluted generators}\label{monotonicitygap}



We start by considering the following densely defined, symmetric operator defined on $\mathscr{F}_{\sigma_\beta}$:


\begin{align}
L_{\sigma_E,\widehat{f},H}(\lambda x+\mu\,\sqrt{\sigma_\beta}):=-\lambda \sum_{\substack{\alpha\in\mathcal{A}\\ \nu_1,\nu_2\in B(H)}}&e^{-\frac{(\nu_1-\nu_2)^2}{8\sigma_E^2}}\,\frac{\overline{\widehat{f}(\nu_1)}\widehat{f}(\nu_2)\,e^{\beta(\nu_1+\nu_2)/4}}{2\cosh((\nu_1-\nu_2)\beta/4)}\, (\delta^{\alpha}_{\nu_1})^\dagger\delta^\alpha_{\nu_2}(x)
\end{align}
with associated form
\begin{align}\label{calEGibbs1}
\mathcal{E}_{\sigma_E,\widehat{f},H}(\lambda\,x+\mu\sqrt{\sigma_\beta}):=|\lambda|^2\sum_{\substack{\alpha\in\mathcal{A}\\ \nu_1,\nu_2\in B(H)}}&e^{-\frac{(\nu_1-\nu_2)^2}{8\sigma_E^2}}\,\frac{\overline{\widehat{f}(\nu_1)}\widehat{f}(\nu_2)\,e^{\beta(\nu_1+\nu_2)/4}}{2\cosh((\nu_1-\nu_2)\beta/4)}\, \langle \delta^\alpha_{\nu_1}(x),\delta^\alpha_{\nu_2}(x)\rangle.
\end{align}
By Fourier integration, we notice that for all $x\in\mathscr{F}$
\begin{align}
\mathcal{E}_{\sigma_E,\widehat{f},H}(x) &=\frac{\sigma_E}{\sqrt{2\pi}}\sum_{\substack{\alpha\\\nu_1,\nu_2}}\int_{-\infty}^\infty \!\!\!e^{-2\sigma_E^2s^2}  e^{-is (\nu_1-\nu_2)} \, \frac{\overline{\widehat{f}(\nu_1)}\widehat{f}(\nu_2)\,e^{\beta(\nu_1+\nu_2)/4}}{\cosh((\nu_1-\nu_2)\beta/4)}\,\langle \delta^\alpha_{\nu_1}(x),\delta^\alpha_{\nu_2}(x)\rangle\,ds\nonumber \\
&=\frac{\sigma_E}{\sqrt{2\pi}}\sum_{\substack{\alpha\\\nu_1,\nu_2}}\int_{-\infty}^\infty \!\!\!e^{-2\sigma_E^2s^2}  \frac{\overline{\widehat{f}(\nu_1)}\widehat{f}(\nu_2)\,e^{\beta(\nu_1+\nu_2)/4}}{\cosh((\nu_1-\nu_2)\beta/4)}\,\langle \delta^\alpha_{\nu_1}(x(s)),\delta^\alpha_{\nu_2}(x(s))\rangle\,ds\nonumber \\
&=\sigma_E{\sqrt{\frac{2}{\pi}}}\,\int_{-\infty}^\infty e^{-2\sigma_E^2s^2}\,\mathcal{E}_{\widehat{f},H}(x(-s))\,ds\label{integralrepLse}
\end{align}
for $x(t):=e^{iHt}x e^{-iHt}$, where we also used that $e^{it H}A^\alpha_\nu e^{-itH}=e^{it\nu}A^\alpha_\nu$. This directly shows that $L_{\smash{\sigma_E,\widehat{f},H}}$ is negative. By Friedrichs extension, $\mathcal{E}_{\smash{\sigma_E,\widehat{f},H}}$ is closable and defines a self-adjoint extension of $L_{\smash{\sigma_E,\widehat{f},H}}$. Both the close form and operator are denoted as $\mathcal{E}_{\smash{\sigma_E,\widehat{f},H}}$ and $L_{\smash{\sigma_E,\widehat{f},H}}$, by slight abuse of notations. 
\begin{proposition}\label{prop.QMSforLinsigmaEappendix}
The closed quadratic form $\mathcal{E}_{\smash{\sigma_E,\widehat{f},H}}$ is completely Dirichlet. Therefore, the strongly continuous semigroup $\{e^{tL_{\sigma_E,\widehat{f},H}}\}_{t\ge 0}$ is completely Markov with respect to the state $\sigma_\beta$.
\end{proposition}





\begin{proof}

By \Cref{GibbsDirichlet}, we know that $\mathcal{E}_{\widehat{f},H}$ is completely Dirichlet. Therefore, for any $\mathbf{x}\in\mathscr{M}_n\otimes  \mathscr{F}_{\sigma_\beta}$, we have that $\mathcal{E}^{(n)}_{\widehat{f},H}(\mathbf{x})\ge \mathcal{E}_{\widehat{f},H}^{(n)}(\mathbf{x}_+),\mathcal{E}_{\widehat{f},H}^{(n)}(\mathbf{x}_\wedge)$. The result follows since $\mathbf{x}_+(s)=(\mathbf{x}(s))_+$, $\mathbf{x}_\wedge(s)=(\mathbf{x}(s))_\wedge$, the integral representation \eqref{integralrepLse} and the form core property of $\mathscr{F}_{\sigma_\beta}$.

\end{proof}


\noindent Next, we aim at deriving an expression for the generator $\cL_{\sigma_E,\widehat{f},H}$ of the strongly continuous semigroup induced on $\mathscr{T}_1(\cH)$ by extension of the reasoning of \Cref{T1semigroupgeneration}. In place of the operator $G$ defined in \eqref{eq:GinGenerationTheore}, we consider the densely defined negative semidefinite symmetric operator $G_{\sigma_E}$ on $\mathcal{F}$ as
\begin{align}\label{GsigmaEdeffirst}
 G_{\sigma_E}
&=-\sum_{\substack{\alpha\in\mathcal{A}\\ \nu_1,\nu_2}}\!e^{-\frac{(\nu_1-\nu_2)^2}{8\sigma_E^2}}\!\!\! \frac{1}{1+e^{\frac{\beta(\nu_2-\nu_1)}{2}}}\overline{\widehat{f}(\nu_1)}\widehat{f}(\nu_2)(A^\alpha_{\nu_1})^\dagger A^\alpha_{\nu_2}.
\end{align}
and the map $\Phi_{\sigma_E,\widehat{f},H}$ defined on $\mathscr{F}$ as
\begin{align*}
\Phi_{\sigma_E,\widehat{f},H}(\ket{E_i}\bra{E_j})&:=\sum_{\substack{\alpha\in\mathcal{A}\\ \nu_1,\nu_2}}e^{-\frac{(\nu_1-\nu_2)^2}{8\sigma_E^2}}\, \overline{\widehat{f}(\nu_1)}\widehat{f}(\nu_2)\, A^\alpha_{\nu_2} \ket{E_i} (A^\alpha_{\nu_1}\ket{E_j})^\dagger.
\end{align*}
\noindent Indeed, that $G_{\sigma_E}$, resp.~$\Phi_{\sigma_E,\widehat{f},H}$, is well-defined on $\mathcal{F}$, resp.~on $\mathscr{F}$, follows directly from the estimates in the proof of \Cref{lem:definL} in the limit $\sigma_E\to \infty$ and the fact that $|e^{\smash{-\frac{(\nu_1-\nu_2)^2}{8\sigma_E^2}}}|\le 1$ uniformly over $\nu_1,\nu_2$. Moreover, by direct computation we have that, as argued in the proof of \Cref{propDirichlettoSchro}, 
\begin{align*}
\cL_{\sigma_E,\widehat{f},H}(\ket{E_i}\bra{E_j})=\sum_{\alpha\in\mathcal{A}}\, G_{\sigma_E}\ket{E_i}\bra{E_j}+\ket{E_i}(G_{\sigma_E}\ket{E_j})^\dagger +\Phi_{\sigma_E,\widehat{f},H}(\ket{E_i}\bra{E_j}).
\end{align*}
% \textcolor{red}{suggestion: remove these, extend to the right domains after together with integral reps
% By a simple adaptation of Lemma \ref{lem.LalphaGwelldefinedF}, we can also show that, under \Cref{eq:condAalphas}, $G_{\sigma_E}$ is well defined on the dense domain $D(\widetilde{H}^{\zeta+\gamma})$ and relatively $\widetilde{H}^{\zeta+\gamma}$-bounded in $\cH$, and $\Phi_{\sigma_E,\widetilde{f},H}$ is well defined on $D(\mathcal{W}_H^{\mu+\gamma,\mu+\gamma})$ and relatively $\mathcal{W}_H^{\kappa_1,\kappa_1}$-bounded in $\mathscr{T}_1(\cH)$. Finally, by adapting the proof of \Cref{prop.generationtheorem}, we arrive at the following: 
% \begin{proposition}\label{prop.generationtheoremsigmaE} Under \Cref{eq:condAalphas}, $D(\mathcal{W}_H^{\mu+\gamma,\mu+\gamma})\cap D(\mathcal{W}_H^{\zeta+\gamma,0})\cap D(\mathcal{W}_H^{0,\zeta+\gamma})\subseteq D(\mathcal{L}_{\sigma_E,\widehat{f},H})$, and for any $x\in D(\mathcal{W}_H^{\mu+\gamma,\mu+\gamma})\cap D(\mathcal{W}_H^{\zeta+\gamma,0})\cap D(\mathcal{W}_H^{0,\zeta+\gamma})$,
% \begin{align}\label{eq.LindbladformL2}
% \mathcal{L}_{\sigma_E,\widehat{f},H}(x)=\Phi_{\sigma_E,\widehat{f},H}(x)+G_{\sigma_E}\cdot x+x\cdot G_{\sigma_E}^\dagger.
% \end{align}
% \end{proposition}}



\noindent Next, we make the important observation that the spectral gap of $L_{\sigma_E,\widehat{f},H}$ increases as $\sigma_E$ decreases (see also \cite{slezak2026polynomial} for a finite-dimensional analogue of the result):

\begin{proposition}\label{lem.gapmonoton} In the notations of the previous paragraph, for any $0< \sigma_E\le \sigma'_E\le \infty$,
\begin{align*}
\operatorname{gap}(L_{\sigma_E,\widehat{f},H})\ge \operatorname{gap}(L_{\sigma'_E,\widehat{f},H}).
\end{align*}
\end{proposition}
\begin{proof}
We make use of the variational formulation of the gap:
\begin{align*}
\operatorname{gap}({L}_{\sigma_E,\widehat{f},H})=\inf_{x\in D(\mathcal{E}_{\sigma_E,\widehat{f},H})\backslash \mathbb{C}\sigma_\beta^{\frac{1}{2}}}\frac{\mathcal{E}_{\sigma_E,\widehat{f},H}(x)}{\Big\|x-\langle \sigma_\beta^{\frac{1}{2}},x\rangle\,\sigma_\beta^{\frac{1}{2}}\Big\|_{2}^2}
\end{align*}
where $\mathcal{E}_{\smash{\sigma_E,\widehat{f},H}}$ is the Dirichlet form associated with $\mathcal{L}_{\smash{\sigma_E,\widehat{f},H}}$. Therefore, denoting the probability density $g_{\sigma_E}(s):=\sigma_E\sqrt{2/\pi}e^{-2\sigma_E^2s^2}$ and using \eqref{integralrepLse} as well as the form core property of $\mathscr{F}_{\sigma_\beta}$, for all $x\in D(\mathcal{E}_{\sigma_E,\widehat{f},H})\backslash \mathbb{C}\sigma_\beta^{\smash{\frac{1}{2}}}$,  there exists a sequence of elements $x_n\in\mathscr{F}_{\sigma_\beta}\backslash \mathbb{C}\sigma^{\smash{\frac{1}{2}}}$ such that $\|x-x_n\|_{D(\mathcal{E}_{\sigma_E,\widehat{f},H})}\to 0$ as $n\to\infty$ and
\begin{align}
 \tfrac{\mathcal{E}_{\sigma_E,\widehat{f},H}(x)}{\Big\|x-\langle \sigma_\beta^{\frac{1}{2}},x\rangle\,\sigma_\beta^{\frac{1}{2}}\Big\|_{2}^2}=\lim_{n\to\infty}\tfrac{\mathcal{E}_{\sigma_E,\widehat{f},H}(x_n)}{\Big\|x_n-\langle \sigma_\beta^{\frac{1}{2}},x_n\rangle\,\sigma_\beta^{\frac{1}{2}}\Big\|_{2}^2}=\lim_{n\to\infty} \int_{-\infty}^{\infty}  \tfrac{g_{\sigma_E}(s)\mathcal{E}_{\widehat{f},H}(x_n(-s))}{\Big\|x_n(-s)-\langle \sigma_\beta^{\frac{1}{2}},x_n(-s)\rangle\,\sigma_\beta^{\frac{1}{2}}\Big\|_{2}^2}\,ds \nonumber   
\end{align}
where we also used the invariance of the denominator under the unitaries $e^{isH}$. From this, we directly get that $\operatorname{gap}({L}_{\sigma_E,\widehat{f},H})\ge \operatorname{gap}({L}_{\widehat{f},H}) $. Now, for any such $x$, denoting $\mathcal{E}_n(s):=\mathcal{E}_{\smash{\widehat{f},H}}(x_n(s))$, the above equations show that $\mathcal{E}_{\sigma_E,\widehat{f},H}(x_n)= \mathcal{E}_n\ast g_{\sigma_E}(0)$. More generally, since $g_{\sigma_E}=g_{\sigma_E'}\ast  g_{\sigma_E^*}$ with $(\sigma_E^*)^2:=\sigma_E^2(\sigma_E')^2/((\sigma_E')^2-\sigma_E^2)$, and by associativity of the convolution, we get
\begin{align*}
\mathcal{E}_{\sigma_E,\widehat{f},H}(x_n)&=\mathcal{E}_n\ast g_{\sigma_E}(0)=\mathcal{E}_n\ast g_{\sigma_E'}\ast g_{\sigma_E^*}(0)=(s\mapsto \mathcal{E}_{\sigma_E',\widehat{f},H}(x_n(s)))\ast g_{\sigma_E^*}(0).
\end{align*}
Thus,
\begin{align*}
 \tfrac{\mathcal{E}_{\sigma_E,\widehat{f},H}(x_n)}{\Big\|x_n-\langle \sigma_\beta^{\frac{1}{2}},x_n\rangle\,\sigma_\beta^{\frac{1}{2}}\Big\|_{2}^2}=\int_{-\infty}^\infty \,\tfrac{g_{\sigma_E^*}(s)\mathcal{E}_{\sigma_E',\widehat{f},H}(x_n(-s))}{\Big\|x_n(-s)-\langle \sigma_\beta^{\frac{1}{2}},x_n(-s)\rangle\,\sigma_\beta^{\frac{1}{2}}\Big\|_{2}^2}\,ds
\end{align*}
from which we directly read off that $\operatorname{gap}({L}_{\sigma_E,\widehat{f},H})\ge \operatorname{gap}({L}_{\sigma_E',\widehat{f},H})$.
\end{proof}


\subsection{Integral representation}\label{integralrepfschwartz}

Since we proved in the last section that the spectral gap can only improve with smaller $\sigma_E$, combining this fact with the results of Section \ref{sec:spectralgaps} for $\sigma_E=\infty$, we directly obtain a set of examples for which the Gibbs samplers at any $\sigma_E\in \mathbb{R}_+$ are gapped. Next, we provide an integral formulation of the generator $\mathcal{L}_{\sigma_E,\widehat{f},H}$ for any $\sigma_E\in(0,\infty)$ and sufficiently nice functions $\widehat{f}$. This form will play a crucial role when considering an implementation of the associated evolution on a discrete variable quantum platform. In what follows, we denote by $\mathcal{S}(\mathbb{R})$ the space of Schwartz functions 

\[
\mathcal{S}(\mathbb{R})
\!=\!\left\{
f\in C^{\infty}(\mathbb{R})\!:\!
\forall m\in\mathbb{N}_0^2,
\sup_{x\in\mathbb{R}}
\big|x^{m_1}f^{(m_2)}(x)\,\big|\!<\!\infty\!
\right\}.
\]

\noindent
Equivalently, \(f\in\mathcal{S}(\mathbb{R})\) if and only if for every pair of non-negative integers \(m,n\),
\[
\lim_{|x|\to\infty} |x|^m |f^{(n)}(x)| = 0.
\]



% In what follows, we denote
% \begin{align*}
% F_3(E):=\sum_{E',E''\in\operatorname{Sp}(H)} |\widehat{f}(E'-E'')\widehat{f}(E'-E)|
% \end{align*}
% \begin{framed}
% \begin{conditionC}\label{eq:CondC}
% Given $f\in
% \mathcal{S}(\mathbb{R})$, there exists $\kappa_1>\mu$ such that 
% \begin{align}
% \label{eq:condC}
% \sum_{E\in\operatorname{Sp}(H)}F_3(E)\,(E+h_0+1)^{-(\kappa_1-\mu)}<\infty\,.
% \end{align}
% \end{conditionC}
% \end{framed}

We start by providing integral representations for each constituent of the generator $\mathcal{L}_{\sigma_E,\widehat{f},H}$.



\begin{proposition}\label{integralrep}
Assume Condition \ref{eq:condAalphas} is satisfied for some $f\in\mathcal{S}(\mathbb{R})$ with Fourier transform $\widehat{f}(\nu):=\int_{-\infty}^\infty f(s)e^{is\nu}ds$. Then for any $\alpha\in\mathcal{A}$, $(L^\alpha,\mathcal{F})$ extends to an operator in $D(\widetilde{H}^{\gamma}),$ which is relatively $\widetilde H^\gamma$-bounded, and for any $\ket{\psi}\in D(\widetilde{H}^{\gamma})$, 
\begin{align*}
L^\alpha\ket{\psi}=\int f(s)\,e^{isH}A^\alpha e^{-isH}\ket{\psi}\,ds=\sum_{E,E'\in\operatorname{Sp}(H)}\widehat{f}(E'-E)P_{E'}A^\alpha P_{E}\ket{\psi}.
\end{align*}
where the latter expression is to be understood as a limit of the finite sums $\sum_{\smash{E\le M,E'\le M'}}$ in $\mathcal{H}$. Moreover, $\widetilde{H}^{\gamma }L^\alpha \widetilde{H}^{-\mu}$ extends to a bounded operator with integral representation
\begin{align*}
\widetilde{H}^{\gamma }L^\alpha \widetilde{H}^{-\mu}:=\int f(s) e^{isH}\widetilde{H}^{\gamma } A^\alpha \widetilde{H}^{-\mu}e^{-isH}\,ds.
\end{align*}
Next, $(G_{\sigma_E},\mathcal{F})$ extends to an operator on $D(\widetilde{H}^{\mu}),$ which is relatively $\widetilde H^\mu$-bounded, and for any {$\sigma_E\in(0,\infty)$,} $\ket{\psi}\in D(\widetilde{H}^{\mu})$, $\ket{\psi}\in D((L^\alpha)^\dagger L^\alpha )\cap D(G_{\sigma_E})$ and
\begin{align*}
G_{\sigma_E}\ket{\psi}&=-\sum_{\alpha\in\mathcal{A}}\int_{-\infty}^\infty g(t)\,e^{itH}(L^\alpha)^\dagger L^\alpha e^{-itH}\ket{\psi}\,dt\\
&=-\!\!\!\!\!\!\!\!\!\!\!\sum_{\substack{\alpha\in\mathcal{A}\\E,E',E''\in\operatorname{Sp}(H)}}\!\!\!\!\!\!\!\!\!\!\widehat g(E'-E'')\overline{\widehat{f}(E-E')}\widehat{f}(E-E'')P_{E'}(A^\alpha)^\dagger P_{E}A^\alpha P_{E''}\ket{\psi},
\end{align*}
with $g(t)=\frac{1}{2\pi}\int_{-\infty}^\infty \frac{e^{-\nu^2/8\sigma_E^2}}{1+e^{\beta\nu/2}} e^{-i\nu t}d\nu$. Similarly, $\Phi_{\sigma_E,\widehat{f},H}$ extends to an operator on $D(\mathcal{W}_H^{\gamma,\gamma})$, which is relatively $\mathcal{W}^{\gamma,\gamma}_H$ bounded, and for any $x\in D(\mathcal{W}_H^{\gamma,\gamma})$ 
\begin{align}
\Phi_{\sigma_E,\widehat{f},H}(x)=\sigma_E\sqrt{\frac{2}{\pi}}\sum_{\alpha}\int_{\mathbb{R}}e^{-2\sigma_E^2 s^2}\, \, X^{\alpha}_s \cdot x\cdot  (X_s^\alpha)^\dagger \,ds
\end{align}
with $X^{\alpha}_s:=e^{isH}L^\alpha e^{-isH}$ and where we recall that $X_s^\alpha\cdot x\cdot (X_s^\alpha)^\dagger \equiv X_s^\alpha\widetilde{H}^{-\gamma}(X^\alpha_s\widetilde{H}^{-\gamma}\mathcal{W}_H^{\gamma,\gamma}(x))^\dagger)^\dagger$.
\end{proposition}
\begin{proof}
For $\ket{\psi}\in \mathcal{F}$ the operator $L^\alpha$ defined in \eqref{eq:LalphainGenerationTheory} satisfies 
\begin{align*}
L^\alpha\ket{\psi}&=\sum_{E,E'\in\spec(H) }\widehat{f}(E'-E)P_{E'}A^\alpha P_{E}\ket{\psi}=\sum_{E,E'\in \operatorname{Sp}(H)}\int_{-\infty}^\infty f(s)e^{is(E'-E)}ds P_{E'}A^\alpha P_E\ket{\psi}\\
&= \lim_{M\to \infty} \lim_{M'\to \infty}\int_{-\infty}^\infty f(s) \sum_{\substack{E,E'\in \operatorname{Sp}(H)\\E\le M,\,E'\le M'}}e^{is(E'-E)}ds P_{E'}A^\alpha P_E\ket{\psi} \\
&= \lim_{M\to \infty} \lim_{M'\to \infty} \int_{-\infty}^\infty f(s)\, P_{H\le M'}e^{isH}A^\alpha e^{-isH} P_{H\le M}\ket{\psi}\,ds  
\end{align*}
where we denoted for $M\in\N_0$ the projection $ P_{H\le M} =  \sum_{\substack{E\in \operatorname{Sp}(H)\\E\le M}}P_E.$ Note that
\begin{align*}
   \| P_{H\le M'}e^{isH}A^\alpha e^{-isH} P_{H\le M}\ket{\psi}\| \le \|A^\alpha\widetilde H^{-\gamma}\| \|\widetilde H^\gamma\ket{\psi}\|
\end{align*}
for all $M,M'\in\N$ and therefore, using $f\in L^1(\R),$ we can use dominated convergence to exchange the limit $M,M'\to\infty$ with the integral and obtain from the above\begin{align*}
L^\alpha\ket{\psi}=\int_{-\infty}^\infty\,f(s)\sum_{E,E'\in\operatorname{Sp}(H)}e^{is(E'-E)}P_{E'}A^\alpha P_E\ket{\psi}\,ds=\int_{-\infty}^\infty f(s)e^{isH}A^\alpha e^{-isH}\ket{\psi}\,ds\,.
\end{align*}
Clearly, the above swapping of sums and integral extend to $\ket{\psi}\in D(\widetilde{H}^{\gamma})$, . We argue similarly for $G_{\sigma_E}$ as defined in \eqref{GsigmaEdeffirst}: for any $\ket{\psi}\in \mathcal{F}$,
\begin{align*}
&G_{\sigma_E}\ket{\psi}\\
&\, =-\sum_{\alpha\in\mathcal{A}}\lim_{M,M'\to\infty}\int_{-\infty}^\infty\sum_{\substack{E,E',E''\in\operatorname{Sp}(H)\\E\le M,E'\le M'}}\!\!\!g(t)\overline{\widehat{f}(E-E')}\widehat{f}(E-E_i)e^{itH}P_{E'}(A^\alpha)^\dagger P_{E}A^\alpha e^{-itH}P_{E''}\ket{\psi}dt\\
&\, =\!-\!\sum_{\alpha\in\mathcal{A}}\lim_{M,M'\to\infty}\int_{-\infty}^\infty \!\!P_{H\le M'}\!\!\!\!\!\!\!\!\!\!\!\!\sum_{E,E',E''\in\operatorname{Sp}(H)}\!\!\!\!\!\!\!\!\!\!g(t)\overline{\widehat{f}(E-E')}\widehat{f}(E-E'')e^{itH}P_{E'}(A^\alpha)^\dagger P_{E}P_{H\le M}A^\alpha e^{-itH}P_{E''}\ket{\psi}dt\\
&\,=\!-\!\sum_{\alpha\in\mathcal{A}}\lim_{M,M'\to\infty}\int_{-\infty}^\infty \!\,g(t)e^{itH}P_{H\le M'}(L^\alpha)^\dagger P_{H\le M}L^\alpha e^{-itH}\ket{\psi}dt.
\end{align*}
Moreover, the term above is uniformly integrable, since $g\in L^1(\mathbb{R})$ and 
\begin{align*}
\|P_{H\le M'}(L^\alpha)^\dagger P_{H\le M} L^\alpha e^{-itH}\ket{\psi}\|\le  \|\widetilde{H}^{-\gamma}L^\alpha\|\,\|\widetilde{H}^\gamma L^\alpha \widetilde{H}^{-\mu} \,\|\widetilde{H}^\mu\ket{\psi}\|<\infty
\end{align*}
% by the summability of
% \begin{align*}
% &\sum_{E,E',E_i\in\operatorname{Sp}(H)}|\widehat{f}(E-E')\widehat{f}(E-E_i)|\,\|P_{E'}(A^\alpha)^{\dagger}P_E A^\alpha e^{-itH}P_{E_i}\ket{\psi}\|\\
% &\qquad\le\! \!\!\sum_{E,E',E_i\in\operatorname{Sp}(H)}\!|\widehat{f}(E-E')\widehat{f}(E-E_i)|\,\|P_{E'}(A^\alpha)^{\dagger}P_E A^\alpha \widetilde{H}^{-\mu}\| \,\!\|\widetilde{H}^{\kappa_1} P_{E_i}\ket{\psi}\|(E_i+h_0+1)^{-(\kappa_1-\mu)}\\
% &\qquad \le \!\!\! \sum_{E,E',E_i\in\operatorname{Sp}(H)}\!\!|\widehat{f}(E-E')\widehat{f}(E-E_i)|\,\|\widetilde{H}^{-\gamma} A^\alpha\|\,\|\widetilde H^\gamma A^\alpha \widetilde{H}^{-\mu}\| \,\|\widetilde{H}^{\kappa_1} P_{E_i}\ket{\psi}\|(E_i+h_0+1)^{-(\kappa_1-\mu)}\\
% &\qquad <\infty
% \end{align*}
% where the final boundedness arises from \Cref{eq:CondC}.
Therefore, using the dominated convergence theorem to exchange the limits $M,M'\to \infty$ with the integral gives
 \begin{align*}
G_{\sigma_E}\ket{\psi}&=-\sum_{\alpha}\int_{-\infty}^{\infty}g(t)\,\sum_{E,E',E''\in\operatorname{Sp}(H)}\overline{\widehat{f}(E-E')}\widehat{f}(E-E_i)e^{itH}P_{E'}(A^\alpha)^\dagger P_E A^\alpha e^{-itH}P_{E_i}\ket{\psi}\\
&=-\sum_{\alpha}\int_{-\infty}^{\infty}g(t)\,e^{itH}(L^\alpha)^\dagger  L^\alpha e^{-itH}\ket{\psi}
\end{align*}
where in the last line we used that $L^\alpha\ket{\psi}\in D(\widetilde{H}^{\gamma})$. 
Once again, the above swapping of sums and integrals extend to $\ket{\psi}\in D(\widetilde{H}^{\mu})$.

For the CP-term $\Phi_{\sigma_E,\widehat f, H}$ we apply a similar dominated convergence argument as for the $G_{\sigma_E}$ operator to exchange the Bohr frequency series with the integral: given $x\in \mathscr{F}$
\begin{align*}
&\Phi_{\sigma_E,\widehat{f},H}(x)\\
&\quad =\sum_{\substack{\alpha\in\mathcal{A}\\ E_1,E_1',E_2,E_2'\in\operatorname{Sp}(H)}}e^{-\frac{(E_1-E_1'-E_2+E_2')^2}{8\sigma_E^2}}\, \overline{\widehat{f}(E_1-E_1')}\widehat{f}(E_2-E_2')\, P_{E_2} A^\alpha P_{E_2'} x P_{E_1'}(A^\alpha)^\dagger P_{E_1}\\
&\quad =\sigma_E\sqrt{\frac{2}{\pi}}\!\!\!\!\!\!\!\!\!\!\!\!\!\!\!\sum_{\substack{\alpha\in\mathcal{A}\\ E_1,E_1',E_2,E_2'\in\operatorname{Sp}(H)}}\!\!\!\!\!\!\!\!\!\!\!\!\!\!\!\int_{-\infty}^\infty e^{-2\sigma_E^2 s^2}  e^{is (-E_1+E_1'+E_2-E_2')} \overline{\widehat{f}(E_1-E_1')}\widehat{f}(E_2-E_2')\, P_{E_2} A^\alpha P_{E_2'} x P_{E_1'}(A^\alpha)^\dagger P_{E_1}ds\\
&\quad =\sigma_E\sqrt{\frac{2}{\pi}}\!\!\!\!\!\!\!\!\!\!\!\!\!\!\!\sum_{\substack{\alpha\in\mathcal{A}\\ E_1,E_1',E_2,E_2'\in\operatorname{Sp}(H)}}\!\!\!\!\!\!\!\!\!\!\!\!\!\!\!\!\int_{-\infty}^\infty \!\!\!\!\!e^{-2\sigma_E^2 s^2}   \overline{\widehat{f}(E_1-E_1')}\widehat{f}(E_2-E_2')\, P_{E_2} e^{isH} A^\alpha e^{-isH}P_{E_2'} x e^{isH}P_{E_1'}(A^\alpha)^\dagger e^{-isH}P_{E_1}ds.
\end{align*}
Arguing by truncation as above, it suffices to show that 
\begin{align*}
&P_{H\le M}\sum_{E_1,E_2,E_1',E_2'\in\operatorname{Sp}(H)} \overline{\widehat{f}(E_1-E_1')}\widehat{f}(E_2-E_2')\, P_{E_2} e^{isH} A^\alpha e^{-isH}P_{E_2'} x e^{isH}P_{E_1'}(A^\alpha)^\dagger e^{-isH}P_{E_1} P_{H\le M'}\\
&\qquad =P_{H\le M}\,\sum_{E_2',E_1'\in\operatorname{Sp}(H)}\,e^{isH}L^\alpha  e^{-isH}P_{E_2'}x P_{E_1'}e^{isH}(L^\alpha)^\dagger e^{-isH}P_{H\le M'}\\
&\qquad =P_{H\le M}\,\,e^{isH}L^\alpha  e^{-isH}x e^{isH}(L^\alpha)^\dagger e^{-isH}P_{H\le M'}
\end{align*}
with norm
\begin{align*}
\|P_{H\le M}\,\,e^{isH}L^\alpha  e^{-isH}x e^{isH}(L^\alpha) e^{-isH}P_{H\le M'}\|_1\le \|x\|_{\mathcal{W}_H^{\gamma,\gamma}}\|L^\alpha \widetilde{H}^{-\gamma}\|^2<\infty.
\end{align*}
Thus, once again, by the dominated convergence theorem, we conclude that 
\begin{align*}
\Phi_{\sigma_E,\widehat{f},H}(x)&=\sigma_E\sqrt{\frac{2}{\pi}}\sum_{\alpha\in\mathcal{A}}\int_{-\infty}^\infty e^{-2\sigma_E^2 s^2}   \,e^{isH}L^\alpha  e^{-isH}\cdot x\cdot  e^{isH}(L^\alpha)^\dagger e^{-isH}ds\\
&=\sigma_E\sqrt{\frac{2}{\pi}}\sum_{\alpha}\int_{\mathbb{R}}e^{-2\sigma_E^2 s^2}\, \, X_s^{\alpha} \cdot x\cdot  (X_s^\alpha)^\dagger \,ds.
\end{align*}
Again, the same swap holds and extends to $x\in \mathcal{W}_H^{\gamma,\gamma}$. 


\end{proof}

\begin{corollary}
 Assume that \Cref{eq:condAalphas} is satisfied for some $f\in\mathcal{S}(\mathbb{R})$ with Fourier transform $\widehat{f}(\nu):=\int_{-\infty}^\infty f(s)e^{is\nu}ds$. Then, $D(\mathcal{W}_H^{\gamma,\gamma})\cap D(\mathcal{W}_H^{0,\mu})\cap D(\mathcal{W}_H^{\mu,0})\subseteq D(\mathcal{L}_{\sigma_E,\widehat{f},H})$, and for any $x\in D(\mathcal{W}_H^{\gamma,\gamma})\cap D(\mathcal{W}_H^{0,\mu})\cap D(\mathcal{W}_H^{\mu,0})$,
 \begin{align*}
\mathcal{L}_{\sigma_E,\widehat{f},H}(x):=\Phi_{\sigma_E,\widehat{f},H}(x)+G_{\sigma_E}\cdot x+x\cdot G_{\sigma_E}^\dagger.
 \end{align*}
\end{corollary}
\begin{proof}
This follows the same lines as in \Cref{prop.generationtheorem}, by invoking the closedness of $\mathcal{L}_{\sigma_E,\widehat{f},H}$ and density of $\mathscr{F}$ in $D(\mathcal{W}_H^{\gamma,\gamma})\cap D(\mathcal{W}_H^{0,\mu})\cap D(\mathcal{W}_H^{\mu,0})$:
given  $x\in D(\mathcal{W}_H^{\gamma,\gamma})\cap D(\mathcal{W}_H^{0,\mu})\cap D(\mathcal{W}_H^{\mu,0})$, $x_n:=P_{H\le n}x P_{H\le n}\in \mathcal{W}_{H}^{\mu,0}$ with $x_n\to x$ in $\mathscr{T}_1(\cH)$ as $n\to\infty$. 
Then, by \Cref{integralrep},
\begin{align*}
G_{\sigma_E}\cdot x_n=-\sum_{\alpha\in\mathcal{A}} \int_{-\infty}^\infty g(t) e^{itH}(L^\alpha)^\dagger L^\alpha e^{-itH}x_n dt\to -\sum_{\alpha\in \mathcal{A}}\int_{-\infty}^\infty g(t) e^{itH}(L^\alpha)^\dagger L^\alpha e^{-itH}x dt\equiv G_{\sigma_E} \cdot x
\end{align*}
by dominated convergence theorem. Similarly, we show  that $x_n\cdot  G_{\sigma_E}^\dagger \to  x\cdot G_{\sigma_E}^\dagger$ and 
\begin{align*}
\Phi_{\sigma_E,\widehat{f},H}(x_n)=\sigma_E\sqrt{\frac{2}{\pi}}\sum_{\alpha\in\mathcal{A}}\int_{\mathbb{R}}e^{-2\sigma_E^2s^2} X_s^\alpha \cdot x_n\cdot  (X_s^\alpha)^\dagger \,ds\to \sigma_E\sqrt{\frac{2}{\pi}}\sum_{\alpha\in\mathcal{A}}\int_{\mathbb{R}}e^{-2\sigma_E^2s^2} X_s^\alpha \cdot x\cdot  (X_s^\alpha)^\dagger \,ds.
\end{align*}
The result follows.
\end{proof}











% \red{Remove Proposition below}
% \begin{proposition}\label{integralrep}
% Let $f\in\mathcal{S}(\mathbb{R})$ with Fourier transform $\widehat{f}(\nu):=\int_{-\infty}^\infty f(s)e^{is\nu}ds$, and assume Condition \ref{eq:condAalphas} is satisfied. Then for any $\alpha\in\mathcal{A}$ and $\ket{\psi}\in D(\widetilde{H}^{\mu+\gamma})$, $\ket{\psi}\in D(A^\alpha)\cap  D(L^\alpha)$ and 
% \begin{align*}
% L^\alpha\ket{\psi}=\int f(s)\,e^{isH}A^\alpha e^{-isH}\ket{\psi}\,ds.
% \end{align*}
% Moreover, for any \robi{$\sigma_E\in(0,\infty)$,} $\ket{\psi}\in D(\widetilde{H}^{\zeta+\gamma})$, $\ket{\psi}\in D((L^\alpha)^\dagger L^\alpha )\cap D(G_{\sigma_E})$ and
% \begin{align*}
% G_{\sigma_E}\ket{\psi}=-\sum_{\alpha\in\mathcal{A}}\int_{-\infty}^\infty g(t)\,e^{itH}(L^\alpha)^\dagger L^\alpha e^{\textcolor{red}{-}itH}\ket{\psi}\,dt\,.
% \end{align*}
% with $g(t)=\frac{1}{2\pi}\int_{-\infty}^\infty \frac{e^{-\nu^2/8\sigma_E^2}}{1+e^{\beta\nu/2}} e^{-i\nu t}d\nu$. Finally, for any $x\in D(\mathcal{W}_H^{\mu+\gamma,\mu+\gamma})$, $x\in D(\Phi_{\sigma_E,\widehat{f},H})$ and 
% \begin{align}
% \Phi_{\sigma_E,\widehat{f},H}(x)=\sigma_E\sqrt{\frac{2}{\pi}}\sum_{\alpha}\int_{\mathbb{R}}e^{-2\sigma_E^2 s^2}\, \, X^{\alpha}_s \cdot x\cdot  (X_s^\alpha)^\dagger \,ds
% \end{align}
% with $X^{\alpha}_s:=e^{isH}L^\alpha e^{-isH}$ and where we recall that $X_s^\alpha\cdot x\cdot (X_s^\alpha)^\dagger \equiv X_s^\alpha\widetilde{H}^{-\kappa_1}(X^\alpha_s\widetilde{H}^{-\kappa_1}\mathcal{W}_H^{\kappa_1,\kappa_1}(x))^\dagger)^\dagger$.
% \end{proposition}



% \begin{proof}

% Fist, we know by Lemma \ref{lem.LalphaGwelldefinedF}, for $\ket{\psi}\in D(\widetilde{H}^{\mu+\gamma})$, $\ket{\psi}\in D(L^\alpha)\cap D(A^\alpha)$ and 
% \begin{align}
% L^\alpha\ket{\psi}&:=\sum_{E,E'\in\spec(H) }\widehat{f}(E'-E)P_{E'}A^\alpha P_{E}\ket{\psi}=\sum_{E,E'\in \operatorname{Sp}(H)}\int_{-\infty}^\infty f(s)e^{is(E'-E)}ds P_{E'}A^\alpha P_E\ket{\psi}
% \end{align}
% From \Cref{eq:condAalphas}, we further get that
% \begin{align*}
% &\int_{-\infty}^\infty\,|f(s)|\,ds \,\sum_{E,E'\in\operatorname{Sp}(H)}\|P_{E'}A^\alpha P_E\ket{\psi}\|\\
% &\qquad \le \|\widetilde H^{\gamma}A^\alpha\widetilde H^{-\mu}\|\,\int_{-\infty}^{\infty}|f(s)|\,ds\,\left(\sum_{E\in\spec(H)}\,(E+h_0+1)^{-\gamma}\right)^2 \|\widetilde{H}^{\mu+\gamma}\ket{\psi}\|\,<\infty
% \end{align*}
% Thus, by Fubini theorem, we can write
% \begin{align*}
% L^\alpha\ket{\psi}=\int_{-\infty}^\infty\,f(s)\sum_{E,E'\in\operatorname{Sp}(H)}e^{is(E'-E)}P_{E'}A^\alpha P_E\ket{\psi}\,ds=\int_{-\infty}^\infty f(s)e^{isH}A^\alpha e^{-isH}\ket{\psi}\,ds\,.
% \end{align*}
% We argue similarly for $G_{\sigma_E}$: for any $\ket{\psi}\in D(\widetilde{H}^{\zeta+\gamma})$,
% \begin{align*}
% G_{\sigma_E}\ket{\psi}&=-\sum_{\substack{\alpha\in\mathcal{A}\\\nu_1,\nu_2}}e^{-\frac{(\nu_1-\nu_2)^2}{8\sigma_E^2}}\,\frac{1}{1+e^{\frac{\beta(\textcolor{red}{\nu_2-\nu_1})}{2}}}\overline{\widehat{f}(\nu_1)}\widehat{f}(\nu_2)\,(A^\alpha_{\nu_1})^\dagger\,A^\alpha_{\nu_2}\ket{\psi}\\
% &=-\sum_{\substack{\alpha\in\mathcal{A}\\\nu_1,\nu_2}}\int_{-\infty}^{\infty}g(t) e^{it(\textcolor{red}{\nu_2-\nu_1})}\overline{\widehat{f}(\nu_1)}\widehat{f}(\nu_2)\,(A^\alpha_{\nu_1})^\dagger\,A^\alpha_{\nu_2}\ket{\psi}\,dt.
% \end{align*}
% Therefore, by \Cref{eq:condAalphas} we have for all $\ket{\psi}\in D(\widetilde{H}^{\zeta+\gamma})$,
% \begin{align*}
% &\|G_{\sigma_E}\ket{\psi}\|\\
% &\quad\le \int_{-\infty}^{\infty}
% |g(t)|\,dt\sum_{\alpha\in\mathcal{A}}\sum_{E,E',E''}\,|\widehat{f}(E'-E)|\, \|P_{E''}(A^\alpha)^{\dagger} P_{E'}A^\alpha P_E\ket{\psi}\|\\
%  &\quad \lesssim \|\widetilde H^{-\mu}A^{\alpha} \widetilde H^{\gamma}\|\|\widetilde H^{\mu+\gamma}A^\alpha \widetilde H^{-\zeta}\| \|\widetilde{H}^{\zeta+\gamma}\ket{\psi}\|\!\int_{-\infty}^{\infty}\!
% |g(t)|dt\left(\sum_{E\in\spec(H)}\!(E+h_0+1)^{-\gamma}\right)^3\\
% &\quad <\infty.
% \end{align*}
% Therefore, by Fubini's theorem, we have that
% \begin{align*}
% G_{\sigma_E}\ket{\psi}&=-\sum_{\alpha}\int_{-\infty}^{\infty}g(t)\,\sum_{\nu_1,\nu_2}e^{it(\textcolor{red}{\nu_2-\nu_1})}\overline{\widehat{f}(\nu_1)}\widehat{f}(\nu_2)(A^\alpha_{\nu_1})^\dagger A^\alpha_{\nu_2}\ket{\psi}\,dt\\
% &=-\sum_{\alpha}\int_{-\infty}^{\infty}g(t)\,\sum_{\nu_1,\nu_2}\overline{\widehat{f}(\nu_1)}\widehat{f}(\nu_2)e^{iHt}(A^\alpha_{\nu_1})^\dagger A^\alpha_{\nu_2}e^{\textcolor{red}{-}iHt}\ket{\psi}\,dt\\
% &=-\sum_{\alpha}\int_{-\infty}^\infty g(t)e^{iHt}(L^\alpha)^\dagger L^\alpha e^{\textcolor{red}{-}iHt}\ket{\psi}dt
% \end{align*}
% Finally, for any $x\in D(\mathcal{W}_H^{\mu+\gamma,\mu+\gamma})$,
% \begin{align*}
% \Phi_{\sigma_E,\widehat{f},H}(x)&=\sum_{\substack{\alpha\in\mathcal{A}\\ \nu_1,\nu_2}}e^{-\frac{(\nu_1-\nu_2)^2}{8\sigma_E^2}}\, \overline{\widehat{f}(\nu_1)}\widehat{f}(\nu_2)\, A^\alpha_{\nu_2} x (A^\alpha_{\nu_1})^\dagger\\
% &=\sigma_E\sqrt{\frac{2}{\pi}}\sum_{\substack{\alpha\in\mathcal{A}\\ \nu_1,\nu_2}}\int_{\mathbb{R}}e^{-2\sigma_E^2 s^2}\, e^{-i(\nu_1-\nu_2)s}\overline{\widehat{f}(\nu_1)}\widehat{f}(\nu_2)A^\alpha_{\nu_2} x (A^\alpha_{\nu_1})^\dagger ds\\
% &=\sigma_E\sqrt{\frac{2}{\pi}}\sum_{\substack{\alpha\in\mathcal{A}\\ \nu_1,\nu_2}}\int_{\mathbb{R}}e^{-2\sigma_E^2 s^2}\, \overline{\widehat{f}(\nu_1)}\widehat{f}(\nu_2) e^{isH}A^\alpha_{\nu_2}e^{-isH} x e^{isH}(A^\alpha_{\nu_1})^\dagger e^{-isH}\,ds.
% \end{align*}
% Since
% \begin{align*}
% &\sum_{\nu_1,\nu_2}\int_{\mathbb{R}}e^{-2\sigma_E^2 s^2}|\widehat{f}(\nu_1)\widehat{f}(\nu_2)|\|A^\alpha_{\nu_2}e^{-isH}x e^{is H}(A^\alpha_{\nu_1})^\dagger \|_1\,ds\\
% &\qquad \le \sum_{\nu_1,\nu_2}\int_{\mathbb{R}}e^{-2\sigma_E^2 s^2}ds|\widehat{f}(\nu_1)\widehat{f}(\nu_2)|\|A^\alpha_{\nu_2}\widetilde{H}^{-\mu+\gamma}\|\,\|\widetilde{H}^{-\mu+\gamma} (A^\alpha_{\nu_1})^\dagger \|\,\|x\|_{\mathcal{W}_H^{\mu+\gamma,\mu+\gamma}}\\
% &\qquad \lesssim \sum_{\nu_1,\nu_2}|\widehat{f}(\nu_1)\widehat{f}(\nu_2)|\|A^\alpha_{\nu_2}\widetilde{H}^{-\mu+\gamma}\|\,\|\widetilde{H}^{-\mu+\gamma} (A^\alpha_{\nu_1})^\dagger \|\,\|x\|_{\mathcal{W}_H^{\mu+\gamma,\mu+\gamma}}\\
% &\qquad \le  \sum_{\substack{E_1,E_1'\\E_2,E_2'}}|\widehat{f}(E_1'-E_1)\widehat{f}(E_2'-E_2)|\|P_{E_2'}A^\alpha P_{E_2}\widetilde{H}^{-\mu+\gamma}\|\,\|\widetilde{H}^{-\mu+\gamma} P_{E_1}(A^\alpha)^\dagger P_{E_1'}\|\,\|x\|_{\mathcal{W}_H^{\mu+\gamma,\mu+\gamma}}\\
% &\qquad \lesssim \|\widetilde H^{\gamma}A^\alpha\widetilde H^{-\mu}\| \left(\sum_{E\in\spec(H)}(E+h_0+1)^{-\gamma}\right)^4 \|x\|_{\mathcal{W}_H^{\mu+\gamma,\mu+\gamma}}\\
% &\qquad <\infty,
% \end{align*}
% we have by one last use of Fubini theorem that
% \begin{align*}
% \Phi_{\sigma_E,\widehat{f},H}(x)&=\sigma_E\sqrt{\frac{2}{\pi}}\sum_{\alpha\in\mathcal{A}}\int_{\mathbb{R}}\sum_{\nu_1,\nu_2}e^{-2\sigma_E^2 s^2}\, \overline{\widehat{f}(\nu_1)}\widehat{f}(\nu_2) e^{isH}A^\alpha_{\nu_2}e^{-isH} x e^{isH}(A^\alpha_{\nu_1})^\dagger e^{-isH}\,ds\\
% &=\sigma_E\sqrt{\frac{2}{\pi}}\sum_{\alpha}\int_{\mathbb{R}}e^{-2\sigma_E^2 s^2}\, \, X_s^{\alpha} x (X_s^\alpha)^\dagger \,ds.
% \end{align*}

% \end{proof}
\subsection{Finite-dimensional truncations}
\label{sec:FiniteTruncBareJumps}
In this section we consider a finite-dimensional\footnote{Technically, the generator  is not finite-dimensional but only finite rank as it still acts on infinite-dimensional space. But we drop this distinction here for simplicity.} truncation of the unbounded generator $\mathcal{L}_{\sigma_E,\widehat{f},H}$ which for truncation parameter $M\in\N$ is denoted by $\mathcal{L}^{\le M}_{\sigma_E,\widehat{f},H_{\le M}}.$
We show in the following that $\mathcal{L}_{\sigma_E,\widehat{f},H}$ is well-approximated by the generator $\mathcal{L}^{\le M}_{\sigma_E,\widehat{f},H_{\le M}}$
 on certain energy constraint input states and for truncation parameter $M$ large enough. 

 In Section~\ref{sec:TruncateJumps}, we first introduce the bounded generator $\mathcal{L}^{\le M}_{\sigma_E,\widehat{f},H}$ by truncating the bare jumps $A^\alpha$ within the unbounded generator and continue to show closeness of this bounded generator to $\mathcal{L}_{\sigma_E,\widehat{f},H}.$ Then in Section~\ref{sec:FinitDimGenerator} we further truncate the Hamiltonian to define the finite-dimensional generator $\mathcal{L}^{\le M}_{\sigma_E,\widehat{f},H_{\le M}}$ and show closeness to $\mathcal{L}^{\le M}_{\sigma_E,\widehat{f},H}.$ Lastly, we show in Section~\ref{sec:FiniteDimPreperationSchwartz} that the dynamics $e^{t\mathcal{L}_{\sigma_E,\widehat{f},H}}$ is well approximated by the finite-dimensional one $e^{t\mathcal{L}^{\le M}_{\sigma_E,\widehat{f},H_{\le M}}}.$

   
\subsubsection{Truncating the bare jumps}
\label{sec:TruncateJumps}
In the following section we focus on approximating bare jumps $\{A^\alpha\}$ by finite-dimensional\footnote{Technically, the truncated bare jumps considered here are finite rank operators defined on an infinite-dimensional space. But we drop this distinction here for simplicity.} bare jumps which will be useful for implementing the Lindlabian dynamics considered in the previous sections on finite-dimensional hardware.

To illustrate this, let us focus first on bare jumps being the  annihilation and creation operators. In particular for truncation parameter $M\in\N,$ we consider $\pi_M:= \sum_{n=0}^M\kb{n}$ and
\begin{align}
\label{eq:DefTruncAnnCrea}
    a^{\le M} :=  \pi_Ma\pi_M=  \sum_{n=1}^M \sqrt{n}\ket{n-1}\!\bra{n},\quad\text{and}\quad\left(a^{\le M}\right)^\dagger = \pi_Ma^\dagger \pi_M=\sum_{n=0}^{M-1}\sqrt{n+1} \ket{n+1}\!\bra{n}.
\end{align}
More generally, for $k\in\N,$ we can also consider truncations of higher order bare jumps like $a^k$ and $(a^k)^\dagger$ which are defined by
\begin{align*}
    \left(a^k\right)^{\le M} &:= \pi_Ma^k\pi_M=\sum_{n=k}^M \sqrt{\frac{n!}{(n-k)!}}\ket{n-k}\!\bra{n},\quad\text{and}
\\\left((a^k)^{\le M}\right)^\dagger &=  \pi_M (a^\dagger)^k\pi_M= \sum_{n=0}^{M-k}\sqrt{\frac{(n+k)!}{n!}}\ket{n+k}\!\bra{n}.
\end{align*}

In the following lemma we see that these finite-dimensional truncations approximate the untruncated bare jumps well when applied on energy constraint input states.

\bin{\begin{lemma}
\label{lem:FiniteTruncBareJumps;}
Let $\kappa\in(0,1/2],$ $k\in\N$ and $M\in\N$ satisfying $M\ge \left(\frac{k}{2\kappa}\right)^{1/\kappa}.$ \bin{$r\ge M^{-1/2}.$} Then for $\ket{\psi}\in D(e^{N^\kappa})$ we have
\begin{align*}
    \|(a - a^{\le M}) \ket{\psi}\| &\le \sqrt{M}\,e^{-M^{\kappa}}\| e^{N^\kappa} \ket{\psi}\|,\\
     \|(a^\dagger - \left(a^{\le M}\right)^\dagger) \ket{\psi}\| &\le \sqrt{M+1}
    \,e^{-M^{\kappa}}\| e^{N^{\kappa}} \ket{\psi}\|.
\end{align*}
Furthermore, for $\psi\in D(e^{2N^{\kappa}})$ we have
\begin{align*}    \left\|e^{N^{\kappa}}(a - a^{\le M}) \ket{\psi}\right\| &\le \sqrt{M}\,e^{-M^{\kappa}}\| e^{2N^{\kappa}} \ket{\psi}\|,\\
     \left\|e^{N^{\kappa}}(a^\dagger - \left(a^{\le M}\right)^\dagger) \ket{\psi}\right\| &\lesssim \sqrt{M}
    \,e^{-M^{\kappa}}\| e^{2N^{\kappa}} \ket{\psi}\|.
\end{align*}
\end{lemma}
\begin{proof}
The first inequality follows from noting \begin{align*}
    \left\|(a- a^{\le M}) e^{-N^{\kappa}}\right\| = \sup_{n\ge M+1} \sqrt{n} e^{-n^\kappa}\le \sqrt{M}e^{-M^{\kappa}},
\end{align*}
where for the second inequality we have used that $M\ge \left(\frac{1}{2\kappa}\right)^{1/\kappa}$ and the fact that the function $\sqrt{x}e^{-x^\kappa}$ is non-increasing for $x\ge \left(\frac{1}{2\kappa}\right)^{1/\kappa}.$
\bin{where for the second inequality we have used $r\ge M^{-1/2}$ together with the fact that the function $xe^{-rx}$ is non-increasing for $x\ge 1/r.$} Similarly, we see
\begin{align*}
\left\|\left(a^\dagger - \left(a^{\le M}\right)^\dagger\right)e^{-\sqrt{N}}\right\|&\le \sup_{n\ge M}\sqrt{n+1}e^{-n^{\kappa}} \le  \sqrt{M+1}e^{-M^{\kappa}},
\end{align*}
where for the second inequality we have used that the function $\sqrt{x+1}e^{-\sqrt{x}}$ is non-increasing for $x\ge \left(\frac{1}{2\kappa}\right)^{1/\kappa}.$ 


 
\end{proof}}

\begin{lemma}
\label{lem:FiniteTruncBareJumps}
Let $\kappa\in(0,1/2],$ $k\in\N$ and $M\in\N$ be such that  $M\ge \left(\frac{k}{2\kappa}\right)^{1/\kappa}+k.$ \bin{$r\ge M^{-1/2}.$} Then for $\ket{\psi}\in D(e^{N^\kappa})$ we have
\begin{align*}
    \left\|\left(a^k - (a^k)^{\le M}\right) \ket{\psi}\right\| &\le M^{k/2}\,e^{-M^{\kappa}}\| e^{N^\kappa} \ket{\psi}\|, \\\left\|\left((a^k)^\dagger - \left((a^k)^{\le M}\right)^\dagger\right) \ket{\psi}\right\| &\lesssim M^{k/2}
    \,e^{-M^{\kappa}}\| e^{N^{\kappa}} \ket{\psi}\|,
\end{align*}
Furthermore, for $\psi\in D(e^{2N^{\kappa}})$ we have
\begin{align*}    \left\|e^{N^{\kappa}}(a^k - (a^k)^{\le M}) \ket{\psi}\right\| &\le M^{k/2}\,e^{-M^{\kappa}}\| e^{2N^{\kappa}} \ket{\psi}\|,\\
     \left\|e^{N^{\kappa}}\left((a^k)^\dagger - \left((a^k)^{\le M}\right)^\dagger\right) \ket{\psi}\right\| &\lesssim M^{k/2}
    \,e^{-M^{\kappa}}\| e^{2N^{\kappa}} \ket{\psi}\|.
\end{align*}
 The constants hidden in the $\lesssim$-notation in the above inequalities depend on $k$ and $\kappa$ but on no other variables.
\end{lemma}

\begin{proof}
The first inequality follows from noting
    \begin{align*}
    \left\|(a^k- (a^k)^{\le M}) e^{-N^{\kappa}}\right\| = \sup_{n\ge \max\{M+1,k\}} \sqrt{\frac{n!}{(n-k)!}} e^{-n^\kappa}\le \sup_{n\ge \max\{M+1,k\}} n^{k/2} e^{-n^\kappa}\le M^{k/2}e^{-M^{\kappa}},
\end{align*}
where for the last inequality we have used that $M\ge \left(\frac{k}{2\kappa}\right)^{1/\kappa}\ge k$ and  that the function $ x^{k/2}e^{-x^{\kappa}}$ is non-increasing for $x\ge \left(\frac{k}{2\kappa}\right)^{1/\kappa}.$

Similarly, we see
\begin{align*}
\left\|\left((a^k)^\dagger - \left((a^k)^{\le M}\right)^\dagger\right)e^{-N^\kappa}\right\|&\le \sup_{n\ge M-k+1}\sqrt{\frac{(n+k)!}{n!}}e^{-n^{\kappa}} \le \sup_{n\ge M-k+1}(n+k)^{k/2}e^{-n^{\kappa}}   \\&\le e^{\kappa(k-1)}(M+1)^{k/2}e^{-M^{\kappa}},
\end{align*}
where for the last inequality we have used that the function $(x+k)^{k/2}e^{-x^\kappa}$ is non-increasing for $x\ge \left(\frac{k}{2\kappa}\right)^{1/\kappa}$ and furthermore that $e^{M^{\kappa}-(M-k+1)^\kappa}\le e^{\kappa(k-1)}.$

The third and fourth inequality follow similarly while noting for the fourth inequality that $e^{(n+k)^{\kappa}-n^\kappa}\le e^{k\kappa}.$
\end{proof}

Going beyond, we want to consider a finite-dimensional truncation scheme for a set of general bare jumps $\{A^\alpha\}_{\alpha\in\mathcal{A}}$ on some separable Hilbert space $\cH:$ \bin{For that let $\{N_\alpha\}_{\alpha \in\cA}$  be a set of self-adjoint operators with $N_\alpha\ge 0$ and having discrete spectrum and eigenvectors $\left(\ket{\nu^\alpha_n}\right)_{n\in \N_0}.$ For $M\in \N$ consider the spectral projection   $\pi^{\alpha}_M =\sum_{n=0}^M \kb{\nu^{\alpha}_n}$ and, assuming that $\pi^{\alpha}_M\cH\subseteq D(A^{\alpha}),$  the truncated bare jump $\left(A^{\alpha}\right)^{\le M} = \pi^{\alpha}_M A^{\alpha}\pi^{\alpha}_M.$ We also assume that we have good control on the operator norm of the truncated jump, i.e. precisely}
 For $M\in \N$ and $\alpha\in\cA$ we consider finite rank projections $\pi^{\alpha}_M$ with $\operatorname{rank}(\pi^{\alpha}_M) =M+1. $ Assuming that $\pi^{\alpha}_M\cH\subseteq D(A^{\alpha})$ we then define the truncated bare jump $\left(A^{\alpha}\right)^{\le M} := \pi^{\alpha}_M A^{\alpha}\pi^{\alpha}_M.$ We assume that we have good control on the operator norm of the truncated jump, i.e. precisely 
\begin{align}
\label{eq:NormTruncatedJump}
    \left\|(A^\alpha)^{\le M}\right\| \le q'(M),
\end{align}
for some polynomially bounded function $q'(M).$ 

    For the remaining section we assume that $A^{\alpha}$ is well-approximated by the truncated bare jump on energy constrained inputs: Precisely, to measure energy we consider a self-adjoint and positive semidefinite operator, $\NA,$ and assume that for $l=1,2$ and some $\kappa\in(0,1/2)$ we have
\begin{align}
\label{eq:AbstractFiniteTruncBareJumps}  
    \left\|e^{(l-1)\NA^{\kappa}}(A^{\alpha}-(A^{\alpha})^{\le M})e^{-l\NA^{\kappa}}\right\| \le q(M)e^{-M^{\kappa}},\bin{ \quad  \left\|e^{(k-1)\NA}((A^{\alpha})^\dagger-\left((A^{\alpha})^{\dagger}\right)^{\le M})e^{-k\NA^{\kappa}}\right\| \le q(M)e^{-M^{\kappa}}, }
\end{align}
where $q(M)$ is some polynomially bounded function. \bin{ $N_{\mathcal{A}} =\frac{1}{2} \sum_{\alpha\in\mathcal{A}} N_\alpha.$} For certain parts of the proofs of Section~\ref{sec:FinitDimGenerator} we also assume that
\begin{align}
\label{eq:NormTruncatedJump2}
\left\|e^{\NA^{\kappa}}(A^{\alpha})^{\le M} e^{-\NA^{\kappa}}\right\| \le q'(M)\quad\text{and}\quad\left\|e^{-\NA^{\kappa}}(A^{\alpha})^{\le M} e^{\NA^{\kappa}}\right\| \le q'(M)
\end{align}
for the polynomially bounded function $q'(M)$ which also appeared in \eqref{eq:NormTruncatedJump}.

To illustrate the above truncation scheme, we consider the example of an $m$-mode bosonic system on the Hilbert space $\left(L^2(\R)\right)^{\otimes m}$: In this case, we can consider multi-indices $\alpha$ taken from the set $\cA=\{(i,+),(i,-)\}_{i=1}^m$ and bare jumps $\{A^\alpha\}_{\alpha\in \cA}=\left\{a^k_i,(a_i^k)^\dagger\right\}_{i=1}^m $ for some $k\in \N.$ Furthermore, for $\alpha= (i,+)$ or $\alpha=(i,-)$, we consider the specific choice of rank-$(M+1)$ projections given by local truncations in the Fock basis, i.e., $\pi^{\alpha}_M \equiv \pi^i_M = \1_{1,\cdots,i-1}\otimes \pi_M\otimes \1_{i+1,\cdots,m} $, where $\pi_M = \sum_{n=0}^{M} \kb{n}.$ 
 In this case, we clearly have that \eqref{eq:NormTruncatedJump} is satisfied with $q'(M) = M^{k/2}.$ 
Lemma~\ref{lem:FiniteTruncBareJumps} ensures for this particular choice of bare jumps and truncations that the relation~\eqref{eq:AbstractFiniteTruncBareJumps} holds, where in this case we take $\NA = \Ntot = \sum_{i=1}^m N_i.$ Alternatively, in \cite{BeckerRouzeSalzmannSchroedinger} we consider $\NA$ to be a finite rank perturbation of $\Ntot$ and verify relation~\eqref{eq:AbstractFiniteTruncBareJumps} for this choice as well. 


Next, we consider the generator $\mathcal{L}_{\sigma_E,\widehat{f},H}$ with bare jumps $\{A^{\alpha}\}_{\alpha\in\cA}$, for which Proposition~\ref{integralrep} provided an explicit integral representation. We see in the following that $\mathcal{L}_{\sigma_E,\widehat{f},H}$ is close to the generator $\mathcal{L}^{\le M}_{\sigma_E,\widehat{f},H}$ which is defined by replacing $A^\alpha$ by $\left(A^{\alpha}\right)^{\le M}.$  For this we assume that energy measured with respect $e^{2\NA^\kappa}$ can only increase subexponentially with respect to time when evolved with the unitary dynamics generated by $H$: More precisely we assume for $k=2,4$ and $\kappa\in(0,1/2)$ as above that there exists $r\ge 0$ such that for all $t\in\R$ we have\footnote{The reason for the choice of range of $\kappa,$ in particular the constraint $\kappa<1/2$, is that this enables us to prove approximation of $\mathcal{L}_{\sigma_E,\widehat{f},H}$ by $\mathcal{L}^{\le M}_{\sigma_E,\widehat{f},H}$ for all $\beta>0.$ The reason for this is that the function $e^{r|t|^{2\kappa}/2}$ needs to be dominated by the function $g(t),$ which is featured in the integral representations of the generators stated in Proposition~\ref{integralrep}, and it can be shown that $|g(t)|
\sim
e^{-\tfrac{\pi}{\beta}|t|}$ for large $t.$ }
\begin{align}\label{eq:ExpBound}
e^{-itH}e^{k\NA^\kappa} e^{itH}\le e^{r|t|^{2\kappa}}\,e^{k\NA^\kappa}.
\end{align}
Such a condition is natural and holds, e.g., for the choice $\NA=\Ntot = \sum_{i=1}^m N_i$ and the Bose-Hubbard model and for all $\kappa\in(0,1/2)$ and $r=0$ since the corresponding Hamiltonian commutes with $\Ntot$, and, as we see in Lemma~\ref{lem:mfBHEnergyLimited} below, for the mean field Bose-Hubbard Hamiltonian, which is studied in the companion paper \cite{BeckerRouzeSalzmannBose}, for all $\kappa\in[0,1/2]$ and some $r>0$ depending only on $\kappa$ and the interaction strength $\psi.$ 

Another natural choice of $\NA$ is given by $\NA = \widetilde H$ where $\widetilde H = H +h_0+1$ and $H \ge -h_0.$ For this choice \eqref{eq:ExpBound} is trivially satisfied for $r=0.$\footnote{In this case, as $r=0,$ we could even consider $\kappa\in(0,1]$.} On the other hand, conditions~\eqref{eq:AbstractFiniteTruncBareJumps} and \eqref{eq:NormTruncatedJump2} are non-trivial and need to be verified explicitly in this case.

We now proceed to show that under conditions \eqref{eq:AbstractFiniteTruncBareJumps} and \eqref{eq:ExpBound}, the generators $\mathcal{L}^{\le M}_{\sigma_E,\widehat{f},H}$ and $\mathcal{L}_{\sigma_E,\widehat{f},H}$ are close. For this, we first prove the following technical Lemma. In what follows, we denote for some function $f\in\mathcal{S}(\mathbb{R})$
\begin{align*}L^\alpha:=\int f(s)e^{isH}A^\alpha e^{-isH}ds,\qquad 
\bin{L^{\alpha}^d:=\int f(s)e^{isH}a_i^\dagger e^{-isH}ds}
L^{\alpha,\le M}&:=\int f(s)e^{isH}(A^{\alpha})^{\le M} e^{-isH}ds.\bin{\qquad L_{i}^{\le M,d}:=\int f(s)e^{isH}(a^{\le M})^\dagger_i e^{-isH}ds.}
\end{align*}

\begin{lemma}\label{lem:integralrepapproxTrunc}
Let $f\in\mathcal{S}(\mathbb{R})$ and set of bare jumps $\{A^\alpha\}_{\alpha\in\cA}$ and assume  \Cref{eq:condAalphas} and, for some self-adjoint and positive semidefinite $\NA$ and $\kappa\in(0,1/2),$  \eqref{eq:AbstractFiniteTruncBareJumps}  are satisfied and further that $e^{\NA^\kappa}A^{\alpha}e^{-2\NA^{\kappa}},$ $e^{-\NA^{\kappa}}A^{\alpha}$ and $A^{\alpha}e^{-\NA^{\kappa}}$ are bounded.
Furthermore, assume that \eqref{eq:ExpBound} holds for $r\ge 0$ and $\kappa\in(0,1/2)$ as above and that
\begin{align}
\label{eq:C_f}
    C_f := \int_{-\infty}^{\infty} e^{r|t|^{2\kappa}} |f(t)|dt < \infty.
\end{align}
Then for $M\in \N$ we have
\begin{align*}
&\|(L^{\alpha}\!-\!L^{\alpha,\le M})e^{-\NA^\kappa}\|\!\le\! C_f\,q(M) e^{-M^{\kappa}},\\
\bin{&\|(L^{d}_{i}\!-\!L^{\le M,d}_{i})e^{-\NA^{\kappa}}\|\!\lesssim  C_f\,\sqrt{M}e^{-\sqrt{M}},\, \,\\}
&\|((L^\alpha)^\dagger L^{\alpha}-\left(L^{\alpha,\le M}\right)^\dagger L^{\alpha,\le M})e^{-2\NA^{\kappa}}\|\le C^2_f C_{A}\,q(M)e^{-M^\kappa},\, \\
\end{align*}
where we denoted $C_A := \max_{\alpha\in\cA}\left(\|e^{\NA^{\kappa}}A^{\alpha}e^{-2\NA^{\kappa}}\|+\|A^{\alpha}e^{-\NA^{\kappa}}\|+\|e^{-\NA^{\kappa}}A^{\alpha}\|\right)<\infty.$
\end{lemma}

\begin{proof}

\bin{First of all note that for $i=1,\cdots, m$ we have that $N_i$ and $\NA$ commute from which we can easily see that
\begin{align}
\label{eq:NtotdomNi}
\left\|e^{\sqrt{N_i}}e^{-\NA^{\kappa}}\right\| = 1.
\end{align}}
 We use \eqref{eq:AbstractFiniteTruncBareJumps} together with \eqref{eq:ExpBound} to see for $\ket{\psi}\in D\left(e^{\NA^{\kappa}}\right)$ that we have
\begin{align}
\label{eq:NLLeqTrunc}
\nn\|(L^{\alpha}-L^{\alpha,\le M}_{i})\ket{\psi}\|
 \nn&\le \int |f(s)|\|(A^{\alpha}-(A^{\alpha})^{\le M})e^{-isH}\ket{\psi}\|\,ds\\&\nn\le  \int |f(s)| \left\|(A^{\alpha}-(A^{\alpha})^{\le M})e^{-\NA^{\kappa}}\right\|\left\|e^{\NA^{\kappa}}e^{-isH}\ket{\psi}\right\| ds\\&\nn\le  \int |f(s)| e^{\tfrac{r|s|^{2\kappa}}{2}}\left\|(A^{\alpha}-(A^{\alpha})^{\le M})e^{-\NA^{\kappa}}\right\|\left\|e^{\NA^{\kappa}}\ket{\psi}\right\| ds\\&\le C_f\,q(M)e^{-M^{\kappa}}\left\|e^{\NA^{\kappa}}\ket{\psi}\right\| ,
\end{align}
which shows the first inequality in Lemma~\ref{lem:integralrepapproxTrunc}. 


\bin{First, we use the estimates derived in Lemma \ref{lem.displaceblocks} to see for  $\ket{\psi}\in D(\NA+I)$, denoting $N_i:=a_i^\dagger a_i$, 
\begin{align*}
\|(\NA+I)\ket{\psi}\|^2&=\bra{\psi}\left(\sum_{j=1}^m N_j+I\right)^2\ket{\psi}\\
&\le m\sum_{j\in[m]}\bra{\psi}(N_j+I)^2\ket{\psi}\\
&=m\sum_{j\in[m]}\|(N_j+I)\ket{\psi}\|^2.
\end{align*}
Applying this bound to $\ket{\psi}=(a_i-(a^D_\alpha)_i)\ket{\varphi}$, we get
\begin{align*}
&\|(\NA+I)(a_i-(a^D_\alpha)_i)\ket{\varphi}\|^2\\
&\quad \le m\,\sum_{j\ne i}4|\alpha|^2\|(N_j+I)(N_i+I)\ket{\varphi}\|^2\\
&\quad + 25m |\alpha|^2\|(N_i+I)^2\ket{\varphi}\|^2\\
&\quad \le 25m^2|\alpha|^2\sum_{j\in[m]}\,\|(N_j+I)^2\ket{\varphi}\|^2\\
&\quad \le 25m^3|\alpha|^2\|(\NA+I)^2\ket{\varphi}\|^2.
\end{align*}



Thus, we have derived that
\begin{align}
\|(\NA\!+\!1)(L_{i}\!-\!L_{\alpha,i})\!\ket{\varphi}\!\|\!\le\! 5C_{2,f}m^{\frac{3}{2}}|\alpha|\|(\NA+I)^2\!\ket{\varphi}\!\|.
\end{align}}

\bin{Similarly, denoting $L^{\alpha}^d:=\int f(s)e^{isH}a_i^\dagger e^{-isH}ds$ and $L_{\alpha,i}^d:=\int f(s)e^{isH}(a_{\alpha}^{Dd})_i e^{-isH}ds$, we get
\begin{align}\label{eq:NLLeqd}
\|(\NA\!+\!1)(L^d_{i}\!-\!L^d_{\alpha,i})\!\ket{\varphi}\!\|\!\le\! 5C_{2,f}m^{\frac{3}{2}}|\alpha|\|(\NA+I)^2\!\ket{\varphi}\!\|.
\end{align}}
Next, we control squares of the $L$ operators for $\ket{\varphi}\in D(e^{2\NA^{\kappa}})$:
\begin{align}\label{eq:twotermstoboundTrunc}
\|((L^\alpha)^\dagger L^{\alpha}-\left(L^{\alpha,\le M}\right)^\dagger L^{\alpha,\le M})\ket{\varphi}\|&\le \|(L^\alpha)^\dagger (L^{\alpha}-L^{\alpha,\le M})\ket{\varphi}\|+\|((L^\alpha)^\dagger-\left(L^{\alpha,\le M}\right)^\dagger)L^{\alpha,\le M}\ket{\varphi}\|.
\end{align}
The first term above can be controlled by realizing that by assumption we have for $\ket{\psi}\in D(e^{\NA^{\kappa}})$ that $\ket{\psi}\in D((L^\alpha)^\dagger)$ with
\begin{align}\label{eq:LiboundTrunc}
\|(L^\alpha)^\dagger\ket{\psi}\|&\le \int |f(s)|\,\|(A^{\alpha})^\dagger e^{-isH}\ket{\psi}\|\,ds\le \int |f(s)|\,\|(A^{\alpha})^\dagger e^{-\NA^{\kappa}}\|\| e^{\NA^{\kappa}}e^{-isH}\ket{\psi}\|\,ds\nn\\&
\le C_f \|e^{-\NA^{\kappa}}A^{\alpha}\| \|e^{\NA^{\kappa}}\ket{\psi}\|. 
\end{align}
Hence, combining with an analogous argument to \eqref{eq:NLLeqTrunc}, we get
for $\ket{\phi}\in D(e^{2\NA^{\kappa}})$
\begin{align}\label{eq:firstboundLLTrunc}
\|(L^\alpha)^\dagger\, (L^{\alpha}-L^{\alpha,\le M})\ket{\varphi}\|\!\le \! C^2_f\,q(M)e^{-M^\kappa}\left\|e^{-\NA^{\kappa}}A^{\alpha}\right\|\|e^{2\NA^{\kappa}}\ket{\varphi}\!\|.
\end{align}
To control the second term in \eqref{eq:twotermstoboundTrunc}, we note that for $\ket{\phi}\in D(e^{2\NA^{\kappa}})$ we have
\begin{align*}
    \left\|e^{\NA^{\kappa}}L^{\alpha,\le M} \ket{\phi}\right\| &\le  \int |f(s)|e^{\frac{r|s|^{2\kappa}}{2}} \left\|e^{\NA^{\kappa}}e^{isH}(A^{\alpha})^{\le M}e^{-2\NA^{\kappa}}\right\|ds\,\left\|e^{2\NA^{\kappa}}\ket{\phi}\right\| \\ &\le \int |f(s)| e^{r|s|^{2\kappa}} ds \,\left\|e^{\NA^{\kappa}}A^{\alpha}e^{-2\NA^{\kappa}}\right\|\left\|e^{2\NA^{\kappa}}\ket{\phi}\right\|= C_f  \left\|e^{\NA^{\kappa}}A^{\alpha}e^{-2\NA^{\kappa}}\right\|\left\|e^{2\NA^{\kappa}}\ket{\phi}\right\|,
\end{align*}
\bin{where in the second inequality we have used that $\|e^{\NA^{\kappa}}a^{\le M}_ie^{-2\NA^{\kappa}}\|\le 1$ and that by \eqref{eq:ExpBound} we have $\|e^{\NA^{\kappa}}e^{isH}e^{-\NA^{\kappa}}\| \le e^{r|s|/2}.$ } From this 
and \eqref{eq:NLLeqTrunc} we bound the second term in \eqref{eq:twotermstoboundTrunc} as 
\begin{align*}
    \|((L^\alpha)^\dagger-\left(L^{\alpha,\le M}\right)^\dagger)L^{\alpha,\le M}\ket{\varphi}\| \le C^2_f\, q(M)e^{-M^{\kappa}} \left\|e^{\NA^{\kappa}}A^{\alpha}e^{-2\NA^{\kappa}}\right\|\left\|e^{2\NA^{\kappa}}\ket{\varphi}\right\|,
\end{align*}
which finishes the proof.

\end{proof}

Next, we show that the generators $\mathcal{L}^{\le M}_{\sigma_E,\widehat{f},H}$ and $\mathcal{L}_{\sigma_E,\widehat{f},H}$ are close when evaluated on states $\rho$ satisfying the superpolynomial energy constraint 
\begin{align*}
    \Tr\left(e^{4\NA^{\kappa}}\rho\right) <\infty.
\end{align*}

\begin{proposition}\label{prop:generatorboundTrunc}
Let $f\in\mathcal{S}(\mathbb{R})$ and set of bare jumps $\{A^\alpha\}_{\alpha\in\cA}$ and assume  \Cref{eq:condAalphas} and, for some self-adjoint and positive semidefinite $\NA$ and $\kappa\in(0,1/2),$  \eqref{eq:AbstractFiniteTruncBareJumps}  are satisfied and further that $e^{\NA^\kappa}A^{\alpha}e^{-2\NA^{\kappa}},$ $e^{-\NA^{\kappa}}A^{\alpha}$ and $A^{\alpha}e^{-\NA^{\kappa}}$ are bounded. Furthermore, assume that \eqref{eq:ExpBound} holds for $r\ge 0$ and $\kappa$ as above and that
\begin{align*}
    C_f := \int_{-\infty}^{\infty} e^{r|t|^{2\kappa}}  |f(t)|dt < \infty.
\end{align*}
Then for all $\beta,\,\sigma_E>0,$ $M\in\N$ and states $\rho$ satisfying $E_k:=\Tr(e^{k\NA^{\kappa}}\rho)<\infty$, $k\in\{2,4\}$, we have
\begin{align}
\left\|(\mathcal{L}_{\sigma_E,\widehat{f},H}-\mathcal{L}^{\le M}_{\sigma_E,\widehat{f},H})(\rho)\right\|_1\le C_{\beta,r,\kappa}C^2_f\, C_A \,|\cA| \,q(M)e^{-M^{\kappa}}\left(\sigma_E
\exp\!\left(\tfrac{\pi^2}{8\beta^2\sigma_E^2}\right)\sqrt{E_4}+\exp\left(r+\tfrac{r^2}{8\sigma^2_E}\right)\,E_2\right),
\end{align}
for some constant $C_{\beta,r,\kappa}\ge 0$ and with $C_A$ being defined in Lemma~\ref{lem:integralrepapproxTrunc}.
\end{proposition}

\begin{proof}


We make use of the integral representation of Proposition \ref{integralrep}. First, we denote the operator 
\begin{align*}
G^{\le M}_{\sigma_E}\!:=\!-\sum_{\alpha\in\cA}\!\int_{-\infty}^\infty\! \!\!\!\!g(t)e^{itH} (L^{\alpha,\le M})^\dagger L^{\alpha,\le M}e^{-itH}\!dt.
\end{align*}
Denoting $\Delta (L^{\alpha})^2:=(L^{\alpha,\le M})^\dagger L^{\alpha,\le M}-(L^{\alpha})^\dagger L^{\alpha},$ we directly obtain from Lemma \ref{lem:integralrepapproxTrunc} that for any state $\rho$ for which  $\Tr(e^{4\NA^{\kappa}}\rho)\le E_4<\infty$ we have 
\begin{align*}
\|(G^{\le M}_{\sigma_E}-G_{\sigma_E})\rho\|_1
  &\le \sum_{\alpha\in\cA} \int |g(t)|\, \|\Delta (L^{\alpha})^2 e^{itH}\rho\|_1 \le \sum_{\alpha\in\cA}\|\Delta (L^{\alpha})^2 e^{-2\NA^{\kappa}}\|\int |g(t)|\left\|e^{2\NA^{\kappa}}e^{-itH}\rho\right\|_1dt\\
&\le C^2_fC_A\,|\cA| q(M)e^{-M^{\kappa}}\sqrt{E_4}\int |g(t)|e^{r|t|^{2\kappa}/2}dt\\& \le 
C_{\beta,r,\kappa}C^2_fC_A\,|\cA|\,\sigma_E
\exp\!\left(\tfrac{\pi^2}{8\beta^2\sigma_E^2}\right)q(M)e^{-M^{\kappa}}\sqrt{E_4},
\end{align*}
where in the last inequality we have used\footnote{Recall the definition $g(t)=\frac{1}{2\pi}\int_{-\infty}^\infty \frac{e^{-\nu^2/8\sigma_E^2}}{1+e^{\beta\nu/2}} e^{-i\nu t}d\nu$, and note $\tfrac{e^{-\nu^2/8\sigma_E^2}}{1+e^{\beta\nu/2}}$ is analytic for $|\Im{\nu}| <\frac{2\pi}{\beta}$. Hence, by shifting the Contour of integration in the definition of $g(t)$ to $\nu\mapsto \nu -i\tfrac{\pi}{\beta}$ for $t>0$ and to $\nu\mapsto \nu +i\tfrac{\pi}{\beta}$ for $t>0$ a straigtforward calculation yields $|g(t)|
\le
\frac{\sigma_E}{\sqrt{2\pi}}
\exp\!\left(\frac{\pi^2}{8\beta^2\sigma_E^2}\right)
e^{-\tfrac{\pi}{\beta}|t|}.$ From this and $\kappa\in(0,1/2),$ another straightforward calculation yields \eqref{eq:gkappaIntegralBound}. Note that the constant $C_{\beta,r,\kappa}$ is finite for all $\beta,r\ge 0$ and scales exponentially for large $\beta.$}
\begin{align}
\label{eq:gkappaIntegralBound}
\int_{\mathbb{R}} |g(t)|\, e^{r |t|^{2\kappa}/2}\, dt
\;\le C_{\beta,r,\kappa} \,
\sigma_E\,\,
\exp\!\left(\frac{\pi^2}{8\beta^2\sigma_E^2}\right),\,
\bin{\exp\!\left[
\frac r2(1-2\kappa)\left(\frac{2\kappa\, r\,\beta}{\pi}\right)^{\frac{2\kappa}{1-2\kappa}}
\right]}
\end{align}
for some finite and $\beta,r$ and $\kappa$ dependent constant $C_{\beta,r,\kappa}\ge 0.$
Similarly, we obtain that 
\begin{align*}
    &\|\rho(G^{\dagger}_{\sigma_E}-\left(G^{\le M}_{\sigma_E}\right)^\dagger)\|_1\le C_{\beta,r,\kappa}C^2_fC_A\,|\cA|\,\sigma_E
\exp\!\left(\tfrac{\pi^2}{8\beta^2\sigma_E^2}\right)q(M)e^{-M^{\kappa}}\sqrt{E_4}.\end{align*}
Next, for 
\begin{align}
 X^\alpha_s:=e^{isH}L^{\alpha} e^{-isH}
,\quad  X^{\alpha,\le M}_s:=e^{isH}L^{\alpha,\le M} e^{-isH}\label{eq:XsaiTrunc}
\end{align}
we get for any $s\in\mathbb{R}$, using an analogous argument to \eqref{eq:LiboundTrunc} and the fact that $|t+s|^{2\kappa}\le |t|^{2\kappa}+|s|^{2\kappa}$ that 
\begin{align*}
\|X^{\alpha}_s\sqrt{\rho}\|_2 &\le C_A\sqrt{E_2} \int |f(t)| e^{r|t+s|^{2\kappa}/2}dt \le C_A C_f\sqrt{E_2} \,e^{r|s|^{2\kappa}/2}\\
    \|X^{\alpha,\le M}_s\sqrt{\rho}\|_2 &\le C_A\sqrt{E_2} \int |f(t)| e^{r|t+s|^{2\kappa}/2}dt\le C_AC_f\sqrt{E_2}\,e^{r|s|^{2\kappa}/2}
\end{align*}
and furthermore, by an analogous argument to \eqref{eq:NLLeqTrunc} that 
\begin{align*}
 &\left\|(X^{\alpha}_s-X^{\alpha,\le M}_s)\sqrt{\rho}\right\|_2\le C_f\,q(M)e^{-M^{\kappa}} \sqrt{E_2} \,e^{r|s|/2}
\end{align*}
which in total gives
\begin{align*}
&\left\|X^{\alpha}_s\rho (X^{\alpha})_s^\dagger-X^{\alpha,\le M}_s\rho(X^{\alpha,\le M}_s)^\dagger\right\|_1
\\&\qquad\le \|((X^{\alpha}_s-X^{\alpha,\le M}_s)\sqrt{\rho}\|_2\|\sqrt{\rho}(X^{\alpha}_s)^\dagger\|_2+ \|X^{\alpha,\le M}_s\sqrt{\rho}\|_2\|\sqrt{\rho}((X^{\alpha}_s)^\dagger-(X^{\alpha,\le M}_{s})^\dagger)\|_2
\\&\qquad\le 2C_AC^2_f\,E_2 e^{r|s|}q(M)e^{-M^{\kappa}}\,
\end{align*}
\bin{Now, by definition, the $X$ operators above correspond to operators $L$ with a shifted function $f_s(u)=f(u-s)$. Hence, the bounds derived in Lemma \ref{integralrepapprox} still hold, up to replacing all integrals of the form $\int |f(u)| C(|u|)du$ by
\begin{align*}
\sup_{s\in\mathbb{R}}\int |f(u-s)|C(|u|)du.
\end{align*}
Thus, we get a constant $C_f'$ such that, given $\Tr(\rho (\NA+I)^2)\le E_2$
\begin{align*}
&\|(X_{i})_s\rho (X_{i})_s^\dagger-(X_{\alpha,i})_s\rho(X_{\alpha,i})^\dagger_s\|_1\\
& \qquad\le C_f'|\alpha|\,\|(\NA+I)^2\sqrt{\rho}\|_2\|\sqrt{\rho}(\NA+I)\|_2\\
&\qquad \le C_f'|\alpha| \,\sqrt{E_2\cdot E_4},
\end{align*}
where we also used \eqref{Liboundeq} in the bound above. A similar bound can also be derived for the operators $X_{i}^d$ and $X_{\alpha,i}^d$.}
Therefore, denoting the map
\begin{align}\label{phialphasigmaEmapTrunc}
\Phi^{\le M}_{\sigma_E,\widehat f,H}(\rho):=\sigma_E\sqrt{\frac{2}{\pi}}\sum_{\alpha\in\cA}\int_{\mathbb{R}}e^{-2\sigma_E^2 s^2}X^{\alpha,\le M}_s \rho (X^{\alpha,\le M}_s)^\dagger ds,
\end{align}
 we get that
\begin{align*}
\|(\Phi^{\le M}_{\sigma_E,\widehat f,H}-\Phi_{\sigma_E,\widehat{f},H})(\rho)\|_1\le 4 |\cA| C_AC^2_f\,E_2 q(M)e^{-M^{\kappa}}\exp\left(\tfrac{r^2}{8\sigma^2_E}\right), 
\end{align*}
where we used that  
\begin{align*}
  \sqrt{\frac{2}{\pi}}\,\sigma_E\int e^{-2\sigma^2_Es^2} e^{r|s|^{2\kappa}}ds\le e^{r} \sqrt{\frac{2}{\pi}}\,\sigma_E\int e^{-2\sigma^2_Es^2} e^{r|s|}ds\le 2\exp\left(r+\tfrac{r^2}{8\sigma^2_E}\right).
\end{align*}

\end{proof}
To end this section, we show that the mean field Bose Hubbard Hamiltonian, $$H_{\text{mfBH}}= -\mu N +\tfrac{U}{2} \tfrac{N(N-1)}{2} -\overline{\psi} a - \psi a^\dagger  + \vert \psi \vert^2,$$ satisfies the assumption \eqref{eq:ExpBound}. More generally, we consider one-mode Hamiltonians of the form
\begin{align*}
    H= h(N) + \overline{\psi} a + \psi a^\dagger
\end{align*}
for some function $h: \N_0 \to \R$ and $\psi\in\C.$  In the following Lemma we find that for $\kappa \in[0,1/2]$ and $k\ge 0$ these Hamiltonians are energy limited with respect to  $e^{kN^{\kappa}}.$  
\begin{lemma}
\label{lem:mfBHEnergyLimited}
Let $\psi \in\C$ and $H = h(N)+ \overline{\psi} a + \psi a^\dagger.$
Then for $0\le \kappa\le 1/2$ and $k\ge 0$ and $t\in\R$ we have
\begin{align}
    e^{itH} e^{kN^{\kappa}}e^{-itH} \le e^{r|t|^{2\kappa}}  e^{kN^\kappa}
\end{align}
for some $r\ge 0$ depending only on $k,\kappa$ and $\psi.$
\end{lemma}
\begin{proof}
We consider $G(t):=e^{itH} e^{k(N+1+t^2)^\kappa}e^{-itH}$ and differentiate 
\begin{align}
\label{eq:GoperatorDeriv}
    \frac{d}{dt}G(t) &=2k\kappa te^{itH} (N+1+t^2)^{\kappa-1} e^{k(N+1+t^2)^\kappa}e^{-itH} + ie^{itH} \left[H,e^{k(N+1+t^2)^\kappa}\right]e^{-itH}.
\end{align}
In the following we estimate both operators appearing on the right hand side respectively. For the first term we use $\kappa -1\le 0$ which gives
\begin{align}
\label{eq:FirstTermBound}
   \nn 2k\kappa|t|\langle \psi,e^{itH} (N+1+t^2)^{\kappa-1} e^{k(N+1+t^2)^\kappa}e^{-itH}\psi\rangle &\le 2k |t|^{2\kappa -1}\langle \psi,e^{itH}  e^{k(N+1+t^2)^\kappa}e^{-itH}\psi\rangle\nn\\&\le 2k |t|^{2\kappa -1}\langle \psi,G(t)\psi\rangle
\end{align}
For the second term we use
\begin{align}
\label{eq:CommBound}\left[H,e^{k(N+1+t^2)^\kappa}\right] &= \overline{\psi} \left(e^{k(N+2+t^2)^\kappa} -e^{k(N+1+t^2)^\kappa} \right)a + \psi \left(e^{k(N+t^2)^\kappa} -e^{k(N+1+t^2)^\kappa} \right)a^\dagger
\end{align}
and bound each term individually. For the first term we expand $\ket{\varphi} = \sum_{n=0}^{\infty} \varphi_n\ket{n}$ and note
\begin{align*}
    &|\bra{\varphi} \left(e^{k(N+2+t^2)^{\kappa}} -e^{k(N+1+t^2)^{\kappa}} \right)a\ket{\varphi}| \le  \sum_{n=1}^{\infty} |\varphi_{n}||\varphi_{n-1}| \sqrt{n}\left|e^{k(n+1+t^2)^{\kappa}} -e^{k(n+t^2)^{\kappa}} \right| \\ &\le k\kappa\sum_{n=1}^{\infty} |\varphi_{n}||\varphi_{n-1}| \sqrt{n} (n+t^2)^{\kappa-1}e^{k(n+1+t^2)^{\kappa}} \le k\kappa |t|^{2\kappa -1} \sum_{n=1}^{\infty} |\varphi_{n}||\varphi_{n-1}| e^{k(n+1+t^2)^{\kappa}}
    \\&\le C |t|^{2\kappa-1}\bra{\varphi} e^{k(N+1+t^2)^\kappa}\ket{\varphi}, 
\end{align*}
 for some $C\ge 0$ which depends on $k$ and $\kappa$ but is independent of all other variables and where we have used in the second inequality we have used the mean value theorem, in the third the fact that $\kappa \le 1/2$ and in the fourth the Cauchy-Schwarz inequality. The second term in \eqref{eq:CommBound} can be treated similarly giving the bound
 \begin{align*}
     |\bra{\varphi} \left(e^{k(N+1+t^2)^{\kappa}} -e^{k(N+t^2)^{\kappa}} \right)a^\dagger\ket{\varphi}| \le C|t|^{2\kappa-1} \bra{\varphi} e^{k(N+1+t^2)^\kappa}\ket{\varphi}.
 \end{align*}
 Combining these estimates with \eqref{eq:GoperatorDeriv} and \eqref{eq:FirstTermBound} we get
 \begin{align*}
     \frac{d}{dt} G(t) \le C'|t|^{2\kappa-1} G(t) 
 \end{align*}
 for some constant $C'\ge 0$ depending only on $k,\kappa$ and $\psi.$ Using Gr\"onwall's inequality in the case $t\ge0$ on the interval $[0,t]$ or for the case $t\le 0$ on the interval $[0,-t]$ proves 
 \begin{align*}
     G(t) \le e^{r |t|^{2\kappa}} G(0) = e^{r |t|^{2\kappa}} e^{kN^\kappa}
 \end{align*}
 for some $r\ge 0$ depending again only on $k,\kappa$ and $\psi.$  Using $e^{kN^{\kappa}} \le e^{k(N+1+t^2)^{\kappa}}$ 
 finishes the proof.
\end{proof}










% \textcolor{red}{need more dicussion here, how to encode $U$?}
\subsubsection{Truncating the Hamiltonian}
\label{sec:FinitDimGenerator}

In the following we provide a finite rank Lindbladian, denoted by $\mathcal{L}^{\le M}_{\sigma_E,\widehat{f},H_{\le M}},$ which is good approximation of the generator $\mathcal{L}_{\sigma_E,\widehat{f},H}$ studied in Section~\ref{monotonicitygap} and \ref{integralrepfschwartz}. For that we consider the finite-dimensional approximations, $\{(A^{\alpha})^{\le M}\}_{\alpha\in\cA},$ of the bare jumps $\{A^\alpha\}_{\alpha\in\cA}$ analysed in Section~\ref{sec:FiniteTruncBareJumps} and replace the Hamiltonian $H$ within the corresponding generator $\mathcal{L}^{\le M}_{\sigma_E,\widehat{f},H}$ by its finite-dimensional truncation, $H_{\le M},$ to obtain $\mathcal{L}^{\le M}_{\sigma_E,\widehat{f},H_{\le M}}.$ As in Section~\ref{sec:FiniteTruncBareJumps} we have already seen that $\mathcal{L}^{\le M}_{\sigma_E,\widehat{f},H}$ is a good approximation of $\mathcal{L}_{\sigma_E,\widehat{f},H}$ on certain energy constraint input states, we focus in this section on the approximation of $\mathcal{L}^{\le M}_{\sigma_E,\widehat{f},H}$ by the fully finite-dimensional\footnote{Technically, $\mathcal{L}^{\le M}_{\sigma_E,\widehat{f},H_{\le M}}$ is a finite rank generator on an infinite-dimensional space. But we drop this distinction here for simplicity.} generator $\mathcal{L}^{\le M}_{\sigma_E,\widehat{f},H_{\le M}}.$


For that, we start with defining the truncated, finite rank Hamiltonian $H_{\le M}$ for some self-adjoint, possibly unbounded, Hamiltonian $H$ on some separable Hilbert space $\cH:$ More precisely we consider a parametrised family of self-adjoint and finite rank projections  $\left(P_M\right)_{M\in\N}$ on $\cH$ such that $P_M\cH \subseteq D(H)$ and with that the finite rank truncation 
\begin{align}\label{HleMpolybound}
    H_{\le M} := P_MHP_M\qquad \text{ with }\qquad  \|H_{\le M}\|\le p_5(|\mathcal{A}|,M)
\end{align}
for some polynomial $p_5$ of the number, $|\mathcal{A}|,$ of jumps and truncation $M$. As described in the Introduction, we refer to the finite dimensional space $\operatorname{im}(P_M)$ as the \emph{system register}. In many-body or multi-mode systems, one usually considers system registers whose dimensions satisfies
$
\log\left( \dim\bigl(\operatorname{im}(P_M)\bigr)\right)
= \mathcal{O}\bigl(|\mathcal{A}|\log(M)\bigr),
$ but we are in principle free to leave the dimension of the system register unspecified at this stage. 

We assume in the following that $H$ is well-approximated by $H_{\le M}$ on energy constraint input states; specifically, we have for some $\kappa\in(0,1/2)$
\begin{align}
 \label{eq:HamiltonianLeakTrunc}  
 \left\|\left(H-H_{\le M}\right)e^{-\NA^\kappa}\right\| \le p(|\cA|,M) e^{-M^\kappa},
\end{align}
where $p(|\cA|,M)$ is some polynomially bounded function and $\NA$ is some self-adjoint and positive semidefinite operator. \bin{Section~\ref{sec:FiniteTruncBareJumps} around equation \eqref{eq:AbstractFiniteTruncBareJumps}.}

Let us illustrate the above with the example of a $m$-mode bosonic system on the Hilbert space $\left(L^2(\R)\right)^{\otimes m},$ with multi-index $\alpha$ taken from the set $\cA=\{(i,+),(i,-)\}_{i=1}^m$ and bare jumps $\{A^\alpha\}_{\alpha\in \cA}=\left\{a^k_i,(a_i^k)^\dagger\right\}_{i=1}^m$ for some $k\in\N:$
In this case, we take $P_M = \pi^{\otimes m}_M$ with $\pi_M = \sum_{n=0}^M\kb{n}$ being the local truncations onto the first $M+1$ Fock states of a fixed mode. Furthermore, for this multimode case, we consider $\NA\equiv\Ntot:=\sum_{i=1}^m a^\dagger_ia_i$ such that \eqref{eq:HamiltonianLeakTrunc} is naturally satisfied for Hamiltonians $H$ being polynomials in $a_i$ and $a^\dagger_i$ with degree $\mathcal{O}(1)$ in the number of modes $m$, e.g., the Bose-Hubbard Hamiltonian. 

Additionally to \eqref{eq:HamiltonianLeakTrunc}, we assume as in Section~\ref{sec:FiniteTruncBareJumps} that $H$ satisfies \eqref{eq:ExpBound}, i.e.~that for $k=2,4$ and $\kappa\in(0,1/2)$ as above there exists $r\ge 0$ such that for all $t\in\R$ we have\footnote{The reason for the choice of range of $\kappa,$ in particular the constraint $\kappa<1/2$, is that this enables us to prove approximation of $\mathcal{L}^{\le M}_{\sigma_E,\widehat{f},H}$ by $\mathcal{L}^{\le M}_{\sigma_E,\widehat{f},H_{\le M}}$ for all $\beta>0.$ The reason for this is that the function $e^{r|t|^{2\kappa}/2}$ needs to be dominated by the function $g(t),$ which is featured in the integral representations of the generators stated in Proposition~\ref{integralrep}, and it can be shown that $|g(t)|
\sim
e^{-\tfrac{\pi}{\beta}|t|}$ for large $t.$ }
\begin{align}\label{eq:ExpBound2}
e^{-itH}e^{k\NA^\kappa} e^{itH}\le e^{r|t|^{2\kappa}}\,e^{k\NA^\kappa}.
\end{align}



\bin{\begin{align*}
    \left\|Q_Me^{-\NA^{\kappa}}\right\| \le \sup_{\substack{n_1,\cdots,n_m\in\N_0\\\exists n_i\ge M}} e^{-\sqrt{\sum_i n_i}} \le e^{-M^{\kappa}}.
\end{align*}}

Under both of these assumptions, we see in the following lemma that the unitary evolutions of the unbounded and truncated Hamiltonians are close on energy constraint inputs:
\begin{lemma}
\label{lem:HamEvolTrunc}
Let $H$ such that \eqref{eq:HamiltonianLeakTrunc} and \eqref{eq:ExpBound2} are satisfied for some $\kappa\in(0,1/2)$ and $r\ge 0$. Then we have
\begin{align*}
    \left\|\left(e^{-itH} - e^{-itH_{\le M}}\right)e^{-\NA^{\kappa}}\,\right\| \le |t|e^{\tfrac{r|t|^{2\kappa}}{2}}p(|\cA|,M)e^{-M^{\kappa}}.
\end{align*}

\end{lemma}
\begin{proof}
By Duhamel's principle we have
\begin{align*}
   &e^{-itH} - e^{-itH_{\le M}}= \int^t_0 e^{-isH_{\le M}}\left(H-H_{\le M}\right)e^{-isH} ds.
\end{align*}
Using \eqref{eq:ExpBound2}, we hence see
\begin{align*}
     &\left\|\left(e^{-itH} - e^{-itH_{\le M}}\right)e^{-\NA^{\kappa}}\,\right\| \le  \int_0^{|t|} e^{\tfrac{r|s|^{2\kappa}}{2}}ds\left\|(H-H_{\le M})e^{-\NA^{\kappa}}\right\|\le |t|e^{\tfrac{r|t|^{2\kappa}}{2}} p(|\cA|,M)e^{-M^{\kappa}},
\end{align*}
where we used \eqref{eq:HamiltonianLeakTrunc} in the last inequality.\end{proof}

\bin{In Section~\ref{sec:FiniteTruncBareJumps} we introduced the finite-dimensional truncations $a^{\le M}_i$ and $\left(a_i^{\le M}\right)^\dagger$ of the bare jumps $a$ and $a^\dagger$ and the generator $\mathcal{L}^{\le M}_{\sigma_E,\widehat{f},H}$ where we replaced the bare jumps in terms of the usual annihilation and creation operators in $\cL_{\sigma_E}$ with their finite-dimensional truncations.} 

Following the notation of Section~\ref{sec:FiniteTruncBareJumps}, we show in the following that under the assumptions \eqref{eq:ExpBound2} and \eqref{eq:HamiltonianLeakTrunc} the finite rank generator, $\mathcal{L}^{\le M}_{\sigma_E,\widehat{f},H_{\le M}},$ is close to $\mathcal{L}^{\le M}_{\sigma_E,\widehat{f},H}.$
For this, we first prove the following technical Lemma. In what follows, we denote for some function $f\in\mathcal{S}(\mathbb{R})$ the integrated jump operators 
\begin{align}
\label{eq:TruncatedL}
L^{\alpha,\le M}:=\int f(s)e^{isH}(A^{\alpha})^{\le M}e^{-isH}ds\quad\text{and}\quad
L^{\alpha,\le M}_{\le M}:=\int f(s)e^{isH_{\le M}}(A^{\alpha})^{\le M} e^{-isH_{\le M}}ds.
\end{align}
By construction, we have that
\begin{align}\label{eqLalphaMMbound}
\|L^{\alpha,\le M}_{\le M}\|\le \|f\|_{L^1(\mathbb{R})}\|(A^\alpha)^{\le M}\|\le \|f\|_{L^1(\mathbb{R})}\, q'(M),
\end{align}
where the polynomial $q'(M)$ is defined in \eqref{eq:NormTruncatedJump}.



\begin{lemma}\label{lem:integralrepapproxFullyFinite}
Let $f\in\mathcal{S}(\mathbb{R})$ and assume that \Cref{eq:condAalphas}, \eqref{eq:NormTruncatedJump} and \eqref{eq:NormTruncatedJump2} are satisfied for the set of bare jumps $\{A^\alpha\}_{\alpha\in\cA}$. Furthermore, for $\NA$ self-adjoint and positive semidefinite assume that \eqref{eq:HamiltonianLeakTrunc} and  \eqref{eq:ExpBound2} holds for some $\kappa\in(0,1/2)$ and $r\ge0$ and that
\begin{align*}
    C'_f := \int_{-\infty}^{\infty} |t|e^{r|t|^{2\kappa}} |f(t)|dt < \infty.
\end{align*}
Then for $M\in \N$ we have
\begin{align*}
\|(L^{\alpha,\le M}\!-\!L^{\alpha,\le M}_{\le M})e^{-\NA^{\kappa}}\|\!&\le\! C'_fp_1(|\cA|,M)e^{-M^{\kappa}},\\
\|((L^{\alpha,\le M})^\dagger L^{\alpha,\le M}-\left(L^{\alpha,\le M}_{\le M}\right)^\dagger L^{\alpha,\le M}_{\le M})e^{-\NA^{\kappa}}\|&\le C_fC'_fp_2(|\cA|,M)e^{-M^{\kappa}},\, 
\end{align*}
where $p_1(|\cA|,M)=2q'(M)p(|\cA|,M)$ and $p_2(|\cA|,M) =4(q'(M))^2p(|\cA|,M)$ with $q'(M)$ being the polynomially bounded function appearing in \eqref{eq:NormTruncatedJump} and \eqref{eq:NormTruncatedJump2} and $C_f$ being defined in \eqref{eq:C_f}.


\end{lemma}

\begin{proof}


For the first inequality in Lemma~\ref{lem:integralrepapproxFullyFinite} we use \eqref{eq:ExpBound2} and Lemma~\ref{lem:HamEvolTrunc} to see for $\ket{\psi}\in D\left(e^{\NA^{\kappa}}\right)$ that we have
\begin{align}
\label{eq:NLLeqFullyFinite}
\nn&\|(L^{\alpha,\le M}-L^{\alpha,\le M}_{\le M})\ket{\psi}\|
\\&\quad\le \int |f(s)|\Big(\left\|(A^{\alpha})^{\le M}\left(e^{-isH} -e^{-isH_{\le M}}\right)\ket{\psi}\right\|\nn+ \left\|\left(e^{isH} -e^{isH_{\le M}}\right)(A^{\alpha})^{\le M} e^{-isH}\ket{\psi}\right\|\Big)ds
\\&\quad\le 2q'(M)p(|\cA|,M)e^{-M^{\kappa}}
\int|f(s)| |s|e^{r|s|^{2\kappa}}ds\|e^{\NA^{\kappa}}\ket{\psi}\| = 2C'_fq'(M)p(|\cA|,M)e^{-M^{\kappa}}
\|e^{\NA^{\kappa}}\ket{\psi}\|,
\end{align}
where we used \eqref{eq:NormTruncatedJump} and \eqref{eq:NormTruncatedJump2} in the second inequality.



\bin{First, we use the estimates derived in Lemma \ref{lem.displaceblocks} to see for  $\ket{\psi}\in D(\Ntot+I)$, denoting $N_i:=a_i^\dagger a_i$, 
\begin{align*}
\|(\Ntot+I)\ket{\psi}\|^2&=\bra{\psi}\left(\sum_{j=1}^m N_j+I\right)^2\ket{\psi}\\
&\le m\sum_{j\in[m]}\bra{\psi}(N_j+I)^2\ket{\psi}\\
&=m\sum_{j\in[m]}\|(N_j+I)\ket{\psi}\|^2.
\end{align*}
Applying this bound to $\ket{\psi}=(a_i-(a^D_\alpha)_i)\ket{\varphi}$, we get
\begin{align*}
&\|(\Ntot+I)(a_i-(a^D_\alpha)_i)\ket{\varphi}\|^2\\
&\quad \le m\,\sum_{j\ne i}4|\alpha|^2\|(N_j+I)(N_i+I)\ket{\varphi}\|^2\\
&\quad + 25m |\alpha|^2\|(N_i+I)^2\ket{\varphi}\|^2\\
&\quad \le 25m^2|\alpha|^2\sum_{j\in[m]}\,\|(N_j+I)^2\ket{\varphi}\|^2\\
&\quad \le 25m^3|\alpha|^2\|(\Ntot+I)^2\ket{\varphi}\|^2.
\end{align*}



Thus, we have derived that
\begin{align}
\|(\Ntot\!+\!1)(L_{i}\!-\!L_{\alpha,i})\!\ket{\varphi}\!\|\!\le\! 5C_{2,f}m^{\frac{3}{2}}|\alpha|\|(\Ntot+I)^2\!\ket{\varphi}\!\|.
\end{align}}

\bin{Similarly, denoting $L_i^d:=\int f(s)e^{isH}a_i^\dagger e^{-isH}ds$ and $L_{\alpha,i}^d:=\int f(s)e^{isH}(a_{\alpha}^{Dd})_i e^{-isH}ds$, we get
\begin{align}\label{eq:NLLeqd}
\|(\Ntot\!+\!1)(L^d_{i}\!-\!L^d_{\alpha,i})\!\ket{\varphi}\!\|\!\le\! 5C_{2,f}m^{\frac{3}{2}}|\alpha|\|(\Ntot+I)^2\!\ket{\varphi}\!\|.
\end{align}}
Next, we control squares of the $L$ operators for $\ket{\varphi}\in D(e^{\NA^{\kappa}})$:
\begin{align}\label{eq:twotermstoboundTrunc2}
&\nn\left\|\left(\left(L^{\alpha,\le M}\right)^\dagger L^{\alpha,\le M}-\left(L^{\alpha,\le M}_{\le M}\right)^\dagger L^{\alpha,\le M}_{\le M}\right)\ket{\varphi}\right\|\\&\le \left\|\left(L^{\alpha,\le M}_{\le M}\right)^\dagger (L^{\alpha,\le M}-L^{\alpha,\le M}_{\le M})\ket{\varphi}\right\|+\left\|\left(\left(L^{\alpha,\le M}\right)^\dagger-\left(L^{\alpha,\le M}_{\le M}\right)^\dagger\right)L^{\alpha,\le M}\ket{\varphi}\right\|.
\end{align}
The first term above can be controlled by using that by \eqref{eq:NormTruncatedJump} we have $\left\|\left(L_{\le M}^{\alpha,\le M}\right)^\dagger\right\| \le q'(M)\|f\|_1\le q'(M)C_f$ and therefore using \eqref{eq:NLLeqFullyFinite} we have
\begin{align*}&\left\|\left(L^{\alpha,\le M}_{\le M}\right)^\dagger (L^{\alpha,\le M}-L^{\alpha,\le M}_{\le M})\ket{\varphi}\right\| \le 2C_fC'_f (q'(M))^2p(|\cA|,M)e^{-M^{\kappa}}
\|e^{\NA^{\kappa}}\ket{\psi}\|.
\end{align*}
To control the second term in \eqref{eq:twotermstoboundTrunc2}, we note that by \eqref{eq:ExpBound2} and \eqref{eq:NormTruncatedJump2} we have 
\begin{align}
\label{eq:ThereAreLoadsOfInequalitiesHere}
\nn&\left\|e^{\NA^{\kappa}}L^{\alpha,\le M} e^{-\NA^{\kappa}}\right\| \le  \int |f(s)|e^{r|s|^{2\kappa}/2} \left\|e^{\NA^{\kappa}}e^{isH}(A^\alpha)^{\le M}e^{-\NA^{\kappa}}\right\|ds\,\\ &\le q'(M)\int |f(s)| e^{r|s|^{2\kappa}} ds = q'(M)C_f,
\end{align}
where $C_f$ is defined in \eqref{eq:C_f}. From this 
and \eqref{eq:NLLeqTrunc} we bound the second term in \eqref{eq:twotermstoboundTrunc2} as 
\begin{align*}
    &\left\|\left(\left(L^{\alpha,\le M}\right)^\dagger-\left(L^{\alpha,\le M}_{\le M}\right)^\dagger\right)L^{\alpha,\le M}\ket{\varphi}\right\| \le 2 C_fC'_f (q'(M))^2p(|\cA|,M)e^{-M^{\kappa}} \left\|e^{\NA^{\kappa}}\ket{\varphi}\right\|,
\end{align*}
which finishes the proof.
\end{proof}
\noindent  We are now ready to show in the following proposition that $\mathcal{L}^{\le M}_{\sigma_E,\widehat{f},H}$ can be well approximated by the finite rank generator $\mathcal{L}^{\le M}_{\sigma_E,\widehat{f},H_{\le M}}.$ 
\begin{proposition}\label{prop:generatorboungFullyFinite}
Let $f\in\mathcal{S}(\mathbb{R})$ and assume that \Cref{eq:condAalphas}, \eqref{eq:NormTruncatedJump} and \eqref{eq:NormTruncatedJump2} are satisfied for the set of bare jumps $\{A^\alpha\}_{\alpha\in\cA}$. Furthermore, for $\NA$ self-adjoint and positive semidefinite assume that \eqref{eq:HamiltonianLeakTrunc} and  \eqref{eq:ExpBound2} holds for some $\kappa\in(0,1/2)$ and $r\ge0$ and that
\begin{align*}
    C'_f := \int_{-\infty}^{\infty} |t|e^{r|t|^{2\kappa}} |f(t)|dt < \infty.
\end{align*}
Then, for all $\beta,\,\sigma_E>0,$ $M\in\N$ and state $\rho$ satisfying $E_2:=\Tr(e^{2\NA^{\kappa}}\rho)<\infty$, we have
\begin{align}
&\nn\left\|(\mathcal{L}^{\le M}_{\sigma_E,\widehat{f},H}-\mathcal{L}^{\le M}_{\sigma_E,\widehat{f},H_{\le M}})(\rho)\right\|_1\le  C_{\beta,r,\kappa} \widetilde C_f \sqrt{E_2}\,p_3(|\cA|,M)e^{-M^{\kappa}}\nn\left(\sigma_E
\exp\!\left(\tfrac{\pi^2}{8\beta^2\sigma_E^2}\right)+\exp\left(r+\tfrac{(r+1)^2}{8\sigma^2_E}\right)\right),
\end{align}
for some constant $C_{\beta,r,\kappa}\ge 0$ and 
where $p_3(|\cA|,M):=(q'(M))^2|\cA|\,p(|\cA|,M)$ is  polynomially bounded, $\widetilde C_f := C_f\max\{C_f,C'_f\}$ and $C_f$ is defined in \eqref{eq:C_f}.
\end{proposition}
\begin{proof}
We make use of the integral representation of Proposition \ref{integralrep}. First, we denote the operators
\begin{align*}
G^{\le M}_{\sigma_E}\!&:=\!-\sum_{\alpha\in\cA}\!\int_{-\infty}^\infty\! \!\!\!\!g(t)e^{itH} (L^{\alpha,\le M})^\dagger L^{\alpha,\le M}e^{-itH}\!dt\\G^{\le M}_{\sigma_E,\le M}\!&:=\!-\sum_{\alpha\in\cA}\!\int_{-\infty}^\infty\! \!\!\!\!g(t)e^{itH_{\le M}} (L_{\le M}^{\alpha,\le M})^\dagger L_{\le M}^{\alpha,\le M}e^{-itH_{\le M}}\!dt,
\end{align*}
    where $G^{\le M}_{\sigma_E}$ appears in the definition of $\mathcal{L}^{\le M}_{\sigma_E,\widehat{f},H}$ and $G^{\le M}_{\sigma_E,\le M}$ appears in the definition of $\mathcal{L}^{\le M}_{\sigma_E,\widehat{f},H_{\le M}}.$
\bin{Following exactly the steps as in the proof of Proposition~\ref{prop:generatorboundTrunc} while utilising Lemma~\ref{lem:integralrepapproxFullyFinite} instead of Lemma~\ref{lem:integralrepapproxTrunc}, we see 
\begin{align}
    \nn&\left\|\left(G^{\le M}_{\sigma_E}+\sum_{\alpha\in\cA}\!\int_{-\infty}^\infty\! \!\!\!\!g(t)e^{-itH} (L_{\le M}^{\alpha,\le M})^\dagger L_{\le M}^{\alpha,\le M}e^{itH}\!dt\right)\rho\right\|_1 \\&\le C_{\beta,r,\kappa}C_fC'_f\,\sigma_E
\exp\!\left(\tfrac{\pi^2}{8\beta^2\sigma_E^2}\right)|\cA|p_2(|\cA|,M)e^{-M^{\kappa}}\sqrt{E_4},
\end{align}
for some constant $C_{\beta,r,\kappa}\ge 0.$
On the other hand, we see
\begin{align*}
  \nn&\left\|\left(G^{\le M}_{\sigma_E,\le M}+\sum_{\alpha\in\cA}\!\int_{-\infty}^\infty\! \!\!\!\!g(t)e^{-itH} (L_{\le M}^{\alpha,\le M})^\dagger L_{\le M}^{\alpha,\le M}e^{itH}\!dt\right)\rho\right\|_1  \\&\nn\le \sum_{\alpha\in\cA}\int |g(t)| \left\|\left(e^{-itH}-e^{-itH_{\le M}}\right) (L_{\le M}^{\alpha,\le M})^\dagger L_{\le M}^{\alpha,\le M} e^{itH}dt\rho \right\|_1\\&\qquad+\sum_{\alpha\in\cA}\int |g(t)| \left\|(L_{\le M}^{\alpha,\le M})^\dagger L_{\le M}^{\alpha,\le M} \left(e^{itH}-e^{itH_{\le M}}\right) dt\rho\right\|_1
\end{align*}}
We see
\begin{align}
\label{eq:MoreGApproxiii}
\nn\left\|\left(G^{\le M}_{\sigma_E}-G^{\le M}_{\sigma_E,\le M}\right)\rho\right\|_1 &\le \sum_{\alpha\in\cA }\int |g(t)| \left\|\left(e^{itH}-e^{itH_{\le M}}\right) (L^{\alpha,\le M})^\dagger L^{\alpha,\le M} e^{-itH}dt\,\rho \right\|_1\\&\quad\nn+\sum_{\alpha\in\cA }\int |g(t)| \left\|\left((L^{\alpha,\le M})^\dagger L^{\alpha,\le M} -(L^{\alpha,\le M}_{\le M})^\dagger L^{\alpha,\le M}_{\le M}\right)e^{-itH}dt\,\rho \right\|_1\\&\quad+\sum_{\alpha\in\cA }\int |g(t)| \left\|(L^{\alpha,\le M}_{\le M})^\dagger L^{\alpha,\le M}_{\le M}\left(e^{-itH}-e^{-itH_{\le M}}\right)dt\,\rho \right\|_1.
\end{align}
For the first term we note using \eqref{eq:ExpBound2} and \eqref{eq:NormTruncatedJump2} 
 together with an analogous argument as for \eqref{eq:ThereAreLoadsOfInequalitiesHere} that 
\begin{align*}
    \left\|e^{\NA^{\kappa}}(L^{\alpha,\le M})^\dagger L^{\alpha,\le M}e^{-\NA^{\kappa}}\right\| \le  \left\|e^{\NA^{\kappa}}(L^{\alpha,\le M})^\dagger e^{-\NA^{\kappa}}\right\|\left\|e^{\NA^{\kappa}} L^{\alpha,\le M}e^{-\NA^{\kappa}}\right\| \le C^2_f (q'(M))^2 
\end{align*}
and therefore using Lemma~\ref{lem:HamEvolTrunc}
\begin{align*}
&\sum_{\alpha\in\cA }\int |g(t)| \left\|\left(e^{itH}-e^{itH_{\le M}}\right) (L^{\alpha,\le M})^\dagger L^{\alpha,\le M} e^{-itH}dt\,\rho \right\|_1\\&\le C^2_f|\cA|(q'(M))^2\sqrt{E_2} \int |g(t)| e^{r|t|^{2\kappa}/2} \left\|\left(e^{itH}-e^{itH_{\le M}}\right)e^{-\NA^{\kappa}}\right\| dt\\&\le C^2_f|\cA|(q'(M))^2\sqrt{E_2} \int |g(t)| |t|e^{r|t|^{2\kappa}}dt\\& \le C_{\beta,r,\kappa}C^2_f\sqrt{E_2} \,
\sigma_E\,\,
\exp\!\left(\frac{\pi^2}{8\beta^2\sigma_E^2}\right) |\cA|(q'(M))^2p(|\cA|,M)e^{-M^{\kappa}}
\end{align*}
for some constant $C_{\beta,r,\kappa}\ge0$ and
where in the last inequality we used \eqref{eq:gkappaIntegralBound} with a larger value of $r$ in there. The third term in \eqref{eq:MoreGApproxiii} can be estimated analogously and the second as well utilising additionally Lemma~\ref{lem:integralrepapproxFullyFinite} which in total gives
\begin{align*}
    \left\|\left(G^{\le M}_{\sigma_E}-G^{\le M}_{\sigma_E,\le M}\right)\rho\right\|_1\le C_{\beta,r,\kappa}\widetilde C_f\sqrt{E_2} \,
\sigma_E\,\,
\exp\!\left(\frac{\pi^2}{8\beta^2\sigma_E^2}\right) |\cA|p_2(|\cA|,M)e^{-M^{\kappa}},
\end{align*}
where we increased the constant $C_{\beta,r,\kappa}$ by a factor 3 compared to the above and, furthermore, denoted $\widetilde C_f := C_f\max\{C_f,C'_f\}.$

Next, for 
\begin{align*}
 X^{\alpha,\le M}_s:=e^{isH}L^{\alpha,\le M} e^{-isH}
,\quad  X^{\alpha,\le M}_{s,\le M}:=e^{isH_{\le M}}L^{\alpha,\le M}_{\le M} e^{-is{H_{\le M}}}
\end{align*}
using that by \eqref{eq:NormTruncatedJump} we have $\|X^{\alpha,\le M}_s\|,\, \|X^{\alpha,\le M}_s\|\le \|f\|_1 q'(M)\le C_f q'(M),$
we see 
\begin{align*}
    &\left\| X^{\alpha,\le M}_s\rho X^{\alpha,\le M}_s- X^{\alpha,\le M}_{s,\le M}\rho X^{\alpha,\le M}_{s,\le M}\right\|_1 \\&\le C_f\,q'(M)\left( \left\|\left(X^{\alpha,\le M}_s-  X^{\alpha,\le M}_{s,\le M}\right)\rho\right\|_1+\left\|\rho\left(X^{\alpha,\le M}_s-  X^{\alpha,\le M}_{s,\le M}\right)\right\|_1\right).
\end{align*}
We focus on the term as the second term can be estimated the same way:
Using \eqref{eq:NormTruncatedJump}
and Lemma~\ref{lem:HamEvolTrunc} we see
\begin{align*}
    \left\|\left(X^{\alpha,\le M}_s-  X^{\alpha,\le M}_{s,\le M}\right)\rho\right\|_1 &\le \int |f(t)|\Big(\left\|\left(e^{i(t+s)H}-e^{i(t+s)H_{\le M}}\right)(A^{\alpha})^{\le M} e^{-i(t+s)H}\rho\right\|_1 \\&\qquad+\left\|e^{i(t+s)H_{\le M}}(A^{\alpha})^{\le M} \left(e^{-i(t+s)H}-e^{-i(t+s)H_{\le M}}\right)\rho\right\|_1 \Big)dt\\&\le 2 q'(M)\sqrt{E_2}\,p(|\cA|,M)e^{-M^{\kappa}}\int |f(t)| |t+s| e^{r|t+s|^{2\kappa}}dt \\&\le2\max\{C_f,C'_f\} \sqrt{E_2}\, q'(M) p(|\cA|,M)e^{-M^\kappa} (1+|s|)e^{r|s|^{2\kappa}}.
\end{align*}
Using this and denoting the map
\begin{align}\label{phialphasigmaEmapTruncDouble}
\Phi^{\le M}_{\sigma_E,\widehat f,H_{\le M}}(\rho):=\sigma_E\sqrt{\frac{2}{\pi}}\sum_{\alpha\in\cA}\int_{\mathbb{R}}e^{-2\sigma_E^2 s^2}X^{\alpha,\le M}_{s,\le M} \rho (X^{\alpha,\le M}_{s,\le M})^\dagger ds,
\end{align}
we see that
\begin{align*}
\left\|\left(\Phi^{\le M}_{\sigma_E,\widehat f,H}-\Phi^{\le M}_{\sigma_E,\widehat f,H_{\le M}}\right)(\rho)\right\|_1\lesssim  \,\widetilde C_f\,\sqrt{E_2} |\cA|(q'(M))^2p(|\cA|,M)e^{-M^{\kappa}} \exp\left(r+\tfrac{(r+1)^2}{8\sigma^2_E}\right)  , 
\end{align*}
where we used that  
\begin{align*}
\,\sigma_E\int e^{-2\sigma^2_Es^2}(1+|s|) e^{r|s|^{2\kappa}}ds&\le  e^{r} \,\sigma_E\int e^{-2\sigma^2_Es^2}(1+|s|) e^{r|s|}ds\\&\lesssim e^{r} \,\sigma_E\int e^{-2\sigma^2_Es^2} e^{(1+r)|s|}ds \lesssim \exp\left(r+\tfrac{(r+1)^2}{8\sigma^2_E}\right)   
\end{align*}
and where we only hide constants independent of all parameters with the $\lesssim$-notation.


\end{proof}

\subsubsection{Finite-dimensional Lindblad dynamics for Schwartz filter function}
\label{sec:FiniteDimPreperationSchwartz}

In this section we combine the results of the previous two sections, in particular Propositions~\ref{prop:generatorboundTrunc} and~\ref{prop:generatorboungFullyFinite}, to see in Theorem~\ref{thm:SchwartzGibbsDynamics} that the dynamics $e^{t\cL_{\sigma_E,\widehat f,H}}$ for Schwartz filter function is well approximated by $e^{t\cL^{\le M}_{\sigma_E,\widehat f,H_{\le M}}}.$ 
Moreover,  under the assumption of spectral gap of the corresponding unbounded generator on the Hilbert-Schmidt space,  we then apply Theorem~\ref{thm:SchwartzGibbsDynamics} to show efficient finite-dimensional preparation of the Gibbs state of $H$ in Corollary~\ref{cor:SchwartzGibbs}. 

We first state both results and then give their proofs at the end of the section.

\begin{theorem}[finite-dimensional approximation of $e^{t\cL_{\sigma_E,\widehat f,H}}$]
\label{thm:SchwartzGibbsDynamics}
Let $f\in\mathcal{S}(\mathbb{R})$ and set of bare jumps $\{A^\alpha\}_{\alpha\in\cA}$ and assume  \Cref{eq:condAalphas} is satisfied. 
Let further $\NA$ self-adjoint and positive semidefinite and $\kappa\in(0,1/2)$ and assume that \eqref{eq:NormTruncatedJump}, \eqref{eq:AbstractFiniteTruncBareJumps}, \eqref{eq:NormTruncatedJump2} are  satisfied and that $C_A := \max_{\alpha\in\cA}\left(\|e^{\NA^{\kappa}}A^{\alpha}e^{-2\NA^{\kappa}}\|+\|A^{\alpha}e^{-\NA^{\kappa}}\|+\|e^{-\NA^{\kappa}}A^{\alpha}\|\right)<\infty.$ 
Furthermore, assume that for $\kappa\in(0,1/2)$ as above \eqref{eq:HamiltonianLeakTrunc} holds and, for some $r\ge 0$ the condition \eqref{eq:ExpBound} is satisfied and that $\int_{-\infty}^{\infty} |t|e^{r|t|^{2\kappa}} |f(t)|dt <\infty.$ 
For $\beta>0$ we assume that the Gibbs state of $H,$ $\sigma_\beta,$ satisfies
\begin{align}E_{\operatorname{Gibbs}}:=\Tr\left(e^{4\NA^\kappa}\, \sigma_\beta\right)<\infty
\end{align}
and consider state $\rho$ such that 
\begin{align}
\label{eq:InputDmaxConstraint}
    \rho \le \mathfrak{c}\,\sigma_\beta
\end{align}
for some $\mathfrak{c}\ge 1.$ Then for all $\sigma_E\in(0,\infty),$ $t\ge 0$ and $M\in\N$ we have
    \begin{align}
\left\|\left(e^{t\cL_{\sigma_E,\widehat f,H}}-e^{t\cL^{\le M}_{\sigma_E,\widehat f,H_{\le M}}}\right)(\rho)\right\|_1\le t\,C_1(\beta,r,\kappa,\sigma_E,f,C_A)\,\mathfrak{c} \,p_4(|\mathcal{A}|,M)E_{\operatorname{Gibbs}}\, e^{-M^\kappa} , 
    \end{align}
    where $C_1(\beta,r,\kappa,\sigma_E,f,C_A)\ge 0$ is some constant depending only on the displayed parameters and $p_4(|\mathcal{A}|,M)$ is some polynomial bounded function. Therefore, for $\eps>0,$ we can achieve
    \begin{align}\label{eq:Metal}
       &\left\|\left(e^{t\cL_{\sigma_E,\widehat f,H}}-e^{t\cL^{\le M}_{\sigma_E,\widehat f,H_{\le M}}}\right)(\rho)\right\|_1\le \eps,\qquad\text{with}\qquad M = \widetilde{\mathcal{O}} \left(\left(\log\left(\frac{t\,\mathfrak{c}\,E_{\operatorname{Gibbs}}|\cA|}{ \eps}\right)\right)^{1/\kappa}\right),
    \end{align}
    where the $\widetilde{\mathcal{O}}$ notation hides constants independent of the displayed parameters and additionally suppresses subdominant $\operatorname{poly}\log\log$ factors.
\end{theorem}
\begin{corollary}[Gibbs state preparation for Schwartz filter function]
\label{cor:SchwartzGibbs}
Under the same assumptions as in Theorem~\ref{thm:SchwartzGibbsDynamics} and assuming additionally that  $L_{\sigma_E,\widehat{f},H},$ the self-adjoint generator on $\mathscr{T}_2(\cH)$ associated to the Lindbladian $\cL_{\sigma_E,\widehat f,H},$ has a positive spectral gap 
$\lambda_2 \equiv \operatorname{gap}\left(L_{\sigma_E,\widehat f,H}\right) >0,$  
we can  achieve for all $\eps>0$
\begin{align}
\nn\left\|e^{t\cL^{\le M}_{\sigma_E,\widehat f,H_{\le M}}}(\rho) - \sigma_\beta\right\|_1 \le \eps\quad \text{with}\quad
t=\mathcal{O}\left(\frac{1}{\lambda_2}\log\left(\frac{\mathfrak{c}}{\eps}\right)\right) \quad\text{and}\quad M = \widetilde{\mathcal{O} }\left(\left(\log\left(\frac{\mathfrak{c}\,E_{\operatorname{Gibbs}}|\cA|}{\lambda_2 \eps}\right)\right)^{1/\kappa}\right).
\end{align}
where the $\mathcal{O}$ and $\widetilde{\mathcal{O}}$ notations hide constants independent of the displayed parameters and $\widetilde{\mathcal{O}}$ additionally suppresses subdominant $\operatorname{poly}\log\log$ factors.
\end{corollary}

\begin{proof}[Proof of Theorem~\ref{thm:SchwartzGibbsDynamics}]
We use that 
\begin{align*}
    \left(e^{t\cL^{\le M}_{\sigma_E,\widehat f,H_{\le M}}} -e^{t\cL_{\sigma_E,\widehat f,H}}\right)(\rho) = \int_0^t e^{(t-s)\cL^{\le M}_{\sigma_E,\widehat f,H_{\le M}}}\left(\cL^{\le M}_{\sigma_E,\widehat f,H_{\le M}} -\cL_{\sigma_E,\widehat f,H}\right)e^{s\cL_{\sigma_E,\widehat f,H}}(\rho)ds
\end{align*}
and therefore, denoting $\rho(s) :=e^{s\cL_{\sigma_E,\widehat f,H}}(\rho),$ using Proposition~\ref{prop:generatorboundTrunc} and~\ref{prop:generatorboungFullyFinite} we see
\begin{align*}
    &\left\|\left(e^{t\cL^{\le M}_{\sigma_E,\widehat f,H_{\le M}}} -e^{t\cL_{\sigma_E,\widehat f,H}}\right)(\rho)\right\|_1 \\
    & \qquad \le   t\sup_{s\in[0,t]}\left\|\left(\cL^{\le M}_{\sigma_E,\widehat f,H_{\le M}}-\cL_{\sigma_E,\widehat f,H}\right)(\rho(s))\right\|_1 \\
 &    \qquad\le t\sup_{s\in[0,t]}\left\|\left(\cL^{\le M}_{\sigma_E,\widehat f,H_{\le M}}-\cL^{\le M}_{\sigma_E,\widehat f,H}\right)(\rho(s))\right\|_1 +t\sup_{s\in[0,t]}\left\|\left(\cL^{\le M}_{\sigma_E,\widehat f,H}-\cL_{\sigma_E,\widehat f,H}\right)(\rho(s))\right\|_1\\
    &\qquad \le t\,C_1(\beta,r,\kappa,\sigma_E,f,C_A) p_4(|\cA|,M)e^{-M^\kappa} \sup_{s\in[0,t]}\max\left\{\sqrt{\Tr\left(e^{4\NA^\kappa}\rho(s)\right)}\,,\, \Tr\left(e^{2\NA^\kappa}\rho(s)\right) \right\}\\
   &  \qquad\le t\,\mathfrak{c}\,C_1(\beta,r,\kappa,\sigma_E,f,C_A) p_4(|\mathcal{A}|,M) e^{-M^\kappa} E_{\operatorname{Gibbs}},
\end{align*}
where $C_1(\beta,r,\kappa,\sigma_E,f,C_A)\ge 0$ is some constant depending only on the displayed parameters, $p_4(|\mathcal{A}|,M)$ is some polynomially bounded function and where we used in the last inequality that by \eqref{eq:InputDmaxConstraint}, positivity of the map $e^{s\cL_{\sigma_E,\widehat f,H}}$ and the fact that $\cL_{\sigma_E,\widehat f,H}$ is KMS-symmetric we have
\begin{align*}
   \rho(s) \le \mathfrak{c}\, e^{s\cL_{\sigma_E,\widehat f,H}}(\sigma_\beta ) = \mathfrak{c}\,\sigma_\beta. 
\end{align*}
\end{proof}

\begin{proof}[Proof of Corollary~\ref{cor:SchwartzGibbs}]
We split
\begin{align}
\label{eq:MajorTriangle}
   \left\|e^{t\cL^{\le M}_{\sigma_E,\widehat f,H_{\le M}}}(\rho) - \sigma_\beta\right\|_1  &\le \left\|e^{t\cL_{\sigma_E,\widehat f,H}}(\rho) - \sigma_\beta\right\|_1+\left\|\left(e^{t\cL^{\le M}_{\sigma_E,\widehat f,H_{\le M}}} -e^{t\cL_{\sigma_E,\widehat f,H}}\right)(\rho)\right\|_1  
\end{align}
For the first term, we argue as in Section~\ref{sec.fastconv} using positivity of spectral gap of $L_{\sigma_E,\widehat f,H}$ together with \eqref{eq:InputDmaxConstraint} to get
\begin{align*} \left\|e^{t\cL_{\sigma_E,\widehat{f},H}}(\rho)- \sigma_\beta\right\|_1 &\le  e^{-\lambda_2 t} \left\|\sigma_\beta^{-1/4}\rho\sigma^{-1/4}_\beta\right\|_2 = e^{-\lambda_2 t} \sqrt{\Tr\left(\sigma^{-1/2}_\beta\rho\sigma^{-1/2}_\beta\rho\right)} \\&\le \sqrt{\mathfrak{c}}\,e^{-\lambda_2 t} \sqrt{\Tr(\rho)}= \sqrt{\mathfrak{c}}\,e^{-\lambda_2 t},
\end{align*}
which gives the desired bound on the mixing time $t.$ The result then follows by Theorem~\ref{thm:SchwartzGibbsDynamics}.
\end{proof}




\subsection{Smoothly approximating generators with singular filter functions}\label{sec:singular}


\bin{Up to this point, we have developed a framework for implementing generators of Gibbs samplers associated with Schwartz functions $f$. However, in infinite dimensions, such generators typically fail to exhibit a spectral gap. In this section, we extend the framework by developing a perturbative approach that enables us to construct generators corresponding to singular potentials, in analogy with~\eqref{eq:Function}.}


The previous sections on efficient implementation of the Gibbs samplers, in particular Section~\ref{integralrepfschwartz}, heavily relied on strong regularity and fast decay of the filter function $f(t).$ More precisely, throughout all of these sections we assumed that $f,$ and hence also its Fourier transform $\widehat f,$ are Schwartz functions. This enabled us, among other things, to provide integral formulations of the generator $\cL_{\sigma_E,\widehat{f},H}$ in Section~\ref{integralrepfschwartz}, which we then utilised for the proceeding finite-dimensional approximation steps in Section~\ref{sec:FiniteTruncBareJumps}.

On the other hand, in Section~\ref{sec:AbsenceOfGap}, we have seen that fast decay of $\widehat f$ leads to the absence of spectral gap and, hence, no mixing time guarantees of the Gibbs samplers for Hamiltonians of the form $H=h(N)$ with superlinear functions $h$. In Section~\ref{sec:GeneralNumberviaBirthDeath}, we have solved this issue by proving that  for such Hamiltonians we get a positive spectral gap if we take instead the  Metropolis-type filter function, which in Fourier space is given by 
\begin{align}
\label{eq:Linsfunction}
\hatfM(\nu) = \exp\left(-\frac{\sqrt{1+(\beta\nu)^2} +\beta\nu}{4}\right)
.\end{align}
Crucially, $\hatfM$ does not decay for large negative Bohr frequencies as $\lim_{\nu\to-\infty}\hatfM(\nu)=1,$ enabling the existence of a positive spectral gap of the corresponding Gibbs sampler. On the flip side, we, however, clearly have that $
\hatfM(\nu),$ and therefore also $\fM(t),$ cannot be Schwartz functions. In fact $\fM$ is only defined as tempered distribution which formally diverges as $t\to 0.$ 

To bridge this gap between Section~\ref{sec:GeneralNumberviaBirthDeath}, which provides convergence guarantees of the Gibbs samplers, with Section~\ref{sec:implementations}, which provides efficient implementation, we establish in Proposition~\ref{prop:DeltaTo0} that generators relying on the filter function \eqref{eq:Linsfunction} can be approximated by generators with Schwartz filter functions. We then apply this result in Section~\ref{sec:FiniteDimPreperationSingular} to show, analogously as in Theorem~\ref{thm:SchwartzGibbsDynamics}, that the dynamics $e^{t\cL_{\sigma_E,\hatfM,H}}$ is well-approximated by a certain finite-dimensional Lindblad dynamics.

For that we consider in the following  a smooth function 
\begin{align*}
\R_{\ge 0}\times \R &\to [0,1]\quad\text{ with }\quad
(\delta,\nu) \mapsto \widehat \phi_\delta(\nu),
\end{align*}
such that $\widehat\phi_0(\nu) = 1$ for all $\nu\in\R$ and further that $\delta\mapsto \widehat\phi_\delta(\nu)$ is non-increasing and the partial derivative in the first variable, $\delta\mapsto \partial_\delta\widehat \phi_\delta(\nu),$ is non-decreasing. With that we define
\begin{align} 
\label{eq:Defeta}
\eta(\nu):=-\partial_\delta\Big|_{\delta=0}\widehat\phi_\delta(\nu) = -\inf_{\delta\ge 0}\partial_\delta\widehat\phi_\delta(\nu)\ge 0.
\end{align}
Furthermore, we define
\begin{align}
\label{eq:Deltahatf}
    \hatfMdelta(\nu) := \hatfM(\nu) \widehat \phi_\delta(\nu),  
\end{align}
which for $\delta>0$ and suitable choice $\widehat \phi_\delta(\nu)$ turns out to be  a Schwartz function as desired as we see below.

We consider in the following Hamiltonians  $H$ satisfying $H\ge -h_0$ for some $h_0\ge 0$ and having discrete spectrum $\spec(H)$ with corresponding spectral projections being denoted by $P_{E}$ for $E\in\spec(H)$. Further, define the functions
\begin{align}
\label{eq:WeirdG}
   \nn F_\eta(E) &:= \sum_{E'\in\spec(H)}  \eta(E'-E)|\hatfM(E'-E)|,\qquad\ F_1(E) := \sum_{E'\in\spec(H)} |\hatfM(E'-E)|,\\
   F_{\eta,\sigma_E,1}(E) &:=\sum_{E'\in \spec(H)} e^{-\frac{(E'-E)^2}{8\sigma^2_E}} F_{\eta}(E'),\qquad F_{\eta,\sigma_E,2}(E) := F_{\eta}(E)\sum_{E'\in \spec(H)} e^{-\frac{(E'-E)^2}{8\sigma^2_E}}.
\end{align}
for $E\in\spec(H).$ In the following proposition, we see that the generator $\cL_{\sigma_E,\hatfM,H}$ defined in \eqref{eq:DefLsigma_E}, i.e.
\begin{align*}
\mathcal{L}_{\sigma_E,\hatfM,H}(\rho)\!&:= \!\sum_{\alpha\in\cA}\sum_{\nu_1,\nu_2\in B(H)}e^{-\frac{(\nu_1-\nu_2)^2}{8\sigma_E^2}}\,\overline{\hatfM(\nu_1)}\hatfM(\nu_2)A^\alpha_{\nu_2}\rho (A^\alpha_{\nu_1})^\dagger+G_{\sigma_E}\rho+\rho G_{\sigma_E}^\dagger\,,\quad\text{where}\\
G_{\sigma_E}&:= \sum_{\substack{\alpha\in\mathcal{A}\\ \nu_1,\nu_2\in B(H)}}\!e^{-\frac{(\nu_1-\nu_2)^2}{8\sigma_E^2}}\!\!\! \frac{1}{1+e^{\frac{\beta(\nu_2-\nu_1)}{2}}}\overline{\hatfM(\nu_1)}\hatfM(\nu_2)(A^\alpha_{\nu_1})^\dagger A^\alpha_{\nu_2},
\end{align*}
and $\hatfM$ being the Metropolis type filter function \eqref{eq:Linsfunction}, is close to $\cL_{\sigma_E,\hatfMdelta,H}$ when evaluated on input states satisfying certain energy constraints depending on the growth of the functions \eqref{eq:WeirdG}.
\begin{proposition}
\label{prop:DeltaTo0}
Assume \Cref{eq:condAalphas} for $\mu\ge \gamma\ge 0$. Let $\sigma_E\in(0,\infty)$ and $H$ be such that $F_\eta(E),\, F_{\eta,\sigma_E}(E)<\infty$ for all $E\in\spec(H)$, and define the self-adjoint operators
$\widetilde F_\eta(H) = \sum_{E\in\spec(H)}  \widetilde F_\eta(E) P_{E},$  $\widetilde F_\eta(H) = \sum_{E\in\spec(H)}  \widetilde F_\eta(E) P_{E}$ and, for $i=1,2,$ $\widetilde F_{\eta,\sigma_E,i}(H) = \sum_{E\in\spec(H)}  \widetilde F_{\eta,\sigma_E,i}(E) P_{E}$ for some functions $\widetilde F_\eta(E),\, \widetilde F_1(E),\widetilde F_{\eta,\sigma_E,i}(E) > 0$
satisfying
\begin{align*}
    \sum_{E\in \spec(H)}\frac{F_\eta(E)}{ \widetilde F_\eta(E)}<\infty,\qquad  \sum_{E\in \spec(H)}\frac{F_1(E)}{ \widetilde F_1(E)}<\infty\qquad\text{and}\qquad    \sum_{E\in \spec(H)}\frac{F_{\eta,\sigma_E,i}(E)}{ \widetilde F_{\eta,\sigma_E,i}(E)}<\infty.
\end{align*}
Then for state $\rho$ we have 
\begin{align}
\label{eq:Deltato0}
&\left\|\left( \mathcal{L}_{\sigma_E,\hatfM,H} - \mathcal{L}_{\sigma_E,\hatfMdelta,H}\right)(\rho) \right\|_1\nn
\\&\le C_{H,\eta,\sigma_E}\,\delta 
    \sum_{\alpha\in\cA} \Big(\|A^\alpha \widetilde  H^{-\gamma}\|^2 \,
    \| \widetilde F_1(H) \widetilde H^{\gamma} \rho \widetilde H^{\gamma}\widetilde F_\eta(H) \|_1\nn
   \\&\qquad\qquad\qquad\qquad+ \|\widetilde H^{-\gamma}A^\alpha\|
      \|\widetilde H^{\gamma}A^\alpha \widetilde H^{-\mu}\|\,
      \left(\|\widetilde F_{\eta,\sigma_E,1}(H) \widetilde H^{\mu}\rho\|_1 +\|\widetilde F_{\eta,\sigma_E,2}(H) \widetilde H^{\mu}\rho\|_1\right)  
\Big) \nn\\& \le C_{H,\eta,\sigma_E}\,\delta 
    \sum_{\alpha\in\cA} \Big(\|A^\alpha \widetilde  H^{-\gamma}\|^2 \,
   \sqrt{ \Tr\left( \widetilde F^2_1(H) \widetilde H^{2\gamma}\rho\right)  \Tr\left( \widetilde F^2_\eta(H) \widetilde H^{2\gamma}\rho\right)} \nn
   \\&\qquad\qquad\qquad\qquad+ \|\widetilde H^{-\gamma}A^\alpha\|
      \|\widetilde H^{\gamma}A^\alpha \widetilde H^{-\mu}\|\,
      \sqrt{\Tr\left(\widetilde F^2_{\eta,\sigma_E,1}(H) \widetilde H^{2\mu}\rho\right)} +\sqrt{\Tr\left(\widetilde F^2_{\eta,\sigma_E,2}(H) \widetilde H^{2\mu}\rho\right)}  
\Big)
\end{align}
where $C_{H,\eta,\sigma_E}:=\max_{i=1,2}\left\{\left(\sum_{E\in\spec(H)} \tfrac{F_\eta(E)}{ \widetilde F_\eta(E)} \right)\left(\sum_{E''\in\spec(H)}\tfrac{F_1(E)}{\widetilde F_1(E)}\right)\,,\ 2\sum_{\substack{E\in\spec(H)}}\!\tfrac{F_{\eta,\sigma_E,i}(E)}{\widetilde F_{\eta,\sigma_E,i}(E)}\right\}.$ \bin{Furthermore, for $H$ having simple spectrum, the estimate \eqref{eq:Deltato0} can be improved by replacing all trace norms involving $\rho$ with operator norms.}
\end{proposition}
\begin{remark}
Proposition \ref{prop:DeltaTo0} can be extended to the case $\sigma_E=\infty$ by using a slightly different strategy to bound the difference between the $G_{\sigma_E}$ terms: In particular in the current proof, we upper bound the Fermi weight $(1+e^{\frac{\beta(\nu_2-\nu_1)}{2}})^{-1}$ by 1 and explicitly use the Gaussian $e^{-\frac{(\nu_1-\nu_2)^2}{8\sigma_E}}$ to bound one of the appearing sum over energies of $H.$ In the case $\sigma_E =\infty$ we would, however, keep the Fermi weight and bound the respective sum over energies using this. 

As in this paper we usually take $\sigma_E \in(0,\infty)$ as a fixed constant, we only explicitly state and proof the current version of Proposition \ref{prop:DeltaTo0} which applies in this case.
\end{remark}
Before giving the proof of Proposition~\ref{prop:DeltaTo0}, we discuss in the following lemma and remark the growth of the functions in \eqref{eq:WeirdG} and therefore the required energy constraints on the input state in Proposition~\ref{prop:DeltaTo0}.  We then continue to discuss specific choices $\widehat \phi_\delta(\nu)$ and $\eta(\nu)$ and then give the proof of Proposition~\ref{prop:DeltaTo0} and Lemma~\ref{lem:GrowthF_eta} at the end of the section.



\begin{lemma}[Growth of $F_\eta(E)$, $F_1(E),$ $F_{\eta,\sigma_E,1}(E)$ and $F_{\eta,\sigma_E,2}(E)$]
\label{lem:GrowthF_eta}
Let $\eta$ be non-decreasing and symmetric, i.e. $\eta(\nu)=\eta(-\nu).$
Denoting \begin{align*}
   N(E):=|\{E'\in\spec(H)\big| E'< E\}|\qquad\text{and}\,\qquad B_j(E) := \left|\left\{ E'\in \spec(H): E'\in [E+j,E+(j+1))\right\}\right|,
\end{align*} for $j\in\N_0,$ 
we have
\begin{align}
\label{eq:F_etaF_1Growth}
   \nn F_1(E) &\le N(E) + \sum_{j=0}^{\infty} B_j(E)\, e^{-\beta j}, \\F_\eta(E) &\le N(E) \eta(E) + \sum_{j=0}^{\infty} B_j(E)\, e^{-\beta j}\eta(j+1). \end{align}
Moreover, we have
   \begin{align}
\label{eq:F_sigma_EGrowth}\nn F_{\eta,\sigma_E,1}(E)&\le  N^2(E)\eta(E) +N(E)\sum_{j=0}^{\infty} N(E+j+1)e^{-\beta j}\eta(j+1) \\&\nn+\sum_{l=0}^{\infty} \Big(B_l(E) N(E+l+1)\eta(E+l+1) + B_l(E)\sum_{j=0}^{\infty} N(E+j+l+2) e^{-\beta j}\eta(j+1)\Big)e^{-\frac{l^2}{8\sigma_E}}\\ 
    F_{\eta,\sigma_E,2}(E) &\le   \left(N(E) \eta(E) + \sum_{j=0}^{\infty} B_j(E) e^{-\beta j/2} \eta(j+1)\right)\left(N(E)+\sum_{j=0}^\infty e^{-\frac{j^2}{8\sigma^2_E}} B_j(E)\right)
\end{align}


\bin{Labelling the spectrum of $H$ in increasing order as $\spec(H) = \{E^{(k)}\}_{k\in\N_0}$ with $E^{(k+1)}\ge E^{(k)},$ we further assume that there exists $s\in\N$ and $\Delta>0$ such that
\begin{align}
\label{eq:GrowthConditionEigenEnergies}   E^{(k+s)} - E^{(k)} \ge \Delta
\end{align}
for all $k\in\N_0.$ Then we have
\begin{align}
\label{eq:GrowthFeta}
    F_\eta(E) \lesssim (E+h_0)\,\eta(E) + \sum_{j=0}^{\infty}e^{-\beta j}\eta(j+1),
\end{align}
where the $\lesssim$-notation hides a multiplicative constant which only depends on $s$ and $\Delta.$}
\end{lemma}
\begin{remark} [Growth of $F_\eta(E)$, $F_1(E),$ $F_{\eta,\sigma_E,1}(E)$ and $F_{\eta,\sigma_E,2}(E)$ assuming Weyl-asymptotics]
For many models of interest, for example the Bose-Hubbard model with repulsive on-site interactions or trapped particles interacting via Coulomb potentials, it can be shown that the functions $N(E)$ and $B_j(E)= N(E+j+1)-N(E+j)$ satisfy \emph{Weyl-type asymptotics}, namely
\begin{align*}
    N(E) = \mathcal{O}( E^k), \qquad\text{and }\qquad B_j(E) = \mathcal{O}((E+j)^{k-1}).
\end{align*}
for some $k\ge 0.$ In this case, we see for  $\eta(E)$ growing subexponentially that the second terms in the upper bounds on $F_{\eta}(E)$ and $F_1(E)$ in \eqref{eq:F_etaF_1Growth} are constant and therefore
\begin{align*}
    F_{\eta}(E) = \mathcal{O}\left(N(E)\eta(E)\right) = \mathcal{O}\left(E^k\eta(E)\right)\qquad\text{and}\qquad F_{1}(E) =  \mathcal{O}\left(E^k\right).
\end{align*}
Similarly, we see under the same assumptions applied to the bounds in \eqref{eq:F_sigma_EGrowth} that
\begin{align*}
    F_{\eta,\sigma_E,i}(E)  = \mathcal{O}\left(E^{2k}\eta(E)\right)
\end{align*}
for $i=1,2.$ Hence, up to the polynomial corrections, the required energy constraints on the input states in Proposition~\ref{prop:DeltaTo0} are determined by the growth of the function $\eta(E).$
\end{remark}

\medskip 

 \noindent In the following we consider specific choices of $\widehat \phi_\delta(\nu)$ and $\eta(\nu).$ In particular we consider
\begin{align*}
    \widehat \phi_{\delta}(\nu) = e^{-\delta \eta(\nu)}
\end{align*}
for some smooth and non-negative function $\eta(\nu).$ Using this construction, it is easy to see that the assumptions around \eqref{eq:Defeta} are clearly satisfied. We want to use this for the specific choices
\begin{align*}
    \eta_1(\nu) := (\beta\nu)^2\qquad\text{or}\qquad \eta_{2,\theta}(\nu) := e^{(1+(\beta\nu)^2)^\theta}
\end{align*}
for some $\theta\in(0,1/2),$ which satisfy the assumptions of Lemma~\ref{lem:GrowthF_eta} and are subexponentially growing. Given that $\hatfM(\nu)$ is bounded, smooth, and generally well-behaved, we can convince ourselves that $\hatfMdelta = \hatfM \,\widehat\phi_\delta$ for $\delta>0$ is a Schwartz function for  both choices $\eta_1$ and $\eta_{2,\theta}.$





\begin{proof}[Proof of Proposition~\ref{prop:DeltaTo0}]


We denote
\begin{align*}
    \Phi_{\sigma_E,\hatfM,H}(\rho)&:=\sum_{\substack{\alpha\in\mathcal{A}\\ \nu_1,\nu_2}}e^{-\frac{(\nu_1-\nu_2)^2}{8\sigma_E^2}}\, \overline{\hatfM(\nu_1)}\hatfM(\nu_2)\, A^\alpha_{\nu_2} \rho (A^\alpha_{\nu_1})^\dagger,\\   \Phi_{\sigma_E,\hatfMdelta,H}(\rho)&:=\sum_{\substack{\alpha\in\mathcal{A}\\ \nu_1,\nu_2}}e^{-\frac{(\nu_1-\nu_2)^2}{8\sigma_E^2}}\, \overline{\hatfMdelta(\nu_1)}\hatfMdelta(\nu_2)\, A^\alpha_{\nu_2} \rho (A^\alpha_{\nu_1})^\dagger 
\end{align*}
and therefore

\begin{align}
\label{eq:DifferencePhiDelta0}
    \nn&\Phi_{\sigma_E,\hatfM,H}(\rho) - \Phi_{\sigma_E,\hatfMdelta,H}(\rho)\\&=\sum_{\substack{\alpha\in\cA\\\nu_1,\nu_2\in B(H)}}e^{-\frac{(\nu_1-\nu_2)^2}{8\sigma_E^2}}\overline{\hatfM(\nu_1)}\hatfM(\nu_2)\left(1-\widehat\phi_\delta(\nu_1) +\widehat\phi_\delta(\nu_1)(1-\widehat\phi_\delta(\nu_2)\right)\,A^\alpha_{\nu_2} \rho (A^\alpha_{\nu_1})^\dagger.
\end{align}
Taking the trace norm, the first term involving $1-\widehat\phi_\delta(\nu_1)$  can be bounded as
\begin{align*}
&\left\|\sum_{\substack{\nu_1,\nu_2\in B(H)}}e^{-\frac{(\nu_1-\nu_2)^2}{8\sigma_E^2}}\,\overline{\hatfM(\nu_1)}\hatfM(\nu_2)\left(1-\widehat\phi_\delta(\nu_1)\right)\,A^\alpha_{\nu_2} \rho (A^\alpha_{\nu_1})^\dagger\right\|_1 \\&\le \delta \|A^{\alpha} \widetilde H^{-\gamma}\|^2 \sum_{E,E',E'',E'''\in\spec(H)}  \eta(E'-E)|\hatfM(E'-E)||\hatfM(E'''-E'')| \left\|\widetilde H^{\gamma} P_{E''}\rho P_{E}\widetilde H^{\gamma}\right\|_1
\\&= \delta \|A^{\alpha} \widetilde H^{-\gamma}\|^2 \sum_{E,E''\in\spec(H)}  F_\eta(E) F_1(E'')\left\|\widetilde H^{\gamma} P_{E''}\rho P_{E}\widetilde H^{\gamma}\right\|_1 \\&\le 
\delta\|A^\alpha \widetilde H^{-\gamma}\|^2 \| \widetilde F_1(H)\widetilde H^{\gamma} \rho \widetilde H^{\gamma}\widetilde F_\eta(H)\|_1 \left(\sum_{E\in\spec(H)} \tfrac{F_\eta(E)}{ \widetilde F_\eta(E)}\right)\left( \sum_{E''\in\spec(H)} \tfrac{F_1(E'')}{ \widetilde F_1(E'')}\right).
\end{align*}
where in the first inequality we upper bounded the Gaussian\footnote{Keeping the Gaussian would usually lead to slightly weaker requirements on the energy of the input state as $\sum_{E',E'''\in\spec(H)} e^{-\frac{(E'+E''- E -E''')^2}{8\sigma^2_E}} \eta(E'-E)|\hatfM(E'-E)||\hatfM(E'''-E'')|$ in many cases of interest does not increase in $E''$ whereas $\sum_{E',E'''\in\spec(H)} \eta(E'-E)|\hatfM(E'-E)||\hatfM(E'''-E'')|$ does linearly. However, in the interest of an easier analysis we bound the Gaussian by 1.} by 1 and further used that
\begin{align*}
    0\le1-\widehat\phi_\delta(\nu) = \widehat\phi_0(\nu)-\widehat\phi_\delta(\nu)\le \delta\sup_{\delta\ge 0}(-\partial_\delta\widehat\phi_\delta(\nu))=  \delta\,\eta(\nu)
\end{align*} by the mean value theorem and the definition of $\eta(\nu)$ in \eqref{eq:Defeta}. 
 \bin{Furthermore, for $H$ having simple spectrum, we could improve the above inequality by using $\|P_{E} \widetilde H^{\gamma}\rho \widetilde H^{\gamma}P_{E''}\|_1 = \|P_{E} \widetilde H^{\gamma}\rho \widetilde H^{\gamma}P_{E''}\|$ in this case which results in all trace norms involving $\rho$ being replaced by operator norms.}
The second summand in \eqref{eq:DifferencePhiDelta0} including $\widehat\phi_\delta(\nu_1)(1-\widehat\phi_\delta(\nu_2))$ can be estimated the same way using $|\widehat\phi_\delta(\nu_1)|\le 1.$ 
Next denote
\begin{align*}
   G_{\sigma_E}&:= -\sum_{\substack{\alpha\in\mathcal{A}\\ \nu_1,\nu_2}}\!e^{-\frac{(\nu_1-\nu_2)^2}{8\sigma_E^2}}\!\!\! \frac{1}{1+e^{\frac{\beta(\nu_2-\nu_1)}{2}}}\overline{\hatfM(\nu_1)}\hatfM(\nu_2)(A^\alpha_{\nu_1})^\dagger A^\alpha_{\nu_2}\quad\text{and}\\G_{\sigma_E,\delta}&:= -\sum_{\substack{\alpha\in\mathcal{A}\\ \nu_1,\nu_2}}\!e^{-\frac{(\nu_1-\nu_2)^2}{8\sigma_E^2}}\!\!\! \frac{1}{1+e^{\frac{\beta(\nu_2-\nu_1)}{2}}}\overline{\hatfMdelta(\nu_1)}\hatfMdelta(\nu_2)(A^\alpha_{\nu_1})^\dagger A^\alpha_{\nu_2}
\end{align*}
and therefore
\begin{align}
\label{eq:DifferenceGDelta0}
G_{\sigma_E}-G_{\sigma_E,\delta}&:= -\sum_{\substack{\alpha\in\mathcal{A}\\ \nu_1,\nu_2}}\!e^{-\frac{(\nu_1-\nu_2)^2}{8\sigma_E^2}}\!\!\! \frac{1}{1+e^{\frac{\beta(\nu_2-\nu_1)}{2}}}\overline{\hatfM(\nu_1)}\hatfM(\nu_2)\left(1-\widehat\phi_\delta(\nu_1) +\widehat\phi_\delta(\nu_1)(1-\widehat\phi_\delta(\nu_2)\right)(A^\alpha_{\nu_1})^\dagger A^\alpha_{\nu_2}.
\end{align}
We focus on the first term involving $1-\widehat\phi_\delta(\nu_1)$ and see by taking its trace norm that
\bin{\footnote{In the first inequality we use $e^{-\frac{(E''-E)^2}{8\sigma_E^2}}\le 1$ and $\frac{1}{1+e^{\frac{\beta(E''-E)}{2}}}\le 1.$ However, keeping these to regularise the sums over $E$ and $E''$ would lead to slightly weaker requirements on the energy of the input state. However, in the interest of an easier analysis we proceed as below.}}
\begin{align*}
&\left\|\sum_{\nu_1,\nu_2}\!e^{-\frac{(\nu_1-\nu_2)^2}{8\sigma_E^2}}\!\!\! \frac{1}{1+e^{\frac{\beta(\nu_2-\nu_1)}{2}}}\overline{\hatfM(\nu_1)}\hatfM(\nu_2)\left(1-\widehat\phi_\delta(\nu_1)\right)(A^\alpha_{\nu_1})^\dagger A^\alpha_{\nu_2}\,\rho\right\|_1\\&\le \delta\sum_{E,E',E''\in\spec(H)}\!e^{-\frac{(E''-E)^2}{8\sigma_E^2}}|\hatfM(E'-E)|\,\eta(E'-E)\left\|P_E(A^\alpha)^\dagger P_{E'}A^\alpha P_{E''}\,\rho\right\|_1 
    \\&= \delta\sum_{E''\in\spec(H)} F_{\eta,\sigma_E,1}(E'') \left\|P_E(A^\alpha)^\dagger P_{E'}A^\alpha P_{E''}\,\rho\right\|_1  \\&\le \delta\|\widetilde H^{-\gamma}A^{\alpha} \|\|\widetilde H^{\gamma}A^\alpha \widetilde H^{-\mu}\| \|\widetilde F_{\eta,\sigma_E,1}(H)\widetilde{H}^{\mu}\rho\|_1\!\sum_{\substack{E''\in\spec(H)}}\!\tfrac{F_{\eta,\sigma_E,1}(E'')}{\widetilde F_{\eta,\sigma_E,1}(E'')}.
\end{align*}
For the second summand in \eqref{eq:DifferenceGDelta0} including $\widehat\phi_\delta(\nu_1)(1-\widehat\phi_\delta(\nu_2))$ we use $|\widehat\phi_\delta(\nu_1)|\le 1$ and see 
\begin{align*}
&\left\|\sum_{\nu_1,\nu_2}\!e^{-\frac{(\nu_1-\nu_2)^2}{8\sigma_E^2}}\!\!\! \frac{1}{1+e^{\frac{\beta(\nu_2-\nu_1)}{2}}}\overline{\hatfM(\nu_1)}\hatfM(\nu_2)\widehat \phi_\delta(\nu_1)\left(1-\widehat\phi_\delta(\nu_2)\right)(A^\alpha_{\nu_1})^\dagger A^\alpha_{\nu_2}\,\rho\right\|_1\\&\le \delta\sum_{E,E',E''\in\spec(H)}\!e^{-\frac{(E''-E)^2}{8\sigma_E^2}}|\hatfM(E'-E'')|\,\eta(E'-E'')\left\|P_E(A^\alpha)^\dagger P_{E'}A^\alpha P_{E''}\,\rho\right\|_1 
    \\&= \delta\sum_{E'',E\in\spec(H)} e^{-\frac{(E''-E)^2}{8\sigma^2_E}}F_{\eta}(E'') \left\|P_E(A^\alpha)^\dagger P_{E'}A^\alpha P_{E''}\,\rho\right\|_1  \\&\le \delta \|\widetilde H^{-\gamma}A^{\alpha} \|\|\widetilde H^{\gamma}A^\alpha \widetilde H^{-\mu}\| \|\widetilde F_{\eta,\sigma_E,2}(H)\widetilde{H}^{\mu}\rho\|_1\!\sum_{\substack{E''\in\spec(H)}}\!\tfrac{F_{\eta,\sigma_E,2}(E'')}{\widetilde F_{\eta,\sigma_E,2}(E'')}.
\end{align*}
\end{proof}

\begin{proof}[Proof of Lemma~\ref{lem:GrowthF_eta}]
Since 
\begin{align*}
    \hatfM (E'-E) &= e^{ -\frac{\sqrt{1 + (\beta (E'-E))^2} + \beta (E'-E)}{4}} \le e^{-\frac{\beta|E'-E| + \beta(E'-E)}{4}} =
    \begin{cases} 1 \quad&\text{ for }\, E'< E\\
e^{-\frac{\beta(E'-E)}{2}} \quad &\text{ else},
\end{cases} 
\end{align*}
 and using that $\eta$ is non-decreasing and symmetric we can bound $F_\eta$ in \eqref{eq:WeirdG} as follows
\begin{align}
\label{eq:GUpperbound}
    \nn F_\eta(E) &\le \left|\left\{E'\in\spec(H)\big| E'< E\right\}\right|\eta(E) + \sum_{\substack{E'\in\spec(H)\\E'\ge E}} e^{-\beta(E'-E)/2}\,\eta(E'-E)\\&\le 
      N(E) \eta(E) + \sum_{j=0}^{\infty} B_j(E) e^{-\beta j/2} \eta(j+1),
\end{align}
where for the second inequality we split the sum over $E'\ge E$ in a a sum over $j\in\N_0$ and over $\{E'\in\spec(H):E'\in[E+j,E+(j+1))\}$ and estimate the exponential and $\eta$ by the upper and lower end points of the energy intervals respectively.

Furthermore, we can bound the Gaussian sum as 
\begin{align*}
\sum_{E'\in\spec(H)} e^{-\frac{(E'-E)^2}{8\sigma^2_E}} = \sum_{\substack{E'\in\spec(H)\\E'< E}} e^{-\frac{(E'-E)^2}{8\sigma^2_E}} + \sum_{\substack{E'\in\spec(H)\\E'\ge E}} e^{-\frac{(E'-E)^2}{8\sigma^2_E}} \le N(E) + \sum_{j=0}^\infty e^{-\frac{j^2}{8\sigma^2_E}} B_j(E).
\end{align*}
Using this we obtain 
\begin{align*}
    F_{\eta,\sigma_E,2}(E) &= F_\eta(E) \sum_{E'\in\spec(H)} e^{-\frac{(E'-E)^2}{8\sigma^2_E}} \\&\le \left(N(E) \eta(E) + \sum_{j=0}^{\infty} B_j(E) e^{-\beta j/2} \eta(j+1)\right)\left(N(E)+\sum_{j=0}^\infty e^{-\frac{j^2}{8\sigma^2_E}} B_j(E)\right)
\end{align*}
Lastly for $F_{\eta,\sigma_E,1}(E)$ we argue similarly as 
\begin{align*}
F_{\eta,\sigma_E,1}(E) &= \sum_{E'\in\spec(H)} F_\eta(E')e^{-\frac{(E'-E)^2}{8\sigma_E}} \le    \sum_{E'\in\spec(H)}\left(N(E') \eta(E') + \sum_{j=0}^{\infty} B_j(E')\, e^{-\beta j}\eta(j+1)\right)e^{-\frac{(E'-E)^2}{8\sigma_E}} \\&= N(E)\left(N(E) \eta(E) + \sum_{j=0}^{\infty} \max_{E'<E}B_j(E')\, e^{-\beta j}\eta(j+1)\right) \\&\qquad+\sum_{E'\ge E} \left(N(E') \eta(E') +\sum_{j=0}^{\infty} B_j(E')\, e^{-\beta j}\eta(j+1)\right)e^{-\frac{(E'-E)^2}{8\sigma_E}} \\&\le N(E)\left(N(E)\eta(E) +\sum_{j=0}^{\infty} \max_{E'<E}B_j(E')e^{-\beta j}\eta(j+1) \right)\\&\qquad+\sum_{l=0}^{\infty} \Big(B_l(E) N(E+l+1)\eta(E+l+1) \\&\qquad\qquad+ B_l(E)\sum_{j=0}^{\infty} \max_{E'\in [E+l,E+(l+1))\cap\spec(H)} B_j(E') e^{-\beta j}\eta(j+1)\Big)e^{-\frac{l^2}{8\sigma_E}} \\ &\le N^2(E)\eta(E) +N(E)\sum_{j=0}^{\infty} N(E+j+1)e^{-\beta j}\eta(j+1) \\&\qquad+\sum_{l=0}^{\infty} \Big(B_l(E) N(E+l+1)\eta(E+l+1) \\&\qquad\qquad+ B_l(E)\sum_{j=0}^{\infty} N(E+j+l+2) e^{-\beta j}\eta(j+1)\Big)e^{-\frac{l^2}{8\sigma_E}}
\end{align*}
where we consistently used that $E\mapsto N(E)$ and $E\mapsto \eta(E)$ are non-decreasing and, furthermore, that by definition $B_j(E) \le N(E+j+1).$
\end{proof}
\subsubsection{Finite-dimensional Lindblad dynamics for singular filter functions}
\label{sec:FiniteDimPreperationSingular}
In this section we combine the results on finite-dimensional truncations for Schwartz filter functions, in particular Propositions~\ref{prop:generatorboundTrunc} and~\ref{prop:generatorboungFullyFinite}, with the approximation result Proposition~\ref{prop:DeltaTo0} for the Metropolis-type filter function to see in Theorem~\ref{thm:SingularGibbsdynamics} that the dynamics $e^{t\cL_{\sigma_E,\hatfM,H}}$ is well approximated by the finite-dimensional dynamics $e^{t\cL^{\le M}_{\sigma_E,\hatfMdelta,H_{\le M}}}.$ 
Moreover,  under the assumption of spectral gap of the corresponding unbounded generator on the Hilbert-Schmidt space,  we then apply Theorem~\ref{thm:SingularGibbsdynamics} to show efficient finite-dimensional preparation of the Gibbs state of $H$ in Corollary~\ref{cor:SingularGibbs}. 

We first start with an supporting technical lemma on the scaling of the $L^1(\R)$ norm and the $C_f$ and $C'_f$ constants from Sections~\ref{sec:TruncateJumps} and \ref{sec:FinitDimGenerator} for the function $\fMdelta.$
We then continue to 
state Theorem~\ref{thm:SingularGibbsdynamics} and Corollary~\ref{cor:SingularGibbs} and give their proofs at the end of the section.
\begin{lemma}
\label{lem:C_fMdeltaBounds}
Let $\beta>0$ and $\hatfM$ defined in \eqref{eq:filterFunction} and for $\delta\in(0,1]$ and $\theta\in(0,1/2)$ the Schwartz function
\begin{align}\label{hatfMdeltaeq}
    \hatfMdelta(\nu) := \hatfM(\nu)e^{-\delta\eta_{2,\theta}(\nu)}
\end{align}
with $\eta_{2,\theta}(\nu) = e^{(1+(\beta\nu)^2)^\theta}.$  Then for $t\in\R$ we have that the Fourier transform of $\hatfMdelta,$ i.e. $\fMdelta(t) = \frac{1}{2\pi}\int_\R \hatfMdelta(\nu) e^{-it\nu} d\nu,$ satisfies the pointwise estimate 
\begin{align}
\label{eq:fMdeltaPointWise}
    |\fMdelta(t)| \le \frac{C_{\theta}\,e^{-\frac{|t|}{2\beta}}}{\beta}\left(\left(\log(1/\delta)\right)^{\tfrac{1}{2\theta}} +1\right),
\end{align}
where $C_{\theta}\ge 0$ denotes some constant depending only on $\theta$. In particular for $r\ge 0$ and $\kappa\in(0,1/2)$ this shows 
\begin{align}
\label{eq:CfMdeltaBound}
\int_\R |\fMdelta(t)|e^{r|t|^{2\kappa}}dt \le C_{\beta,r,\kappa,\theta}\, \left(\left(\log(1/\delta)\right)^{\tfrac{1}{2\theta}} +1\right)
\end{align}
for some constant $C_{\beta,r,\kappa,\theta}\ge 0$ depending only on $\beta,r,\kappa$ and $\theta$.

\end{lemma}
We prove Lemma~\ref{lem:C_fMdeltaBounds} in Appendix~\ref{sec:fmDeltaProofs}.
\begin{remark}
    A similar result can be obtained for the alternative choice considered in Section~\ref{sec:singular}, namely $\hatfMdelta (\nu) =\hatfM(\nu) e^{-\delta\eta_1(\nu)}$ with $\eta_1(\nu) = (\beta\nu)^2.$ In this case it can be shown that the Fourier transform satisfies
    \begin{align*}
      |\fMdelta(t)| \le \frac{C_{\theta}\,e^{-\frac{|t|}{2\beta}}}{\beta}\left(\frac{1}{\sqrt{\delta}} +1\right)\quad\text{and hence}\quad  \int_\R |\fMdelta(t)|e^{r|t|^{2\kappa}}dt \le C_{\beta,r,\kappa,\theta}\, \left(\frac{1}{\sqrt{\delta}} +1\right)
    \end{align*}
    under the same parameter choices as in Lemma~\ref{lem:C_fMdeltaBounds}. As this scaling as $\delta\to 0$ is  much worse in terms of $\delta$ compared to the one obtained in Lemma~\ref{lem:C_fMdeltaBounds}, we focus on the choice $\eta_{2,\theta}$ in the following. On the other hand, it should be noted that choosing $\eta_1$ instead of $\eta_{2,\theta}$ leads to much weaker requirements on the energy of the input state for the approximation of $\cL_{\sigma_E,\hatfM,H}$ by $\cL_{\sigma_E,\hatfMdelta,H}$ in Proposition~\ref{prop:DeltaTo0}.
\end{remark}

\begin{theorem}[finite-dimensional approximation of $e^{t\cL_{\sigma_E,\hatfM,H}}$]
\label{thm:SingularGibbsdynamics}
For inverse temperature $\beta>0$ let $\hatfM$ defined in \eqref{eq:filterFunction} and consider a set of bare jumps $\{A^\alpha\}_{\alpha\in\cA}$ such that  \Cref{eq:condAalphas} is satisfied for some $\mu\ge \gamma\ge 0$. 
Let further $\NA$ self-adjoint and positive semidefinite and $\kappa\in(0,1/2)$ and assume that \eqref{eq:NormTruncatedJump}, \eqref{eq:AbstractFiniteTruncBareJumps}, \eqref{eq:NormTruncatedJump2} are satisfied and denote \begin{align*}
    C'_{A} := \max_{\alpha\in\cA}\left\{\|e^{\NA^{\kappa}}A^{\alpha}e^{-2\NA^{\kappa}}\|,\|A^{\alpha}e^{-\NA^{\kappa}}\|,\|e^{-\NA^{\kappa}}A^{\alpha}\|, \|A^\alpha\widetilde H^{-\gamma}\|,\|\widetilde H^{\gamma}A^\alpha\widetilde H^{-\mu}\|\right\}<\infty
\end{align*} 
for $\widetilde H= H + (h_0+1)I.$ Furthermore, assume that for $\kappa\in(0,1/2)$ as above \eqref{eq:HamiltonianLeakTrunc} holds and, for some $r\ge 0,$ the condition \eqref{eq:ExpBound} is satisfied.  For $\theta\in(0,1/2)$ consider $\eta_{2,\theta}(\nu) := e^{(1+(\beta\nu)^2)^\theta}$ and assume that the Gibbs state of $H,$ $\sigma_\beta,$ satisfies
\begin{align*}
E'_{\operatorname{Gibbs}}&:=\max\Big\{\,\Tr\left(e^{4\NA^\kappa}\, \sigma_\beta\right)\,,\,\Tr\left( \widetilde F^2_{\eta_{2,\theta}}(H) \widetilde H^{2\gamma}\sigma_\beta\right)\,,\, \sqrt{\Tr\left(\widetilde F^2_{\eta_{2,\theta},\sigma_E,1}(H) \widetilde H^{2\mu}\sigma_\beta\right)},\\&\qquad\qquad \sqrt{\Tr\left(\widetilde F^2_{\eta_{2,\theta},\sigma_E,2}(H) \widetilde H^{2\mu}\sigma_\beta\right)}\,\Big\}<\infty,
\end{align*}
where the functions $\widetilde F_{\eta_{2,\theta}},\widetilde F_{\eta_{2,\theta},\sigma_E,1}$ and $\widetilde F_{\eta_{2,\theta},\sigma_E,2}$ are defined in Proposition~\ref{prop:DeltaTo0}.
Let $\rho$ be a state such that 
\begin{align}
\label{eq:InputDmaxConstraint2}
    \rho \le \mathfrak{c}\,\sigma_\beta
\end{align}
for some $\mathfrak{c}\ge 1.$ 
Consider Metropolis-type filter function  $\hatfM$ defined in \eqref{eq:filterFunction} and for $\delta\in(0,1]$  the Schwartz function
 \begin{align*}
    \hatfMdelta(\nu) := \hatfM(\nu)e^{-\delta\eta_{2,\theta}(\nu)}.
\end{align*}
Then for all $\sigma_E\in(0,\infty),$ $t\ge 0$ and $M\in\N$ we have
    \begin{align}
\nn&\left\|\left(e^{t\cL_{\sigma_E,\hatfM,H}}-e^{t\cL^{\le M}_{\sigma_E,\hatfMdelta,H_{\le M}}}\right)(\rho)\right\|_1\nn\le  t\, \delta \,\mathfrak{c}\, |\cA|\,C_2\!\left(C_{H,\eta,\sigma_E},C_A\right) \, E'_{\operatorname{Gibss}} \\&\qquad\qquad\quad +t\,\mathfrak{c}\,C_3(\beta,r,\kappa,\sigma_E,C_A,\theta) \left((\log(1/\delta))^{\frac{1}{2\theta}} +1\right)^2 p_4(|\mathcal{A}|,M) e^{-M^\kappa} E'_{\operatorname{Gibbs}}
    \end{align}
    where $C_{H,\eta,\sigma_E}$ is the explicit constant defined in Proposition~\ref{prop:DeltaTo0}  and $C_2\!\left(C_{H,\eta,\sigma_E},C_A\right)\ge 0$ and $C_3(\beta,r,\kappa,\sigma_E,C_A)\ge 0$ are some constants depending only on the displayed parameters and $p_4(|\mathcal{A}|,M)$ is some polynomially bounded function. Therefore, we can achieve
    \begin{align}
\nn&\left\|\left(e^{t\cL_{\sigma_E,\hatfM,H}}-e^{t\cL^{\le M}_{\sigma_E,\hatfMdelta,H_{\le M}}}\right)(\rho)\right\|_1 \le \eps\qquad\text{with}\\& \delta=\Omega\left(\frac{\eps}{\mathfrak{c}\,E'_{\operatorname{Gibbs}}\,|\cA|\,t} \right) \quad\text{and}\quad M = \widetilde{\mathcal{O}}\left(\left(\log\left(\frac{t\,\mathfrak{c}\,E'_{\operatorname{Gibbs}}|\cA|}{ \eps}\right)\right)^{1/\kappa}\right),
   \end{align}
    where the $\Omega$ and $\widetilde{\mathcal{O}}$ notations hide constants independent of the displayed parameters and the $\widetilde \Omega$ notation suppresses  $\operatorname{poly}\log\log$ factors.
\end{theorem}
\begin{corollary}[Gibbs state preparation for singular filter function]\label{cor:SingularGibbs}

Under the same assumptions as in Theorem~\ref{thm:SingularGibbsdynamics} and assuming additionally that $L_{\sigma_E,\hatfM,H},$ the self-adjoint generator on $\mathscr{T}_2(\cH)$ associated to the Lindbladian $\cL_{\sigma_E,\hatfM,H},$ has a positive spectral gap 
$\lambda_2 \equiv \operatorname{gap}\left(L_{\sigma_E,\hatfM,H}\right) >0,$ we can achieve for all $\eps>0$
\begin{align}\nn&\left\|e^{t\cL^{\le M}_{\sigma_E,\hatfMdelta,H_{\le M}}}(\rho) - \sigma_\beta\right\|_1 \le \eps\qquad
\text{with} \qquad t=\mathcal{O}\left(\frac{1}{\lambda_2}\log\left(\frac{\mathfrak{c}}{\eps}\right)\right) \\&\delta=\Omega\left(\frac{\lambda_2\,\eps}{\mathfrak{c}\,E'_{\operatorname{Gibbs}}|\cA|\log\left(\frac{\mathfrak{c}}{\eps}\right)} \right) \qquad\text{and}\qquad M = \widetilde{\mathcal{O}}\left(\left(\log\left(\frac{\mathfrak{c}\,E'_{\operatorname{Gibbs}}|\cA|}{\lambda_2 \eps}\right)\right)^{1/\kappa}\right),
\end{align}
where the $\Omega,$ $\mathcal{O}$ and $\widetilde{\mathcal{O}}$ notations hide constants independent of the displayed parameters and the $\widetilde \Omega$ notation suppresses subdominant $\operatorname{poly}\log\log$ factors.


\end{corollary}
\begin{proof}[Proof of Theorem~\ref{thm:SingularGibbsdynamics}]

We use
\begin{align*}
   & \left(e^{t\cL^{\le M}_{\sigma_E,\hatfMdelta,H_{\le M}}} -e^{t\cL_{\sigma_E,\hatfM,H}}\right)(\rho) \\&= \int_0^t e^{(t-s)\cL^{\le M}_{\sigma_E,\hatfMdelta,H_{\le M}}}\left(\cL^{\le M}_{\sigma_E,\hatfMdelta,H_{\le M}} -\cL_{\sigma_E,\hatfM,H}\right)e^{s\cL_{\sigma_E,\hatfM,H}}(\rho)ds
\end{align*}
and therefore, denoting $\rho(s) :=e^{s\cL_{\sigma_E,\hatfM,H}}(\rho)$ we split
\begin{align}
\label{eq:MoreTriangle}
\nn&\left\|\left(e^{t\cL^{\le M}_{\sigma_E,\hatfMdelta,H_{\le M}}} -e^{t\cL_{\sigma_E,\hatfM,H}}\right)(\rho)\right\|_1  \le   t\sup_{s\in[0,t]}\left\|\left(\cL^{\le M}_{\sigma_E,\hatfMdelta,H_{\le M}}-\cL_{\sigma_E,\hatfM,H}\right)(\rho(s))\right\|_1 \\&\nn\le t\sup_{s\in[0,t]}\left\|\left(\cL^{\le M}_{\sigma_E,\hatfMdelta,H_{\le M}}-\cL^{\le M}_{\sigma_E,\hatfMdelta,H}\right)(\rho(s))\right\|_1 +t\sup_{s\in[0,t]}\left\|\left(\cL^{\le M}_{\sigma_E,\hatfMdelta,H}-\cL_{\sigma_E,\hatfMdelta,H}\right)(\rho(s))\right\|_1\\&\quad+t\sup_{s\in[0,t]}\left\|\left(\cL_{\sigma_E,\hatfMdelta,H}-\cL_{\sigma_E,\hatfM,H}\right)(\rho(s))\right\|_1.
\end{align}
For the first two terms, we use Proposition~\ref{prop:generatorboundTrunc} and~\ref{prop:generatorboungFullyFinite} which gives 
\begin{align*}
    &t\sup_{s\in[0,t]}\left\|\left(\cL^{\le M}_{\sigma_E,\hatfMdelta,H_{\le M}}-\cL^{\le M}_{\sigma_E,\hatfMdelta,H}\right)(\rho(s))\right\|_1 +t\sup_{s\in[0,t]}\left\|\left(\cL^{\le M}_{\sigma_E,\hatfMdelta,H}-\cL_{\sigma_E,\hatfMdelta,H}\right)(\rho(s))\right\|_1\\&\le t \,C_{\fMdelta}\max\{C_{\fMdelta},C'_{\fMdelta}\}\,C_4(\beta,r,\kappa,\sigma_E,C_A) p_4(|\cA|,M)e^{-M^\kappa} \sup_{s\in[0,t]}\max\left\{\sqrt{\Tr\left(e^{4\NA^\kappa}\rho(s)\right)}\,,\, \Tr\left(e^{2\NA^\kappa}\rho(s)\right) \right\}\\&\le t C_3(\beta,r,\kappa,\sigma_E,C_A,\theta) \left((\log(1/\delta))^{\frac{1}{2\theta}} +1\right)^2 p_4(|\cA|,M)e^{-M^\kappa} \sup_{s\in[0,t]}\max\left\{\sqrt{\Tr\left(e^{4\NA^\kappa}\rho(s)\right)}\,,\, \Tr\left(e^{2\NA^\kappa}\rho(s)\right) \right\}
     \\&\le t\,\mathfrak{c}\,C_3(\beta,r,\kappa,\sigma_E,C_A,\theta) \left((\log(1/\delta))^{\frac{1}{2\theta}} +1\right)^2 p_4(|\mathcal{A}|,M) e^{-M^\kappa} E'_{\operatorname{Gibbs}}.
\end{align*}
Here, $C_4(\beta,r,\kappa,\sigma_E,C_A)\ge 0$ and $C_2(\beta,r,\kappa,\sigma_E,C_A,\theta)\ge 0$ are some constants depending only on the displayed parameters, $p_4(|\mathcal{A}|,M)$ is some polynomially bounded function and we denoted
\begin{align*}
    C_{\fMdelta} := \int_{-\infty}^{\infty} e^{r|t|^{2\kappa}} |\fMdelta(t)| dt\quad\text{and}\quad C'_{\hatfMdelta} := \int_{-\infty}^{\infty} |t|e^{r|t|^{2\kappa}} |\fMdelta(t)| dt
\end{align*}
and used in the second inequality that by Lemma~\ref{lem:C_fMdeltaBounds} we have
\begin{align*}
    \max\left\{C_{\fMdelta}\,,\, C'_{\fMdelta}\right\}\le C_{\beta,r,\kappa,\theta}\left((\log(1/\delta))^{\frac{1}{2\theta}} +1\right)
\end{align*}
 and for the last inequality that by \eqref{eq:InputDmaxConstraint}, positivity of the map $e^{s\cL_{\sigma_E,\hatfM,H}}$ and the fact that $\cL_{\sigma_E,\hatfM,H}$ is KMS-symmetric we have
\begin{align}
\label{eq:BoundEverythingByGibbs}
   \rho(s) \le \mathfrak{c}\, e^{s\cL_{\sigma_E,\hatfM,H}}(\sigma_\beta ) = \mathfrak{c}\,\sigma_\beta. 
\end{align}
Lastly, for the third term in \eqref{eq:MoreTriangle} we use Proposition~\ref{prop:DeltaTo0} which yields
\begin{align*}
&t\sup_{s\in[0,t]}\left\|\left(\cL_{\sigma_E,\hatfMdelta,H}-\cL_{\sigma_E,\hatfM,H}\right)(\rho(s))\right\|_1 \le t\, \delta \,|\cA|\,C_2\!\left(C_{H,\eta,\sigma_E},C_A\right) \times\\&\times\!\!\sup_{s\in[0,t]}\max\Big\{\Tr\left( \widetilde F^2_{\eta_{2,\theta}}(H) \widetilde H^{2\gamma}\rho(s)\right)\,,\, \sqrt{\Tr\left(\widetilde F^2_{\eta_{2,\theta}\sigma_E,1}(H) \widetilde H^{2\mu}\rho(s)\right)}, \sqrt{\Tr\left(\widetilde F^2_{\eta_{2,\theta},\sigma_E,2}(H) \widetilde H^{2\mu}\rho(s)\right)}\,\Big\} \\&\le t\,\delta \,\mathfrak{c}\, |\cA|\,C_2\!\left(C_{H,\eta,\sigma_E},C_A\right) \, E'_{\operatorname{Gibss}}.
\end{align*}
Here, $C_{H,\eta,\sigma_E}$ is the explicit constant defined in Proposition~\ref{prop:DeltaTo0} and
$C_2\!\left(C_{H,\eta,\sigma_E},C_A\right)\ge 0$  is some constant depending only on the displayed parameters and for the last inequality we have used  \eqref{eq:BoundEverythingByGibbs}.


\end{proof}
\begin{proof}[Proof of Corollary~\ref{cor:SingularGibbs}]
We split
\begin{align}
\label{eq:MajorTriangle2}
   \left\|e^{t\cL^{\le M}_{\sigma_E,\hatfMdelta,H_{\le M}}}(\rho) - \sigma_\beta\right\|_1  &\le \left\|e^{t\cL_{\sigma_E,\hatfM,H}}(\rho) - \sigma_\beta\right\|_1+\left\|\left(e^{t\cL^{\le M}_{\sigma_E,\hatfMdelta,H_{\le M}}} -e^{t\cL_{\sigma_E,\hatfM,H}}\right)(\rho)\right\|_1  
\end{align}
For the first term, we argue as in Section~\ref{sec.fastconv} using positivity of spectral gap of $L_{\sigma_E,\hatfM,H}$ together with \eqref{eq:InputDmaxConstraint2} to get
\begin{align*} \left\|e^{t\cL_{\sigma_E,\hatfM,H}}(\rho)- \sigma_\beta\right\|_1 &= e^{-\lambda_2 t} \left\|\sigma_\beta^{-1/4}\rho\sigma^{-1/4}_\beta\right\|_2 = e^{-\lambda_2 t} \sqrt{\Tr\left(\sigma^{-1/2}_\beta\rho\sigma^{-1/2}_\beta\rho\right)} \\&\le \sqrt{\mathfrak{c}}\,e^{-\lambda_2 t} \sqrt{\Tr(\rho)}= \sqrt{\mathfrak{c}}\,e^{-\lambda_2 t},
\end{align*}
which gives the desired bound on the mixing time $t.$ The result then follows by Theorem~\ref{thm:SingularGibbsdynamics}.
\end{proof}













\subsection{Finite-dimensional circuit implementation}
\label{sec:FullFiniteDimPipeline}


\noindent In this final implementation section we show how the dynamics resulting from the generator $\cL_{\smash{\sigma_E,\widehat{f},H_{\le M}}}^{\le M}$ can be implemented by a quantum circuit through discretizations of the integral representations of its jump and coherent parts. By the integral representation in Proposition~\ref{integralrep}, we note that the generator $\cL_{\smash{\sigma_E,\widehat{f},H_{\le M}}}^{\le M}$ is constructed from a continuous family of jump operators. In Section~\ref{sec:IntegralDiscretization} below we, therefore, first establish that this generator can be written as a Gaussian integral over certain generators of the form \eqref{eq:GKLS} which can then be approximated by a Lindblad generator with a finite number of jumps using Gauss-Hermite quadratures, c.f.~Proposition ~\ref{prop:DiscreteGauss}. After this discretization, and assuming the filter function satisfies \cite[Assumption 13]{ding2025efficient}, as recalled in \Cref{cond:Window} and \ref{condfgevrey}, we can invoke \cite[Theorem 18]{ding2025efficient} to obtain an efficient circuit implementation of the finite-dimensional Lindblad dynamics.

Finally, combining this with Theorem~\ref{thm:SchwartzGibbsDynamics} yields an efficient finite-dimensional circuit implementation of the infinite-dimensional dynamics \(e^{t\cL_{\sigma_E,\widehat f,H}}\) for Schwartz filter functions \(\widehat f\) in Theorem~\ref{thm:CircuitDynamicsSchwartz}, while combining it with Theorem~\ref{thm:SingularGibbsdynamics} yields the corresponding result for \(e^{t\cL_{\sigma_E,\hatfM,H}}\) in Theorem~\ref{thm:CircuitDynamicsSingular}, with \(\hatfM\) given by \eqref{eq:filterFunction}. Given positive spectral gap of the associated generators on the space of Hilbert-Schmidt operators this gives efficient Gibbs state preparation in Corollary~\ref{cor:GibbsCircuitSchwartz} and \ref{cor:GibbsCircuitSingular} respectively.


\subsubsection{Integral discretizations}
\label{sec:IntegralDiscretization}

We start by recalling \cite[Assumption 13]{ding2025efficient} for the cut-off part of the filter function.
\begin{definition}[Gevrey function]
 A complex-valued $C^\infty$ function
$h : \mathbb{R} \to \mathbb{C}$ is called a \emph{Gevrey function of order} $s \ge 0$ if there exist
constants $C_1, C_2 > 0$ such that, for every nonnegative integer
$n $, the derivatives of $h$ satisfy
\begin{equation}
\label{eq:gevrey-def}
\| h^{(n)}\|_{L^\infty(\mathbb{R})}
\le
C_1 C_2^{n} n^{n s}.
\end{equation}
For fixed constants $C_1, C_2, s$, the set of such Gevrey functions is denoted by $\mathcal{G}^s_{C_1,C_2}$. Furthermore, we define
\[
\mathcal{G}^s = \bigcup_{C_1,C_2>0} \mathcal{G}^s_{C_1,C_2}.
\]


\end{definition}


\begin{condition}
\label{cond:Window}
 In what follows, we also consider a cut-off function $w\in \mathcal{G}^s_{\xi_q,\xi_w}$ for some $\xi_w\ge 1$, with $w(\nu) = \overline{w(-\nu)}$ for all $\nu\in\mathbb{R}$ and such that 
    \[
    w(\nu) = 1 \quad \text{when } |\nu| \le \tfrac12,
    \qquad
    w(\nu) = 0 \quad \text{when } |\nu| \ge 1.
    \]
    Finally, we denote by $\kappa(\nu)=w(\nu/S)$ for some $S>0$ which we will choose later.
\end{condition}






\noindent Next, for filter function $\widehat f\in\mathcal{S}(\R)$ we aim at showing that the generator $\cL_{\smash{\sigma_E,\widehat{f},H_{\le M}}}^{\le M},$ which was defined in Section~\ref{sec:FinitDimGenerator} and studied in Sections~\ref{sec:FiniteDimPreperationSingular} and \ref{sec:FiniteDimPreperationSingular},  can be approximated by a generator made of a constant number of jumps by integral discretization. We recall that 
\begin{align*}
\cL_{\smash{\sigma_E,\widehat{f},H_{\le M}}}^{\le M}(\rho)=\sum_{\alpha\in\cA} \int_\R \gamma(t) X_{t,\le M}^{\alpha,\le M}\rho \left(X_{t,\le M}^{\alpha,\le M}\right)^\dagger dt\, + G^{\le M}_{\sigma_E,\le M}\rho +\rho \left(G^{\le M}_{\sigma_E,\le M}\right)^\dagger 
\end{align*}
with 
\begin{align*}
G^{\le M}_{\sigma_E,\le M} &=- \int_\R g(t) (X_{t,\le M}^{\alpha,\le M})^\dagger X_{t,\le M}^{\alpha,\le M}dt,\quad &
X_{t,\le M}^{\alpha,\le M}:=e^{itH_{\le M}}L_{\le M}^{\alpha,\le M}e^{-itH_{\le M}},\\\gamma(t)&:=\sigma_E\sqrt{\frac{2}{\pi}}e^{-2\sigma_E^2 t^2}\qquad\qquad\qquad\text{and}&g(t):=\frac{1}{2\pi}\int _\R \frac{e^{-\nu^2/8\sigma_E^2}}{1+e^{\beta\nu/2}}e^{-i\nu t}d\nu
\end{align*}
and $L^{\alpha,\le M}_{\le M}$ has been defined in \eqref{eq:TruncatedL}.
  By decomposing the function $g$ into real and imaginary part as
\begin{align*}
g(t)&=\frac{1}{2\pi}\int \frac{e^{-\nu^2/8\sigma_E^2}}{1+e^{\beta\nu/2}}e^{-i\nu t}d\nu=\frac{1}{2}\gamma(t)+\frac{1}{2\pi}\int {e^{-\nu^2/8\sigma_E^2}}e^{-i\nu t}  \frac{1-e^{\beta\nu/2}}{2(1+e^{\beta\nu/2})} d\nu \\&=\frac{1}{2}\gamma(t)+\frac{1}{4\pi}\int {e^{-\nu^2/8\sigma_E^2}}e^{-i\nu t} \tanh(-\beta \nu/4) d\nu,
\end{align*}
we get that 
\begin{align*}
\cL_{\smash{\sigma_E,\widehat{f},H_{\le M}}}^{\le M}(\rho)&=\sum_\alpha \int \gamma(t)\left( X_{t,\le M}^{\alpha,\le M}\rho \left(X_t^{\alpha,\le M}\right)^\dagger-\frac{1}{2}\left\{\left(X_{t,\le M}^{\alpha,\le M}\right)^\dagger X_{t,\le M}^{\alpha,\le M},\rho\right\}\right) dt-i\big[B^{\le M}_{\sigma_E},\rho\big],
\end{align*}
with
\begin{align*}
B_{\sigma_E}^{\le M}:=\sum_{\alpha\in\mathcal{A}}\frac{(-i)}{4\pi}\int\int {e^{-\nu^2/8\sigma_E^2}}e^{-i\nu t}  \tanh(-\beta\nu/4)  e^{itH_{\le M}}(L_{\le M}^{\alpha,\le M})^\dagger L_{\le M}^{\alpha,\le M} e^{-itH_{\le M}}\,d\nu dt=(B_{\sigma_E}^{\le M})^\dagger,
\end{align*}
where the self-adjointness follows from the antisymmetry of $\tanh$. Decomposing into the eigenbasis of $H_{\le M} = \sum_{E\in\spec(H_{\le M})} E \,P_E$ and denoting $\big[(L_{\le M}^{\alpha,\le M})^\dagger L_{\le M}^{\alpha,\le M}\big]_\mu=\sum_{E'-E=\mu}P_{E'}(L_{\le M}^{\alpha,\le M})^\dagger L_{\le M}^{\alpha,\le M}P_{E}$, we get
\begin{align*}
B_{\sigma_E}^{\le M}&=\sum_{\alpha\in\mathcal{A}}\sum_{\mu\in B(H_{\le M})}\frac{(-i)}{4\pi}\int\int {e^{-\nu^2/8\sigma_E^2}}e^{-i(\nu-\mu) t}  \tanh(-\beta\nu/4)  \big[(L_{\le M}^{\alpha,\le M})^\dagger L_{\le M}^{\alpha,\le M} \big]_\mu\,d\nu dt\\
&=\sum_{\alpha\in\mathcal{A}}\sum_{\mu\in B(H_{\le M})}\frac{(-i)}{2} {e^{-\mu^2/8\sigma_E^2}}  \tanh(-\beta\mu/4)  \big[(L_{\le M}^{\alpha,\le M})^\dagger L_{\le M}^{\alpha,\le M} \big]_\mu.
\end{align*}
Using the notation from \Cref{condfgevrey}, we define the Schwartz function $\widehat{t}_\kappa(\mu):=-\frac{i}{2}\tanh(-\beta\mu/4)\kappa(\mu)$ for $S\ge 4\|H_{\le M}\|$, and using that for all Bohr frequencies of the truncated Hamiltonian $\mu\in B(H_{\le M})$ we have $|\mu|\le 2\|H_{\le M}\| \le S/2$ and therefore $\widehat t_\kappa(\mu) = -\frac{i}{2}\tanh(-\beta \mu/4),$ we see
\begin{align*}
B_{\sigma_E}^{\le M}&=\sigma_E\sqrt{\frac{2}{\pi}}\sum_{\alpha\in\mathcal{A}}\sum_{\mu\in B(H_{\le M})}\int e^{-i\mu t-2\sigma_E^2 t^2 }   \widehat{t}_\kappa(\mu)  \big[(L_{\le M}^{\alpha,\le M})^\dagger L_{\le M}^{\alpha,\le M} \big]_\mu dt\\
&=\sigma_E\sqrt{\frac{2}{\pi}}\sum_{\alpha\in\mathcal{A}}\sum_{\mu\in B(H_{\le M})}\int e^{-2\sigma_E^2 t^2 }   \widehat{t}_\kappa(\mu) e^{-itH} \big[(L_{\le M}^{\alpha,\le M})^\dagger L_{\le M}^{\alpha,\le M} \big]_\mu e^{itH}\,dt\\
&=\sum_{\alpha\in\mathcal{A}}\int \gamma(t)   {t}_\kappa(s) e^{i(s-t)H_{\le M}} (L_{\le M}^{\alpha,\le M})^\dagger L_{\le M}^{\alpha,\le M}  e^{-i(s-t)H_{\le M}}\,ds\,dt\\
&= \int \gamma(t)\,e^{-it H_{\le M}}B^{\le M}e^{itH_{\le M}}dt ,
\end{align*}
where in the last equality we have denoted 
\begin{align*}
   B^{\le M} := \sum_{\alpha\in\cA} \int_\R t_\kappa(s) \,e^{isH_{\le M}}\left(L^{\alpha,\le M}_{\le M}\right)^\dagger L^{\alpha,\le M}_{\le M}e^{-isH_{\le M}} \,ds. 
\end{align*}
All in all, we have found that 
\begin{align}\label{convcomb}
\mathcal{L}_{\sigma_E,\widehat{f},H_{\le M}}^{\le M}(\rho)=\int \gamma(t) \,\mathcal{L}_{\widehat{f},H_{\le M}}^{\le M,\,t}(\rho)\, dt,
\end{align}
where $\mathcal{L}_{\widehat{f},H_{\le M}}^{\le M,\,t}$ coincides with the generator $\mathcal{L}_{\widehat{f},H_{\le M}}$ defined in \eqref{eq:GKLS} with filter function $\widehat{f}$, jumps $\left(A^{\alpha}\right)^{\le M}_t:= e^{itH_{\le M}} (A^{\alpha})^{\le M} e^{-itH_{\le M}}$ and Hamiltonian $H_{\le M}$. 


Next, we need to argue that, without loss of generality, the filter function $\widehat{f}$ can be replaced by  $\widehat{f}_{\kappa}(\nu):=\widehat{f}(\nu)\kappa(\nu)$ with $\kappa$ satisfying \Cref{cond:Window} for some choice of parameter $S$. This is done by observing that
\begin{align}
\label{eq:LMkappaVsNoKappa}
\nn L_{\le M}^{\alpha,\le M}=\int f(s)\,e^{isH_{\le M}}(A^{\alpha})^{\le M}e^{-isH_{\le M}}\,ds&=\sum_{\nu\in B(H_{\le M})}\widehat{f}(\nu) ((A^{\alpha})^{\le M})_{\nu}\\\nn
&=\sum_{\nu\in B(H_{\le M})}\widehat{f}_\kappa(\nu) ((A^{\alpha})^{\le M})_{\nu}\\
&=\int f_\kappa(s)\,e^{is H_{\le M}}(A^{\alpha})^{\le M}e^{-is H_{\le M}} ds,
\end{align}
for $S\ge 4\|H_{\le M}\|$. Thus, without loss of generality, we can replace $\mathcal{L}_{\widehat{f},H_{\le M}}^{\le M,\,t}$ by $\mathcal{L}_{\widehat{f}_\kappa,H_{\le M}}^{\le M,\,t}$ in \Cref{convcomb}.

We invoke Gaussian-Hermite quadratures, recalled in the following supporting lemma, to find a discretization of the Gaussian integral in \eqref{convcomb} in Proposition~\ref{prop:DiscreteGauss} below.




\begin{lemma}
\label{lem:GaussHermite}
Let $X$ be a Banach space, $:\mathbb R\to X$ be $C^\infty$ and assume that
\[
\|F^{(m)}(t)\|_X \le K\,C^m
\qquad
\text{for all }m\in\mathbb N_0,\ t\in\mathbb R,
\]
for some constants $C,\,K\ge 0$. Let $\{(x_k,w_k)\}_{k=1}^n \subset \R\times \R$ be the $n$-point Gauss--Hermite nodes and weights with $\max_k|x_k|\le 2\sqrt{n}$ and $\max_k w_k=\mathcal{O}(n^{-1/2})$, and define
\[
Q_n(f)
:=
\frac{1}{\sqrt{\pi}}
\sum_{k=1}^n
w_k\,F\!\left(\frac{x_k}{\sqrt2\,\sigma_E}\right).
\]
Then, 
\[
n \ge \frac{C^2}{2\sigma_E^2} + \log_2\!\left(\frac{K}{\varepsilon}\right)\quad \Rightarrow\quad 
\left\|
\int_{\mathbb R} \gamma(t)\,F(t)\,dt
-
Q_n(F)
\right\|_X
\le \varepsilon.
\]
\end{lemma}

\begin{proof}
Set 
\[
I(F):=\int_{\mathbb R} \gamma(t)\,F(t)\,dt.
\]
With the change of variables
$
x=\sqrt2\,\sigma_E t,
$ we get
\[
I(F)
=
\frac{1}{\sqrt{\pi}}
\int_{\mathbb R}
e^{-x^2}
F\!\left(\frac{x}{\sqrt2\,\sigma_E}\right)\,dx.
\]
For an element  of $X'$, the dual of $X$, i.e. a bounded linear functional $\varphi: X\to \C$, define
\(
g_\varphi(x):=
\varphi\!\left(F\!\left(\frac{x}{\sqrt2\,\sigma_E}\right)\right).
\)
Then
\[
\varphi(I(F)-Q_n(F))
=
\frac{1}{\sqrt{\pi}}
\left(
\int e^{-x^2} g_\varphi(x)\,dx
-
\sum_{k=1}^n w_k g_\varphi(x_k)
\right).
\]
By the scalar Gauss-Hermite error formula,
\[
\left|
\int e^{-x^2} g_\varphi(x)\,dx
-
\sum_{k=1}^n w_k g_\varphi(x_k)
\right|
\le
\frac{\sqrt{\pi}\,n!}{2^n(2n)!}
\sup_{x\in\R} |g_\varphi^{(2n)}(x)|.
\]
By the chain rule we have,
\[
g_\varphi^{(2n)}(x)
=
\left(\frac{1}{\sqrt2\,\sigma_E}\right)^{2n}
\varphi\!\left(F^{(2n)}\!\left(\frac{x}{\sqrt2\,\sigma_E}\right)\right),
\]
and therefore
\[
|g_\varphi^{(2n)}(x)|
\le
\|\varphi\|
\left(\frac{1}{2\sigma_E^2}\right)^n
\left\|
F^{(2n)}\!\left(\frac{x}{\sqrt2\,\sigma_E}\right)
\right\|_X.
\]
Using the assumption,
$
\left\|
F^{(2n)}(t)
\right\|_X
\le K\,
C^{2n},$
 we obtain
\[
\sup_{x\in\R}|g_\varphi^{(2n)}(x)|
\le K
\|\varphi\|
\left(\frac{C^2}{2\sigma_E^2}\right)^n.
\]
\bin{Since the Hermite nodes satisfy $|x_k|\lesssim \sqrt{n}$, we bound
\[
\sup_x |g_\varphi^{(2n)}(x)|
\lesssim K
\|\varphi\|
\left(\frac{C^2 n}{\sigma_E^2}\right)^n.
\]}
Combining the estimates and using the dual expression of the norm on $X$ gives
\[
\|I(F)-Q_n(F)\|_X = \sup_{\substack{\phi\in X'\\\|\phi\|\le 1}} \left|\phi\left(I(F)) - Q_n(F)\right)\right|
\le  K
\frac{n!}{(2n)!}
\left(\frac{C }{2\sigma_E}\right)^{2n}.
\]
Using $
\frac{n!}{(2n)!}\le \frac{1}{n^n}$
we obtain
\[
\|I(F)-Q_n(F)\|_X
\le K
\left(\frac{C^2}{4
\sigma_E^2 n}\right)^n.
\]
Hence, for $n\ge \max\left\{\frac{C^2}{2\sigma^2_E},\log_2\left(\frac{K}{\eps}\right)\right\}$ this gives $\|I(F)-Q_n(F)\|_X\le \eps.$
\end{proof}

\begin{proposition}
\label{prop:DiscreteGauss}
 Let $H$ be such that the truncated Hamiltonian $H_{\le M}$, defined in Section~\ref{sec:FinitDimGenerator}, satisfies $\|H_{\le M}\| \le  p_5(\cA,M)$ for some polynomially bounded function $p_5.$ Moreover, let $\{\cA^\alpha\}_{\alpha\in\cA}$ be a set of bare jumps satisfying \eqref{eq:NormTruncatedJump}, i.e. $\|\left(A^\alpha\right)^{\le M}\|\le q'(M)$ for some polynomially bounded function $q'(M).$ Let $\kappa$ satisfy  \Cref{cond:Window} for $s\ge 1$ and $\xi_q,\xi_w\ge 0$ and  
 $S=4p_5(\cA|,M).$ 
 Let furthermore $\beta>0$ and $\widehat f\in\mathcal{S}(\R)$ and denote 
\begin{align*}
    K:=|\mathcal{A}|\xi_q\,(1+\log\left(((\beta+\xi_w/S)\max\{S,1\}\right))   \,\left(q'(M)\right)^2 \min\left\{\|f\|^2_{L^1(\R)},\|f_\kappa\|^2_{L^1(\R)}\right\}.
\end{align*} 
Then for $\beta,\eps>0$ and $t\ge 0$ we can achieve
\begin{align}
\label{eq:DiscreteGaussEvolution}
 \Big\|e^{t\mathcal{L}^{\le M}_{\sigma_E,\widehat{f},H_{\le M}}}-e^{t\mathcal{L}^{\le M,\mathfrak{n}}_{\sigma_E,\widehat{f},H_{\le M}}}\Big\|_{1\to 1}\le \epsilon \qquad \text{with}\qquad \mathfrak{n}=\mathcal{O}\left(  \frac{S^2}{\sigma_E^2}+\log\Big(\frac{K t}{\epsilon}\Big)\right)
\end{align}
with, given the $\mathfrak{n}$-point Gauss-Hermite nodes and weights $\{(x_k,w_k)\}_{k=1}^{\mathfrak{n}}$ with $\max_{k}w_k=\mathcal{O}(\mathfrak{n}^{-1/2})$ and $\max_k|x_k|\le \mathfrak{n}^{1/2}$, 
\begin{align}
\label{eq:ImplementedGenerator}
\mathcal{L}^{\le M,\mathfrak{n}}_{\sigma_E,\widehat{f},H_{\le M}}:=\frac{1}{\sqrt{\pi}}
\sum_{k=1}^{\mathfrak{n}}
w_k\,\mathcal{L}_{\widehat{f},H_{\le M}}^{\le M,x_k/(\sqrt{2}\sigma_E)}
\end{align}
\end{proposition}

\begin{proof}
We use
\begin{align*}
e^{t\mathcal{L}^{\le M}_{\sigma_E,\widehat{f},H_{\le M}}}-e^{t\mathcal{L}^{\le M,\mathfrak{n}}_{\sigma_E,\widehat{f},H_{\le M}}} = \int^t_0 e^{s\mathcal{L}^{\le M}_{\sigma_E,\widehat{f},H_{\le M}}}\left(\mathcal{L}^{\le M}_{\sigma_E,\widehat{f},H_{\le M}} - \mathcal{L}^{\le M,\mathfrak{n}}_{\sigma_E,\widehat{f},H_{\le M}}\right) e^{(1-s)\mathcal{L}^{\le M,\mathfrak{n}}_{\sigma_E,\widehat{f},H_{\le M}}} \,dt
\end{align*}
and, therefore, as both dynamics are contractive in trace norm and the $1\to 1$ norm is submultiplicative, we have
\begin{align*}
\Big\|e^{t\mathcal{L}^{\le M}_{\sigma_E,\widehat{f},H_{\le M}}}-e^{t\mathcal{L}^{\le M,\mathfrak{n}}_{\sigma_E,\widehat{f},H_{\le M}}}\Big\|_{1\to 1} \le \left\|\mathcal{L}^{\le M}_{\sigma_E,\widehat{f},H_{\le M}} - \mathcal{L}^{\le M,\mathfrak{n}}_{\sigma_E,\widehat{f},H_{\le M}}\right\|_{1\to 1}
\end{align*}
Hence, using Lemma~\ref{lem:GaussHermite} for the function $t\mapsto F(t):=\mathcal{L}_{\widehat{f},H_{\le M}}^{\le M,t}$, it suffices to compute estimate the  derivatives as  
\begin{align*}
\|F^{(m)}(t)\|_{1\to 1}&\le 2|\mathcal{A}|\, (4\|H_{\le M}\|)^m\max_{\alpha\in\mathcal{A}}\left\|L_{\le M}^{\alpha,\le M}\right\|^2+2^m\|H_{\le M}\|^m \|B^{\le M}\|\\
&\le 2|\mathcal{A}|\, (4\|H_{\le M}\|)^m\max_{\alpha\in\mathcal{A}}\left\|L_{\le M}^{\alpha,\le M}\right\|^2(1+\|t_\kappa\|_{L^1(\mathbb{R})})\\
&\le 2|\mathcal{A}|\, (4\|H_{\le M}\|)^m\min\left\{\|f\|^2_{L^1(\R)},\|f_\kappa\|^2_{L^1(\R)}\right\} \max_{\alpha\in\mathcal{A}}\left\|\left(A^\alpha\right)^{\le M}\right\|^2(1+\|t_\kappa\|_{L^1(\mathbb{R})})\\ &\le 2|\mathcal{A}|\, (4\|H_{\le M}\|)^m\min\left\{\|f\|^2_{L^1(\R)},\|f_\kappa\|^2_{L^1(\R)}\right\} (q'(M))^2(1+\|t_\kappa\|_{L^1(\mathbb{R})}),
\end{align*}
where in the second to last inequality we have used \eqref{eq:LMkappaVsNoKappa} to bound the operator norm of $L_{\le M}^{\alpha,\le M}.$
Using now \Cref{cond:Window} combined with \cite[Lemma 30]{ding2025efficient}, we see
\begin{align*}
\|t_\kappa\|_{L^1(\mathbb{R})} = \mathcal{O}\big(\xi_q\,(1+\log((\beta+\xi_w/S)\max\{S,1\})),
\end{align*}
which finishes the proof.
\end{proof}



We can now employ the \cite[Theorem 18]{ding2025efficient} to provide a circuit implementation of the finite-dimensional dynamics. For that we first recall the condition on the filter function in \cite[Assumption 13]{ding2025efficient}.



\begin{condition} \label{condfgevrey}
For $\beta > 0$, we consider a filter function $\widehat{f}$ such that $\nu\mapsto \widehat{f}(\nu/\beta)\in   \mathcal{G}^s_{\smash{\xi_q,\xi_u}}$ for some constants $\xi_q,\xi_u \ge 1$ and $s \ge 1$. In addition, we assume
$
\frac{d}{d\nu}\widehat{f}(\nu) \in L^1(\mathbb{R}),
$
and denote
\[
C_{1,f}
:=
\left\|
\frac{d}{d\nu}\widehat{f}(\nu/\beta)
\right\|_{L^1(\mathbb{R})}.
\]
\end{condition} 


\noindent The next result shows that \Cref{condfgevrey} is satisfied for the function $\hatfMdelta$ defined in \eqref{eq:hatfMdeltaeqIntro} and analysed in Sections~\ref{sec:singular} and~\ref{sec:FiniteDimPreperationSingular}.

\begin{lemma}
\label{LemmGevreycond}
Let $\beta>0,$ $\delta\ge 0$ and $\theta\in [0,1/2).$ and denote
$
h_{\delta,\theta}(\nu)
:=
\hatfMdelta\!\left(\nu/\beta\right),$ where the function $\hatfMdelta$ is defined in \eqref{eq:hatfMdeltaeqIntro}. Then \(h_{\delta,\theta}\) satisfies \Cref{condfgevrey}. More precisely, there exists $\xi_{\theta},\, C_{1,\theta}\ge 0$ only depending on $\theta$ but independent of $\delta$ such that \(h_{\delta,\theta}\in \mathcal G^1_{\xi_\theta,4}\) and 
\begin{align*}
\left\|\frac{d}{d\nu} h_{\theta,\delta}\right\|_1 \le C_{1}
\end{align*}
for some $C_1\ge 0$ independent of $\theta$ and $\delta.$
\end{lemma}
We prove Lemma~\ref{LemmGevreycond} in Appendix~\ref{sec:fmDeltaProofs}.





\noindent Following the notation of Proposition~\ref{prop:DiscreteGauss}, it remains to argue that the evolution generated by the Lindbladian $\mathcal{L}^{\le M,\mathfrak{n}}_{\smash{\sigma_E,\widehat{f},H_{\le M}}}$ can be efficiently implemented on a quantum computer. For this, we assume access to the same oracles as in \cite{ding2025efficient}, with renormalized bare jumps $ (A^{\alpha})^{\le M}/q'(M)$ and Hamiltonian simulation of $H_{\le M}$. Therefore, using the fact that the Gauss-Hermite nodes satisfy $\max_{k\in\{1,\cdots,\mathfrak{n}\}} |x_k|\le \mathfrak{n}^{1/2},$  for $\alpha\in\cA$ and $k\in\{1,\cdots,\mathfrak{n}\}$, the time-evolved jump $(A^{\alpha})^{\le M}_{x_k/(\sqrt{2}\sigma_E)}/q'(M)$ can be prepared via total Hamiltonian simulation time of order $  \mathfrak{n}^{1/2}/\sigma_E$. Then, since each query to a block encoding of the jump operators of $\mathcal{L}^{\le M,\mathfrak{n}}_{\smash{\sigma_E,\widehat{f},H_{\le M}}}$ requires a single query to a block encoding of the bare jumps, and each query to the block encoding of coherent part requires two queries of it, we can bound the total Hamiltonian simulation time using \cite[Theorem 18]{ding2025efficient} as follows. 

\begin{theorem}[{\cite[Theorem 18]{ding2025efficient}}]
\label{thm:Lin}
Let $\beta>0$ and $M,\,\mathfrak{n}\in\N.$ Moreover, let $\widehat f$ satisfy \Cref{condfgevrey} and $\kappa$ satisfy \Cref{cond:Window} for some $s,S\ge 1.$  For filter functions $\widehat f_\kappa(\nu) = \widehat f(\nu)\kappa(\nu)$ and $\widehat t_\kappa(\nu) := -\frac{i}{2}\tanh(-\beta\nu/4)\kappa(\nu)$
assume oracle access to the respective Fourier transforms ${f}_\kappa$ and  $t_\kappa$, block encodings of the subnormalised bare jumps $(A^{\alpha})^{\le M}/q'(M)$, and controlled Hamiltonian simulation $U_{H_{\le M}}$. The Lindbladian evolution generated by $\mathcal{L}^{\le M,\mathfrak{n}}_{\smash{\sigma_E,\widehat{f},H_{\le M}}},$ defined via \eqref{eq:ImplementedGenerator} with $\max_{k\in\{1,\cdots,\mathfrak{n}\}} |x_k|\le \mathfrak{n}^{1/2}$ can be simulated up to time $t$ within $\epsilon$-diamond distance with total Hamiltonian simulation time
\begin{align*}
\widetilde {\mathcal{O}}\Big(\sigma_E^{-1} (\beta+1)\left(\log^{2+s}(S)+1\right)\mathfrak{n}^{5/2}|\cA|^2 \ t\ \operatorname{poly}(M)\log^{1+s}(1/\eps)\Big).
\end{align*}
In addition the algorithm requires 
\begin{align*}
\widetilde{\mathcal{O}}\Big(\log\left(S^2 +S\|H_{\le M}\|\right) + \log^2\left(t|\cA|/\eps\right) + \log\left(\beta+1\right)\Big)
\end{align*}
many ancilla qubits.
Here the $\widetilde{\mathcal{O}}$ absorbs subdominant polylogarithmic dependencies on all the parameters. 

\end{theorem}




\begin{remark}[Circuit implementation of \(a_i^{\le M}\) and \((a_i^{\le M})^\dagger\)]
\label{rem:CircuitTruncJumps}
For an $m$-mode continuous variable system with Hilbert space $(L^2(\R))^\otimes \cong L^2(\R^m),$ we usually consider bare jumps being the annihilation and creation operators on each site, i.e.~$\{A^\alpha\}_{\alpha\in\cA} \equiv \{a_i,a^\dagger_i\}_{i=1}^{m}.$ Furthermore, we then consider their Fock basis truncations $a_i^{\le M}$ and $\left(a_i^{\le M}\right)^\dagger$ (c.f.~\eqref{eq:DefTruncAnnCrea}), whose operator norms satisfy $\|a^{\le M}_i\|\,,\, \|(a^{\le M}_i)^\dagger\|\le \sqrt{M}.$
For fixed mode $i$ these truncated operators can be seen as acting on a $(M+1)$-dimensional space and we, hence, argue in the following how to implement them on a \(q=\lceil \log_2(M+1)\rceil\) qubit device. In binary encoding, the  subnormalised operators
\(a_i^{\le M}/\sqrt{M}\) and \((a_i^{\le M})^\dagger/\sqrt{M}\) are weighted shift operators. A naive
implementation via a flat lookup of the \(M+1\) coefficients has depth
\(\mathcal O(M)\), up to precision overhead. 

However, this can be drastically improved
by noting that the shift part \(\ket{n}_i \mapsto \ket{n\pm1}_i\) is simply reversible
increment/decrement on the \(q=\lceil \log_2(M+1)\rceil\) qubits, while the corresponding
matrix elements are given by the efficiently computable functions
\(n\mapsto \sqrt{ n/M}\) and \(n\mapsto \sqrt{(n+1)/M}\). Using reversible arithmetic to compute
these coefficients to accuracy \(\varepsilon\) and loading them into an ancilla rotation
yields a block-encoding of the subnormalized operators
\(a_i^{\le M}/\sqrt M\) and \((a_i^{\le M})^\dagger/\sqrt M\) with operator-norm error at
most \(\varepsilon\), and with $$\mathcal{O}\left(\operatorname{poly}(\log(M),\log(1/\eps))\right)$$ circuit depth.
\end{remark}

\begin{remark}[Circuit implementation of $e^{itH_{\le M}}$] 
For an $m$-mode continuous variable system with Hilbert space $(L^2(\R))^\otimes \cong L^2(\R^m),$ we usually consider the projection onto the system register given by $P_M = \pi^{\otimes m}_M$ with local projections $\pi_M = \sum_{n=0}^{M} \kb{n}$. Due to the tensor product structure of $P_M$, we see for $H$ being an $\mathcal{O}(1)$ degree polynomial in the creation and annihilation operators $a_i$ and $a^\dagger_i,$ that also the truncated Hamiltonian $H_{\le M} = P_MHP_M$ is a sum of local terms. In particular, encoding the system register $\operatorname{im}(P_M)$ by $q=\ceil{m\log_2(M+1)}$ many qubits, we see that the locality of $H_M$ is of order $\log(M).$ Therefore, the unitary evolution $e^{itH_{\le M}}$ can be implemented efficiently using QSVT techniques in $$\mathcal{O}\left(\operatorname{poly}(M,\log(1/\eps))\right)$$ circuit depth
\label{rem:CircuitTruncHam}
\cite{Low_2017,low2017hamiltoniansimulationuniformspectral,Low2019hamiltonian,PhysRevX.6.041067,10.1145/3313276.3316366}.
\end{remark}


\subsubsection{Circuit implementation for Schwartz filter functions}
Combining Proposition~\ref{prop:DiscreteGauss} and Theorem~\ref{thm:Lin} with Theorem~\ref{thm:SchwartzGibbsDynamics} yields the following result on efficient implementation of the Lindblad dynamics for Schwartz filter functions.  Recall that the image of the finite rank projection $P_M$ involved in the truncation of the Hamiltonian in Section~\ref{sec:FinitDimGenerator}, i.e. $H_{\le M} = P_MHP_M,$ is referred to as the system register, which governs the size of the quantum computer on which can provide an circuit implementation of the Lindblad dynamics. To make the following results more explicit we consider the scaling $\log(\operatorname{im}(P_M)) = \mathcal{O}(|\cA|\log(M))$ which is the typical scaling required for quantum many-body or multi-mode continuous variable systems.We note, however, that Theorems~\ref{thm:SchwartzGibbsDynamics} and~\ref{thm:CircuitDynamicsSingular}, as well as Corollaries~\ref{cor:GibbsCircuitSchwartz} and~\ref{cor:GibbsCircuitSingular}, remain valid for more general projections $P_M$, provided the remaining assumptions are satisfied. The corresponding number of required qubits is then simply 
\(
\log_2\left( \dim\bigl(\operatorname{im}(P_M)\bigr)\right).
\)

\begin{theorem}[Circuit implementation of $e^{t\cL_{\sigma_E,\widehat f,H}}$ for $\widehat f$ Schwartz]
\label{thm:CircuitDynamicsSchwartz}
Let $\beta>0,$ $t\ge 0$ and $\widehat f\in\mathcal{S}(\R).$
 Under the assumptions of Theorem~\ref{thm:SchwartzGibbsDynamics},~\ref{thm:Lin} and  Proposition~\ref{prop:DiscreteGauss} and assuming oracle access to ${f}_\kappa$ and  $t_\kappa$, block encodings of the subnormalised bare jumps $(A^{\alpha})^{\le M}/q'(M)$, controlled Hamiltonian simulation $U_{H_{\le M}}$, where 
\begin{align}
\label{eq:MDimension}
 M =\widetilde\Theta\left(\operatorname{poly}\left(\log\left(\frac{t\,\mathfrak{c}\,E_{\operatorname{Gibbs}}|\cA|}{ \eps}\right)\right)\right)
    \end{align}
    and assume oracle access to a state preparation circuit for input state $\rho$ satisfying
    \begin{align*}
        \rho \le \mathfrak{c}\, \sigma_\beta.
    \end{align*}
Then for $\sigma_E\in(0,\infty)$ the state $e^{t\cL_{\sigma_E,\widehat f, H}}(\rho)$ can be prepared within $\eps$-trace distance on a quantum computer with $\mathcal{O}\left(|\cA|\log(M)\right)$ many qubits with total Hamiltonian simulation time 
    \begin{align*}
\widetilde{\mathcal{O}}\left( \,t\,\operatorname{poly}\left(|\mathcal{A}|\,,\, \log\left(\frac{\mathfrak{c}\,E_{\operatorname{Gibbs}}}{\epsilon}\right)\right)\right).
    \end{align*} 
Here, $\widetilde{\mathcal{O}}$ and $\widetilde\Theta$ treat all non-displayed parameters as constant
and further absorb  polylogarithmic dependencies in the leading order.
\end{theorem}
\begin{remark}
To illustrate the above theorem, let us consider an $m$-mode continuous variable system on the Hilbert space $(L^2(\R))^{\otimes m} \cong L^2(\R^m),$ choice of bare jumps being $\{A^\alpha\}_{\alpha\in\cA} \equiv \{a_i,a^\dagger_i\}_{i=1}^m$ and with local truncations in the Fock basis $a^{\le M}_i$ and $(a^{\le M}_i)^{\le M}$ being defined in \eqref{eq:DefTruncAnnCrea}. In this case all assumptions on the jump operators in the above theorem are naturally satisfied for the choice $\NA = N_{\operatorname{tot}} = \sum_{i=1}^m a^\dagger_ia_i$ at the relevant places in Section~\ref{sec:TruncateJumps}.
Furthermore, as noted in Remark~\ref{rem:CircuitTruncJumps}, the truncated jumps $a^{\le M}_i$ and $(a^{\le M}_i)^{\le M}$ can be implemented efficiently by a quantum circuit.

Moreover, for same choice and Hamiltonian $H$ being a bounded degree polynomial in the annihilation and creation operators, e.g. the Bose-Hubbard Hamiltonian, all relevant assumptions in the above theorem are satisfied as well as we have seen in Section~\ref{sec:FinitDimGenerator}. As seen in Remark~\ref{rem:CircuitTruncHam}, Hamiltonian simulation with respect to the truncated Hamiltonian $H_{\le M}$ can also be efficiently implemented on a quantum circuit.
\end{remark}

\begin{proof}[Proof of Theorem~\ref{thm:CircuitDynamicsSchwartz}]
By Theorem~\ref{thm:SchwartzGibbsDynamics} we know for $\sigma_E\in(0,\infty)$ fixed that
\begin{align*}
\left\|\left(e^{t\mathcal{L}_{\sigma_E,\widehat{f},H}}-e^{t\mathcal{L}^{\le M}_{\sigma_E,\widehat{f},H_{\le M}}}\right)(\rho)\right\|_1  \le \eps/3 
\end{align*}
for some $M$ satisfying \eqref{eq:MDimension}. 

Furthermore, by Proposition~\ref{prop:DiscreteGauss} we have 
\begin{align*}
\Big\|e^{t\mathcal{L}^{\le M}_{\sigma_E,\widehat{f},H_{\le M}}}-e^{t\mathcal{L}^{\le M,\mathfrak{n}}_{\sigma_E,\widehat{f},H_{\le M}}}\Big\|_{1\to 1}\le \epsilon/3
\end{align*}
for  $$\mathfrak{n} = \widetilde\Theta\left(\operatorname{poly}(M)+ \log\left(\frac{|\cA|t}{\eps}\right) \right) =\widetilde \Theta\left( \operatorname{poly}\left(\log\left(\frac{t\,\mathfrak{c}\,E_{\operatorname{Gibbs}}|\cA|}{\eps}\right)\right)\right).$$ The result follows by Theorem~\ref{thm:Lin}.

\end{proof}
As a direct consequence of Theorem~\ref{thm:CircuitDynamicsSchwartz} and Corollary~\ref{cor:SchwartzGibbs}, we find the following result on efficient Gibbs state preparation under the assumption of positive spectral gap.
\begin{corollary} [Gibbs state preparation for Schwartz filter functions]
\label{cor:GibbsCircuitSchwartz}

Under the same assumptions as in Theorem~\ref{thm:CircuitDynamicsSchwartz} and assuming additionally that  $L_{\sigma_E,\widehat{f},H},$ the self-adjoint generator on $\mathscr{T}_2(\cH)$ associated to the Lindbladian $\cL_{\sigma_E,\widehat f,H},$ has a positive spectral gap 
$\lambda_2 \equiv \operatorname{gap}\left(L_{\sigma_E,\widehat f,H}\right) >0,$ 
the Gibbs state of the Hamiltonian can be prepared within $\eps$-trace distance on a quantum computer with $\mathcal{O}(|\cA|\log(M))$ many qubits, for some 
\begin{align*}    
M =\widetilde\Theta\left(\operatorname{poly}\left(\log\left(\frac{\mathfrak{c}\,E_{\operatorname{Gibbs}}|\cA|}{\lambda_2 \eps}\right)\right)\right), 
\end{align*}
with total Hamiltonian simulation time 
\begin{align*}
\widetilde{\mathcal{O}}\left(\frac{1}{\lambda_2}\operatorname{poly}\left(|\mathcal{A}|\,,\,\log\left(\frac{\mathfrak{c}E_{\operatorname{Gibbs}}}{\epsilon}\right)\right)\right).
    \end{align*} 
Here, $\widetilde{\mathcal{O}}$ and $\widetilde\Theta$ treat all non-displayed parameters as constant
and further absorb  polylogarithmic dependencies in the leading order.
\end{corollary}

\subsubsection{Circuit implementation for singular filter functions}
Combining Proposition~\ref{prop:DiscreteGauss} and Theorem~\ref{thm:Lin} with Theorem~\ref{thm:SingularGibbsdynamics} yields the following result on efficient implementation of the Lindblad dynamics for Metropolis-type filter function \eqref{eq:filterFunction}. 
\begin{theorem}[Circuit implementation of $e^{t\cL_{\sigma_E,\hatfM,H}}$]
\label{thm:CircuitDynamicsSingular}
Let $\beta>0,$ $t\ge 0$ and $\hatfM$ be the Metropolis-type filter function defined in \eqref{eq:filterFunction} and $\hatfMdelta$ from \eqref{eq:hatfMdeltaeqIntro} with
\begin{align}
\label{eq:deltCircuit}
    \delta = \Theta\left(\frac{\eps}{\mathfrak{c}\,E'_{\operatorname{Gibbs}}\,|\cA|\,t} \right).
\end{align}
 Under the assumptions of Theorem~\ref{thm:SingularGibbsdynamics},~\ref{thm:Lin} and  Proposition~\ref{prop:DiscreteGauss} and assuming oracle access to\footnote{Analogously as before, $\fMdeltakappa$ is defined as the Fourier transform of the function $\nu\mapsto\hatfMdelta(\nu)\kappa(\nu)$ where $\kappa$ satisfies the assumptions of Proposition~\ref{prop:DiscreteGauss}.} $\fMdeltakappa$  and  $t_\kappa$, block encodings of the subnormalised bare jumps $(A^{\alpha})^{\le M}/q'(M)$, controlled Hamiltonian simulation $U_{H_{\le M}}$, where 
\begin{align}
\label{eq:MDimension1}
 M =\widetilde\Theta\left(\operatorname{poly}\left(\log\left(\frac{t\,\mathfrak{c}\,E'_{\operatorname{Gibbs}}|\cA|}{ \eps}\right)\right)\right)
    \end{align}
    and assume oracle access to a state preparation circuit for input state $\rho$ satisfying
    \begin{align*}
        \rho \le \mathfrak{c}\, \sigma_\beta.
    \end{align*}
Then for $\sigma_E\in(0,\infty)$ the state $e^{t\cL_{\sigma_E,\hatfM, H}}(\rho)$ can be prepared within $\eps$-trace distance on a quantum computer with $\mathcal{O}\left(|\cA|\log(M)\right)$ many qubits with total Hamiltonian simulation time 
    \begin{align*}
\widetilde{\mathcal{O}}\left( \,t\,\operatorname{poly}\left(|\mathcal{A}|\,,\, \log\left(\frac{\mathfrak{c}\,E'_{\operatorname{Gibbs}}}{\epsilon}\right)\right)\right).
    \end{align*} 
Here, $\widetilde{\mathcal{O}}$ and $\widetilde\Theta$ treat all non-displayed parameters as constant
and further absorb  polylogarithmic dependencies in the leading order.
\end{theorem}

\begin{proof}
By Theorem~\ref{thm:SingularGibbsdynamics} we know for $\sigma_E\in(0,\infty)$ fixed that
\begin{align*}
\left\|\left(e^{t\mathcal{L}_{\sigma_E,\hatfM,H}}-e^{t\mathcal{L}^{\le M}_{\sigma_E,\hatfMdelta,H_{\le M}}}\right)(\rho)\right\|_1  \le \eps/3 
\end{align*}
for some $\delta$ and $M$ satisfying \eqref{eq:deltCircuit} and \eqref{eq:MDimension1} respectively.
From Lemma~\ref{LemmGevreycond}, we know that the function $\hatfMdelta$ satisfies \Cref{condfgevrey} with $\xi_q,\xi_u,s,C_{1,\hatfMdelta}$ being constant in the relevant free parameters and, furthermore, by Lemma~\ref{lem:C_fMdeltaBounds} we see
\begin{align*}
\|\hatfMdelta\|_1 = \widetilde{\mathcal{O}}\left(\operatorname{poly}(\log(1/\delta)\right)
\end{align*}
Hence, we can apply  Proposition~\ref{prop:DiscreteGauss} which gives
\begin{align*}
\Big\|e^{t\mathcal{L}^{\le M}_{\sigma_E,\hatfMdelta,H_{\le M}}}-e^{t\mathcal{L}^{\le M,\mathfrak{n}}_{\sigma_E,\hatfMdelta,H_{\le M}}}\Big\|_{1\to 1}\le \epsilon/3
\end{align*}
for  $$\mathfrak{n} = \widetilde\Theta\left(\operatorname{poly}(M)+ \log\left(\frac{|\cA|t\  \operatorname{poly}\!\left(\log(1/\delta)\right)}{\eps}\right) \right) =\widetilde \Theta\left( \operatorname{poly}\left(\log\left(\frac{t\,\mathfrak{c}\,E'_{\operatorname{Gibbs}}|\cA|}{\eps}\right)\right)\right).$$ The result follows by Theorem~\ref{thm:Lin}.

\end{proof}

As a direct consequence of Theorem~\ref{thm:CircuitDynamicsSingular}, we find the following result on efficient Gibbs state preparation under the assumption of positive spectral gap.


\begin{corollary} [Gibbs state preparation for singular filter functions]
\label{cor:GibbsCircuitSingular}

Under the same assumptions as in Theorem~\ref{thm:CircuitDynamicsSingular} and assuming additionally that  $L_{\sigma_E,\hatfM,H},$ the self-adjoint generator on $\mathscr{T}_2(\cH)$ associated to the Lindbladian $\cL_{\sigma_E,\hatfM,H},$ has a positive spectral gap 
$\lambda_2 \equiv \operatorname{gap}\left(L_{\sigma_E,\hatfM,H}\right) >0,$ 
the Gibbs state of the Hamiltonian can be prepared within $\eps$-trace distance on a quantum computer with $\widetilde{\mathcal{O}}(|\cA|\log(M))$ many qubits, for some 
\begin{align*}    
M =\widetilde\Theta\left(\operatorname{poly}\left(\log\left(\frac{\mathfrak{c}\,E'_{\operatorname{Gibbs}}|\cA|}{\lambda_2 \eps}\right)\right)\right), 
\end{align*}
with total Hamiltonian simulation time 
\begin{align*}
\widetilde{\mathcal{O}}\left(\frac{1}{\lambda_2}\operatorname{poly}\left(|\mathcal{A}|\,,\,\log\left(\frac{\mathfrak{c}E'_{\operatorname{Gibbs}}}{\epsilon}\right)\right)\right).
    \end{align*} 
Here, $\widetilde{\mathcal{O}}$ and $\widetilde\Theta$ treat all non-displayed parameters as constant
and further absorb  polylogarithmic dependencies in the leading order.
\end{corollary}



\appendix



\section{\Cref{eq:condAalphas} for Schr\"odinger operators}
\label{appendixC}
In this section we prove Theorem \ref{thm:Schroedinger_operators}. To verify Condition \theconditionA \ for Schr\"odinger operators, we start by defining the generalized Sobolev spaces
\begin{definition}
Let $s,\sigma \ge 0$. We define
\[
\mathcal{H}^{s,\sigma}(\mathbb{R}^d)
:=
\Bigl\{
f\in L^2(\mathbb{R}^d)
:\;
\langle x\rangle^\sigma f \in L^2(\mathbb{R}^d),
\quad
\langle \xi\rangle^s \widehat{f}(\xi)\in L^2(\mathbb{R}^d)
\Bigr\},
\]
where $\langle x\rangle := (1+|x|^2)^{1/2}$ and $\widehat{f}$ denotes the Fourier transform of $f$.
\end{definition}

\begin{proposition}
\label{prop:interp}
Let $s,\sigma \ge 0$ and $0<\theta<1$. Then
\[
\bigl[L^2(\mathbb{R}^d),\mathcal{H}^{s,\sigma}(\mathbb{R}^d)\bigr]_\theta
=
\mathcal{H}^{\theta s,\theta\sigma}(\mathbb{R}^d)
\]
with equivalence of norms.
\end{proposition}

\begin{proof}
We write
\[
\mathcal{H}^{s,\sigma}
=
\{f\in L^2:\ \langle x\rangle^\sigma f\in L^2,\ \langle D\rangle^s f\in L^2\}
\]
with norm
\[
\|f\|_{\mathcal{H}^{s,\sigma}}^2
=
\|\langle x\rangle^\sigma f\|_{L^2}^2
+
\|\langle D\rangle^s f\|_{L^2}^2.
\]
Thus
\[
\mathcal{H}^{s,\sigma}
=
L^2(\langle x\rangle^{2\sigma}dx) \cap H^s(\mathbb{R}^d)
\]
with equivalent norms.

We use that complex interpolation preserves intersections of compatible Hilbert couples \cite{Peetre1974}
\[
[L^2(\mathbb R^d), \mathcal H^{s,0}(\mathbb R^d)\cap \mathcal H^{0,\sigma}(\mathbb R^d)]_\theta
=
[L^2(\mathbb R^d),\mathcal H^{s,0}(\mathbb R^d)]_\theta \cap [L^2(\mathbb R^d),\mathcal H^{0,\sigma}(\mathbb R^d)]_\theta.
\]

It is standard \cite[Thm. 5.5.3]{BerghLofstrom1976} that
\[
[L^2(\mathbb{R}^d), \mathcal H^{0,\sigma}(\mathbb R^d)]_\theta
=
\mathcal H^{0,\theta \sigma}(\mathbb R^d),
\]
and
\[
[L^2(\mathbb{R}^d),H^{s}(\mathbb{R}^d)]_\theta
=
H^{\theta s}(\mathbb{R}^d).
\]

Combining these two identities yields the result.
\end{proof}

\noindent We illustrate these assumptions, especially \eqref{eq:BoundBareJumpsWithHam}, for Schr\"odinger operators and $A^{\alpha}$ the standard creation and annihilation operators in the following Proposition. 

For positive self-adjoint operators one has \cite[Theo.4.17]{Lunardi2018} 
\begin{theorem}[Interpolation of domains for positive self-adjoint operators]
\label{thm:interpolation}
Let $\mathcal H$ be a Hilbert space and let $A \ge 0$ be a positive self-adjoint operator on $H$. 
For $s \ge 0$, define
\[
D(A^s) := \left\{ u \in \mathcal H : \int_0^\infty \lambda^{2s} \, d\|E_\lambda u\|^2 < \infty \right\},
\]
where $(E_\lambda)$ is the spectral resolution of $A$.

Then for $0 \le \alpha < \beta$ and $\theta \in (0,1)$,
\[
\bigl[ D(A^\alpha), \, D(A^\beta) \bigr]_\theta 
= D\bigl(A^{(1-\theta)\alpha + \theta \beta}\bigr),
\]
with equivalence of norms.
\end{theorem}

Let $D(\tilde H) = \mathcal H^{2,2}(\mathbb R^d).$ Then, we have by Theorem \ref{thm:interpolation} and Proposition \ref{prop:interp} that $D(\tilde H^{1/2}) = \mathcal H^{1,1}(\mathbb R^d).$ On the other hand, the creation and annihilation operators are continuous linear operators $a_j,a_j^{\dagger}:\mathcal H^{k,k}(\mathbb R^d) \to  \mathcal H^{k-1,k-1}(\mathbb R^d)$ for $k \in \mathbb N.$






\begin{lemma}[Domain of the quantum harmonic oscillator]
Let
\[
H_0:=-\Delta+|x|^2
\]
initially on Schwartz space \(\mathcal S(\mathbb R^d)\subset L^2(\mathbb R^d)\). Then \(H_0\)
is essentially self-adjoint, and its self-adjoint realization satisfies
\[
D(H_0)
=
\{u\in H^2(\mathbb R^d): |x|^2u\in L^2(\mathbb R^d)\}.
\]
Equivalently,
\[
D(H_0)
=
\{u\in L^2(\mathbb R^d): -\Delta u\in L^2(\mathbb R^d),\ |x|^2u\in L^2(\mathbb R^d)\}.
\]
Moreover, the graph norm of \(H_0\) is equivalent to
\[
u\mapsto \|u\|+\|\Delta u\|+\||x|^2u\|.
\]
\end{lemma}

\begin{proof}
For \(u\in \mathcal S(\mathbb R^d)\), the harmonic oscillator is symmetric and hence closable. We compute
\[
\|H_0u\|^2
=
\|-\Delta u+|x|^2u\|^2
=
\|\Delta u\|^2+\||x|^2u\|^2-2\Re\langle \Delta u, |x|^2u\rangle.
\]
By integration by parts,
\[
-\langle \Delta u, |x|^2u\rangle
=
\sum_{j=1}^d \langle \partial_j u,\partial_j(|x|^2u)\rangle
=
\sum_{j=1}^d \langle \partial_j u,2x_ju+|x|^2\partial_j u\rangle.
\]
Taking real parts gives
\[
-\Re\langle \Delta u, |x|^2u\rangle
=
\||x|\nabla u\|^2
+
2\Re\sum_{j=1}^d \langle \partial_j u,x_ju\rangle.
\]
Now
\[
2\Re\langle \partial_j u,x_ju\rangle
=
\int_{\mathbb R^d} x_j\,\partial_j(|u|^2)\,dx
=
-\int_{\mathbb R^d}|u|^2\,dx
=
-\|u\|^2,
\]
hence
\[
2\Re\sum_{j=1}^d \langle \partial_j u,x_ju\rangle
=
-d\|u\|^2.
\]
Therefore
\[
-\Re\langle \Delta u, |x|^2u\rangle
=
\||x|\nabla u\|^2-d\|u\|^2,
\]
and so
\[
\|H_0u\|^2
=
\|\Delta u\|^2+\||x|^2u\|^2+2\||x|\nabla u\|^2-2d\|u\|^2.
\]
In particular,
\[
\|\Delta u\|^2+\||x|^2u\|^2
\le
\|H_0u\|^2+2d\|u\|^2.
\]
Since also
\[
\|H_0u\|\le \|\Delta u\|+\||x|^2u\|,
\]
we obtain 
\[
\|u\|+\|H_0u\|
\text{ is equivalent to }
\|u\|+\|\Delta u\|+\||x|^2u\|
\text{ for }u\in\mathcal S(\mathbb R^d).
\]

Now let \(H_0\) denote the closure of the operator on \(\mathcal S(\mathbb R^d)\).
By the graph norm equivalence, \(u\in D(H_0)\) if and only if there exists a
sequence \(u_n\in\mathcal S(\mathbb R^d)\) such that
\[
u_n\to u,\qquad \Delta u_n\to v,\qquad |x|^2u_n\to w
\quad\text{in }L^2(\mathbb R^d)
\]
for some \(v,w\in L^2(\mathbb R^d)\). Passing to distributions shows that
\(v=\Delta u\) and \(w=|x|^2u\). Thus
\[
D(H_0)\subset
\{u\in L^2(\mathbb R^d): \Delta u\in L^2(\mathbb R^d),\ |x|^2u\in L^2(\mathbb R^d)\}.
\]

Conversely, if \(u\in L^2(\mathbb R^d)\) satisfies
\[
\Delta u\in L^2(\mathbb R^d),\qquad |x|^2u\in L^2(\mathbb R^d),
\]
then \(u\in H^2(\mathbb R^d)\). Choosing \(u_n\in \mathcal S(\mathbb R^d)\) with
\[
u_n\to u,\qquad \Delta u_n\to \Delta u,\qquad |x|^2u_n\to |x|^2u
\quad\text{in }L^2(\mathbb R^d),
\]
we obtain
\[
H_0u_n=-\Delta u_n+|x|^2u_n\to -\Delta u+|x|^2u
\quad\text{in }L^2(\mathbb R^d).
\]
Hence \(u\in D(H_0)\). Therefore
\[
D(H_0)=
\{u\in L^2(\mathbb R^d): \Delta u\in L^2(\mathbb R^d),\ |x|^2u\in L^2(\mathbb R^d)\}.
\]
The minimal operator $H_0$ is equal to the maximal operator, thus $H_0$ is the self-adjoint realization. 
Finally, since \(u,\Delta u\in L^2\) implies \(u\in H^2(\mathbb R^d)\), this is
equivalent to
\[
D(H_0)=
\{u\in H^2(\mathbb R^d): |x|^2u\in L^2(\mathbb R^d)\}.
\]

\end{proof}

It is well-known that 
\[ V \in L^{\max\{2,d/2\}}(\mathbb R^d)+ \langle x \rangle^{\alpha} L^{\infty}(\mathbb R^d)\text{ with } \alpha<2\]
is relatively zero-bounded by \cite[Theorem X.15]{reed1975ii} and \cite[Theorem 4.28]{Zworski2012} with respect to the harmonic oscillator $H_0$. Thus, Kato-Rellich's theorem \cite[Theorem X.12]{reed1975ii} shows the self-adjointness of $H_0+V$ on $\mathcal H^{2,2}(\mathbb R^d).$


\begin{proposition}
Let
\[
H=-\Delta+V(x)
\quad\text{on }L^2(\mathbb R^d),
\]
where $V\in C^\infty(\mathbb R^d)$ is real-valued and satisfies
\[
V(x)\ge c\langle x\rangle^r-C_0,
\qquad c>0,\quad r\ge 1,
\]
as well as
\[
|\partial_x^\alpha V(x)|\le C_\alpha \langle x\rangle^{r-|\alpha|},
\qquad \alpha\in\mathbb N^d.
\]
For $h_0 \ge C_0$, we have 
\[
\widetilde H:=H+(h_0+1)I.
\]
Let $a_j,a_j^{\dagger}$ be the annihilation and creation operators.
Then for every $n\in\mathbb N$,
\[
\widetilde H^n a_j \widetilde H^{-n-1},\widetilde H^n a_j^{\dagger} \widetilde H^{-n-1}\in \mathcal B(L^2(\mathbb R^d)).
\]
\end{proposition}

\begin{proof}
We work in the scattering calculus $\Psi^{m,l}_{\mathrm{sc}}(\mathbb R^d)$. Recall \cite[(1.3)]{HassellJiaSussman2025} that a smooth function $a(x,\xi)\in S^{m,l}_{\mathrm{sc}}$ if
\[
|D_x^\alpha D_\xi^\beta a(x,\xi)|
\le C_{\alpha\beta}\langle x\rangle^{l-|\alpha|}\langle \xi\rangle^{m-|\beta|}
\]
for all multi-indices $\alpha,\beta.$
Its quantization \cite[(1.1)]{HassellJiaSussman2025} defines $\Psi^{m,l}_{\mathrm{sc}}$.

We first show that
\[
\widetilde H\in \Psi^{2,r}_{\mathrm{sc}}.
\]
Indeed, its symbol is
\[
\widetilde{p}(x,\xi)=|\xi|^2+V(x)+\lambda.
\]
The term $|\xi|^2$ satisfies
\[
|D_x^\alpha D_\xi^\beta |\xi|^2|
\le C_{\alpha\beta}\langle \xi\rangle^{2-|\beta|}
\le C_{\alpha\beta}\langle x\rangle^{r-|\alpha|}\langle \xi\rangle^{2-|\beta|},
\]
while for $V(x)+\lambda$ we use the assumption on derivatives of $V$. 
By the composition property of the scattering calculus \cite[(2.1)]{HassellJiaSussman2025},
\[
\Psi^{m_1,l_1}_{\mathrm{sc}}\circ \Psi^{m_2,l_2}_{\mathrm{sc}}
\subset \Psi^{m_1+m_2,\;l_1+l_2}_{\mathrm{sc}}.
\]
Since $\widetilde H\in \Psi^{2,r}_{\mathrm{sc}}$, an induction gives
\begin{equation}\label{eq:Hn}
\widetilde H^n\in \Psi^{2n,rn}_{\mathrm{sc}}
\qquad\text{for all }n\in\mathbb N.
\end{equation}

Moreover, $\widetilde H$ (totally) elliptic, hence we have \cite[Prop. 2.1]{HassellJiaSussman2025},
\begin{equation}\label{eq:Hneg}
\widetilde H^{-n-1}\in \Psi^{-2n-2,-r(n+1)}_{\mathrm{sc}}.
\end{equation}

The operator
\[
a_j=\frac1{\sqrt2}(x_j+\partial_{x_j}) \text{ and }a_j^{\dagger}=\frac1{\sqrt2}(x_j-\partial_{x_j})
\]
have symbols
\[
b_j(x,\xi)=\frac1{\sqrt2}(x_j+i\xi_j) \text{ and }\overline{b_j}(x,\xi)=\frac1{\sqrt2}(x_j-i\xi_j).
\]
A direct inspection shows that for all $\alpha,\beta$,
\[
|D_x^\alpha D_\xi^\beta b_j(x,\xi)|
\le C_{\alpha\beta}\langle x\rangle^{1-|\alpha|}\langle \xi\rangle^{1-|\beta|},
\]
hence
\begin{equation}\label{eq:aj}
a_j,a_j^{\dagger}\in \Psi^{1,1}_{\mathrm{sc}}.
\end{equation}
Combining \eqref{eq:Hn}, \eqref{eq:aj}, and \eqref{eq:Hneg}, we obtain by the composition rule \cite[(2.1)]{HassellJiaSussman2025}
\[
\widetilde H^n a_j \widetilde H^{-n-1},\widetilde H^n a_j^{\dagger}
 \widetilde H^{-n-1}
\in
\Psi^{2n,rn}_{\mathrm{sc}}
\circ
\Psi^{1,1}_{\mathrm{sc}}
\circ
\Psi^{-2n-2,-r(n+1)}_{\mathrm{sc}} \subset \Psi^{-1,\;1-r}_{\mathrm{sc}}.
\]
Since $-1\le 0$ and $1-r\le 0$ (because $r\ge 1$), the standard boundedness
result for the scattering calculus \cite[Prop. 3.6]{HassellJiaSussman2025} implies that every operator in
$\Psi^{m,l}_{\mathrm{sc}}$ with $m\le 0$ and $l\le 0$ is bounded on $L^2(\mathbb R^d)$.
Therefore
\[
\widetilde H^n a_j \widetilde H^{-n-1},\widetilde H^n a_j^{\dagger} \widetilde H^{-n-1}\in \mathcal B(L^2(\mathbb R^d)).
\]

\end{proof}


\section{Auxiliary results to study general one-mode Hamiltonians}

\noindent The following Lemma is key in establishing Theorem \ref{thm:SpectralGapGeneralh(N)}:

 
\begin{lemma}
\label{lem:EnEquivalences}
Let $n_0\in\N_0$ and $(E_n)_{n\in\N_0}$ be a sequence of real numbers which is non-decreasing for all $n\ge n_0$. The following are equivalent:
\begin{enumerate}
     \item  \label{it:LowerBoundEnDiff}There exists $\delta>0$ and  $s\in\N$ such that for all $m\ge n_0$ we have $E_{m+s} - E_m\ge \delta.$
\item
\label{it:LowerBoundEnDiffstrong} For all $\delta>0$ there exists $s\in\N$ such that for all $m\ge n_0$ we have $E_{m+s} - E_m\ge \delta$.
    \item \label{it:FiniteSum}For all $\beta>0$ we have\,\,$\sup_{m} \sum_{j\ge m} e^{-\beta(E_{j}-E_m)} <\infty.$
    \item \label{it:FiniteSumGamma} For all $\beta>0$ there exists $\gamma \in(0,1)$ such that\,\, $\sup_{m} \sum_{j\ge m} \gamma^{-(j-m)}e^{-\beta(E_{j}-E_m)} <\infty.$
\end{enumerate}
Furthermore, if condition \eqref{it:LowerBoundEnDiff} holds true, we have for all $\gamma\in (e^{-\beta\delta/s},1]$ that
\begin{align}
\label{eq:ExplicitBoundGamma}
    \sup_{m,k\ge 0}\sum_{j\ge m} \gamma^{-(j-m)}e^{-\beta(E_{j+k} - E_{m+k}+E_{j} - E_m)/2}  \le 
(n_0+1)\gamma^{-n_0} e^{\beta \Delta_E} \frac{e^{\beta\delta}}{1-q}  <\infty,
\end{align}
where we denoted $q:=\gamma^{-1}e^{-\beta\delta/s}<1$ and $\Delta_E :=\max_{0\le j,m\le 2n_0} |E_j-E_m|.$

\end{lemma}

\begin{proof}
It is obvious that  \eqref{it:LowerBoundEnDiffstrong} $\implies$ \eqref{it:LowerBoundEnDiff} and also that \eqref{it:FiniteSumGamma} $\implies$ \eqref{it:FiniteSum}. Hence, we focus in the following on the non-trivial directions. First, we assume \eqref{it:LowerBoundEnDiff} and show that this implies \eqref{it:LowerBoundEnDiffstrong}: Since the $E_n$ are non-decreasing for all $n\ge n_0$, we have for all $m\in\N_0$ such that $m\ge n_0$ and $j\ge m$ that
\begin{align*}
    E_{j}- E_m &\ge E_{m+s\lfloor(j-m)/s}\rfloor- E_m =  \sum_{k=1}^{\lfloor\frac{j-m}{s}\rfloor} E_{m+ks} - E_{m+(k-1)s} \ge \lfloor\frac{j-m}{s}\rfloor \delta.
\end{align*}
This already shows \eqref{it:LowerBoundEnDiffstrong}: for any $\delta'>0$, choose $s'$ such that $\lfloor s'/s\rfloor \delta\ge \delta'$, and hence for all $m$, the above shows that $E_{m+s'}-E_m\ge \delta'$.

Next, we prove \eqref{eq:ExplicitBoundGamma} from condition \eqref{it:LowerBoundEnDiff}:
For $\gamma \in (e^{-\beta\delta/s},1]$ we denote $q = \gamma^{-1} e^{-\beta\delta/s} <1$ and see by the above that 
\begin{align}
\label{eq:GammaSumBound}
   &\sup_{\substack{m\ge n_0\\k\ge 0}} \sum_{j\ge m} \gamma^{-(j-m)}e^{-\beta(E_{j+k} -E_{m+k} +E_j- E_m)/2}  \le \sum_{j=0}^\infty \gamma^{-j}e^{-\beta\lfloor\frac{j}{s}\rfloor\delta}  \le 
    e^{\beta\delta}\sum_{j=0}^\infty q^j = \frac{e^{\beta\delta}}{1-q} <\infty,
\end{align}
where for the second inequality, we have used that $\lfloor\frac{j}{s}\rfloor \ge \frac{j}{s} -1$.

Furthermore, for $m<n_0$ and $k\ge0$, we have
\begin{align}
\label{eq:FiniteSumGamma}
    \sum_{j=m}^{n_0-1} \gamma^{-(j-m)}e^{-\beta(E_{j+k}-E_{m+k}+E_j-E_m)/2} &\le n_0\max_{m\le j\le n_0-1}\gamma^{-(j-m)}e^{-\beta(E_{j+k}-E_{m+k}+E_j-E_m)/2} \\
    \nn&
   \le n_0\gamma^{-n_0}\max_{0\le j,m\le 2n_0}\left\{e^{-\beta(E_{j}-E_{m})},1\right\} 
\end{align}
where we have used the fact that 
\begin{equation}
\begin{split}
\label{eq:Maaan}
\sup_{k\ge 0}e^{-\beta(E_{j +k}-E_{m+k})/2} \le \max_{0\le k\le n_0}\left\{e^{-\beta(E_{j +k}-E_{m+k})/2},1\right\} \le \max_{0\le j',m'\le 2n_0}\left\{e^{-\beta(E_{j}-E_{m})/2},1\right\}
\end{split}
\end{equation}
where the first inequality follows from the fact that for $k\ge n_0$ we have $E_{j+k}\ge E_{m+k}$ as $j\ge m$, and furthermore that $j,m\le n_0.$
Moreover, we see 
\begin{align}
\label{eq:SeriesGamma}
    \sum_{j\ge n_0} \gamma^{-(j-m)}   e^{-\beta(E_{j+k}-E_{m+k}+E_j-E_m)/2}  &
      = e^{-\beta(E_{n_0+k}-E_{m+k}+E_{n_0}-E_m)/2}\gamma^{-(n_0-m)}\nn
      \\
      &\quad\nn\times\sum_{j\ge n_0}   \gamma^{-(j-n_0)}e^{-\beta(E_{j+k}-E_{n_0+k}+E_j-E_{n_0})/2} \\
      &\le \nn e^{-\beta(E_{n_0+k}-E_{m+k}+E_{n_0}-E_m)/2}\gamma^{-(n_0-m)}\frac{e^{\beta\delta}}{1-q}\\&\le \gamma^{-n_0}\max_{0\le j,m\le 2n_0}\left\{e^{-\beta(E_{j}-E_{m})},1\right\}\frac{e^{\beta\delta}}{1-q},
\end{align}
where, for the second to last inequality, we have used \eqref{eq:GammaSumBound} and for the last \eqref{eq:Maaan} for $j=n_0.$

Combining \eqref{eq:GammaSumBound}, \eqref{eq:FiniteSumGamma}, and \eqref{eq:SeriesGamma}, we obtain
\begin{align*}
     \sup_{m,k\ge 0}\sum_{j\ge m} \gamma^{-(j-m)}e^{-\beta(E_{j+k} - E_{m+k}+E_{j} - E_m)/2} &\le 
(n_0+1)\gamma^{-n_0} \max_{0\le j,m\le 2n_0}\left\{ e^{-\beta(E_j-E_m)} ,1\right\}\frac{e^{\beta\delta}}{1-q}  
\\&\le (n_0+1)\gamma^{-n_0}  e^{\beta\Delta_E} \frac{e^{\beta\delta}}{1-q}
<\infty,  
\end{align*}
which shows that \eqref{eq:ExplicitBoundGamma} 
holds true. Furthermore, \eqref{it:FiniteSumGamma} immediately follows by restricting to $k=0.$


We finish the proof by showing that \eqref{it:FiniteSum} implies \eqref{it:LowerBoundEnDiffstrong}: Assume \eqref{it:LowerBoundEnDiffstrong} is not satisfied, i.e., that there exists $\delta>0$ such that for all $s\in\N$ there exists $m_s\ge n_0$ such that $E_{m_s+s}- E_{m_s}<\delta.$ \bin{For $s$ large enough we can without loss of generality\footnote{To see this assume that there exists $\delta>$ such that for all $s\in\N$ there exists $m_s<n_0$ satisfying $E_{m_s+s}-E_{m_s}<\delta.$ But this implies that the subsequence $\left(E_{m_s+s}\right)_{s\in \N}$ is uniformly bounded as $E_{m_s+s}<\delta +\max_{m<n_0}E_m$. From this we immediately see that \eqref{it:FiniteSum} cannot hold true.} assume that $m_s\ge n_0$ and } Since the sequence $(E_n)_{n\in\N_0}$ is non-decreasing for $n\ge n_0$, we therefore have that $E_{j}- E_{m_s}<\delta$ for all $m_s\le j\le m_s+s$ and consequently for $\beta>0$ that
\begin{align*}
    \sum_{j\ge m_s} e^{-\beta(E_j-E_{m_s})} > s e^{-\beta\delta}.
\end{align*}
Since $s\in \N$ was arbitrary, we see
\begin{align*}
     \sup_{m\in\N}\sum_{j\ge m} e^{-\beta(E_j-E_m)} = \infty,
\end{align*}
which finishes the proof.

\end{proof}





\section{Auxiliary results on the filter function $\hatfMdelta$}
\label{sec:fmDeltaProofs}
In this section we prove Lemma~\ref{lem:C_fMdeltaBounds} and~\ref{LemmGevreycond}. For that we first state and prove two supporting lemmas.
\begin{lemma}
\label{lem:AaaahFm} 
    Denote the closed strip of the complex plane $
\overline S:=\{z\in\mathbb C:\ |\Im z|\le \tfrac12\}$ and the function $h_1(z):= \exp\left(-\frac{\sqrt{1+z^2} +z}{4}\right),$ where $\sqrt{1+z^2}$ denotes the principal square root. Then we have 
\begin{align}
    \sup_{z\in\overline S}|h_1(z)|\le 1.
\end{align}
\end{lemma}
\begin{proof}
Let $z\in\overline S$ and 
$
w:=1+z^2=1+u^2-v^2+2iuv.$
For the principal square root we use the known relation
\begin{align}
\label{eq:PrincSquareRoot}
\Re\sqrt{w}=\sqrt{\frac{|w|+\Re w}{2}}.
\end{align}
From a direct computation we see $
|w|^2-(u^2-1+v^2)^2=4u^2\ge 0,
$ and hence
$
|w|\ge |u^2-1+v^2|,
$
which gives 
$
|w|+\Re w \ge  |u^2-1+v^2| + 1 + u^2 -v^2 \ge 2u^2$
Hence, using \eqref{eq:PrincSquareRoot} we see
\[
\Re(\sqrt{1+z^2}+z)=\Re\sqrt{w}+u\ge |u|+u\ge 0,
\]
and thus
\[
|h_1(z)|
=
\exp\!\left(-\frac{\Re(\sqrt{1+z^2}+z)}{4}\right)
\le 1.
\]
\end{proof}


\begin{lemma}
\label{lem:RealPartBig}
Denoting the closed strip of the complex plane $
\overline S:=\{z\in\mathbb C:\ |\Im z|\le \tfrac12\},
$ we have for all $\theta\in[0,1/2)$ that there exists a constant  $C_\theta\ge 0$ such that for all $z\in \overline S$ with $|\Re z| \ge  C_\theta$ we have
\begin{align}
\label{eq:LowerBoundRealPart}
\Re\left(e^{(1+z^2)^\theta}\right) \ge \frac{1}{2} \exp\left(\frac{|\Re z|^{2\theta}}{2}\right).
\end{align}
\end{lemma}
\begin{proof}
We use
\begin{align}
\label{eq:RealExponential}
\Re\left(e^{(1+z^2)^\theta}\right)  = e^{\Re((1+z^2)^\theta)} \cos(\Im((1+z^2)^\theta)).
\end{align}
and show in the following for $z\in \overline S$ that $\Im((1+z^2)^\theta)$ is small and $\Re((1+z^2))^\theta)$ is large for $|\Re z|$ large.
For that we write 
\begin{align*}
    R(z) &:= |(1+z^2)|, \qquad \phi(z) := \operatorname{arg}(1+z^2) 
\end{align*}
and hence 
\begin{align*}
    (1+z^2)^\theta = R^\theta(\nu) e^{i\theta \phi(z)} = R^{\theta}(z)\Big(\cos(\theta \phi(z)) + i\sin(\theta\phi(z))\Big)
\end{align*}
 We see using $|\Im z|\le 1/2$ and a direct computation that $\Re(1+z^2)\ge (\Re\nu)^2\ge 0$ and $|\Im(1+z^2)| \le |\Re z|$ and therefore
\begin{align}
\label{eq:phi(nu)UpperBound}
    \left|\phi(z)\right| = \left|\arctan\left(\frac{\Im\left(1+z^2\right)}{\Re\left(1+z^2\right)}\right)\right| \le \frac{1}{|\Re z|},
\end{align}
where we used $|\!\arctan(x)|\le |x|.$  By a direct computation we see for $|\Re z|\ge 1$ that \begin{align}
\label{eq:R(nu)bound}
  (\Re z)^2 \le R(z) \le 3(\Re z)^2
\end{align} and using $|\!\sin(x)| \le |x|$ together with \eqref{eq:phi(nu)UpperBound}, we get
\begin{align*}
   \left| \Im\left(((1+z^2)^\theta\right)\right| = R^\theta(z)|\sin(\theta\phi(z))|\le  \theta 3^\theta \, |\Re z|^{2\theta -1} \xrightarrow[|\Re z|\to\infty]{} 0 
\end{align*}
as $\theta<1/2.$ Therefore, we can pick $C_\theta\ge 0$ such that for all $|\Re z|\ge C_\theta$ we have
\begin{align}
\label{eq:CosImLowerBound}
\cos\left(\Im\left((1+z^2)^\theta\right)\right) \ge \frac{1}{2}.
\end{align}
Furthermore, using \eqref{eq:R(nu)bound} and \eqref{eq:phi(nu)UpperBound} again, we see, after possibly increasing $C_\theta,$ that for $|\Re z| \ge C_\theta$ we have
\begin{align}
    \Re((1+z^2)^\theta = R^{\theta}(z) \cos\left(\theta \phi(z)\right) \ge \frac{1}{2}|\Re z|^{2\theta}. 
\end{align}
Combining this with \eqref{eq:RealExponential} and  \eqref{eq:CosImLowerBound} finishes the proof.
\end{proof}

\begin{proof} [Proof of Lemma~\ref{lem:C_fMdeltaBounds}]

We start by realising that $\nu\to \hatfMdelta(\nu)$ is analytic\footnote{Here, we choose the branch such that $\sqrt{1+(\beta\nu)^2}$ and $(1+(\beta\nu)^2)^{\theta}$ are positive for $\nu\in\R.$} on the strip in the complex plane with $|\Im(\nu)|<1/\beta.$ Hence, for $a:=\frac{1}{2\beta}>0$ we can shift the contour of integration in $\fMdelta(t) = \frac{1}{2\pi}\int_\R \hatfMdelta(\nu)e^{-it\nu}d\nu$ for $t>0$ as $\nu\mapsto \nu - ia$ and for $t<0$ as $\nu\mapsto \nu+ia$ yielding 
\begin{align}
\label{eq:ExpDecayfMdelta(t)}
    |\fMdelta(t)| \le e^{-a|t|} \int_\R \left|\hatfMdelta(\nu-i\operatorname{sgn}(t) a)\right|d\nu.
\end{align}
We focus on bounding the integral on the right hand side.
From Lemma~\ref{lem:AaaahFm} we have $|\hatfM(\nu-i\operatorname{sgn}(t) a)| \le 1$
and therefore 
\begin{align}
\label{eq:RemainderIntegral}
   \nn \int_\R \left|\hatfMdelta(\nu-i\operatorname{sgn}(t)a)\right|d\nu &\le \int_\R \exp\left(-\delta\,\Re\left(e^{((1+(\beta(\nu-i\operatorname{sgn}(t)a))^2)^\theta}\right)\right)d\nu\\ &=\frac{1}{\beta}\int_\R \exp\left(-\delta\,\Re\left(e^{((1+(\nu-i\operatorname{sgn}(t)/2))^2)^\theta}\right)\right)d\nu
\end{align}
where we have used that $|e^z| = e^{\Re(z)}.$ 
We use Lemma~\ref{lem:RealPartBig} to see that there exists a constant $\nu_0(\theta)\ge 0$ only depending on $\theta$ such that for all $|\nu|\ge \nu_0(\theta)$ we have
\begin{align}
\label{eq:LowerBoundRealPart2}
\Re\left(e^{(1+(\nu-i\operatorname{sgn}(t)/2)^2)^\theta}\right) \ge \frac{1}{2} \exp\left(\frac{|\nu|^{2\theta}}{2}\right).
\end{align}
 Using \eqref{eq:LowerBoundRealPart2} with \eqref{eq:RemainderIntegral}, we get
\begin{align*}
    &\beta\int_\R \left|\hatfMdelta(\nu-i\operatorname{sgn}(t)/2)\right|d\nu \le 2\nu_0(\theta)+ \int_{|\nu|\ge \nu_0(\theta)} \exp\left(-\frac{\delta}{2}\exp\left(\frac{|\nu|^{2\theta}}{2}\right)\right)d\nu\\&\le 2\nu_0(\theta)+ (2\log\left(1/\delta)\right))^{\tfrac{1}{2\theta}} + \int_{\left\{|\nu|\ge (2\log\left(1/\delta)\right))^{1/(2\theta)}\right\}} \exp\left(-\frac{\delta}{2}\exp\left(\frac{|\nu|^{2\theta}}{2}\right)\right)d\nu.
\end{align*}
For the remainder of the proof, we bound the integral on the right-hand side of the last inequality. 

For $\nu>0$  consider the change of variable $y = \delta\exp\left(\nu^{2\theta}/2\right)$ which gives
\begin{align*}
  &\int_{\left\{|\nu|\ge (2\log\left(1/\delta)\right))^{1/(2\theta)}\right\}} \exp\left(-\frac{\delta}{2}\exp\left(\frac{|\nu|^{2\theta}}{2}\right)\right)d\nu =2\int^{\infty}_{(2\log\left(1/\delta)\right))^{1/(2\theta)}} \exp\left(-\frac{\delta}{2}\exp\left(\frac{\nu^{2\theta}}{2}\right)\right)d\nu\\&=\frac{2}{\theta}\int^{\infty}_{1} \frac{(2\log(y/\delta))^{\tfrac{1-2\theta}{2\theta}}}{y}\,e^{-y/2}dy  \le C_\theta\,\left( (\log(1/\delta))^{\tfrac{1-2\theta}{2\theta}} +1\right)
\end{align*}
for some finite constant $C_\theta\ge 0.$ Using \eqref{eq:ExpDecayfMdelta(t)} this shows \eqref{eq:fMdeltaPointWise}. Furthermore, using $\kappa\in(0,1/2)$ gives \eqref{eq:CfMdeltaBound}.

\end{proof}

\begin{proof}[Proof of Lemma~\ref{LemmGevreycond}]
Denote $$h_{\delta,\theta}(\nu):= \hatfMdelta(\nu/\beta) = \exp\left(-\tfrac{\sqrt{1+\nu^2}+\nu}{4}\right)\,
\exp\!\left(-\delta e^{(1+\nu^2)^\theta}\right). $$
We prove that \(h_{\delta,\theta}\) extends holomorphically and boundedly to a fixed complex strip, and then apply Cauchy's estimate. Set
\[
S_1:=\{z\in\mathbb C:\ |\Im z|<\tfrac 23\}.
\]
Write \(z=u+iv\). Then for \(z\in S\) we have $
\Re(1+z^2)=1+u^2-v^2 > 5/9$. In particular, the function $S_1\to \C,\ z\to 1+z^2$ does not meet the branch cut \((-\infty,0]\) and hence the principal branches
$
\sqrt{1+z^2}$
$(1+z^2)^\theta$
are holomorphic on \(S_1\). Therefore
\[
h_{\delta,\theta}(z)
=
\exp\!\left(-\frac{\sqrt{1+z^2}+z}{4}\right)\,
\exp\!\left(-\delta e^{(1+z^2)^\theta}\right)
\]
is holomorphic on \(S_1\). Next we show that \(h_{\delta,\theta}\) is bounded on the closed strip
\[
\overline S:=\{z\in\mathbb C:\ |\Im z|\le \tfrac12\}.
\]
From Lemma~\ref{lem:AaaahFm} we know that the first factor
\(
h_1(z):=\exp\!\left(-\frac{\sqrt{1+z^2}+z}{4}\right)
\)
satisfies 
$
|h_1(z)|\le 1$ for all $z\in \overline S$.

Now consider the second factor
$
h_2(z):=\exp\!\left(-\delta e^{(1+z^2)^\theta}\right).
$
From Lemma~\ref{lem:RealPartBig} we know that there exists $C_\theta$ such that for all $|\Re z| \ge C_\theta$ we have
\[
|h_2(z)|
=
\exp\!\left(-\delta\,\Re\!\left(e^{(1+z^2)^\theta}\right)\right)
\le
\exp\!\left(-\tfrac{\delta}{2}e^{\frac{|\Re z|^{2\theta}}{2}}\right) \le 1.
\]
As $h_2$ is continuous, we also have that $h_2(z)$ is uniformly bounded for $|\Re z|\le C_\theta$ which in summary gives $\sup_{z\in \overline S} |h_2(z)|<\infty.$ We have hence shown
\[
\xi_{\theta}:=\sup_{z\in \overline S}|h_{\delta,\theta}(z)|<\infty,
\]
where $\xi_{\theta}$ depends on $\theta$ but is independent on $\delta.$
Note that for all \(x\in\mathbb R\) the closed disc
$
B_{1/4}(x) := \left\{z\in\C\, :\, |z|\le 1/4 \right\}$
lies inside the open strip \(S\). Since \(h_{\delta,\theta}\) is holomorphic on \(S\), Cauchy's estimate gives
\[
|h_{\delta,\theta}^{(n)}(x)|
\le
4^n n!
\sup_{|z-x|\le \frac14}|h_{\delta,\theta}(z)|
\le
\xi_{\theta}\,(4n)^n.
\]
Taking the supremum over \(x\in\mathbb R\), we get
\[
\|h_{\delta,\theta}^{(n)}\|_{L^\infty(\mathbb R)}
\le \xi_{\theta}\ (4n)^n.
\]
Therefore, we have shown \(h_{\delta,\theta}\in \mathcal G^1_{\xi_{\theta},4}\), as claimed.

Next, we prove $\frac{d}{d\nu} h_{\delta,\theta} \in L^1(\R).$ We use the product rule $\frac{d}{d\nu} h_{\delta,\theta} = h_2 \frac{d}{d\nu}h_1 + h_1\frac{d}{d\nu}h_2$ and treat term individually. For the first term we use $|h_2(\nu)|\le 1$ for $\nu\in\R$ and furthermore \cite[Lemma 28]{ding2025efficient} which gives 
\begin{align*}
   \left\|h_2 \frac{d}{d\nu}h_1 \right\|_1 \le \left\|\frac{d}{d\nu}h_1 \right\|_1 \le c _1,
\end{align*}
for some $c_1\ge 0,$ which is, just as $h_1,$ independent of $\theta$ and $\delta.$ 
For the second part, assume without loss of generality $\delta>0$ and $\theta>0$ as otherwise $\frac{d}{d\nu} h_2=0,$ and note
\begin{align*}
    \frac{d}{d\nu} h_2(\nu) = -\frac{2\delta \theta \, \nu}{(1+\nu^2)^{1-\theta}} e^{(1+\nu^2)^\theta }\exp\left(-\delta e^{(1+\nu^2)^\theta}\right),
\end{align*}
from which we see
\begin{align*}
  \left| \frac{d}{d\nu} h_2(\nu)\right| = \begin{cases}
   -\frac{d}{d\nu} h_2(\nu),\quad\text{for} \, \nu\ge 0,\\
   \frac{d}{d\nu} h_2(\nu),\qquad\text{for} \, \nu<0.
   \end{cases} 
\end{align*}
Using further $|h_1(\nu)|\le 1,$ we finish the proof by noting
\begin{align*}
    \left\|h_1\frac{d}{d\nu} h_2\right\|_1 \le \left\|\frac{d}{d\nu} h_2\right\|_1 = -\int_0^\infty \frac{d}{d\nu} h_2(\nu) d\nu + \int_{-\infty}^0 \frac{d}{d\nu} h_2(\nu) d\nu = 2 h_2(0) = 2 e^{-\delta e}\le 2,
\end{align*}
where in the second to last equality we have used that $h_2(\nu)\xrightarrow[|\nu|\to \infty]{}0.$
\end{proof}


\bin{By standard arguments\footnote{To see this note $|\fMdelta(t)| \le \frac{1}{2\pi}\|\hatfMdelta\|_1.$  Furthermore, by  partially integrating twice $\int \hatfMdelta(\nu) e^{-it\nu}d\nu = \tfrac{-1}{t^2} \int \hatfMdelta''(\nu) e^{-it\nu}d\nu$, where we used that the boundary terms vanish as $\delta>0,$ and therefore $|\fMdelta(t)| \le \tfrac{1}{t^2}\|\hatfMdelta''\|_1.$ Combining both, we see for $R>0$ that $\|\fMdelta\|_1 \le \frac{1}{2\pi}\left(\int_{|t|\le R} \|\hatfMdelta\|_1 dt + \int_{|t|> R} \tfrac{\|\hatfMdelta''\|_1 }{t^2}dt\right) = \tfrac{1}{\pi}\left(R\|\hatfMdelta\|_1+\tfrac{\|\hatfMdelta''\|_1}{R}\right).$ Choosing $R = \sqrt{\|\hatfMdelta''\|_1/\|\hatfMdelta\|_1}$ yields \eqref{eq:StandardL1Bound}.}
\begin{align}
\label{eq:StandardL1Bound}
\left\|\fMdelta\right\|_1 \le \frac{2}{\pi}\sqrt{\left\|\hatfMdelta\right\|_1,\left\|\hatfMdelta''\right\|_1}
\end{align}}









\bibliographystyle{unsrtnat}
\bibliography{ref}

\end{document}
